
\documentclass[11pt,a4paper]{amsart}

\usepackage{amsmath,amssymb,amsthm,mathtools}
\usepackage[T1]{fontenc}
\usepackage[utf8]{inputenc}
\usepackage{lmodern}
\usepackage{microtype}
\usepackage[hidelinks]{hyperref}
\usepackage{orcidlink}
\usepackage{enumitem}
\usepackage[margin=1in]{geometry}
\raggedbottom

\theoremstyle{plain}
\newtheorem{theorem}{Theorem}[section]
\newtheorem{proposition}[theorem]{Proposition}
\newtheorem{lemma}[theorem]{Lemma}
\newtheorem{corollary}[theorem]{Corollary}

\theoremstyle{definition}
\newtheorem{definition}[theorem]{Definition}
\newtheorem{assumption}[theorem]{Assumption}
\newtheorem{condition}[theorem]{Condition}

\theoremstyle{remark}
\newtheorem{remark}[theorem]{Remark}

\DeclareMathOperator{\Var}{Var}
\DeclareMathOperator{\Cov}{Cov}
\DeclareMathOperator{\spr}{spr}
\DeclareMathOperator{\diag}{diag}
\DeclareMathOperator{\dist}{dist}
\DeclareMathOperator{\He}{He}
\newcommand{\R}{\mathbb{R}}
\newcommand{\N}{\mathbb{N}}
\newcommand{\PP}{\mathbb{P}}
\newcommand{\QQ}{\mathbb{Q}}
\newcommand{\EE}{\mathbb{E}}
\newcommand{\cA}{\mathcal{A}}
\newcommand{\cB}{\mathcal{B}}
\newcommand{\cC}{\mathcal{C}}
\newcommand{\cD}{\mathcal{D}}
\newcommand{\cF}{\mathcal{F}}
\newcommand{\cG}{\mathcal{G}}
\newcommand{\cH}{\mathcal{H}}
\newcommand{\cI}{\mathcal{I}}
\newcommand{\cJ}{\mathcal{J}}
\newcommand{\cK}{\mathcal{K}}
\newcommand{\cE}{\mathcal{E}}
\newcommand{\cL}{\mathcal{L}}
\newcommand{\cM}{\mathcal{M}}
\newcommand{\cP}{\mathcal{P}}
\newcommand{\cQ}{\mathcal{Q}}
\newcommand{\cR}{\mathcal{R}}
\newcommand{\cS}{\mathcal{S}}
\newcommand{\cV}{\mathcal{V}}
\newcommand{\cX}{\mathcal{X}}
\newcommand{\one}{\mathbf{1}}
\newcommand{\dd}{\,d}

% Preload AMS symbol fonts so microtype expansion is configured before first page use.
\AtBeginDocument{\begingroup\setbox0=\hbox{$\square\mathbb{R}$}\endgroup}

\begin{document}

\title[Risk-Neutral Compression]{%
Risk-Neutral Compression and the Geometry of Latent Market Stress}
\author[J.~Vidal Llaurad\'o]{J.~Vidal Llaurad\'o\,\orcidlink{0009-0002-1918-5458}}
\thanks{DOI: \href{https://doi.org/10.2139/ssrn.6680938}{10.2139/ssrn.6680938}.}
\email{vidal.llaurado@gmail.com}
\subjclass[2020]{60G22, 91G20, 91G80, 62M15}
\keywords{risk-neutral compression, latent market stress, rough volatility, incomplete markets, option-surface geometry, state-price densities, market-variance risk, visible hedging}
\date{April 2026}

\begin{abstract}
Public option surfaces can price and quote local surface coefficients without determining the
latent market-stress loading that moves non-public claims.  This paper formulates that boundary as
risk-neutral compression.  For a square-integrable claim or local option-surface summary \(X\), the
\(\QQ\)-priced structural loading is
\(L_X^\QQ=C_{\cH_\varpi}^{\QQ}(X)\).  A public information set
\(\cH\subseteq\cG_T\) fixes only
\(C_{\cH}^{\QQ}(L_X^\QQ)\).  The main theorem identifies this conditional projection as the common
source of public pricing summaries, option-surface transfers, matched signed \(\PP/\QQ\) coefficient
comparisons, nonlinear variance-compression gaps, and visible-hedging residuals.  Calibration of a
surface therefore leaves latent-state recovery and tradeable hedgeability as separate conditions.
Residual components survive as orthogonal projection errors unless the specified public information
or hedge span completes the priced structural loading.  The SPX/NDX
OptionMetrics exercise, rebuilt from raw Oxford--Man realized-volatility data and external option
quotes, estimates the finite public projection: implied-volatility level, total-variance level, and
signed downside total-variance skew align with the rebuilt Perron stress direction in full-sample
and rolling-window designs.  Companion~A supplies the Gaussian source-family estimates used by the
option-surface specialization.
\end{abstract}

\maketitle

\section{Introduction}\label{sec:intro}

Public option surfaces are public price functionals.  They can determine a price, a skew
coefficient, or a local implied-volatility derivative without determining the latent loading that
prices market-variance risk.  The paper treats option surfaces, premium comparisons, variance
claims, and visible hedge classes as public images of one object, the \(\QQ\)-priced structural
loading.

The latent input is the synchronized market-variance state
\((\varpi,\cH_\varpi,\Delta_{Q_M})\).  The market-scale Volterra--Perron construction derives this
state from directed contagion and Perron synchronization \cite{vidal2026market}.  The preceding
smoothing, dynamic, and projected-score papers separate pricing visibility, path-space
detectability, observable screening, and feasible \(\PP\)-side inference
\cite{vidal2026latent,vidal2026dynamic,vidal2026econometric}.  The present paper changes the
measure and changes the question.  It asks which part of the resulting priced loading is selected
by public option and hedge information under \(\QQ\).

The first boundary is measure-theoretic.  The directed Volterra kernels, rough exponents, signed
edge support, and Perron mode of the normalized synchronization matrix are fixed by the primitive
geometry.  Equivalent pricing reweights conditional expectations.  It does not make public
screened edge selections, visibility weights, compressed loadings, or claim-specific projections
measure invariant unless a separate projection-invariance assumption is imposed.  The
Perron-generated reduction from \cite{vidal2026market} supplies the sigma-field and factor on which
risk-neutral compression acts.  Write
\[
        (\varpi,\cH_\varpi,\Delta_{Q_M})
        \quad\longmapsto\quad
        L_X^\QQ=C_{\cH_\varpi}^{\QQ}(X)
        \quad\longmapsto\quad
        C_{\cH}^{\QQ}(L_X^\QQ),
        \qquad
        \cH\subseteq\cG_T.
\]
The market-scale construction of \cite{vidal2026market} derives the synchronized aggregate state.
This paper prices its public images.

The pricing layer is fixed before the empirical layer.  Prices are represented by a
state-price density; visible pricing summaries are \(\QQ\)-measurable public summaries; visible
hedge classes are closed subspaces of \(L^2(\QQ,\cG_T)\).  The martingale-pricing background is the
state-price-density and equivalent-martingale-measure tradition
\cite{harrisonKreps1979,harrisonPliska1981,merton1973,delbaen1994,karatzasShreve1998}.  The
projection step uses the Hilbert-space geometry common to local-risk-minimization,
mean-variance hedging, and incomplete-information hedging
\cite{follmerSondermann1986,follmerSchweizer1991,schweizer1992meanVariance,
schweizer1995minimal,choulliVandaeleVanmaele2010}.  Conditional expectation is the standard
\(L^2\) projection geometry \cite{williams1991,karatzasShreve1991}.  Partial-information
investment and information-based pricing provide boundary references for the visible/latent split,
not inputs to the theorem \cite{lakner1995partial,brodyHughstonMacrina2008}.

The governing operator is \(C_{\cH}^{\QQ}\).  A claim's priced structural loading is visible,
screened while still priced through \(\cH_\varpi\), or absent according to
\(\|C_{\cH}^{\QQ}(L_X^\QQ)\|_{L^2(\QQ)}\) and \(\|L_X^\QQ\|_{L^2(\QQ)}\).  A centered aggregate-factor
component is then separated inside \(L_X^\QQ\) relative to the market-variance benchmark selected
by \(\cH_\varpi\).  Public summaries classify the full priced loading.  The centered factor enters
only when nonlinear compression gaps or visible-hedging residuals expose it.

The same loading then moves through the market objects studied below.  Equivalent pricing preserves
the mode system and the derived sigma-field.  Compression and projection theorems make public
pricing and hedging objects factor through \(L_X^\QQ\).  Option sections turn that projection into
surface coefficients, with signed downside total-variance skew as the tail-facing empirical
coordinate.  The premium section compares matched signed \(\PP/\QQ\) projected coefficients on a
common basis.  The variance and hedging sections record the convexity gaps and projection residuals
left by the same operator.

\subsection*{Main results}
Fix the synchronized Perron state.  A claim or local option summary has under \(\QQ\) the priced
structural loading
\[
        L_X^\QQ=C_{\cH_\varpi}^{\QQ}(X),
\]
and public information \(\cH\subseteq\cG_T\) determines only
\[
        C_{\cH}^{\QQ}(L_X^\QQ).
\]
This is the paper's common object.  The first results prove that, once the Perron-generated
reduction and the equivalent pricing measure are fixed, every treated public pricing or hedging
summary is a compression, projection, or measurable transformation of \(L_X^\QQ\).  The
option-transfer block then states the additional quote-level conditions under which this compressed
loading appears as a skew, price-surface, or implied-volatility coefficient.  Companion~A supplies
the primitive Volterra estimates that reduce the Gaussian option-facing loading to the scalar seam
used by the local option-factor statement.

The other branches are not parallel theorem copies.  The premium branch compares matched signed
\(\PP/\QQ\) projected coefficients on a common source-screened basis and uses the comparison for
support classification.  The variance branch identifies the convexity gaps generated by the
aggregate factor and by recentering defects.  The hedging branch expresses visible incompleteness as
the distance from \(L_X^\QQ\) to a visible hedge space.

The governing results are Theorems~\ref{thm:q-invariance},
\ref{thm:main-risk-neutral-compression}, \ref{thm:option-abstract}, \ref{thm:market-skew},
\ref{thm:calibration-to-premium}, \ref{thm:q-compression}, and~\ref{thm:q-hedging}.  The appendices
supply the endpoint calculations: option-transfer variants in
Appendix~\ref{sec:option-transfer-variants}, premium estimator and support details in
Appendix~\ref{sec:premium-implementation-appendix}, endpoint obstructions in
Appendix~\ref{sec:detailed-scope-gates}, and Gaussian source-family estimates from Companion~A,
summarized in Propositions~\ref{thm:main-benchmark-model-to-seam},
\ref{thm:main-common-channel-headline}, and~\ref{thm:main-gaussian-option-factor}.

The Perron-market construction supplies the latent stress coordinate and visible channel geometry.
The present analysis identifies its risk-neutral public projections.

\subsection*{Organization}
Sections~\ref{sec:model}--\ref{sec:class} define the Perron-generated reduction, prove invariance
of the primitive geometry under equivalent pricing, define \(L_X^\QQ\), and classify public
visibility by the norm of the compressed loading.  Theorem~\ref{thm:main-risk-neutral-compression}
is the common projection identity.

Section~\ref{sec:qsetup} gives the option-facing pricing setup, the abstract first-variation
statement, and the cumulant-class transfer from compressed loading to quoted surface coefficient.
Section~\ref{sec:option-transfer-main} records the option consequence used empirically.  The
Hermite, primitive non-Gaussian, rare-jump, Gaussian-Volterra, Bachelier, and Black calculations are
in Appendices~\ref{sec:option-transfer-variants}, \ref{sec:skew}, and~\ref{sec:black}; Companion~A
supplies the Gaussian source-family verification \cite{vidal2026gaussianseams}.

Section~\ref{sec:empirical} reports the SPX/NDX finite-projection evidence.  Section~\ref{sec:premium}
compares harmonized signed \(\PP/\QQ\) projected coefficients on a cleaned matched derivative
surface.  Sections~\ref{sec:variance} and~\ref{sec:hedging} treat nonlinear variance compression
and visible hedge residuals.  Section~\ref{sec:scope-gates} states the endpoint scope, and
Appendix~\ref{sec:detailed-scope-gates} records the detailed obstructions.  Section~\ref{sec:conclusion}
closes the argument.


\section{Latent market geometry under \texorpdfstring{\(\PP\)}{P} and \texorpdfstring{\(\QQ\)}{Q}}\label{sec:model}

Let \(N\ge2\) be fixed and let \(W=(W^1,\dots,W^N)\) be an \(N\)-dimensional Brownian motion with
independent coordinates on a filtered probability space
\((\Omega,\cF,(\cF_t)_{t\ge0},\PP)\).
The Volterra kernel specification used below follows the rough-volatility and affine-Volterra
lineage \cite{bayer2016pricing,gatheral2018volatility,abijaber2019affine,eleuch2019roughheston},
with broader survey and calibration-context references in
\cite{dinunno2023survey,abijaberLi2025practice}.  The screened Perron construction is taken
from \cite{vidal2026market}; the Perron--Frobenius terminology follows the standard
nonnegative-matrix background \cite{hornJohnson2012,bermanPlemmons1994}.

\begin{assumption}[Directed Volterra market under \(\PP\)]\label{ass:model}
For each \(1\le i,j\le N\), there exist exponent \(H_{ij}\in(0,1)\), amplitude
\(\beta_{ij}\in\R\), and continuous modulation \(a_{ij}\colon [0,T]^2\to\R\) such that for
\(0\le s<t\le T\),
\[
        K_{ij}(t,s)
        =
        \beta_{ij}\,a_{ij}(t,s)\,(t-s)^{H_{ij}-1/2}.
\]
Moreover \(a_{ij}(t,t)>0\) whenever \(\beta_{ij}\neq0\), and the drift \(t\mapsto \mu_i^\PP(t)\)
is deterministic and continuous.  The latent volatility channel of asset \(i\) is
\[
        V_i(t)
        =
        \mu_i^\PP(t)
        +
        \sum_{j=1}^N \int_0^t K_{ij}(t,s)\,\dd W_s^j.
\]
Fix deterministic weights \(w\in(0,\infty)^N\) and define the aggregate market-volatility process
\[
        M(t):=\sum_{i=1}^N w_i V_i(t).
\]
\end{assumption}

The latent contagion geometry is encoded by the signed edge matrix and the screened
synchronization matrix of \cite{vidal2026market}.  Fix a reference time
\(t_0\in(0,T]\), screening weights \(\pi_{ij}\in[0,1]\), and a market-dominant active set
\(\cE_M\subseteq\{1,\dots,N\}^2\).  Define the structural signed edge matrix
\[
        B_{ij}
        :=
        \beta_{ij}a_{ij}(t_0,t_0)\one_{\{(i,j)\in \cE_M\}},
        \qquad
        1\le i,j\le N.
\]
The nonnegative screened synchronization matrix is
\[
        A_{ij}:=|B_{ij}|\,\pi_{ij}.
\]
The weights \(\pi_{ij}\) are normalized visibility weights for the local stress-envelope
experiment.  They are part of the observational projection, not primitive edge signs.
Throughout, work on the irreducible aggregate regime inherited from
\cite{vidal2026market}.  In particular, \(A\) is assumed irreducible; let
\(v\in(0,\infty)^N\) be its Perron right eigenvector and let \(\spr(A)\) denote its spectral
radius.

This separates structural and observational edge objects.  The structural edge set
\[
        E^{\mathrm{str}}
        :=
        \{(i,j):\beta_{ij}a_{ij}(t_0,t_0)\ne0\}
\]
is part of the primitive Volterra geometry.  Physical screened selections
\(E^{\mathrm{vis},\PP}\) and weights \(\pi_{ij}^{\PP}\), and risk-neutral or option-surface
projected selections \(E^{\mathrm{vis},\QQ}\) and weights \(\pi_{ij}^{\QQ}\), are observational
objects.  Equivalent pricing preserves the primitive directed Volterra geometry under the stated
representation assumptions.  It does not, without an additional projection-invariance assumption,
preserve finite-sample screened edge selections or empirical visibility weights.

\begin{assumption}[Equivalent square-integrable pricing measure]\label{ass:q}
There exists a predictable \(\R^N\)-valued process \(\theta\) such that
\[
        \int_0^T \|\theta_s\|^2\,\dd s < \infty
        \qquad \PP\text{-a.s.},
\]
and the stochastic exponential
\[
        Z_t
        :=
        \exp\!\left(
                -\int_0^t \theta_s^\top \dd W_s
                - \frac12 \int_0^t \|\theta_s\|^2\,\dd s
        \right)
\]
is a strictly positive martingale.  Define \(\QQ\) by
\[
        \frac{\dd\QQ}{\dd\PP}=Z_T.
\]
All claims considered below are assumed to belong to \(L^2(\QQ)\).  When the notation is used for
discounted no-arbitrage prices of traded assets, those discounted traded assets are
\(\QQ\)-local martingales.
\end{assumption}

\begin{condition}[Optional bounded-density \(L^2\)-comparability]\label{ass:q-bounds}
The state-price density satisfies
\[
        0<\underline z \le Z_T \le \overline z<\infty
        \qquad \PP\text{-a.s.}
\]
for some deterministic constants \(\underline z,\overline z\).
This optional condition is invoked only in comparability statements that identify \(L^2(\PP)\)
and \(L^2(\QQ)\) setwise.  The main risk-neutral compression results are stated natively on
\(L^2(\QQ)\) and do not require this bounded-density condition.
\end{condition}

\begin{assumption}[Aggregate Perron sufficient-state reduction]\label{ass:aggregate-perron-bridge}
For the aggregate claim families treated here, the market-scale construction of \cite{vidal2026market}
supplies a Perron-projected amplitude \(\varpi\), a visible sigma-field \(\cG_T\), a derived
structural sigma-field \(\cH_\varpi\), and, for integrated market variance, a centered structural
increment \(\Delta_{Q_M}\).  These objects satisfy the aggregate Perron sufficient-state reduction
\[
        (\varpi,\cH_\varpi,\Delta_{Q_M})
        \quad\text{is sufficient for the aggregate claim families considered below.}
\]
The analysis below studies the \(\QQ\)-priced public compressions generated by this reduction.
\end{assumption}

\begin{definition}[Aggregate Perron structural sigma-field]\label{def:lambda-structure}
Assume the aggregate claim families treated here satisfy the Perron-projected
sufficient-state reduction of \cite{vidal2026market}.  Let \(u,v\in(0,\infty)^N\) be the left and
right Perron eigenvectors of the irreducible screened synchronization matrix, normalized so that
\[
        u^\top v=1.
\]
Let \(x\) denote the aggregate latent envelope from the market-scale construction of
\cite{vidal2026market} and define the
canonically projected aggregate amplitude by
\[
        \varpi(t):=u^\top x(t),
        \qquad
        0\le t\le T.
\]
Assume
\[
        \int_0^T \varpi(t)^2\,\dd t \in L^2(\PP,\cF_T)\cap L^2(\QQ,\cF_T).
\]
Set
\[
        \cL_T:=\sigma(\varpi(t):0\le t\le T),
        \qquad
        \cH_\varpi:=\cG_T\vee \cL_T.
\]
Thus \(\cH_\varpi\) is the derived aggregate structural sigma-field inherited from
\cite{vidal2026market} for the aggregate claim classes treated here; for integrated market
variance, the same reduction selects the centered structural increment
\[
        \Delta_{Q_M}
        =
        C_{\cH_\varpi}(Q_M)-C_{\cG_T}(Q_M).
\]
\cite{vidal2026market} derives these objects from primitive aggregate coupling, Perron projection,
and off-Perron aggregate screening, and illustrates the reduction empirically on the Oxford-Man
realized-measure panel \cite{heber2009realized,barndorffNielsen2008realized,shephardSheppard2010heavy}.
\end{definition}


\begin{remark}[Aggregate reduction, not a fresh existence claim]
Definition~\ref{def:lambda-structure} imports the aggregate-first reduction established in
\cite{vidal2026market}: the Perron-projected amplitude \(\varpi=u^\top x\), the derived sigma-field
\(\cH_\varpi\), the market-variance increment \(\Delta_{Q_M}\), and the claim families for which
that state is sufficient.  The present argument starts after that reduction and identifies the
corresponding \(\QQ\)-priced loadings and their public compressions, with no new Perron-factor
existence assumption inside this manuscript.
\end{remark}



\begin{definition}[Pricing, visible summaries, and risk-neutral compression]\label{def:compression}
Let \(\cG_t\subseteq \cF_t\) be the visible market filtration.  For a discounted claim
\(X\in L^2(\QQ,\cF_T)\), define the pricing operator
\[
        \pi_t(X):=\EE_{\QQ}[X\mid \cF_t].
\]
For a date-\(t\) public information set \(\cH_t\subseteq\cG_t\), define the public date-\(t\)
compression by
\[
        C_{t,\cH}^{\QQ}(X):=\EE_{\QQ}[X\mid\cH_t].
\]
For every sub-\(\sigma\)-field \(\cH\subseteq \cG_T\), define the risk-neutral compression of
\(X\) onto \(\cH\) by
\[
        C^\QQ_{\cH}(X):=\EE_{\QQ}[X\mid \cH].
\]
A \emph{visible pricing summary} is any \(\mathfrak S\in L^2(\QQ,\cG_T)\), and a
\emph{visible hedge class} is any closed linear subspace \(\cS\subseteq L^2(\QQ,\cG_T)\).
The time-zero option calculations below use the terminal notation \(C_{\cH}^{\QQ}\) after
suppressing the observation date.  The empirical panel is the repeated-date analogue of the same
local conditional-pricing operation, with each option surface measurable with respect to the public
information available on its observation date.
\end{definition}

\begin{definition}[Derived sufficient structural sigma-field and structural claim spaces]\label{def:mode-space}
Under Definition~\ref{def:lambda-structure}, define the \(\QQ\)-side structural claim space
\[
        \cM_\QQ:=L^2(\QQ,\cH_\varpi),
\]
and the corresponding \(\PP\)-side structural claim space
\[
        \cM_\PP:=L^2(\PP,\cH_\varpi).
\]
\end{definition}

\begin{remark}[Scope of the pricing layer]
The state-price density gives an arbitrage-free pricing operator on discounted claims, but the
analysis remains incomplete-market in scope.  It does not seek a representative-agent
foundation or a market-completion theorem.  This is the same martingale-valuation separation
between pricing measures and market completeness that underlies
\cite{harrisonKreps1979,delbaen1994,karatzasShreve1998}.  The pricing layer identifies which
parts of latent market stress remain publicly visible under \(\QQ\) and which survive only through
compression gaps and residual hedge risk.
\end{remark}

\section{Structural geometry and measure change}\label{sec:geom}

Equivalent pricing preserves the primitive Volterra--Perron geometry but not the conditional
projection operators.  The first theorem fixes that separation before any visibility or hedging
claim is made.

\begin{theorem}[Geometry invariance under equivalent pricing]\label{thm:q-invariance}
Assume Assumptions~\ref{ass:model} and~\ref{ass:q}, together with
Definition~\ref{def:lambda-structure}.  Then there exists an
\(N\)-dimensional Brownian motion \(W^\QQ\) on \((\Omega,\cF,(\cF_t),\QQ)\) such that for every
asset \(i\),
\[
       V_i(t)
       =
       \mu_i^\QQ(t)
       +
       \sum_{j=1}^N \int_0^t K_{ij}(t,s)\,\dd W_s^{\QQ,j},
\]
where the possibly random Volterra drift shift is
\[
       \mu_i^\QQ(t)
       :=
       \mu_i^\PP(t)
       -
       \sum_{j=1}^N \int_0^t K_{ij}(t,s)\,\theta_s^j\,\dd s.
\]
It is deterministic, or of finite variation as a function of \(t\), only under additional
regularity restrictions on the Girsanov kernel and the Volterra kernels; no subsequent structural
statement uses such restrictions.
In particular:
\begin{enumerate}[label=(\roman*)]
\item the kernels \(K_{ij}\), the exponents \(H_{ij}\), the structural active set \(\cE_M\), the
signed edge matrix \(B\), the normalized synchronization matrix \(A\), its spectral radius
\(\spr(A)\), and its Perron mode \(v\) are the same under \(\PP\) and \(\QQ\), subject to the
declared representation and normalization assumptions;
\item the structural mode sigma-field \(\cH_\varpi=\cG_T\vee \cL_T\) is the same under
\(\PP\) and \(\QQ\); consequently the mode classes \(\cM_\PP\) and \(\cM_\QQ\) are built from
the same underlying information, although they need not coincide as sets without additional
integrability or comparability conditions;
\item the structural projections \(C_{\cH_\varpi}^{\PP}\) and \(C_{\cH_\varpi}^{\QQ}\) act on
the same mode-generating sigma-field, but they remain measure-dependent.  Thus equivalent pricing
preserves the latent geometry while potentially changing its public projection.
\end{enumerate}
\end{theorem}

\begin{proof}
By Girsanov's theorem \cite{karatzasShreve1991},
\[
        W_t^\QQ:=W_t+\int_0^t \theta_s\,\dd s
\]
is an \(N\)-dimensional Brownian motion under \(\QQ\).  Substituting
\(\dd W_s=\dd W_s^\QQ-\theta_s\,\dd s\) into the Volterra representation gives
\[
        V_i(t)
        =
        \mu_i^\PP(t)
        +
        \sum_{j=1}^N \int_0^t K_{ij}(t,s)\,\dd W_s^{\QQ,j}
        -
        \sum_{j=1}^N \int_0^t K_{ij}(t,s)\,\theta_s^j\,\dd s,
\]
which is the claimed representation.  Since the kernels themselves are unchanged, so are the
exponents, the active set, the screened operator, and the Perron system built from them.

Because the aggregate Perron amplitude \(\varpi=u^\top x\) is fixed by the underlying
measurable structure rather than by either probability measure, its path sigma-field
\(\cL_T=\sigma(\varpi(t):0\le t\le T)\) is unchanged by the equivalent measure change.  Hence the
sigma-field \(\cH_\varpi=\cG_T\vee\cL_T\) is the same measurable object under \(\PP\) and
\(\QQ\), and the spaces \(\cM_\PP\) and \(\cM_\QQ\) are simply the respective \(L^2\)-classes on
that common sigma-field.  The last statement follows because conditional expectation is determined
by both the sub-\(\sigma\)-field and the measure.  The sub-\(\sigma\)-field is common, but the
\(L^2(\PP)\)- and \(L^2(\QQ)\)-inner products, and therefore the corresponding orthogonal
projections, may differ.
\end{proof}

\begin{proposition}[Optional \(L^2\)-comparability under bounded density]\label{prop:q-bounds}
Under Assumption~\ref{ass:q} and Condition~\ref{ass:q-bounds}, the spaces \(L^2(\PP)\) and \(L^2(\QQ)\)
coincide as sets and have equivalent norms.  Consequently \(\cM_\PP\) and \(\cM_\QQ\) also
coincide as sets and have equivalent norms.
\end{proposition}

\begin{proof}
Let \(X\) be measurable.  If \(X\in L^2(\PP)\), then Condition~\ref{ass:q-bounds} gives
\[
        \EE_{\QQ}[X^2]
        =
        \EE_{\PP}[Z_T X^2]
        \le
        \overline z\,\EE_{\PP}[X^2]
        <\infty.
\]
Conversely, if \(X\in L^2(\QQ)\), then
\[
        \EE_{\PP}[X^2]
        =
        \EE_{\QQ}\!\left[\frac{X^2}{Z_T}\right]
        \le
        \underline z^{-1}\EE_{\QQ}[X^2]
        <\infty.
\]
So \(L^2(\PP)=L^2(\QQ)\) as sets.  The two norms are equivalent because
\[
        \underline z\,\EE_{\PP}[X^2]
        \le
        \EE_{\QQ}[X^2]
        \le
        \overline z\,\EE_{\PP}[X^2]
\]
whenever \(X\in L^2(\PP)=L^2(\QQ)\).  Restricting to \((\cG_T\vee\cL_T)\)-measurable random
variables gives the same conclusion for \(\cM_\PP\) and \(\cM_\QQ\).
\end{proof}

\begin{corollary}[Same geometry, different visible projections]\label{cor:same-geometry}
Equivalent pricing preserves the structural mode system and the mode-generating sigma-field
\(\cH_\varpi\), but not its public visibility.  The core measure-invariant objects are the
kernels, the screened Perron construction, and \(\cH_\varpi\).  The measure-dependent objects are
the visible projections \(C_{\cG_T}^{\PP}\) and \(C_{\cG_T}^{\QQ}\) together with the structural
projections \(C_{\cH_\varpi}^{\PP}\) and \(C_{\cH_\varpi}^{\QQ}\).  If
Condition~\ref{ass:q-bounds} also holds, one may additionally identify \(\cM_\PP\) and
\(\cM_\QQ\) setwise by Proposition~\ref{prop:q-bounds}.
\end{corollary}


\begin{proof}
Theorem~\ref{thm:q-invariance} fixes the measurable sigma-field
\(\cH_\varpi\) and the kernel-Perron objects under the equivalent change of measure.  It does not
identify the conditional-expectation operator.  For any sub-\(\sigma\)-field
\(\cA\subseteq\cH_\varpi\), Bayes' formula gives, whenever the expression is defined,
\[
        C_{\cA}^{\QQ}(Y)
        =
        \frac{\EE_\PP[Z_TY\mid\cA]}{\EE_\PP[Z_T\mid\cA]},
\]
whereas \(C_{\cA}^{\PP}(Y)=\EE_\PP[Y\mid\cA]\).  The range sigma-field is the same, but the
projection may change with the state-price density.  Proposition~\ref{prop:q-bounds} gives the
stronger optional setwise identification of the corresponding \(L^2\)-classes under bounded-density
comparability.
\end{proof}

\begin{remark}[Same latent geometry]
The phrase means the same kernel system, the same screened Perron construction, and the same
mode-generating sigma-field.  Under the optional Condition~\ref{ass:q-bounds}, the \(\PP\)- and
\(\QQ\)-side mode classes may also be identified setwise.  It does not mean that the \(\PP\)- and
\(\QQ\)-conditional projections of a fixed claim agree.  The distinction is exactly that they need
not.
\end{remark}

\section{Risk-neutral compression and projection}\label{sec:compression}

The compression problem begins after the structural sigma-field has been fixed.  The map
\(C_{\cH_\varpi}^{\QQ}\) extracts the priced structural loading, and the smaller map
\(C_{\cH}^{\QQ}\) records the part admitted by public information.

\begin{proposition}[Optional projection stability under bounded-density comparability]\label{prop:proj-stability}
Under Assumption~\ref{ass:q} and Condition~\ref{ass:q-bounds}, fix a sub-\(\sigma\)-field
\(\cH\subseteq \cH_\varpi\).  Then \(L^2(\PP,\cH)=L^2(\QQ,\cH)\) as sets, and the operators
\[
        C_{\cH}^{\PP}\colon \cM_\PP=\cM_\QQ\to L^2(\PP,\cH),
        \qquad
        C_{\cH}^{\QQ}\colon \cM_\PP=\cM_\QQ\to L^2(\QQ,\cH)
\]
are bounded linear projections on the same ambient claim class.  For every
\(X\in \cM_\PP=\cM_\QQ\),
\[
        \|C_{\cH}^{\PP}X\|_{L^2(\QQ)}
        \le
        \Bigl(\frac{\overline z}{\underline z}\Bigr)^{1/2}
        \|X\|_{L^2(\QQ)},
\]
\[
        \|C_{\cH}^{\QQ}X\|_{L^2(\PP)}
        \le
        \Bigl(\frac{\overline z}{\underline z}\Bigr)^{1/2}
        \|X\|_{L^2(\PP)}.
\]
Hence the \(\QQ\)-side visible projection is a continuous reweighting of the same structural span,
not a new geometric object.  This is an optional comparability statement; the compression
theorems use the native \(L^2(\QQ)\) Hilbert space and do not rely on this bounded-density
identification.
\end{proposition}

\begin{proof}
By Proposition~\ref{prop:q-bounds}, \(L^2(\PP)=L^2(\QQ)\) as sets with equivalent norms, and the
same holds after restricting to \(\cH\)-measurable random variables.  Since conditional
expectations are orthogonal projections in their native Hilbert norms,
\[
        \|C_{\cH}^{\PP}X\|_{L^2(\PP)}\le \|X\|_{L^2(\PP)},
        \qquad
        \|C_{\cH}^{\QQ}X\|_{L^2(\QQ)}\le \|X\|_{L^2(\QQ)}.
\]
Applying the norm equivalence from Proposition~\ref{prop:q-bounds} to the projected terms gives the
displayed bounds.  The final statement is the interpretation of the fact that the range
\(L^2(\cH)\) is fixed while the projection operator changes continuously with the weighted norm.
\end{proof}

\begin{proposition}[Supporting structural projection of claims]\label{thm:claim-decomp}
Under Assumptions~\ref{ass:model} and~\ref{ass:q}, together with
Definition~\ref{def:mode-space}, every claim \(X\in L^2(\QQ,\cF_T)\) admits a
unique orthogonal decomposition
\[
        X=L_X^\QQ+U_X^\QQ,
\]
where
\[
        L_X^\QQ:=C_{\cH_\varpi}^\QQ(X)\in \cM_\QQ,
        \qquad
        U_X^\QQ:=X-C_{\cH_\varpi}^\QQ(X)\perp_{L^2(\QQ)} \cM_\QQ.
\]
Equivalently,
\[
        \EE_\QQ[U_X^\QQ\mid \cH_\varpi]=0.
\]
For every public sub-\(\sigma\)-field \(\cH\subseteq \cG_T\),
\[
        C_{\cH}^\QQ(X)=C_{\cH}^\QQ(L_X^\QQ).
\]
If \(L_X^\QQ\neq0\), define
\[
        \alpha_X^\QQ:=\|L_X^\QQ\|_{L^2(\QQ)},
        \qquad
        \Lambda_X:=\frac{L_X^\QQ}{\alpha_X^\QQ};
\]
if \(L_X^\QQ=0\), set \(\alpha_X^\QQ:=0\) and \(\Lambda_X:=0\).  The term \(L_X^\QQ\) is the
claim's \emph{priced structural loading under \(\QQ\)}.
\end{proposition}

The full loading and the residual loadings have different jobs.  For every public
sub-\(\sigma\)-field \(\cH\subseteq\cG_T\) and every closed tradeable span
\(\cS\subseteq L^2(\QQ,\cH)\), write
\[
        R_{X,\cH}^\QQ
        :=
        L_X^\QQ-C_{\cH}^\QQ(L_X^\QQ),
        \qquad
        R_{X,\cS}^\QQ
        :=
        L_X^\QQ-C_{\cS}^\QQ(L_X^\QQ).
\]
Public compression concerns \(C_{\cH}^\QQ(L_X^\QQ)\).  Price intervals and convex residual
charges concern \(R_{X,\cH}^\QQ\).  Visible hedge errors concern \(R_{X,\cS}^\QQ\).

\begin{proof}
Because \(\cM_\QQ=L^2(\QQ,\cH_\varpi)\) is a closed subspace of \(L^2(\QQ,\cF_T)\), the Hilbert
projection theorem yields the unique orthogonal projection \(P_{\cM_\QQ}^\QQ(X)\).  Conditional
expectation onto \(\cH_\varpi\) is exactly this projection, so
\[
        L_X^\QQ=C_{\cH_\varpi}^\QQ(X)
\]
and \(U_X^\QQ=X-L_X^\QQ\) is orthogonal to \(\cM_\QQ\), equivalently
\(\EE_\QQ[U_X^\QQ\mid \cH_\varpi]=0\).  Since every public sigma-field
\(\cH\subseteq \cG_T\subseteq \cH_\varpi\), the tower property gives
\[
        C_{\cH}^\QQ(X)
        =
        C_{\cH}^\QQ\bigl(C_{\cH_\varpi}^\QQ(X)\bigr)
        =
        C_{\cH}^\QQ(L_X^\QQ).
\]
The remaining statements are notation.
\end{proof}

\begin{proposition}[Approximate sufficiency of the Perron-generated aggregate structural sigma-field]\label{prop:approx-suff}
Let \(X\in L^2(\QQ,\cF_T)\), and suppose that
\[
        X=X_\varpi+R_X,
\]
where \(X_\varpi\in L^2(\QQ,\cH_\varpi)\).  Let
\[
        L_X^\QQ:=C_{\cH_\varpi}^\QQ(X)
\]
be the priced structural loading from Proposition~\ref{thm:claim-decomp}.  Then
\[
        L_X^\QQ
        =
        X_\varpi + C_{\cH_\varpi}^\QQ(R_X).
\]
Consequently,
\[
        \|L_X^\QQ-X_\varpi\|_{L^2(\QQ)}
        \le
        \|R_X\|_{L^2(\QQ)}.
\]
Moreover, for every public sub-\(\sigma\)-field \(\cH\subseteq \cG_T\),
\[
        \|C_{\cH}^\QQ(L_X^\QQ)-C_{\cH}^\QQ(X_\varpi)\|_{L^2(\QQ)}
        \le
        \|R_X\|_{L^2(\QQ)},
\]
and for every visible hedge class \(\cS\subseteq L^2(\QQ,\cG_T)\),
\[
        \|P_{\cS}^\QQ(L_X^\QQ)-P_{\cS}^\QQ(X_\varpi)\|_{L^2(\QQ)}
        \le
        \|R_X\|_{L^2(\QQ)}.
\]
Thus, if \(\|R_X\|_{L^2(\QQ)}\le \delta_X\), replacing fuller aggregate structural content by its
\(\cH_\varpi\)-measurable Perron-generated component perturbs structural projections, public
compressions, and visible projections by at most \(\delta_X\).
\end{proposition}

\begin{proof}
Because \(X_\varpi\) is \(\cH_\varpi\)-measurable,
\[
        L_X^\QQ
        =
        C_{\cH_\varpi}^\QQ(X_\varpi+R_X)
        =
        X_\varpi+C_{\cH_\varpi}^\QQ(R_X),
 \]
which gives the first display.  Since conditional expectation is a contraction on \(L^2(\QQ)\),
\[
        \|L_X^\QQ-X_\varpi\|_{L^2(\QQ)}
        =
        \|C_{\cH_\varpi}^\QQ(R_X)\|_{L^2(\QQ)}
        \le
        \|R_X\|_{L^2(\QQ)}.
\]
Applying the contractions \(C_{\cH}^\QQ\) and \(P_\cS^\QQ\) to \(L_X^\QQ-X_\varpi\) yields the
remaining bounds.
\end{proof}

\begin{remark}[Economic reading of \(\cH_\varpi\) for aggregate claim families]\label{rem:approx-suff}
Proposition~\ref{prop:approx-suff} is deterministic.  If the off-Perron aggregate remainder is
small in \(L^2(\QQ)\), then the Perron-generated sigma-field \(\cH_\varpi\) is sufficient for the
structural reduction and approximately sufficient in operator norm for the aggregate claim families
inherited from \cite{vidal2026market}.  Read together with the Oxford-Man appendix of
\cite{vidal2026market}, where the aggregate off-Perron share is empirically reported to be small for
the market-stress proxy treated there
\cite{heber2009realized,barndorffNielsen2008realized,shephardSheppard2010heavy}, the sigma-field
\(\cH_\varpi\) has a specific economic interpretation.  For the aggregate claims studied in the
market-scale analysis of \cite{vidal2026market} and in the present analysis, it is the empirically
dominant enlargement in that reported design.
\end{remark}

\begin{definition}[Canonical centered aggregate factor direction]\label{def:agg-factor}
For the aggregate claim families inherited from \cite{vidal2026market}, let
\[
        \Lambda_M^\star\in L^2(\QQ,\cH_\varpi)
\]
denote the centered aggregate factor direction obtained by normalizing the canonical centered
structural increment of integrated market variance under the present pricing measure.  Thus
\[
        \EE_\QQ[\Lambda_M^\star\mid \cG_T]=0,
        \qquad
        \|\Lambda_M^\star\|_{L^2(\QQ)}=1
\]
whenever the centered increment is nonzero; if that centered increment vanishes, set
\(\Lambda_M^\star:=0\).  Define the centered structural subspace
\[
        \cM_\QQ^\circ
        :=
        \{Y\in L^2(\QQ,\cH_\varpi):\EE_\QQ[Y\mid \cG_T]=0\}.
\]
Then
\[
        \cM_\QQ^\circ
        =
        \cM_\QQ\ominus L^2(\QQ,\cG_T),
\]
and \(\Lambda_M^\star\in \cM_\QQ^\circ\).
\end{definition}

\begin{proposition}[Canonical aggregate factor decomposition inside the priced structural loading]\label{thm:factor-decomp}
Under Proposition~\ref{thm:claim-decomp} and Definition~\ref{def:agg-factor}, every claim
\(X\in L^2(\QQ,\cF_T)\) admits a unique decomposition
\[
        L_X^\QQ
        =
        V_X^\QQ
        +
        \alpha_X^{\QQ,\star}\Lambda_M^\star
        +
        R_X^{\QQ,\star},
\]
where
\[
        V_X^\QQ:=C_{\cG_T}^\QQ(L_X^\QQ)\in L^2(\QQ,\cG_T),
\]
\[
        \alpha_X^{\QQ,\star}\in\R,
        \qquad
        R_X^{\QQ,\star}\in \cM_\QQ^\circ,
        \qquad
        R_X^{\QQ,\star}\perp_{L^2(\QQ)} \Lambda_M^\star.
\]
Equivalently, if
\[
        \widetilde L_X^\QQ:=L_X^\QQ-C_{\cG_T}^\QQ(L_X^\QQ)\in \cM_\QQ^\circ,
\]
then
\[
        \widetilde L_X^\QQ=\alpha_X^{\QQ,\star}\Lambda_M^\star+R_X^{\QQ,\star},
        \qquad
        \alpha_X^{\QQ,\star}
        =
        \langle \widetilde L_X^\QQ,\Lambda_M^\star\rangle_{L^2(\QQ)}.
\]
If \(\Lambda_M^\star=0\), then \(\alpha_X^{\QQ,\star}:=0\) and
\(R_X^{\QQ,\star}:=\widetilde L_X^\QQ\).  A claim is called \emph{aggregate-factor aligned} when
\[
        R_X^{\QQ,\star}=0.
\]
\end{proposition}

\begin{proof}
The space \(\cM_\QQ^\circ\) is a closed subspace of \(L^2(\QQ,\cF_T)\).  Since
\(L_X^\QQ\in\cM_\QQ\) and \(\widetilde L_X^\QQ=L_X^\QQ-\EE_\QQ[L_X^\QQ]\) has zero
\(\QQ\)-mean, \(\widetilde L_X^\QQ\in\cM_\QQ^\circ\).  If \(\Lambda_M^\star=0\), set
\(\alpha_X^{\QQ,\star}:=0\) and \(R_X^{\QQ,\star}:=\widetilde L_X^\QQ\).  Otherwise,
\(\mathrm{span}\{\Lambda_M^\star\}\) is a one-dimensional closed subspace of \(\cM_\QQ^\circ\), so
the Hilbert projection theorem yields the unique orthogonal decomposition
\[
        \widetilde L_X^\QQ
        =
        \langle \widetilde L_X^\QQ,\Lambda_M^\star\rangle_{L^2(\QQ)}\,\Lambda_M^\star
        +
        R_X^{\QQ,\star}
\]
with \(R_X^{\QQ,\star}\perp \Lambda_M^\star\).  Adding back
\(V_X^\QQ=C_{\cG_T}^\QQ(L_X^\QQ)\) gives the stated claim.
\end{proof}

\begin{corollary}[Linear public screening of the centered aggregate factor]\label{cor:factor-screening}
Under Definition~\ref{def:agg-factor}, for every public sub-\(\sigma\)-field \(\cH\subseteq\cG_T\),
\[
        C_{\cH}^\QQ(\Lambda_M^\star)=0.
\]
Consequently, if \(X\) is aggregate-factor aligned so that
\[
        L_X^\QQ=V_X^\QQ+\alpha_X^{\QQ,\star}\Lambda_M^\star,
\]
then
\[
        C_{\cH}^\QQ(L_X^\QQ)=C_{\cH}^\QQ(V_X^\QQ).
\]
Thus the canonical centered aggregate factor is always screened from linear public compression,
even when its coefficient \(\alpha_X^{\QQ,\star}\) is nonzero inside the priced structural loading.
\end{corollary}

\begin{proof}
Because \(\Lambda_M^\star\in \cM_\QQ^\circ\), one has
\(\EE_\QQ[\Lambda_M^\star\mid \cG_T]=0\).  Since \(\cH\subseteq \cG_T\), the tower property gives
\[
        C_{\cH}^\QQ(\Lambda_M^\star)
        =
        \EE_\QQ\bigl[\EE_\QQ[\Lambda_M^\star\mid \cG_T]\mid \cH\bigr]
        =0.
\]
Linearity of \(C_{\cH}^\QQ\) gives
\[
        C_{\cH}^\QQ(L_X^\QQ)
        =
        C_{\cH}^\QQ(V_X^\QQ)
        +
        \alpha_X^{\QQ,\star}C_{\cH}^\QQ(\Lambda_M^\star)
        =
        C_{\cH}^\QQ(V_X^\QQ),
\]
which proves the aligned-claim identity.
\end{proof}

\begin{remark}[Structural objects versus priced claim objects]\label{rem:struct-vs-priced}
The measure-invariant objects are the kernel system, the screened Perron construction, and the
derived sufficient structural sigma-field \(\cH_\varpi\).  By contrast, the priced structural
loading \(L_X^\QQ\) is claim-specific and measure-dependent.  Its scalar weight
\(\alpha_X^\QQ\) is a pricing object, and its normalized direction \(\Lambda_X\) is the normalized
direction of that priced loading for the chosen claim under \(\QQ\).  Thus \(\Lambda_X\) is not
measure-invariant claim by claim.  The invariant content is the structural geometry in which those
claim-specific loadings live.  The aggregate factor benchmark \(\Lambda_M^\star\) is the exception
isolated here.  Its role is tied to the canonical centered market-variance factor construction of
\cite{vidal2026market} rather than to an arbitrary claim.  Even there, the \(\QQ\)-side realization remains
measure-dependent through the conditional projection defining that centered structural increment.
\end{remark}

\begin{remark}[What the factor decomposition does and does not classify]\label{rem:factor-scope}
Proposition~\ref{thm:factor-decomp} is a factor-level decomposition inside the centered structural
loading while leaving the loading itself in place.  The claim-specific object classified by public
visibility remains \(L_X^\QQ=C_{\cH_\varpi}^\QQ(X)\), because visible summaries act on the full
priced structural loading.  The canonical centered aggregate factor \(\Lambda_M^\star\) has a
separate role.  It isolates the aggregate synchronized benchmark direction inside that loading and
contributes an exact orthogonal term to compression and hedging residuals.  Corollary~\ref{cor:factor-screening}
identifies why \(\Lambda_M^\star\) is not directly linearly visible through the option-surface
sigma-field.
\end{remark}

\begin{corollary}[Compression principle for priced structural loading]\label{prop:compression-unifier}
Let \(X\in L^2(\QQ,\cF_T)\) with canonical decomposition
\[
        X=L_X^\QQ+U_X^\QQ.
\]
Then:
\begin{enumerate}[label=(\roman*)]
\item for every sub-\(\sigma\)-field \(\cH\subseteq \cG_T\),
\[
        C_{\cH}^{\QQ}(X)
        =
        C_{\cH}^{\QQ}(L_X^\QQ);
\]
\item for every closed visible hedge class \(\cS\subseteq L^2(\QQ,\cG_T)\),
\[
        P_{\cS}^{\QQ}(X)
        =
        P_{\cS}^{\QQ}(L_X^\QQ),
\]
where \(P_{\cS}^{\QQ}\) is the \(L^2(\QQ)\)-orthogonal projection onto \(\cS\).
\item if \(L_X^\QQ\neq0\), so that \(L_X^\QQ=\alpha_X^\QQ\Lambda_X\), then
\[
        C_{\cH}^{\QQ}(X)=\alpha_X^\QQ C_{\cH}^{\QQ}(\Lambda_X),
        \qquad
        P_{\cS}^{\QQ}(X)=\alpha_X^\QQ P_{\cS}^{\QQ}(\Lambda_X),
\]
for every public \(\cH\subseteq \cG_T\) and every visible hedge class
\(\cS\subseteq L^2(\QQ,\cG_T)\).
\item more generally, if \(\Psi_{\cH}\colon L^2(\QQ,\cH)\to\R^m\) is any Borel measurable
functional on a public claim class \(L^2(\QQ,\cH)\), then
\[
        \Psi_{\cH}(C_{\cH}^{\QQ}(X))
        =
        \Psi_{\cH}(C_{\cH}^{\QQ}(L_X^\QQ)),
\]
and if \(L_X^\QQ\neq0\),
\[
        \Psi_{\cH}(C_{\cH}^{\QQ}(X))
        =
        \Psi_{\cH}(\alpha_X^\QQ C_{\cH}^{\QQ}(\Lambda_X)).
\]
\end{enumerate}
In particular, every public pricing object treated here is obtained by compressing,
projecting, or measurably transforming the same priced structural loading \(L_X^\QQ=\alpha_X^\QQ\Lambda_X\).
\end{corollary}

\begin{proof}
Because \(U_X^\QQ\perp \cM_\QQ\) and \(L^2(\QQ,\cG_T)\subseteq \cM_\QQ\), one has
\(C_{\cH}^{\QQ}(U_X^\QQ)=0\) for every \(\cH\subseteq \cG_T\), which proves part~(i).  Likewise
\(U_X^\QQ\perp\cS\) for every \(\cS\subseteq L^2(\QQ,\cG_T)\), which proves part~(ii).
Part~(iii) is the identity \(L_X^\QQ=\alpha_X^\QQ\Lambda_X\).  Part~(iv) follows by applying the
measurable map \(\Psi_{\cH}\) to part~(i), and then using part~(iii) when \(L_X^\QQ\neq0\).
\end{proof}

\begin{corollary}[Optional same-span reading under bounded-density comparability]\label{cor:compression-same-span}
Under Condition~\ref{ass:q-bounds}, the geometric span is the same under \(\PP\) and \(\QQ\), but
public visibility changes because \(C_{\cH}^{\QQ}\) and \(C_{\cH}^{\PP}\) are different bounded
projections on that common span.  Thus screening under \(\QQ\) is a reweighting of the same priced
structural directions, not a loss of structural support.
\end{corollary}

\begin{proof}
Proposition~\ref{prop:proj-stability} identifies the structural span as a common vector space under
\(\PP\) and \(\QQ\) when the density is bounded above and away from zero.  Corollary~\ref{prop:compression-unifier}
then says that the public pricing object is obtained by applying the risk-neutral projection
\(C_{\cH}^{\QQ}\) to the priced loading in that span.  Thus the support directions are common, while
the visible component may change because the conditional-expectation operator is measure-dependent.
\end{proof}

The risk-neutral compression theorem identifies the Hilbert-space source of public pricing,
visible projection, and residual hedge risk.

\begin{theorem}[Main risk-neutral compression theorem]\label{thm:main-risk-neutral-compression}
Assume Assumptions~\ref{ass:model}, \ref{ass:q}, and~\ref{ass:aggregate-perron-bridge}, together
with Definition~\ref{def:lambda-structure}.  Let
\[
        \cH\subseteq\cG_T\subseteq\cH_\varpi\subseteq\cF_T.
\]
For every \(X\in L^2(\QQ,\cF_T)\), set
\[
        L_X^\QQ=C_{\cH_\varpi}^\QQ(X),
        \qquad
        U_X^\QQ=X-L_X^\QQ.
\]
Then
\[
        X=L_X^\QQ+U_X^\QQ,
        \qquad
        U_X^\QQ\perp_{L^2(\QQ)}L^2(\QQ,\cH_\varpi).
\]
Every public compression factors through the priced structural loading:
\[
        C_{\cH}^\QQ(X)=C_{\cH}^\QQ(L_X^\QQ).
\]
For every closed visible hedge class \(\cS\subseteq L^2(\QQ,\cG_T)\),
\[
        P_{\cS}^\QQ(X)=P_{\cS}^\QQ(L_X^\QQ)
\]
and the residual risk decomposes as
\[
        \inf_{H\in\cS}\EE_\QQ[(X-H)^2]
        =
        \dist_{L^2(\QQ)}(L_X^\QQ,\cS)^2
        +
        \|U_X^\QQ\|_{L^2(\QQ)}^2.
\]
If \(L_X^\QQ\neq0\), then the normalized loading
\[
        \Lambda_X=\frac{L_X^\QQ}{\|L_X^\QQ\|_{L^2(\QQ)}}
\]
determines every public compression and visible projection up to the scalar
\(\alpha_X^\QQ=\|L_X^\QQ\|_{L^2(\QQ)}\).

After centering against \(L^2(\QQ,\cG_T)\), the same loading has the unique aggregate-factor split
\[
        L_X^\QQ
        =
        V_X^\QQ
        +
        \alpha_X^{\QQ,\star}\Lambda_M^\star
        +
        R_X^{\QQ,\star},
        \qquad
        R_X^{\QQ,\star}\perp_{L^2(\QQ)}\Lambda_M^\star.
\]
Finally, whenever \(X=X_\varpi+R_X\) with \(X_\varpi\) \(\cH_\varpi\)-measurable, replacing the
full aggregate structural content by the Perron-generated component \(X_\varpi\) perturbs
structural projection, public compression, and visible projection by at most
\(\|R_X\|_{L^2(\QQ)}\).
\end{theorem}

\begin{proof}
Conditional expectation onto \(\cH_\varpi\) is the \(L^2(\QQ)\)-orthogonal projection onto the
closed subspace \(L^2(\QQ,\cH_\varpi)\).  Hence \(U_X^\QQ\) is orthogonal to
\(L^2(\QQ,\cH_\varpi)\).  Since \(\cH\subseteq\cG_T\subseteq\cH_\varpi\), the tower property gives
\[
        C_{\cH}^\QQ(X)
        =
        C_{\cH}^\QQ\bigl(C_{\cH_\varpi}^\QQ(X)\bigr)
        =
        C_{\cH}^\QQ(L_X^\QQ).
\]
Also \(\cS\subseteq L^2(\QQ,\cG_T)\subseteq L^2(\QQ,\cH_\varpi)\), so
\(U_X^\QQ\perp\cS\).  Therefore the closest element of \(\cS\) to \(X=L_X^\QQ+U_X^\QQ\) is the
closest element of \(\cS\) to \(L_X^\QQ\), and Pythagoras gives
\[
        \|X-H\|_{L^2(\QQ)}^2
        =
        \|L_X^\QQ-H\|_{L^2(\QQ)}^2+\|U_X^\QQ\|_{L^2(\QQ)}^2,
        \qquad H\in\cS.
\]
Taking the infimum over \(H\in\cS\) proves the hedge residual identity.  The normalized-loading
statement follows from \(L_X^\QQ=\alpha_X^\QQ\Lambda_X\).

The aggregate-factor split is the orthogonal decomposition of
\(L_X^\QQ-C_{\cG_T}^\QQ(L_X^\QQ)\) along the one-dimensional span of
\(\Lambda_M^\star\), as in Proposition~\ref{thm:factor-decomp}.  If \(X=X_\varpi+R_X\) with
\(X_\varpi\) \(\cH_\varpi\)-measurable, then
\[
        L_X^\QQ-X_\varpi=C_{\cH_\varpi}^\QQ(R_X),
\]
so the final perturbative bound follows from the \(L^2(\QQ)\)-contraction property of conditional
expectation and orthogonal projection.
\end{proof}

\begin{remark}[Relation to partial-information hedging]\label{rem:partial-information-positioning}
The Hilbert-space step in Theorem~\ref{thm:main-risk-neutral-compression} is the familiar
projection step from reduced-information pricing and hedging.  Conditional expectation onto a
smaller information set is an \(L^2\)-projection, and closed visible hedge spaces enter by
orthogonal projection.  This is the same geometric layer underlying local risk minimization,
mean-variance hedging, the F\"ollmer--Schweizer decomposition, and its relationship with the
Galtchouk--Kunita--Watanabe decomposition
\cite{follmerSondermann1986,follmerSchweizer1991,schweizer1992meanVariance,
schweizer1995minimal,choulliVandaeleVanmaele2010}.  The theorem leaves the Galtchouk--Kunita--Watanabe decomposition and the F\"ollmer--Schweizer
partial-information hedging construction in place.  Its role is to identify the object projected here: the
\(\QQ\)-priced structural loading \(L_X^\QQ=C_{\cH_\varpi}^{\QQ}(X)\) generated by the
Perron-sufficient latent market state.  That object is then the common source of the option,
premium, variance, and hedging endpoints below.
\end{remark}

\section{Risk-neutral visibility and screening}\label{sec:class}

Visibility is a norm statement on the compressed loading.  A public summary sees a claim only
through \(C_{\cH}^{\QQ}(L_X^\QQ)\); the uncompressed loading may still be nonzero.

\begin{definition}[Compression-based \(\QQ\)-visibility of priced structural loading]\label{def:q-visible}
Fix a claim \(X\in L^2(\QQ,\cF_T)\) with priced structural loading \(L_X^\QQ\), and fix a public
sub-\(\sigma\)-field \(\cH\subseteq \cG_T\).  The claim's structural loading is called
\begin{enumerate}[label=(\roman*)]
\item \emph{\(\QQ\)-visible for \(X\) at public scale \(\cH\)} if
\[
        \|C_{\cH}^{\QQ}(L_X^\QQ)\|_{L^2(\QQ)}>0;
\]
\item \emph{\(\QQ\)-screened for \(X\) at public scale \(\cH\)} if
\[
        \|L_X^\QQ\|_{L^2(\QQ)}>0
        \quad\text{and}\quad
        \|C_{\cH}^{\QQ}(L_X^\QQ)\|_{L^2(\QQ)}=0;
\]
\item \emph{structurally absent for \(X\)} if
\[
        \|L_X^\QQ\|_{L^2(\QQ)}=0.
\]
\end{enumerate}
\end{definition}

\begin{proposition}[Operational summary criterion]\label{prop:summary-criterion}
Let \(\{X_h\}_{h>0}\subset L^2(\QQ,\cF_T)\) be a claim family with priced structural loadings
\(L_{X_h}^\QQ\), let \(\cH_h\subseteq\cG_T\) be public sigma-fields, and let
\(\{\mathfrak S_h(\varepsilon)\}_{h>0,|\varepsilon|<\varepsilon_0}\subset L^2(\QQ,\cH_h)\) be
visible summary families.  Suppose that for a positive deterministic normalization \(a_h>0\) there are
remainders \(r_h(\varepsilon)\in L^2(\QQ)\) such that
\[
        \frac{\mathfrak S_h(\varepsilon)-\mathfrak S_h(0)}{\varepsilon}
        =
        C_{\cH_h}^{\QQ}(L_{X_h}^\QQ)
        +
        r_h(\varepsilon),
\]
and
\[
        \lim_{h\downarrow0,\,\varepsilon\to0}
        a_h^{-1}\|r_h(\varepsilon)\|_{L^2(\QQ)}=0.
\]
Assume moreover that the normalized compressed priced structural loadings converge:
\[
        a_h^{-1}C_{\cH_h}^{\QQ}(L_{X_h}^\QQ)
        \longrightarrow
        Y
        \qquad\text{in }L^2(\QQ).
\]
Then the normalized perturbation has the joint \(L^2(\QQ)\)-limit
\[
        D_{\mathrm{load}}^\QQ\mathfrak S
        :=
        \lim_{h\downarrow0,\,\varepsilon\to0}
        \frac{\mathfrak S_h(\varepsilon)-\mathfrak S_h(0)}{\varepsilon a_h}
        =Y,
\]
and consequently
\[
        D_{\mathrm{load}}^\QQ\mathfrak S\neq0
        \quad\Longleftrightarrow\quad
        \lim_{h\downarrow0}
        a_h^{-1}\|C_{\cH_h}^{\QQ}(L_{X_h}^\QQ)\|_{L^2(\QQ)}
        =
        \|Y\|_{L^2(\QQ)}>0.
\]
Thus the derivative-based criterion detects the limiting normalized compressed loading.
The normalization \(a_h\) need not vanish; this convention covers both decaying and exploding
short-maturity quote scales.  Without convergence of the normalized compressed loadings it should
be read only as a scale-wise sufficient test.
\end{proposition}

\begin{proof}
After division by \(a_h\), the difference between the normalized summary perturbation and
\(a_h^{-1}C_{\cH_h}^{\QQ}(L_{X_h}^\QQ)\) is \(a_h^{-1}r_h(\varepsilon)\), which converges to zero
in the stated joint limit.  The normalized compressed loading converges to \(Y\) in \(L^2(\QQ)\).
The normalized summary perturbation therefore has the same limit.  Norm convergence follows from
the reverse triangle inequality, and the final equivalence is exactly \(\|Y\|_{L^2(\QQ)}>0\).
\end{proof}

\begin{proposition}[Risk-neutral visibility classification for priced structural loadings]\label{thm:q-trichotomy}
Fix a claim \(X\in L^2(\QQ,\cF_T)\), let
\[
        X=L_X^\QQ+U_X^\QQ
\]
be the canonical decomposition of Proposition~\ref{thm:claim-decomp}, and fix a public
sub-\(\sigma\)-field \(\cH\subseteq \cG_T\).  Then exactly one of the following three cases holds:
\begin{enumerate}[label=(\roman*)]
\item \emph{visible and structurally priced}:
\[
        \|C_{\cH}^{\QQ}(L_X^\QQ)\|_{L^2(\QQ)}>0;
\]
\item \emph{screened but structurally priced}:
\[
        \|L_X^\QQ\|_{L^2(\QQ)}>0
        \quad\text{and}\quad
        \|C_{\cH}^{\QQ}(L_X^\QQ)\|_{L^2(\QQ)}=0;
\]
\item \emph{structurally absent for the claim}:
\[
        \|L_X^\QQ\|_{L^2(\QQ)}=0.
\]
\end{enumerate}
These conditions are necessary and sufficient.
\[
        \text{If } \alpha_X^\QQ>0,\qquad
        \|C_{\cH}^{\QQ}(L_X^\QQ)\|_{L^2(\QQ)}
        =
        \alpha_X^\QQ \|C_{\cH}^{\QQ}(\Lambda_X)\|_{L^2(\QQ)}.
\]
\end{proposition}

\begin{proof}
Definition~\ref{def:q-visible} partitions the claim by the two norms
\(\|L_X^\QQ\|_{L^2(\QQ)}\) and
\(\|C_{\cH}^{\QQ}(L_X^\QQ)\|_{L^2(\QQ)}\).  If the first norm is zero, the loading is absent.  If it
is positive, the second norm is either positive or zero, giving visibility or screening and no
fourth case.  When \(\alpha_X^\QQ>0\), the scalar factorization
\(
        L_X^\QQ=\alpha_X^\QQ\Lambda_X
\)
and linearity of \(C_{\cH}^{\QQ}\) give the displayed identity.
\end{proof}

\begin{remark}[Why this is the classification theorem]
The classification does not test existence of the structural geometry; Theorem~\ref{thm:q-invariance}
has already fixed that geometry.  It tests whether a chosen public compression selects a nonzero
component of the claim's priced structural loading.  Proposition~\ref{thm:q-trichotomy}
characterizes the distinction by the norm of \(C_{\cH}^{\QQ}(L_X^\QQ)\).
\end{remark}

\section{Market-index price dynamics and option setup under \texorpdfstring{\(\QQ\)}{Q}}\label{sec:qsetup}

The option block separates two maps.  Risk-neutral compression produces the option-facing priced
loading; a quote functional decides which part of that loading appears as an observed surface
coefficient.  The abstract first-variation theorem is Gaussian-free.  The cumulant-class condition
then turns the compressed scalar signal into a short-maturity skew law, and the
Gaussian-Volterra/Bachelier specialization gives kernel-level coefficients.  The quote conventions
and option-surface background come from the Bachelier, Black--Scholes/Black, rational option-pricing,
state-price-density, local-volatility, and skew/moment literatures
\cite{bachelier1900,blackScholes1973,black1976,merton1973,breeden1978,dupire1994,
jarrowRudd1982,corradoSu1996,bakshi2003}.

\begin{assumption}[Regular public option-summary first variation]\label{ass:option-fv}
For every sufficiently small maturity \(\tau>0\), there exist:
\begin{enumerate}[label=(\roman*)]
\item a public option sigma-field \(\cH_{M,\tau}^{\mathrm{opt}}\subseteq \cG_T\);
\item an option first-variation pricing object \(\Xi_{M,\tau}\in L^2(\QQ,\cF_T)\);
\item a continuous linear functional
\[
        \mathfrak D_{M,\tau}\colon L^2(\QQ,\cH_{M,\tau}^{\mathrm{opt}})\to\R;
\]
\item a deterministic scale \(a_{M,\tau}>0\) and a remainder \(r_{M,\tau}\in\R\)
\end{enumerate}
such that the quoted option summary family \(\mathfrak S_M(\tau;\varepsilon)\) satisfies
\[
        \partial_\varepsilon \mathfrak S_M(\tau;\varepsilon)\big|_{\varepsilon=0}
        =
        \mathfrak D_{M,\tau}
        \bigl(C_{\cH_{M,\tau}^{\mathrm{opt}}}^\QQ(\Xi_{M,\tau})\bigr)
        +
        r_{M,\tau},
\]
with
\[
        a_{M,\tau}^{-1}r_{M,\tau}\to0
        \qquad
        \text{as } \tau\downarrow0.
\]
\end{assumption}

\begin{definition}[Option first-variation pricing object and option-facing priced structural loading]\label{def:option-loading}
Under Assumption~\ref{ass:option-fv}, define
\[
        L_{M,\tau}^\QQ:=C_{\cH_\varpi}^\QQ(\Xi_{M,\tau}),
        \qquad
        U_{M,\tau}^\QQ:=\Xi_{M,\tau}-C_{\cH_\varpi}^\QQ(\Xi_{M,\tau}).
\]
Under the standing hypotheses of Proposition~\ref{thm:claim-decomp}, this is the canonical
priced-structural decomposition of the option first-variation claim \(\Xi_{M,\tau}\).
If \(L_{M,\tau}^\QQ\neq0\), set
\[
        \alpha_{M,\tau}^\QQ:=\|L_{M,\tau}^\QQ\|_{L^2(\QQ)},
        \qquad
        \Lambda_{M,\tau}:=\frac{L_{M,\tau}^\QQ}{\alpha_{M,\tau}^\QQ};
\]
otherwise set \(\alpha_{M,\tau}^\QQ:=0\) and \(\Lambda_{M,\tau}:=0\).  The family
\(\{L_{M,\tau}^\QQ\}_{\tau>0}\) is the option-facing priced structural loading.
\end{definition}

The abstract option theorem separates the compressed loading from the scalar quote functional that
may annihilate part of it.

\begin{theorem}[Abstract option-summary compression of priced structural loading]\label{thm:option-abstract}
Under Assumptions~\ref{ass:model}, \ref{ass:q}, and~\ref{ass:option-fv}, together with
Definitions~\ref{def:mode-space} and~\ref{def:option-loading},
\[
        \partial_\varepsilon \mathfrak S_M(\tau;\varepsilon)\big|_{\varepsilon=0}
        =
        \mathfrak D_{M,\tau}
        \bigl(C_{\cH_{M,\tau}^{\mathrm{opt}}}^\QQ(L_{M,\tau}^\QQ)\bigr)
        +
        r_{M,\tau}.
\]
If \(\alpha_{M,\tau}^\QQ>0\), then
\[
        \partial_\varepsilon \mathfrak S_M(\tau;\varepsilon)\big|_{\varepsilon=0}
        =
        \mathfrak D_{M,\tau}
        \bigl(
                \alpha_{M,\tau}^\QQ
                C_{\cH_{M,\tau}^{\mathrm{opt}}}^\QQ(\Lambda_{M,\tau})
        \bigr)
        +
        r_{M,\tau}.
\]
In particular, the first variation of every regular public option summary treated here is
a continuous linear functional of the compressed option-facing priced structural loading.  Hence
the norm of the compressed option-facing priced structural loading governs structural visibility
and screening at the chosen option scale:
\[
        \|C_{\cH_{M,\tau}^{\mathrm{opt}}}^\QQ(L_{M,\tau}^\QQ)\|_{L^2(\QQ)}.
\]
The particular scalar option summary detects that visible loading only through the additional
linear functional \(\mathfrak D_{M,\tau}\); a nonzero compressed loading may be invisible to this
scalar summary if it lies in \(\ker\mathfrak D_{M,\tau}\).
\end{theorem}

\begin{proof}
By Corollary~\ref{prop:compression-unifier},
\[
        C_{\cH_{M,\tau}^{\mathrm{opt}}}^\QQ(\Xi_{M,\tau})
        =
        C_{\cH_{M,\tau}^{\mathrm{opt}}}^\QQ(L_{M,\tau}^\QQ).
\]
Substituting this identity into Assumption~\ref{ass:option-fv} gives the first displayed formula.
The second formula follows from the scalar factorization
\(
        L_{M,\tau}^\QQ=\alpha_{M,\tau}^\QQ\Lambda_{M,\tau}
\)
when \(\alpha_{M,\tau}^\QQ>0\).  The final statement separates two kernels.  Visibility is governed
by the kernel of the projection \(C_{\cH_{M,\tau}^{\mathrm{opt}}}^\QQ\) on the option-facing loading.
Detection by the chosen scalar quote is governed, after that projection, by
\(\ker\mathfrak D_{M,\tau}\).
\end{proof}

\begin{remark}[Scope of the abstract option theorem]
Theorem~\ref{thm:option-abstract} identifies the option-side compression map.  The first variation
of a regular public option summary is generated by compressing the option-facing priced structural
loading and then applying a continuous linear summary functional.  No Gaussian assumption enters
that result.  The explicit non-Gaussian cumulant-class law and the later
Gaussian-Volterra/Bachelier asymptotics are specializations of the same abstract compression
statement.
\end{remark}

\begin{assumption}[Non-Gaussian cumulant-class public skew condition]\label{ass:cumulant-class}
For every sufficiently small maturity \(\tau>0\), let
\[
        v_M(\tau):=\Var_\QQ(R_\tau^0)
\]
denote the baseline option-facing return variance at maturity \(\tau\), and let
\[
        \kappa_{3,M}^\QQ(\tau;\varepsilon)
\]
denote the third cumulant of the perturbed option-facing return at that maturity.  Assume:
\begin{enumerate}[label=(\roman*)]
\item there exists \(\vartheta_M>0\) such that
\[
        v_M(\tau)=\vartheta_M \tau + o(\tau),
        \qquad
        \tau\downarrow0;
\]
\item there exists a finite set of exponents \(\mathcal B_M\subset(0,\infty)\) and coefficients
\(\{A_M(\beta)\}_{\beta\in\mathcal B_M}\subset\R\) such that
\[
        \mathrm{supp}^{\mathrm{comp}}_M
        :=
        \{\beta\in\mathcal B_M:A_M(\beta)\neq0\}
\]
and, whenever \(\mathrm{supp}^{\mathrm{comp}}_M\neq\varnothing\),
\[
        \beta_M^{\mathrm{comp}}
        :=
        \inf \mathrm{supp}^{\mathrm{comp}}_M,
\]
with the convention \(\beta_M^{\mathrm{comp}}:=+\infty\) when
\(\mathrm{supp}^{\mathrm{comp}}_M=\varnothing\), and such that:
\begin{enumerate}[label=(\alph*)]
\item if \(\mathrm{supp}^{\mathrm{comp}}_M\neq\varnothing\), then
\[
        \partial_\varepsilon \kappa_{3,M}^\QQ(\tau;\varepsilon)\big|_{\varepsilon=0}
        =
        6\vartheta_M^{3/2}
        \sum_{\beta\in\mathcal B_M}
        A_M(\beta)\,\tau^{\beta+3/2}
        +
        o\!\left(\tau^{\beta_M^{\mathrm{comp}}+3/2}\right),
        \qquad
        \tau\downarrow0,
\]
\item if \(\mathrm{supp}^{\mathrm{comp}}_M=\varnothing\), then
\[
        \partial_\varepsilon \kappa_{3,M}^\QQ(\tau;\varepsilon)\big|_{\varepsilon=0}
        =
        0
        \qquad
        \text{for every sufficiently small }\tau>0.
\]
\end{enumerate}
\item there exists a deterministic function \(G_M(\tau)\to G_M(0)>0\) such that
\begin{enumerate}[label=(\alph*)]
\item if \(\mathrm{supp}^{\mathrm{comp}}_M\neq\varnothing\), then
\[
        \partial_\varepsilon
        \mathfrak S_M(\tau;\varepsilon)\big|_{\varepsilon=0}
        =
        \frac{G_M(\tau)}{6v_M(\tau)^{3/2}\tau^{1/2}}\,
        \partial_\varepsilon
        \kappa_{3,M}^\QQ(\tau;\varepsilon)\big|_{\varepsilon=0}
        +
        o\!\left(\tau^{\beta_M^{\mathrm{comp}}-1/2}\right).
\]
\item if \(\mathrm{supp}^{\mathrm{comp}}_M=\varnothing\), then
\[
        \partial_\varepsilon
        \mathfrak S_M(\tau;\varepsilon)\big|_{\varepsilon=0}
        =
        0
        \qquad
        \text{for every sufficiently small }\tau>0.
\]
\end{enumerate}
\end{enumerate}
\end{assumption}

\begin{definition}[Cumulant-selected compressed exponent support of the scalar public signal]\label{def:compressed-support}
Under Assumption~\ref{ass:cumulant-class}, define the cumulant-selected compressed exponent
support by
\[
        \mathrm{supp}^{\mathrm{comp}}(\Gamma_{M,\tau}^\QQ)
        :=
        \{\beta\in\mathcal B_M:A_M(\beta)\neq0\},
\]
where
\[
        \Gamma_{M,\tau}^\QQ
        :=
        \mathfrak D_{M,\tau}
        \bigl(C_{\cH_{M,\tau}^{\mathrm{opt}}}^\QQ(L_{M,\tau}^\QQ)\bigr)
\]
is the scalar public first-variation signal generated from the compressed option-facing priced
structural loading.  The effective public leading exponent is
\[
        \beta_M^{\mathrm{comp}}
        :=
        \inf \mathrm{supp}^{\mathrm{comp}}(\Gamma_{M,\tau}^\QQ),
\]
with the convention \(+\infty\) when the compressed support is empty.
The exponent support is selected by the cumulant-class condition for the scalar public signal
generated by the compressed loading.  The full \(L^2\)-valued compressed loading has no such scalar
support notion.
\end{definition}

The cumulant-class theorem characterizes when the scalar public skew coefficient is the leading
power selected by the compressed option-facing loading.

\begin{theorem}[Non-Gaussian cumulant-class market skew as compressed priced structural loading]\label{thm:market-skew}
Assume Assumptions~\ref{ass:model}, \ref{ass:q}, \ref{ass:option-fv},
and~\ref{ass:cumulant-class}, together with
Definitions~\ref{def:mode-space} and~\ref{def:option-loading}.  Let
\[
        \Gamma_{M,\tau}^\QQ
        :=
        \mathfrak D_{M,\tau}
        \bigl(C_{\cH_{M,\tau}^{\mathrm{opt}}}^\QQ(L_{M,\tau}^\QQ)\bigr).
\]
If \(\mathrm{supp}^{\mathrm{comp}}(\Gamma_{M,\tau}^\QQ)=\varnothing\), then the cumulant-class
coefficient system does not select a nontrivial leading public power at the modeled scales.
Assume henceforth that
\[
        \mathrm{supp}^{\mathrm{comp}}(\Gamma_{M,\tau}^\QQ)\neq\varnothing,
        \qquad
        \text{equivalently }
        \beta_M^{\mathrm{comp}}<\infty.
\]
Then, as \(\tau\downarrow0\), the quoted public first variation satisfies
\[
        \partial_\varepsilon
        \mathfrak S_M(\tau;\varepsilon)\big|_{\varepsilon=0}
        =
        G_M(\tau)
        \sum_{\beta\in\mathrm{supp}^{\mathrm{comp}}(\Gamma_{M,\tau}^\QQ)}
        A_M(\beta)\,\tau^{\beta-1/2}
        +
        o\!\left(\tau^{\beta_M^{\mathrm{comp}}-1/2}\right).
\]
Moreover Theorem~\ref{thm:option-abstract} gives
\[
        \partial_\varepsilon
        \mathfrak S_M(\tau;\varepsilon)\big|_{\varepsilon=0}
        =
        \Gamma_{M,\tau}^\QQ
        +
        r_{M,\tau},
        \qquad
        a_{M,\tau}^{-1}r_{M,\tau}\to0.
\]
If the abstract option-summary remainder is negligible at the cumulant leading scale,
\[
        r_{M,\tau}
        =
        o\!\left(\tau^{\beta_M^{\mathrm{comp}}-1/2}\right),
\]
then the scalar public compressed-loading signal itself has the same expansion:
\[
        \Gamma_{M,\tau}^\QQ
        =
        G_M(\tau)
        \sum_{\beta\in\mathrm{supp}^{\mathrm{comp}}(\Gamma_{M,\tau}^\QQ)}
        A_M(\beta)\,\tau^{\beta-1/2}
        +
        o\!\left(\tau^{\beta_M^{\mathrm{comp}}-1/2}\right).
\]
Equivalently, under this scale-compatible remainder condition and when
\(\alpha_{M,\tau}^\QQ>0\),
\[
        \mathfrak D_{M,\tau}
        \bigl(
                \alpha_{M,\tau}^\QQ
                C_{\cH_{M,\tau}^{\mathrm{opt}}}^\QQ(\Lambda_{M,\tau})
        \bigr)
        =
        G_M(\tau)
        \sum_{\beta\in\mathrm{supp}^{\mathrm{comp}}(\Gamma_{M,\tau}^\QQ)}
        A_M(\beta)\,\tau^{\beta-1/2}
        +
        o\!\left(\tau^{\beta_M^{\mathrm{comp}}-1/2}\right).
\]
Thus the leading public skew derivative is generated by the first-order public compression of the
option-facing priced structural loading.  The expansion for the scalar compressed-loading signal
requires the displayed scale compatibility; otherwise the cumulant-class expansion is an expansion
of the quoted derivative rather than of \(\Gamma_{M,\tau}^\QQ\) itself.
\end{theorem}

\begin{proof}
Assume \(\mathrm{supp}^{\mathrm{comp}}(\Gamma_{M,\tau}^\QQ)\neq\varnothing\).  The case of empty
compressed support is exactly the qualitative first sentence of the theorem.
Assumption~\ref{ass:cumulant-class}(iii) expresses the public first variation through the
normalized third-cumulant derivative.  Substituting the asymptotic expansion from
Assumption~\ref{ass:cumulant-class}(ii) and using
\[
        v_M(\tau)^{3/2}
        =
        \vartheta_M^{3/2}\tau^{3/2}(1+o(1))
\]
from Assumption~\ref{ass:cumulant-class}(i) yields
\[
        \partial_\varepsilon
        \mathfrak S_M(\tau;\varepsilon)\big|_{\varepsilon=0}
        =
        G_M(\tau)\sum_{\beta\in\mathcal B_M} A_M(\beta)\tau^{\beta-1/2}
        +
        o\!\left(\tau^{\beta_M^{\mathrm{comp}}-1/2}\right).
\]
Removing the zero coefficients is exactly the passage from \(\mathcal B_M\) to
\(\mathrm{supp}^{\mathrm{comp}}(\Gamma_{M,\tau}^\QQ)\).  The identity with
\(\Gamma_{M,\tau}^\QQ+r_{M,\tau}\) is Theorem~\ref{thm:option-abstract}, written in scalar
notation.  If
\(r_{M,\tau}=o(\tau^{\beta_M^{\mathrm{comp}}-1/2})\), subtracting this remainder from the
derivative expansion gives the displayed expansion for \(\Gamma_{M,\tau}^\QQ\).  The factorized
form follows from Theorem~\ref{thm:option-abstract} when
\(\alpha_{M,\tau}^\QQ>0\).
\end{proof}


The Hermite obstruction, primitive non-Gaussian Volterra estimates, and rare-jump exact price
calculations delimit the option-transfer endpoint.  They are collected in
Appendix~\ref{sec:option-transfer-variants}.  The article uses the compact consequence in
Theorem~\ref{thm:market-skew}.  When a quoted surface coefficient satisfies the stated transfer
condition, it is a public compression of the option-facing priced structural loading.

\begin{assumption}[Gaussian-Volterra market-index dynamics with synchronized perturbation]\label{ass:q-index}
Fix deterministic index weights \(b=(b_1,\dots,b_N)\in(0,\infty)^N\), a continuous baseline normal
volatility \(\sigma_0^M\colon[0,T]\to(0,\infty)\), and an additional \(\QQ\)-Brownian motion
\(B^\QQ\) such that
\[
        \langle B^\QQ,W^{\QQ,j}\rangle_t=\rho_j t,
        \qquad
        \sum_{j=1}^N \rho_j^2\le 1,
        \qquad
        1\le j\le N.
\]
Equivalently, possibly after enlarging the filtered space,
\[
        B_t^\QQ
        =
        \sum_{j=1}^N \rho_j W_t^{\QQ,j}
        +
        \sqrt{1-\|\rho\|^2}\,W_t^{\QQ,0},
\]
where \(W^{\QQ,0}\) is a \(\QQ\)-Brownian motion independent of \(W^\QQ\).
For a local synchronized perturbation parameter \(\varepsilon\), let the discounted market-index
price satisfy the additive dynamics
\[
        S_t^{M,\varepsilon}
        =
        S_0^M
        +
        \int_0^t \sigma_M^\varepsilon(u)\,\dd B_u^\QQ,
        \qquad
        \sigma_M^\varepsilon(u):=\sigma_0^M(u)+\varepsilon Y_M(u),
\]
where
\[
        Y_M(t)
        :=
        \sum_{j=1}^N \int_0^t G_j(t,s)\,\dd W_s^{\QQ,j},
\]
with
\[
        G_j(t,s)
        :=
        \sum_{i=1}^N
        b_i\,\dot\beta_{ij}\,a_{ij}(t,s)\,(t-s)^{H_{ij}-1/2}\,
        \one_{\{(i,j)\in\cE_M\}}.
\]
Here \(\dot\beta_{ij}\in\R\) is the synchronized-mode perturbation loading of the market-visible
channel \((i,j)\).  The synchronized perturbation satisfies
\[
        \EE_\QQ\int_0^T Y_M(u)^2\,\dd u<\infty,
\]
so the stochastic integral defining \(Z_\tau\) below is well defined for \(\tau\le T\).
Set
\[
        X_\tau:=\int_0^\tau \sigma_0^M(u)\,\dd B_u^\QQ,
        \qquad
        Z_\tau:=\int_0^\tau Y_M(u)\,\dd B_u^\QQ,
        \qquad
        R_\tau^\varepsilon:=S_\tau^{M,\varepsilon}-S_0^M=X_\tau+\varepsilon Z_\tau.
\]
Assume further that for every sufficiently small \(\tau>0\),
\[
        X_\tau^2 Z_\tau\in L^2(\QQ).
\]
For \(\tau>0\), let \(C_M(\tau,k;\varepsilon)\) denote the local date-\(t\) price, with \(t\)
suppressed, of a maturity-\(\tau\) call on \(S^{M,\varepsilon}\) with strike \(S_t^M+k\), let
\(\sigma_N^M(\tau,k;\varepsilon)\) be its Bachelier implied volatility, let
\(\cH_{M,\tau}^{\mathrm{opt}}\subseteq \cG_t\) denote the public sigma-field generated by the
maturity-\(\tau\) quote family at that observation date, and let
\[
        \mathfrak S_{M,N}(\tau;\varepsilon)\in \R \subset L^2(\QQ,\cH_{M,\tau}^{\mathrm{opt}})
 \]
denote the associated public ATM normal skew summary.  The displayed local formulas are written
after conditioning on the observation-date public information.  In a single-date time-zero
calculation this sigma-field may be trivial and the quote is identified with a constant element of
\(L^2(\QQ)\); in the empirical panel, the same coefficient is observed repeatedly across dates as a
\(\cG_t\)-measurable option-surface coordinate.  Its local value is the quoted normal slope
\[
        \partial_k \sigma_N^M(\tau,k;\varepsilon)\big|_{k=0}.
\]
\end{assumption}

\begin{definition}[Market-visible option coefficients and exponents]\label{def:market-option-coeff}
Under Assumption~\ref{ass:q-index}, define the return-visible channel set
\[
        \cE_M^{\mathrm{ret}}
        :=
        \{(i,j)\in\cE_M: b_i\,\dot\beta_{ij}\,\rho_j\neq0\}.
\]
For each \((i,j)\in\cE_M^{\mathrm{ret}}\), set
\[
        \ell_{ij}^M
        :=
        \frac{b_i\,\dot\beta_{ij}\,a_{ij}(0,0)\,\rho_j}{H_{ij}+1/2},
        \qquad
        \kappa_{ij}^M
        :=
        \frac{6\,(\sigma_0^M(0))^2\,b_i\,\dot\beta_{ij}\,a_{ij}(0,0)\,\rho_j}
        {(H_{ij}+1/2)(H_{ij}+3/2)},
\]
and define
\[
        c_{ij}^M
        :=
        \frac{\kappa_{ij}^M}{6(\sigma_0^M(0))^3}
        =
        \frac{b_i\,\dot\beta_{ij}\,a_{ij}(0,0)\,\rho_j}
        {\sigma_0^M(0)(H_{ij}+1/2)(H_{ij}+3/2)}.
\]
Set the market-scale leading exponents
\[
        \beta_M
        :=
        \min\{H_{ij}:(i,j)\in\cE_M^{\mathrm{ret}}\},
\]
\[
        \beta_M^{\mathrm{off}}
        :=
        \min\{H_{ij}:i\neq j,\ (i,j)\in\cE_M^{\mathrm{ret}}\},
        \qquad
        \beta_M^{\mathrm{diag}}
        :=
        \min\{H_{ii}:(i,i)\in\cE_M^{\mathrm{ret}}\},
\]
with the convention \(+\infty\) when the corresponding set is empty; the same convention also
applies to \(\beta_M\).  Define the aggregate
leading coefficients
\[
        A_M^{\mathrm{off}}
        :=
        \sum_{i\neq j:\,H_{ij}=\beta_M^{\mathrm{off}}} c_{ij}^M,
        \qquad
        A_M^{\mathrm{diag}}
        :=
        \sum_{i:\,H_{ii}=\beta_M^{\mathrm{diag}}} c_{ii}^M.
\]
When \(\cE_M^{\mathrm{ret}}=\varnothing\), all coefficients above are taken to vanish and the skew
response is trivial; in particular \(\beta_M=+\infty\).
\end{definition}

\begin{remark}[Role of the market-index setup]
Assumption~\ref{ass:q-index} supplies the market-facing input for the option calculation.  It fixes
the pricing measure, return Brownian channel, leverage coefficients, and synchronized perturbation
map.  The option summary is generated by this market-index model, not by an external single-channel
pricing specification.  In the Gaussian--Volterra/Bachelier specialization, the cumulant-class
condition is verified in closed form.
\end{remark}


\section{Option-surface transfer: main statement}\label{sec:option-transfer-main}

The empirical option input is deliberately narrow.  Section~\ref{sec:qsetup} supplies the abstract
first-variation theorem, the cumulant-class transfer, and the local option-facing loading.
Appendices~\ref{sec:skew}--\ref{sec:black} supply the Gaussian-Volterra, Bachelier, and Black
calculations.  Under the source-family hypotheses verified in Companion~A, the local option-surface
coefficient is a public \(\QQ\)-compression of the same priced structural loading that appears in
Theorem~\ref{thm:main-risk-neutral-compression}.  The SPX/NDX exercise estimates that finite
projection by aligning option-surface level, total-variance level, and downside total-variance skew
with the SPX/NDX restriction of the full Oxford--Man Perron direction.


\section{Empirical option-surface evidence}\label{sec:empirical}

The empirical exercise is a finite public-projection diagnostic.  The theory treats a local pricing
problem under a fixed \(\QQ\).  After conditioning on observation-date public information, a single
time-zero quote is a deterministic element of that local price system.  The panel uses the
date-indexed analogue: on date \(d\), the fitted option-surface coordinate is the realized value of
the same local public-summary coefficient map.  The regressions test the panel analogue of the
compression coefficient, not randomness of one time-zero quote.

The Perron state is rebuilt in a separate reproducibility archive from raw Oxford--Man
realized-volatility data and then matched to SPX/NDX OptionMetrics surfaces over the common window
2000--2018.  The \(\PP\)-side object is the latent Perron stress direction rebuilt on the full
eight-index Oxford--Man universe
\[
        \{\text{SPX},\text{NASDAQ},\text{FTSE},\text{DAX},\text{CAC},\text{EuroStoxx 50},
        \text{Nikkei},\text{Hang Seng}\}.
\]
The \(\QQ\)-side object is the SPX/NDX option-surface vector projected onto the SPX/NDX restriction
of the full Perron direction, not onto a Perron vector recomputed from the two-index submatrix.
The test asks whether public option-surface coordinates align with the already rebuilt
\(\PP\)-side direction, rather than calibrating the Perron state to options.

For feature \(f\), maturity bucket \(\tau\), and observation date \(d\), the panel
regression is
\[
        q_{d,\tau}^{(f)}
        =
        \alpha_{\tau}^{(f)}
        +
        \beta_{\tau}^{(f)}p_d
        +
        \varepsilon_{d,\tau}^{(f)}.
\]
Here \(p_d\) is the rebuilt \(\PP\)-side Perron coordinate and \(q_{d,\tau}^{(f)}\) is the
\(\QQ\)-side option-feature vector, centered ticker-by-ticker and projected onto the same
restricted full-universe Perron direction.  The reported slopes use the raw feature units.
Correlations, sign shares, and standardized signed compression-gap
diagnostics use separately standardized \(p_d\) and \(q_{d,\tau}^{(f)}\) within the selected
feature/state-mode/maturity cell.  Newey--West standard errors use the automatic lag rule in the
empirical code,
\[
        \ell=\max\{1,\min(20,\operatorname{round}(4(n/100)^{2/9}))\}.
\]
Regressions are estimated separately by maturity bucket; the reported HAC errors are time-series
Newey--West errors within each bucket rather than cross-sectional date-clustered errors.  Rolling
midpoint rows use one observation per rolling Perron window, while rolling completed-window rows
match each option date to the most recent completed rolling state.  For each matched date, let
\(Q_{d,\tau}^{(f)}\in\R^2\) be the centered SPX/NDX option-feature vector and write its Perron
coordinate as
\[
        q_{d,\tau}^{(f)}
        =
        \langle \widehat u_{\mathrm{opt}},Q_{d,\tau}^{(f)}\rangle,
        \qquad
        \Pi_{\mathrm{Perron}}Q_{d,\tau}^{(f)}
        =
        q_{d,\tau}^{(f)}\widehat v_{\mathrm{opt}},
\]
where \(\widehat u_{\mathrm{opt}}\) and \(\widehat v_{\mathrm{opt}}\) are the SPX/NDX restrictions
of the full left and right Perron vectors, normalized by
\(\|\widehat v_{\mathrm{opt}}\|_2=1\) and
\(\langle \widehat u_{\mathrm{opt}},\widehat v_{\mathrm{opt}}\rangle=1\).  With
\(a_{\mathrm{opt}}=(1/2,1/2)\), the off-Perron aggregate residual is
\[
        o_{d,\tau}^{(f)}
        =
        a_{\mathrm{opt}}^\top
        \bigl(Q_{d,\tau}^{(f)}-q_{d,\tau}^{(f)}\widehat v_{\mathrm{opt}}\bigr),
\]
and the off-Perron statistic reported below is
\[
        \frac{1}{n_\tau}\sum_{d=1}^{n_\tau}|o_{d,\tau}^{(f)}|,
\]
in the raw feature units of \(f\).  The statistic is therefore a residual magnitude diagnostic on
the two-index restriction rather than a scale-free percentage.  Placebo and residualized controls
compare the Perron coordinate with the off-Perron coordinate, rotated-date Perron series, median
random directions in the Perron/off-Perron plane, and regime/aggregate controls.  For the
downside total-variance skew coordinate, the residualized checks also control directly for the
contemporaneous ATM total-variance level.

The surface coordinates are deliberately low-dimensional.  The level coordinates are the fitted ATM
implied volatility and ATM total variance at 30, 60, and 90 days.  The tail-facing coordinate is the
signed downside total-variance skew
\[
        -\partial_k\{\tau\sigma_B(\tau,k)^2\}\big|_{k=0},
\]
estimated by a filtered local fit to the total-variance smile.  Raw implied-volatility skew is
retained as a secondary tail statistic, but the main tail coordinate is total-variance based because
it is the surface jet that matches the theorem-side scaling most directly.  The ingestion keeps
European index-option quotes with positive bid, offer, mid, implied volatility, and forward, valid
expiry dates, \(7\le \mathrm{DTE}\le365\), and strike scaled as \(\texttt{strike\_price}/1000\).
The local fit uses quotes in the maturity buckets \(\{30,60,90\}\) days with a ten-day half-width
and \(|k|\le0.35\).  The certified skew fit requires both SPX and NDX quotes to satisfy local-surface
quality conditions: enough quotes on both sides of ATM, near-ATM support, high direct total-variance
fit quality, and a well-conditioned local quadratic design.

\begin{table}[t]
\centering
\footnotesize
\begin{tabular}{lrrrrrr}
\hline
Feature & Days & Obs. & Corr. & \(R^2\) & NW \(t\) & Off-Perron \\
\hline
ATM IV level & 30 & 3207 & 0.807 & 0.651 & 23.219 & 0.058 \\
ATM IV level & 60 & 2903 & 0.805 & 0.648 & 25.040 & 0.054 \\
ATM IV level & 90 & 1997 & 0.801 & 0.641 & 21.607 & 0.050 \\
ATM total-variance level & 30 & 3207 & 0.708 & 0.501 & 13.338 & 0.004 \\
ATM total-variance level & 60 & 2903 & 0.734 & 0.539 & 15.409 & 0.005 \\
ATM total-variance level & 90 & 1997 & 0.738 & 0.545 & 14.247 & 0.006 \\
Downside total-variance skew & 30 & 3207 & 0.685 & 0.469 & 10.308 & 0.009 \\
Downside total-variance skew & 60 & 2903 & 0.692 & 0.479 & 10.642 & 0.007 \\
Downside total-variance skew & 90 & 1997 & 0.668 & 0.446 & 9.283 & 0.010 \\
\hline
\end{tabular}
\caption{Full-sample SPX/NDX compression evidence.  The \(\QQ\)-side option-surface Perron
coordinate is regressed on the rebuilt \(\PP\)-side Perron coordinate.  The final column
reports the mean absolute \(\QQ\)-side off-Perron component on the same two-index restriction.}
\label{tab:empirical-main-compression}
\end{table}

Table~\ref{tab:empirical-main-compression} is the descriptive benchmark.  ATM implied-volatility
level explains roughly two thirds of the Perron-aligned option movement, and ATM total-variance
level explains about one half.  Filtered downside total-variance skew has lower explanatory power
than level, but it remains a Perron-compressed tail coordinate.  Raw IV skew is weaker and less
stable.

\begin{table}[t]
\centering
\small
\begin{tabular}{lrrrrr}
\hline
Feature & Days & Midpoint $R^2$ & Midpoint $t$ & Completed $R^2$ & Completed $t$ \\
\hline
ATM IV level & 30 & 0.268 & 4.275 & 0.274 & 7.301 \\
ATM IV level & 60 & 0.220 & 4.022 & 0.282 & 7.211 \\
ATM IV level & 90 & 0.255 & 4.629 & 0.238 & 5.687 \\
ATM total variance & 30 & 0.218 & 3.082 & 0.232 & 5.312 \\
ATM total variance & 60 & 0.202 & 3.389 & 0.266 & 5.577 \\
ATM total variance & 90 & 0.248 & 3.902 & 0.233 & 4.674 \\
Downside TV skew & 30 & 0.305 & 4.116 & 0.267 & 6.660 \\
Downside TV skew & 60 & 0.300 & 4.682 & 0.324 & 7.242 \\
Downside TV skew & 90 & 0.330 & 4.664 & 0.279 & 5.948 \\
\hline
\end{tabular}
\caption{Rolling-window compression evidence.  The midpoint column uses one observation per rebuilt rolling
Perron window; the completed-window column matches each option date to the most recent completed
rolling state.}
\label{tab:empirical-rolling-robustness}
\end{table}

Table~\ref{tab:empirical-rolling-robustness} separates the descriptive full-sample benchmark from
rolling-state evidence.  The midpoint column uses one observation per rolling Perron window; the
completed-window column matches each option date to the most recent completed rolling state.  The
level effects fall, but remain statistically visible.  The downside total-variance skew retains a
nontrivial rolling signal, which is the empirical counterpart of the option-tail branch used below.

\begin{table}[t]
\centering
\small
\begin{tabular}{rrrrrr}
\hline
Days & Obs. & Corr. & $R^2$ & NW $t$ & Mean abs. off-Perron \\
\hline
30 & 3087 & 0.684 & 0.468 & 10.002 & 0.009 \\
60 & 2888 & 0.695 & 0.483 & 10.671 & 0.007 \\
90 & 1973 & 0.674 & 0.455 & 9.331 & 0.010 \\
\hline
\end{tabular}
\caption{Filtered downside total-variance skew statistics from the local surface fit.}
\label{tab:empirical-certified-skew}
\end{table}

The SPX/NDX pair supplies a focused two-index setting for the compression diagnostic.  The measured
coordinates match the finite-projection pattern.  The dominant public option coordinates are
compressed versions of the rebuilt latent Perron state, while skew becomes stable only after the
signed total-variance surface jet and the local-surface quality filters are imposed.

\subsection{Signed premium evidence}\label{subsec:empirical-premium-diagnostic}
The premium branch compares matched \(\PP\)- and \(\QQ\)-projected coefficients on a common basis.
The empirical implementation is conservative.  It reports the raw-unit transfer slope from the
regression above and separately records the standardized signed compression gap
\[
        \Delta^z_{M,\tau}
        =
        z\!\left(\text{\(\QQ\)-option Perron coordinate}\right)
        -
        z\!\left(\text{\(\PP\)-Perron stress coordinate}\right).
\]
Because the mean of this standardized signed compression gap is mechanically zero in the full sample, the
signed information comes from the raw transfer slope, the share of dates on which the
\(\PP\)- and \(\QQ\)-coordinates have the same standardized sign, and the difference between the
mean gap in the top and bottom \(\PP\)-stress quintiles.

\begin{table}[t]
\centering
\small
\begin{tabular}{lrrrrrr}
\hline
Feature & Days & Raw slope & \(R^2\) & NW \(t\) & Same sign & High-low gap \\
\hline
ATM IV level & 30 & 0.023 & 0.651 & 23.219 & 0.828 & -0.582 \\
ATM IV level & 60 & 0.022 & 0.648 & 25.040 & 0.832 & -0.599 \\
ATM IV level & 90 & 0.019 & 0.641 & 21.607 & 0.834 & -0.590 \\
ATM total variance & 30 & 0.001 & 0.501 & 13.338 & 0.776 & -0.910 \\
ATM total variance & 60 & 0.002 & 0.539 & 15.409 & 0.787 & -0.832 \\
ATM total variance & 90 & 0.003 & 0.545 & 14.247 & 0.805 & -0.786 \\
Downside TV skew & 30 & 0.002 & 0.468 & 10.002 & 0.770 & -1.009 \\
Downside TV skew & 60 & 0.003 & 0.483 & 10.671 & 0.776 & -1.008 \\
Downside TV skew & 90 & 0.003 & 0.455 & 9.331 & 0.764 & -1.006 \\
\hline
\end{tabular}
\caption{Signed common-basis compression statistics.  Slope is reported in raw feature units.
The high-low gap is the mean of \(\Delta^z_{M,\tau}\) in the top \(\PP\)-stress quintile minus the
corresponding mean in the bottom quintile.}
\label{tab:empirical-premium-diagnostic}
\end{table}

The signed diagnostic contains information that an absolute-gap proxy discards.  The transfer
slopes are positive and statistically separated from zero, and same-sign shares are high across the
level, total-variance, and filtered downside-skew coordinates.  The high-minus-low stress gap is
uniformly negative.  In high \(\PP\)-stress states the standardized option coordinate remains aligned
with the Perron direction but lies below the standardized physical stress coordinate.  The object is
a signed compression and shrinkage diagnostic rather than a direct economic premium magnitude.

\begin{table}[t]
\centering
\footnotesize
\begin{tabular}{llrrrr}
\hline
Feature & State & \(R^2\) range & NW \(t\) range & Same-sign range & High-low range \\
\hline
ATM IV level & Full sample & 0.641--0.651 & 21.607--25.040 & 0.828--0.834 & -0.599-- -0.582 \\
ATM IV level & Rolling daily & 0.238--0.282 & 5.687--7.301 & 0.600--0.624 & -1.546-- -1.427 \\
ATM IV level & Rolling midpoint & 0.220--0.268 & 4.022--4.629 & 0.610--0.657 & -1.543-- -1.320 \\
ATM total variance & Full sample & 0.501--0.545 & 13.338--15.409 & 0.776--0.805 & -0.910-- -0.786 \\
ATM total variance & Rolling daily & 0.232--0.266 & 4.674--5.577 & 0.604--0.622 & -1.654-- -1.514 \\
ATM total variance & Rolling midpoint & 0.202--0.248 & 3.082--3.902 & 0.598--0.635 & -1.554-- -1.469 \\
Downside TV skew & Full sample & 0.455--0.483 & 9.331--10.671 & 0.764--0.776 & -1.009-- -1.006 \\
Downside TV skew & Rolling daily & 0.257--0.324 & 5.952--7.243 & 0.656--0.686 & -1.560-- -1.344 \\
Downside TV skew & Rolling midpoint & 0.244--0.330 & 4.225--4.682 & 0.621--0.697 & -1.563-- -1.461 \\
\hline
\end{tabular}
\caption{Rolling signed common-basis premium statistics.  Ranges are taken
over the 30-, 60-, and 90-day maturity buckets.  Rolling daily uses the completed-window
convention; rolling midpoint uses one observation per rolling Perron window.}
\label{tab:empirical-premium-rolling}
\end{table}

Table~\ref{tab:empirical-premium-rolling} shows that the signed premium comparison survives the
same rebuilt rolling-state discipline used in the compression experiment.  The rolling
statistics are smaller than the full-sample benchmark, but transfer \(R^2\) and same-sign shares
remain positive, while the high-minus-low \(\PP\)-stress gap remains negative across the three
option-surface coordinates.  The comparison is a targeted two-index signed comparison aligned with
the rolling implementation used for the core compression evidence.

\section{Pairwise latent-contagion premium on a matched \texorpdfstring{\(\PP/\QQ\)}{P/Q} surface}\label{sec:premium}

The premium section compares the visible pricing geometry under \(\QQ\) with the projected
physical-measure object supplied by the projected-score construction of \cite{vidal2026econometric}.
The estimand is a signed coefficient difference.  Both sides are projected onto the same
source-screened basis, and the \(\QQ\)-side comes from a cleaned option-implied derivative surface
observed on dates matched to the realized-volatility sample.  Option-implied surface information is
used in the state-price-density and local-surface sense \cite{breeden1978,dupire1994}; the realized
side follows the Oxford-Man/realized-kernel provenance \cite{heber2009realized,barndorffNielsen2008realized,
shephardSheppard2010heavy}.  The object is the
matched \(\PP/\QQ\) projected coefficient, not a standalone roughness comparison.

\begin{definition}[Harmonized pairwise projected basis and premiums]\label{def:premium}
Fix an ordered pair \((i,j)\).  Let \(\varphi_{ij}\) denote the common unit-norm source-screened
projected basis direction supplied by the projected-score geometry of
\cite{vidal2026econometric} for that pair,
normalized identically under \(\PP\) and \(\QQ\).  Let \(D_{ij}^\PP\) and \(D_{ij}^\QQ\) denote the corresponding local
covariance-derivative objects under \(\PP\) and \(\QQ\), represented on the same reduced
experiment.  Let \(L_{ij}^\PP\) and \(L_{ij}^\QQ\) denote the pairwise priced structural loadings of
those reduced experiments on the harmonized basis, so that \(D_{ij}^\PP\) and \(D_{ij}^\QQ\) are
their public projected representatives.  Define the projected pairwise structural-loading coefficients
\[
        \gamma_{ij}^\PP
        :=
        \langle D_{ij}^\PP,\varphi_{ij}\rangle_{\PP,ij},
        \qquad
        \gamma_{ij}^\QQ
        :=
        \langle D_{ij}^\QQ,\varphi_{ij}\rangle_{\QQ,ij},
\]
the pairwise latent-contagion premium
\[
        \Pi_{ij}:=\gamma_{ij}^\QQ-\gamma_{ij}^\PP,
\]
and the companion reduced-form roughness gap
\[
        \Delta \widetilde H_{ij}:=\widetilde H_{ij}^\QQ-\widetilde H_{ij}^\PP,
\]
where \(\widetilde H_{ij}^\PP\) and \(\widetilde H_{ij}^\QQ\) are pairwise reduced-form roughness
exponents estimated from the harmonized physical and option-implied experiments, not the
structural Volterra exponents of Section~\ref{sec:model}.
\end{definition}

\begin{remark}[Why the premium is coefficient-based]
The primary premium object is \(\Pi_{ij}\), with \(\Delta \widetilde H_{ij}\) retained only as a
secondary roughness summary.  A coefficient gap is closer to the projected structural-loading
geometry, uses the same source-screened normalization on both sides, and maps directly to the
projected-score machinery of \cite{vidal2026econometric}.
\end{remark}

\begin{assumption}[Matched cleaned derivative-surface experiment]\label{ass:premium}
For each ordered pair \((i,j)\), there exist matched observation dates
\[
        d=1,\dots,n,
\]
a maturity grid \(\mathcal T_n\subset(0,\bar\tau]\), and a strike or log-moneyness grid
\(\mathcal M_n\subset[-\bar k,\bar k]\) such that:
\begin{enumerate}[label=(\roman*)]
\item on every matched date \(d\), the \(\PP\)-side reduced experiment is represented by a
projected-score object \(\widehat D_{ij,d}^\PP\) built from the realized-volatility geometry of
\cite{vidal2026econometric};
\item on the same matched date \(d\), the \(\QQ\)-side reduced experiment is represented by a
cleaned, arbitrage-regularized option surface
\[
        \widehat C_{i,d}(\tau,k)
        \quad\text{or equivalently}\quad
        \widehat \sigma_{i,d}^{\mathrm{imp}}(\tau,k),
        \qquad
        (\tau,k)\in\mathcal T_n\times \mathcal M_n,
\]
from which a local derivative object \(\widehat D_{ij,d}^\QQ\) is extracted by a linear
surface-functional compatible with the same source-screened nuisance projection and the same
pairwise normalization as on the \(\PP\)-side;
\item both sides are projected onto the same harmonized source-screened basis \(\varphi_{ij}\),
yielding paired coefficient observations
\[
        \widehat g_{ij,d}^\PP
        :=
        \langle \widehat D_{ij,d}^\PP,\varphi_{ij}\rangle_{\PP,ij},
        \qquad
        \widehat g_{ij,d}^\QQ
        :=
        \langle \widehat D_{ij,d}^\QQ,\varphi_{ij}\rangle_{\QQ,ij},
\]
and sample averages
\[
        \widehat\gamma_{ij,n}^\PP
        :=
        \frac1n\sum_{d=1}^n \widehat g_{ij,d}^\PP,
        \qquad
        \widehat\gamma_{ij,n}^\QQ
        :=
        \frac1n\sum_{d=1}^n \widehat g_{ij,d}^\QQ;
\]
\item the paired estimator vector satisfies a joint central limit theorem:
\[
        \sqrt{n}
        \begin{pmatrix}
                \widehat\gamma_{ij,n}^\PP-\gamma_{ij}^\PP\\
                \widehat\gamma_{ij,n}^\QQ-\gamma_{ij}^\QQ
        \end{pmatrix}
        \Rightarrow
        \mathcal N\!\left(
                0,\,
                \begin{pmatrix}
                        \Sigma_{ij}^\PP & \Sigma_{ij}^{\PP\QQ}\\
                        \Sigma_{ij}^{\PP\QQ} & \Sigma_{ij}^\QQ
                \end{pmatrix}
        \right).
\]
\end{enumerate}
\end{assumption}

\begin{definition}[Cleaned option-calibration panel]\label{def:calibration-panel}
For each matched date \(d=1,\dots,n\), asset or index \(i\), and quote index
\(\ell=1,\dots,L_{i,d}\), observe a mid quote or implied-volatility quote
\[
        Y_{i,d,\ell}^{\mathrm{obs}}
\]
at maturity--moneyness coordinate
\[
        x_{i,d,\ell}:=(\tau_{i,d,\ell},k_{i,d,\ell})\in(0,\bar\tau]\times[-\bar k,\bar k].
\]
A cleaned \(\QQ\)-surface is an estimator
\[
        \widehat{\mathsf C}_{i,d}\in\cX
\]
obtained from the quote cloud by weighted least squares, likelihood, or constrained smoothing over
a no-arbitrage admissible class \(\Theta_{\mathrm{arb}}\):
\[
        \widehat\theta_{i,d}
        \in
        \arg\min_{\theta\in\Theta_{\mathrm{arb}}}
        \sum_{\ell=1}^{L_{i,d}}w_{i,d,\ell}
        \bigl(Y_{i,d,\ell}^{\mathrm{obs}}-\mathsf C_\theta(x_{i,d,\ell})\bigr)^2
        +\operatorname{pen}_{i,d}(\theta),
        \qquad
        \widehat{\mathsf C}_{i,d}:=\mathsf C_{\widehat\theta_{i,d}}.
\]
Here \(\cX\) is a Banach space strong enough to control the local surface derivatives used below,
for example a weighted \(C^r\)-norm on the maturity--moneyness window.
\end{definition}

\begin{definition}[Option-surface coefficient extraction]\label{def:surface-coeff-extraction}
Fix an ordered pair \((i,j)\).  Let
\[
        \cA_{ij}:\cX\to\cD_{ij}
\]
be the local linear derivative map that extracts the \(\QQ\)-side pairwise derivative object from a
cleaned option surface.  Examples include local ATM-skew derivatives, local total-variance
curvature derivatives, or a vector of maturity-local derivatives.  Let
\[
        P_{ij}:\cD_{ij}\to\cD_{ij}^{\mathrm{scr}}
\]
be the source-screening nuisance projection, \(N_{ij}>0\) the pairwise normalizer, and
\(\varphi_{ij}\in\cD_{ij}^{\mathrm{scr}}\) the unit harmonized source-screened basis vector.  The
population \(\QQ\)-side coefficient observation on date \(d\) is
\[
        g_{ij,d}^{\QQ}
        :=
        \left\langle
                N_{ij}^{-1}P_{ij}\cA_{ij}(\mathsf C_{i,d}^{\QQ}),
                \varphi_{ij}
        \right\rangle_{ij},
\]
where \(\mathsf C_{i,d}^{\QQ}\) is the latent arbitrage-free option surface on that date.  The empirical
coefficient observation is
\[
        \widehat g_{ij,d}^{\QQ}
        :=
        \left\langle
                \widehat N_{ij,n}^{-1}\widehat P_{ij,n}\cA_{ij}(\widehat{\mathsf C}_{i,d}),
                \widehat\varphi_{ij,n}
        \right\rangle_{ij,n}.
\]
The \(\PP\)-side coefficient \(\widehat g_{ij,d}^{\PP}\) is defined analogously from the
realized-volatility projected-score experiment on the same date.
\end{definition}

\begin{assumption}[Derivative-stable calibration and harmonization]\label{ass:derivative-stable-calibration}
For each ordered pair \((i,j)\), the following hold.
\begin{enumerate}[label=(\roman*)]
\item The derivative extraction map \(\cA_{ij}\) is bounded linear from \(\cX\) to \(\cD_{ij}\),
and \(P_{ij}\) is a bounded projection.
\item The empirical projection, basis, inner product, and normalizer are stable in the sense that
substituting
\[
        (\widehat P_{ij,n},\widehat\varphi_{ij,n},\widehat N_{ij,n},
        \langle\cdot,\cdot\rangle_{ij,n})
\]
for
\[
        (P_{ij},\varphi_{ij},N_{ij},\langle\cdot,\cdot\rangle_{ij})
\]
changes the sample mean of the \(\QQ\)-side coefficient observations by \(o_\PP(n^{-1/2})\).
\item The cleaned surface has the coefficient-level asymptotic linear representation
\[
        \widehat g_{ij,d}^{\QQ}
        =
        g_{ij,d}^{\QQ}
        +
        \zeta_{ij,d}^{\QQ}
        +
        b_{ij,n}^{\QQ}
        +
        r_{ij,d,n}^{\QQ},
\]
where \(\zeta_{ij,d}^{\QQ}\) is the calibration influence contribution,
\(b_{ij,n}^{\QQ}\) is a deterministic smoothing or regularization bias, and
\[
        \sqrt n\, b_{ij,n}^{\QQ}\to b_{ij}^{\QQ},
        \qquad
        \frac1{\sqrt n}\sum_{d=1}^n r_{ij,d,n}^{\QQ}=o_\PP(1).
\]
The same representation holds on the \(\PP\)-side, with influence term
\(\zeta_{ij,d}^{\PP}\) and bias limit \(b_{ij}^{\PP}\).
\item The matched vector
\[
        Y_{ij,d}
        :=
        \begin{pmatrix}
                g_{ij,d}^{\PP}-\gamma_{ij}^{\PP}+\zeta_{ij,d}^{\PP}\\
                g_{ij,d}^{\QQ}-\gamma_{ij}^{\QQ}+\zeta_{ij,d}^{\QQ}
        \end{pmatrix}
\]
is strictly stationary and satisfies a bivariate central limit theorem
\[
        \frac1{\sqrt n}\sum_{d=1}^nY_{ij,d}
        \Rightarrow
        \mathcal N(0,\Omega_{ij}^{\mathrm{cal}}).
\]
\end{enumerate}
\end{assumption}

\begin{theorem}[Calibration-to-premium coefficient transfer]\label{thm:calibration-to-premium}
Under Assumption~\ref{ass:derivative-stable-calibration}, define
\[
        \widehat\gamma_{ij,n}^{\mu}
        :=
        \frac1n\sum_{d=1}^n\widehat g_{ij,d}^{\mu},
        \qquad
        \mu\in\{\PP,\QQ\},
\]
and
\[
        \widehat\Pi_{ij,n}:=\widehat\gamma_{ij,n}^{\QQ}-\widehat\gamma_{ij,n}^{\PP}.
\]
Then
\[
        \sqrt n
        \begin{pmatrix}
                \widehat\gamma_{ij,n}^{\PP}-\gamma_{ij}^{\PP}\\
                \widehat\gamma_{ij,n}^{\QQ}-\gamma_{ij}^{\QQ}
        \end{pmatrix}
        \Rightarrow
        \mathcal N\!\left(
                \begin{pmatrix}
                        b_{ij}^{\PP}\\ b_{ij}^{\QQ}
                \end{pmatrix},
                \Omega_{ij}^{\mathrm{cal}}
        \right).
\]
Consequently,
\[
        \sqrt n\left(\widehat\Pi_{ij,n}-\Pi_{ij}\right)
        \Rightarrow
        \mathcal N\!\left(
                b_{ij}^{\QQ}-b_{ij}^{\PP},
                \begin{pmatrix}-1&1\end{pmatrix}
                \Omega_{ij}^{\mathrm{cal}}
                \begin{pmatrix}-1\\1\end{pmatrix}
        \right).
\]
If the regularization biases are undersmoothed or bias-corrected so that
\(b_{ij}^{\PP}=b_{ij}^{\QQ}=0\), this is exactly the matched premium CLT in
Proposition~\ref{thm:premium-estimator}, with the covariance matrix enlarged to include
option-calibration and realized-volatility estimation noise.
\end{theorem}

\begin{proof}
The coefficient-level representations imply
\[
        \sqrt n
        \begin{pmatrix}
                \widehat\gamma_{ij,n}^{\PP}-\gamma_{ij}^{\PP}\\
                \widehat\gamma_{ij,n}^{\QQ}-\gamma_{ij}^{\QQ}
        \end{pmatrix}
        =
        \frac1{\sqrt n}\sum_{d=1}^nY_{ij,d}
        +
        \begin{pmatrix}
                \sqrt n\, b_{ij,n}^{\PP}\\
                \sqrt n\, b_{ij,n}^{\QQ}
        \end{pmatrix}
        +o_\PP(1).
\]
The assumed CLT and bias limits give the bivariate limit.  The premium estimator is the linear map
\((x_P,x_Q)\mapsto x_Q-x_P\), so the second claim follows by the continuous mapping theorem.
\end{proof}

\begin{proposition}[A simple sufficient rate contract for negligible calibration error]\label{prop:surface-rate-contract}
Suppose the harmonization objects are deterministic or estimated on an auxiliary sample, and
there exists a deterministic sequence \(a_n\downarrow0\) such that
\[
        \max_{1\le d\le n}\|\widehat{\mathsf C}_{i,d}-\mathsf C_{i,d}^{\QQ}\|_{\cX}=O_\PP(a_n),
        \qquad
        \sqrt n\,a_n\to0.
\]
Then the calibration contribution to
\(\widehat\gamma_{ij,n}^{\QQ}\) is \(o_\PP(n^{-1/2})\).  In this undersmoothed regime the cleaned
surface can be treated as providing the oracle coefficient observations
\(g_{ij,d}^{\QQ}\) at the first-order premium scale.
\end{proposition}

\begin{proof}
By boundedness of \(\cA_{ij}\), \(P_{ij}\), and the coefficient functional, there is a constant
\(C_{ij}\) such that
\[
        |\widehat g_{ij,d}^{\QQ}-g_{ij,d}^{\QQ}|
        \le
        C_{ij}\|\widehat{\mathsf C}_{i,d}-\mathsf C_{i,d}^{\QQ}\|_{\cX}
\]
up to the separately controlled harmonization error.  Hence
\[
        \sqrt n\left|
        \frac1n\sum_{d=1}^n(\widehat g_{ij,d}^{\QQ}-g_{ij,d}^{\QQ})
        \right|
        \le
        C_{ij}\sqrt n\max_d\|\widehat{\mathsf C}_{i,d}-\mathsf C_{i,d}^{\QQ}\|_{\cX}
        =o_\PP(1).
\]
\end{proof}

\begin{remark}[Operational calibration path for premium coefficients]\label{rem:premium-calibration-path}
The premium coefficients are no longer abstract observations.  The market data path is
quote cloud \(\to\) no-arbitrage cleaned surface \(\widehat{\mathsf C}_{i,d}\) \(\to\) local derivative
object \(\cA_{ij}(\widehat{\mathsf C}_{i,d})\) \(\to\) source-screened projection and normalization
\(\to\) coefficient observation \(\widehat g_{ij,d}^{\QQ}\) \(\to\) matched premium estimator
\(\widehat\Pi_{ij,n}\).  The transfer theorem identifies where calibration noise enters the
asymptotic variance and when undersmoothing or bias correction removes it at first order.
\end{remark}

\begin{proposition}[Matched pairwise premium estimator from a cleaned \texorpdfstring{\(\PP/\QQ\)}{P/Q} surface]\label{thm:premium-estimator}
Under Assumption~\ref{ass:premium}, define
\[
        \widehat\Pi_{ij,n}:=\widehat\gamma_{ij,n}^\QQ-\widehat\gamma_{ij,n}^\PP.
\]
Then \(\widehat\Pi_{ij,n}\to\Pi_{ij}\) in probability and
\[
        \sqrt{n}\,(\widehat\Pi_{ij,n}-\Pi_{ij})
        \Rightarrow
        \mathcal N\!\left(
                0,\,
                \Sigma_{ij}^\QQ
                +
                \Sigma_{ij}^\PP
                -
                2\,\Sigma_{ij}^{\PP\QQ}
        \right).
\]
In particular, the paired covariance term is the default asymptotic correction when the realized
and option-implied inputs are observed on common dates after option-surface cleaning.
\end{proposition}

\begin{proof}
Consistency follows because \(\widehat\Pi_{ij,n}-\Pi_{ij}\) is the difference of the two consistent
coefficient estimators.  For the asymptotic law, write
\[
        \sqrt{n}\,(\widehat\Pi_{ij,n}-\Pi_{ij})
        =
        \begin{pmatrix}
                -1 & 1
        \end{pmatrix}
        \sqrt{n}
        \begin{pmatrix}
                \widehat\gamma_{ij,n}^\PP-\gamma_{ij}^\PP\\
                \widehat\gamma_{ij,n}^\QQ-\gamma_{ij}^\QQ
        \end{pmatrix}.
\]
The variance is the quadratic form of the joint covariance matrix under the contrast \((-1,1)\):
\[
        (-1,1)
        \begin{pmatrix}
        \Sigma_{ij}^\PP & \Sigma_{ij}^{\PP\QQ}\\
        \Sigma_{ij}^{\PP\QQ} & \Sigma_{ij}^\QQ
        \end{pmatrix}
        \begin{pmatrix}-1\\1\end{pmatrix}
        =
        \Sigma_{ij}^\QQ+\Sigma_{ij}^\PP-2\Sigma_{ij}^{\PP\QQ}.
\]
This is the multivariate delta method for the linear map \((x,y)\mapsto y-x\).
\end{proof}

\begin{remark}[Asymptotic independence as a special case]
If the \(\PP\)- and \(\QQ\)-side inputs are asymptotically independent, or if the paired design is
replaced by asymptotically disjoint samples, then \(\Sigma_{ij}^{\PP\QQ}=0\) and
Proposition~\ref{thm:premium-estimator} simplifies to the familiar variance sum
\[
        \Sigma_{ij}^\QQ+\Sigma_{ij}^\PP.
\]
The paired covariance term is retained because matched dates are the natural design
for comparing realized and option-implied structural loading on the same market states.
\end{remark}

\begin{corollary}[Pairwise \(\PP/\QQ\) classification on a matched cleaned basis]\label{cor:premium-class}
Fix tolerance levels \(\eta_P,\eta_Q,\eta_\Pi>0\).  For each ordered pair \((i,j)\), the common
projected basis \(\varphi_{ij}\) and the cleaned matched \(\QQ\)-surface yield the following
population classification:
\begin{enumerate}[label=(\roman*)]
\item \emph{dynamically observable}: \(|\gamma_{ij}^\PP|\ge\eta_P\) and
\(|\gamma_{ij}^\QQ|\ge\eta_Q\);
\item \emph{pricing-visible only}: \(|\gamma_{ij}^\PP|<\eta_P\) and
\(|\gamma_{ij}^\QQ|\ge\eta_Q\);
\item \emph{path-detectable only}: \(|\gamma_{ij}^\PP|\ge\eta_P\) and
\(|\gamma_{ij}^\QQ|<\eta_Q\);
\item \emph{neutral/low-premium}: \(|\gamma_{ij}^\PP|<\eta_P\),
\(|\gamma_{ij}^\QQ|<\eta_Q\), and \(|\Pi_{ij}|<\eta_\Pi\);
\item \emph{premium-dominant}: \(|\Pi_{ij}|\ge\eta_\Pi\) on the common basis, even after
harmonized source-screening.
\end{enumerate}
Confidence-region analogues follow from Proposition~\ref{thm:premium-estimator}.
\end{corollary}

\begin{definition}[Pairwise premium support]\label{def:premium-support}
For \(\eta_\Pi\ge0\), define the pairwise premium support at level \(\eta_\Pi\) by
\[
        \cE_\Pi(\eta_\Pi)
        :=
        \left\{
                (i,j): |\Pi_{ij}|\ge \eta_\Pi
        \right\}.
\]
When \(\eta_\Pi=0\), write
\[
        \cE_\Pi
        :=
        \left\{
                (i,j): \Pi_{ij}\neq0
        \right\}.
\]
\end{definition}

\begin{proposition}[Separated \(\PP/\QQ\) observability forces premium support]\label{prop:premium-support-separated}
Fix \(\eta_P,\eta_Q>0\) and an ordered pair \((i,j)\).
\begin{enumerate}[label=(\roman*)]
\item If \(\eta_P>\eta_Q\) and
\[
        |\gamma_{ij}^\PP|\ge \eta_P,
        \qquad
        |\gamma_{ij}^\QQ|<\eta_Q,
\]
then
\[
        |\Pi_{ij}|
        \ge
        \eta_P-\eta_Q.
\]
Consequently,
\[
        (i,j)\in \cE_\Pi(\eta_P-\eta_Q)\subseteq \cE_\Pi.
\]
\item If \(\eta_Q>\eta_P\) and
\[
        |\gamma_{ij}^\QQ|\ge \eta_Q,
        \qquad
        |\gamma_{ij}^\PP|<\eta_P,
\]
then
\[
        |\Pi_{ij}|
        \ge
        \eta_Q-\eta_P.
\]
Consequently,
\[
        (i,j)\in \cE_\Pi(\eta_Q-\eta_P)\subseteq \cE_\Pi.
\]
\end{enumerate}
In particular, every pair whose harmonized projected coefficients are separated across measures by
a strict visibility gap on the common source-screened basis belongs to the nontrivial premium
support.
\end{proposition}

\begin{proof}
Since
\[
        \Pi_{ij}=\gamma_{ij}^\QQ-\gamma_{ij}^\PP,
\]
the reverse triangle inequality gives
\[
        |\Pi_{ij}|
        \ge
        \bigl||\gamma_{ij}^\PP|-|\gamma_{ij}^\QQ|\bigr|.
\]
Under the hypotheses of part~(i),
\[
        |\Pi_{ij}|
        >
        \eta_P-\eta_Q.
\]
This proves part~(i).  Part~(ii) is identical after exchanging \(\PP\) and \(\QQ\).
\end{proof}

\begin{corollary}[Structural support inclusion for the pairwise premium]\label{cor:premium-support-structural}
Fix an ordered pair \((i,j)\) and suppose the harmonized projected basis \(\varphi_{ij}\) realizes
the same source-screened structural direction under \(\PP\) and \(\QQ\).
\begin{enumerate}[label=(\roman*)]
\item If there exist \(0<\eta_Q<\eta_P\) such that the \(\PP\)-side pair is separated from zero on
the common basis,
\[
        |\gamma_{ij}^\PP|\ge \eta_P,
\]
while the matched cleaned \(\QQ\)-side surface places the same pair below the pricing-visibility
threshold on that basis,
\[
        |\gamma_{ij}^\QQ|<\eta_Q,
\]
then
\[
        (i,j)\in \cE_\Pi(\eta_P-\eta_Q)\subseteq \cE_\Pi.
\]
\item If there exist \(0<\eta_P<\eta_Q\) such that the matched cleaned \(\QQ\)-side surface is
separated from zero on the common basis,
\[
        |\gamma_{ij}^\QQ|\ge \eta_Q,
\]
while the \(\PP\)-side projected-score geometry places the same pair below the physical
source-screened threshold on that basis,
\[
        |\gamma_{ij}^\PP|<\eta_P,
\]
then
\[
        (i,j)\in \cE_\Pi(\eta_Q-\eta_P)\subseteq \cE_\Pi.
\]
\end{enumerate}
Equivalently, every harmonized pair with path-detectability alone or pricing-visibility alone
must belong to the premium support.
\end{corollary}

\begin{proof}
For part~(i), Proposition~\ref{prop:premium-support-separated}(i) says that a pair with
\(|\gamma_{ij}^{\PP}|\ge\eta_P\) and \(|\gamma_{ij}^{\QQ}|<\eta_Q\) has
\(|\gamma_{ij}^{\QQ}-\gamma_{ij}^{\PP}|\ge \eta_P-\eta_Q\).  Part~(ii) is the same inequality with
\(\PP\) and \(\QQ\) interchanged, using Proposition~\ref{prop:premium-support-separated}(ii).
Hence any support disagreement above the stated thresholds produces a nonzero signed premium
coefficient.
\end{proof}

\begin{remark}[How visibility inputs feed the support inclusion]\label{rem:premium-support-trilogy}
The structural content of Corollary~\ref{cor:premium-support-structural} is that premium support
is now tied to \(\PP/\QQ\) visibility separation rather than left entirely at the level of a fixed
pairwise estimator.  On the \(\PP\)-side, \cite{vidal2026econometric,vidal2026market} identify the
source-screened, physically nondegenerate pairwise channels.  On the \(\QQ\)-side,
\cite{vidal2026latent} identifies when a pairwise contagion channel is pricing-visible or screened
at the tested maturity scale.  Whenever the matched cleaned \(\QQ\)-surface and the projected-score
\(\PP\)-geometry are harmonized on the same source-screened basis \(\varphi_{ij}\), those structural
inputs feed directly into the support inclusion theorem above.
\end{remark}

\begin{proposition}[One-dimensional harmonized pairwise nondegeneracy is coefficient nonvanishing]\label{prop:premium-coeff-nonzero}
Fix an ordered pair \((i,j)\), let \(\mu\in\{\PP,\QQ\}\), and suppose the harmonized reduced
experiment for that pair is one-dimensional on the common source-screened basis \(\varphi_{ij}\),
so that
\[
        D_{ij}^\mu=\gamma_{ij}^\mu \varphi_{ij}.
\]
Then
\[
        D_{ij}^\mu\neq0
        \qquad\Longleftrightarrow\qquad
        \gamma_{ij}^\mu\neq0.
\]
\end{proposition}

\begin{proof}
Since \(\varphi_{ij}\) has unit norm, the displayed representation gives
\[
        \|D_{ij}^\mu\|_{\mu,ij}=|\gamma_{ij}^\mu|.
\]
Thus \(D_{ij}^\mu=0\) if and only if its norm is zero, and this is equivalent to
\(\gamma_{ij}^\mu=0\).
\end{proof}

\begin{corollary}[Exact premium support from harmonized path-only or pricing-only pairs]\label{cor:premium-support-exact}
Fix an ordered pair \((i,j)\) and suppose the harmonized reduced experiments under \(\PP\) and
\(\QQ\) are both one-dimensional on the common source-screened basis \(\varphi_{ij}\).
\begin{enumerate}[label=(\roman*)]
\item If the \(\PP\)-side pair is physically nondegenerate on the harmonized basis,
\[
        D_{ij}^\PP\neq0,
\]
while the matched cleaned \(\QQ\)-side pair is screened on that same basis at the tested scale,
\[
        D_{ij}^\QQ=0,
\]
then
\[
        \Pi_{ij}\neq0,
        \qquad
        (i,j)\in\cE_\Pi.
\]
\item If the matched cleaned \(\QQ\)-side pair is pricing-nondegenerate on the harmonized basis,
\[
        D_{ij}^\QQ\neq0,
\]
while the \(\PP\)-side source-screened reduced loading vanishes on that same basis,
\[
        D_{ij}^\PP=0,
\]
then
\[
        \Pi_{ij}\neq0,
        \qquad
        (i,j)\in\cE_\Pi.
\]
\end{enumerate}
In particular, exact harmonized path-only pairs and exact harmonized pricing-only pairs must belong
to the pairwise premium support.
\end{corollary}

\begin{proof}
Under the one-dimensional reduction,
Proposition~\ref{prop:premium-coeff-nonzero} gives
\[
        D_{ij}^\PP\neq0\Longleftrightarrow \gamma_{ij}^\PP\neq0,
        \qquad
        D_{ij}^\QQ\neq0\Longleftrightarrow \gamma_{ij}^\QQ\neq0.
\]
If \(D_{ij}^\PP\neq0\) and \(D_{ij}^\QQ=0\), then
\[
        \Pi_{ij}
        =
        \gamma_{ij}^\QQ-\gamma_{ij}^\PP
        =
        -\gamma_{ij}^\PP
        \neq0.
\]
This proves part~(i).  Part~(ii) is identical after exchanging \(\PP\) and \(\QQ\).
\end{proof}

\begin{remark}[Exact versus tolerance-separated premium support]\label{rem:premium-support-exact}
Corollary~\ref{cor:premium-support-exact} is the exact-support counterpart to
Proposition~\ref{prop:premium-support-separated}.  The former is the sharp statement when the
harmonized pairwise reduction is genuinely one-dimensional and one side vanishes exactly on the
common source-screened basis.  The latter is the threshold-separated asymptotic statement when only a visibility
gap separated by thresholds \(\eta_P,\eta_Q\) is available.
\end{remark}


The one-dimensional projected-information estimator, covariance-normalizer construction,
finite-dimensional if-and-only-if support law, and unsigned-information counterexamples are
implementation infrastructure.  They are collected in Appendix~\ref{sec:premium-implementation-appendix}.
The body keeps the premium claim at coefficient level.  Signed matched coefficients identify the
priced \(\PP/\QQ\) gap on a common source-screened basis; unsigned information is only a support
diagnostic unless signs and normalizers are tracked.

\section{Variance claims and risk-neutral compression}\label{sec:variance}

Index-option summaries are only one public class.  Nonlinear claims expose the same compression
operator through convexity.  Corollary~\ref{prop:compression-unifier} makes the public
variance-pricing object a compression of the claim's priced structural loading.  The section first
records the abstract compression theorem for square-integrable claims and then specializes it to the
market-wide object
\[
        Q_M:=\int_0^T M(t)^2\,\dd t.
\]
The market-variance motivation is adjacent to the log-contract and volatility-derivatives literature
\cite{neuberger1994,carrlee2009}.  The result is a compression statement, not a variance-swap
replication formula.

\begin{assumption}[Raw visible-plus-synchronized aggregate reduction under \(\QQ\)]\label{ass:q-raw-visible-aggregate}
There exist a visible process \(M^{\mathrm{vis},\QQ}\), the structural synchronized amplitude
process \(\varpi\), and a raw residual process \(\varepsilon^{\QQ,\mathrm{raw}}\) such that
\[
        M(t)=M^{\mathrm{vis},\QQ}(t)+(w^\top v)\varpi(t)+\varepsilon^{\QQ,\mathrm{raw}}(t),
        \qquad
        0\le t\le T,
\]
where \(M^{\mathrm{vis},\QQ}(t)\) is \(\cG_t\)-measurable for every \(t\).  Assume further that
\[
        \int_0^T \bigl(M^{\mathrm{vis},\QQ}(t)\bigr)^2\,\dd t,
        \qquad
        \int_0^T \varpi(t)^2\,\dd t,
        \qquad
        \int_0^T \bigl(\varepsilon^{\QQ,\mathrm{raw}}(t)\bigr)^2\,\dd t
\]
all belong to \(L^2(\QQ,\cF_T)\), and that the conditional Fubini operations below are valid.
\end{assumption}

\begin{definition}[Risk-neutral aggregate-centering defect]\label{def:q-centering-defect}
Under Assumption~\ref{ass:q-raw-visible-aggregate}, set
\[
        \cH_\varpi:=\cG_T\vee\cL_T,
        \qquad
        \zeta(t):=(w^\top v)\varpi(t),
        \qquad
        A^\QQ(t):=M^{\mathrm{vis},\QQ}(t)+\zeta(t).
\]
Define the \(\QQ\)-aggregate-centering defect by
\[
        b_\varpi^\QQ(t)
        :=
        \EE_\QQ[\varepsilon^{\QQ,\mathrm{raw}}(t)\mid \cH_\varpi],
        \qquad
        0\le t\le T,
\]
and the recentered residual by
\[
        \eta_\varpi^{\QQ,\perp}(t)
        :=
        \varepsilon^{\QQ,\mathrm{raw}}(t)-b_\varpi^\QQ(t).
\]
Thus \(b_\varpi^\QQ\) measures the failure of the aggregate residual to remain centered after
the state-price-density reweighting.
\end{definition}

\begin{proposition}[Risk-neutral recentering of the aggregate reduction]\label{prop:q-recentering}
Under Assumption~\ref{ass:q-raw-visible-aggregate} and
Definition~\ref{def:q-centering-defect},
\[
        M(t)=A^\QQ(t)+b_\varpi^\QQ(t)+\eta_\varpi^{\QQ,\perp}(t),
        \qquad
        0\le t\le T,
\]
and
\[
        \EE_\QQ[\eta_\varpi^{\QQ,\perp}(t)\mid\cH_\varpi]=0
        \qquad
        \text{for almost every }t\in[0,T].
\]
Moreover \(b_\varpi^\QQ\) and \(\eta_\varpi^{\QQ,\perp}\) inherit the required square
integrability, and the priced structural loading of integrated market variance is
\[
        L_{Q_M}^\QQ
        =
        \int_0^T \bigl(A^\QQ(t)+b_\varpi^\QQ(t)\bigr)^2\,\dd t
        +
        \int_0^T
        \EE_\QQ\!\left[
                \bigl(\eta_\varpi^{\QQ,\perp}(t)\bigr)^2
                \mid \cH_\varpi
        \right]\dd t.
\]
\end{proposition}

\begin{proof}
The decomposition is the definition of \(b_\varpi^\QQ\) and
\(\eta_\varpi^{\QQ,\perp}\).  Conditional centering follows from
\[
        \EE_\QQ[\eta_\varpi^{\QQ,\perp}(t)\mid\cH_\varpi]
        =
        \EE_\QQ[\varepsilon^{\QQ,\mathrm{raw}}(t)\mid\cH_\varpi]
        -
        b_\varpi^\QQ(t)
        =
        0.
\]
Conditional Jensen and conditional Fubini give
\[
        \int_0^T \bigl(b_\varpi^\QQ(t)\bigr)^2\,\dd t
        \le
        \EE_\QQ\!\left[
                \int_0^T \bigl(\varepsilon^{\QQ,\mathrm{raw}}(t)\bigr)^2\,\dd t
                \mid \cH_\varpi
        \right],
\]
so \(b_\varpi^\QQ\) is square-integrable at the required integrated scale; the same follows for
\(\eta_\varpi^{\QQ,\perp}\) from \((x-y)^2\le2x^2+2y^2\).  Finally,
\[
        Q_M
        =
        \int_0^T
        \bigl(A^\QQ(t)+b_\varpi^\QQ(t)+\eta_\varpi^{\QQ,\perp}(t)\bigr)^2\,\dd t.
\]
Taking \(\EE_\QQ[\cdot\mid\cH_\varpi]\), using that \(A^\QQ+b_\varpi^\QQ\) is
\(\cH_\varpi\)-measurable, and applying the displayed centering removes the cross term and gives
the formula for \(L_{Q_M}^\QQ=C_{\cH_\varpi}^\QQ(Q_M)\).
\end{proof}

\begin{proposition}[State-price-density formula for the centering defect]\label{prop:q-centering-density}
Assume Assumption~\ref{ass:q}.  Under Assumption~\ref{ass:q-raw-visible-aggregate}, for almost
every \(t\),
\[
        b_\varpi^\QQ(t)
        =
        \frac{
                \EE_\PP\!\left[
                        Z_T\varepsilon^{\QQ,\mathrm{raw}}(t)
                        \mid \cH_\varpi
                \right]
        }{
                \EE_\PP[Z_T\mid\cH_\varpi]
        }.
\]
If, in addition, the raw residual is \(\PP\)-centered relative to the same structural
sigma-field,
\[
        \EE_\PP[\varepsilon^{\QQ,\mathrm{raw}}(t)\mid\cH_\varpi]=0,
\]
then
\[
        b_\varpi^\QQ(t)
        =
        \frac{
                \Cov_\PP\!\left(
                        Z_T,\varepsilon^{\QQ,\mathrm{raw}}(t)
                        \mid \cH_\varpi
                \right)
        }{
                \EE_\PP[Z_T\mid\cH_\varpi]
        }.
\]
Consequently the exact zero-defect branch holds whenever the state-price density is conditionally
orthogonal to the raw aggregate residual given \(\cH_\varpi\).
\end{proposition}

\begin{proof}
For every bounded \(\cH_\varpi\)-measurable \(Y\),
\[
        \EE_\QQ[Y\varepsilon^{\QQ,\mathrm{raw}}(t)]
        =
        \EE_\PP[Z_TY\varepsilon^{\QQ,\mathrm{raw}}(t)]
        =
        \EE_\PP\!\left[
                Y\EE_\PP[Z_T\varepsilon^{\QQ,\mathrm{raw}}(t)\mid\cH_\varpi]
        \right].
\]
Likewise
\[
        \EE_\QQ[Y b]
        =
        \EE_\PP[Z_TY b]
        =
        \EE_\PP\!\left[Y b\,\EE_\PP[Z_T\mid\cH_\varpi]\right]
\]
for every \(\cH_\varpi\)-measurable candidate \(b\).  Since \(Z_T>0\) \(\PP\)-a.s. and
\(\QQ\sim\PP\), the denominator \(\EE_\PP[Z_T\mid\cH_\varpi]\) is strictly positive a.s.  The
displayed ratio is therefore the \(\QQ\)-conditional expectation.  If the raw residual is
\(\PP\)-centered, subtract
\(\EE_\PP[Z_T\mid\cH_\varpi]\EE_\PP[\varepsilon^{\QQ,\mathrm{raw}}(t)\mid\cH_\varpi]=0\) from
the numerator to obtain the conditional covariance formula.  The final statement is the case in
which this conditional covariance vanishes.
\end{proof}

\begin{proposition}[Small centering defect gives a perturbative variance-loading formula]\label{prop:q-centering-defect-bound}
Under Assumption~\ref{ass:q-raw-visible-aggregate} and
Definition~\ref{def:q-centering-defect}, define the zero-defect benchmark
\[
        L_{Q_M}^{\QQ,0}
        :=
        \int_0^T \bigl(A^\QQ(t)\bigr)^2\,\dd t
        +
        \int_0^T
        \EE_\QQ\!\left[
                \bigl(\varepsilon^{\QQ,\mathrm{raw}}(t)\bigr)^2
                \mid \cH_\varpi
        \right]\dd t.
\]
Then
\[
        L_{Q_M}^\QQ-L_{Q_M}^{\QQ,0}
        =
        2\int_0^T A^\QQ(t)b_\varpi^\QQ(t)\,\dd t.
\]
Consequently,
\[
        \|L_{Q_M}^\QQ-L_{Q_M}^{\QQ,0}\|_{L^2(\QQ)}
        \le
        2
        \left\|
                \int_0^T \bigl(A^\QQ(t)\bigr)^2\,\dd t
        \right\|_{L^2(\QQ)}^{1/2}
        \left\|
                \int_0^T \bigl(b_\varpi^\QQ(t)\bigr)^2\,\dd t
        \right\|_{L^2(\QQ)}^{1/2}.
\]
In particular, the exact centered branch is recovered when \(b_\varpi^\QQ=0\), and a small
state-price-density centering defect produces only the displayed first-order loading error.
\end{proposition}

\begin{proof}
Since
\[
        \varepsilon^{\QQ,\mathrm{raw}}
        =
        b_\varpi^\QQ+\eta_\varpi^{\QQ,\perp},
\]
conditional orthogonality gives
\[
        \EE_\QQ\!\left[
                \bigl(\varepsilon^{\QQ,\mathrm{raw}}(t)\bigr)^2
                \mid\cH_\varpi
        \right]
        =
        \bigl(b_\varpi^\QQ(t)\bigr)^2
        +
        \EE_\QQ\!\left[
                \bigl(\eta_\varpi^{\QQ,\perp}(t)\bigr)^2
                \mid\cH_\varpi
        \right].
\]
Subtracting this identity from the loading formula in Proposition~\ref{prop:q-recentering}
leaves the claimed difference identity; the \(\int (b_\varpi^\QQ)^2\) terms cancel.  Pathwise
Cauchy--Schwarz gives
\[
        \left|\int_0^T A^\QQ(t)b_\varpi^\QQ(t)\,\dd t\right|^2
        \le
        \left(\int_0^T (A^\QQ(t))^2\,\dd t\right)
        \left(\int_0^T (b_\varpi^\QQ(t))^2\,\dd t\right).
\]
Taking \(L^2(\QQ)\)-norms and applying Cauchy--Schwarz to the product gives the displayed bound.
\end{proof}

\begin{assumption}[Visible aggregate-volatility reduction under \(\QQ\)]\label{ass:q-visible-aggregate}
There exist a visible process \(M^{\mathrm{vis},\QQ}\), the structural synchronized amplitude
process \(\varpi\), and a residual process \(\eta^\QQ\) such that
\[
        M(t)=M^{\mathrm{vis},\QQ}(t)+(w^\top v)\varpi(t)+\eta^\QQ(t),
        \qquad
        0\le t\le T,
\]
where \(M^{\mathrm{vis},\QQ}(t)\) is \(\cG_t\)-measurable for every \(t\) and
\[
        \EE_\QQ[\eta^\QQ(t)\mid \cG_T\vee\cL_T]=0
        \qquad
        \text{for almost every } t\in[0,T].
\]
Assume further that
\[
        \int_0^T \bigl(M^{\mathrm{vis},\QQ}(t)\bigr)^2\,\dd t,
        \qquad
        \int_0^T \varpi(t)^2\,\dd t,
        \qquad
        \int_0^T \bigl(\eta^\QQ(t)\bigr)^2\,\dd t
\]
all belong to \(L^2(\QQ,\cF_T)\).
\end{assumption}

\begin{remark}[What is assumed and what is derived in the variance block]
Assumption~\ref{ass:q-visible-aggregate} is the zero-defect specialization of the raw recentered
reduction in Propositions~\ref{prop:q-recentering}--\ref{prop:q-centering-defect-bound}.  It gives
the minimal centered variance formula.  When exact \(\QQ\)-centering fails, the market-variance
loading remains explicit with the additional defect \(b_\varpi^\QQ\), and the error relative to the
zero-defect benchmark is controlled by Proposition~\ref{prop:q-centering-defect-bound}.
\end{remark}

\begin{proposition}[Canonical \(\QQ\)-decomposition of integrated market variance]\label{prop:q-qm}
Assume Assumption~\ref{ass:q-visible-aggregate} and set
\[
        \cH_\varpi:=\cG_T\vee \cL_T.
\]
Then
\[
        Q_M=L_{Q_M}^\QQ+U_{Q_M}^\QQ,
\]
where
\[
        L_{Q_M}^\QQ:=\EE_\QQ[Q_M\mid \cH_\varpi]\in \cM_\QQ,
\]
and
\[
        U_{Q_M}^\QQ:=Q_M-\EE_\QQ[Q_M\mid \cH_\varpi].
\]
Moreover,
\[
        \EE_\QQ[U_{Q_M}^\QQ\mid \cH_\varpi]=0.
\]
\end{proposition}

\begin{proof}
The decomposition is the tower-property identity
\[
        Q_M
        =
        \EE_\QQ[Q_M\mid \cH_\varpi]
        +
        \bigl(Q_M-\EE_\QQ[Q_M\mid \cH_\varpi]\bigr).
\]
The first term belongs to \(\cM_\QQ\).  The second lies in the residual fiber because
\[
        \EE_\QQ\bigl[Q_M-\EE_\QQ[Q_M\mid \cH_\varpi]\mid \cH_\varpi\bigr]=0.
\]
The normalization is notation.
\end{proof}

\begin{proposition}[Explicit form of the market-variance loading]\label{prop:q-qm-explicit}
Writing \(\zeta(t):=(w^\top v)\varpi(t)\), one has
\[
        L_{Q_M}^\QQ
        =
        \int_0^T \bigl(M^{\mathrm{vis},\QQ}(t)+\zeta(t)\bigr)^2\,\dd t
        +
        \int_0^T \EE_\QQ[\eta^\QQ(t)^2\mid \cH_\varpi]\,\dd t.
\]
\end{proposition}

\begin{proof}
Writing \(\zeta(t):=(w^\top v)\varpi(t)\), one has
\[
        M(t)=M^{\mathrm{vis},\QQ}(t)+\zeta(t)+\eta^\QQ(t),
\]
so
\[
        Q_M
        =
        \int_0^T \bigl(M^{\mathrm{vis},\QQ}(t)+\zeta(t)\bigr)^2\,\dd t
        +
        2\int_0^T \bigl(M^{\mathrm{vis},\QQ}(t)+\zeta(t)\bigr)\eta^\QQ(t)\,\dd t
        +
        \int_0^T \eta^\QQ(t)^2\,\dd t.
\]
Taking \(\EE_\QQ[\cdot\mid \cH_\varpi]\) and using
\[
        \EE_\QQ[\eta^\QQ(t)\mid \cH_\varpi]=0
\]
for almost every \(t\), together with conditional Fubini and the square-integrability
assumptions, eliminates the cross term and yields the displayed formula.
\end{proof}

\begin{remark}[Interpretation]
Proposition~\ref{prop:q-qm-explicit} identifies the market-variance loading as the
sum of the visible-plus-synchronized square and the conditional second moment of the residual
channel.
\end{remark}

\begin{proposition}[Canonical extraction of the aggregate factor from integrated market variance]\label{prop:q-qm-factor}
Assume Assumption~\ref{ass:q-visible-aggregate}.  Let
\[
        V_{Q_M}^\QQ:=C_{\cG_T}^\QQ(L_{Q_M}^\QQ),
        \qquad
        \widetilde L_{Q_M}^\QQ:=L_{Q_M}^\QQ-V_{Q_M}^\QQ.
\]
Then
\[
        \widetilde L_{Q_M}^\QQ
        =
        \alpha_{Q_M}^{\QQ,\star}\Lambda_M^\star,
        \qquad
        R_{Q_M}^{\QQ,\star}=0,
\]
where
\[
        \alpha_{Q_M}^{\QQ,\star}
        :=
        \|\widetilde L_{Q_M}^\QQ\|_{L^2(\QQ)}
\]
whenever \(\widetilde L_{Q_M}^\QQ\neq0\), with the convention
\(\alpha_{Q_M}^{\QQ,\star}:=0\) otherwise.  Equivalently,
\[
        L_{Q_M}^\QQ
        =
        V_{Q_M}^\QQ+\alpha_{Q_M}^{\QQ,\star}\Lambda_M^\star,
\]
so integrated market variance generates the canonical aggregate factor benchmark exactly in the
sense of Proposition~\ref{thm:factor-decomp}.
\end{proposition}

\begin{proof}
By Definition~\ref{def:agg-factor}, \(\Lambda_M^\star\) is obtained by normalizing the canonical
centered structural increment of integrated market variance under the present pricing measure.  In
the notation of Proposition~\ref{prop:q-qm}, that centered increment is exactly
\[
        \widetilde L_{Q_M}^\QQ
        =
        L_{Q_M}^\QQ-C_{\cG_T}^\QQ(L_{Q_M}^\QQ).
\]
If \(\widetilde L_{Q_M}^\QQ=0\), the claim is trivial.  Otherwise,
\[
        \Lambda_M^\star
        =
        \frac{\widetilde L_{Q_M}^\QQ}{\|\widetilde L_{Q_M}^\QQ\|_{L^2(\QQ)}},
\]
which yields the stated representation and shows that the orthogonal residual in
Proposition~\ref{thm:factor-decomp} vanishes for \(Q_M\).
\end{proof}

\begin{proposition}[Canonical \(\QQ\)-decomposition of variance-linked claims]\label{prop:q-claim}
Assume Assumption~\ref{ass:q-visible-aggregate}, let \(X=\Psi(Q_M)\in L^2(\QQ,\cF_T)\), and set
\[
        \cH_\varpi:=\cG_T\vee \cL_T.
\]
Define
\[
        L_X^\QQ:=\EE_\QQ[X\mid \cH_\varpi],
\]
\[
        U_X^\QQ:=X-\EE_\QQ[X\mid \cH_\varpi].
\]
Then
\[
        X=L_X^\QQ+U_X^\QQ,
\]
where \(L_X^\QQ\in \cM_\QQ\) and
\[
        \EE_\QQ[U_X^\QQ\mid \cH_\varpi]=0.
\]
If \(L_X^\QQ\neq0\), define
\[
        \alpha_X^\QQ:=\|L_X^\QQ\|_{L^2(\QQ)},
        \qquad
        \Lambda_X:=\frac{L_X^\QQ}{\alpha_X^\QQ};
\]
otherwise set \(\alpha_X^\QQ:=0\) and \(\Lambda_X:=0\).  In all cases,
\[
        \EE_\QQ[U_X^\QQ\mid \cG_T,\Lambda_X]=0.
\]
\end{proposition}

\begin{proof}
The decomposition is again the tower-property identity
\[
        X
        =
        \EE_\QQ[X\mid \cH_\varpi]
        +
        \bigl(X-\EE_\QQ[X\mid \cH_\varpi]\bigr).
\]
The first term belongs to \(\cM_\QQ\) exactly as in Proposition~\ref{prop:q-qm}.  The definition of
\(U_X^\QQ\) gives \(\EE_\QQ[U_X^\QQ\mid \cH_\varpi]=0\).  Since \(\Lambda_X\) is \(\cH_\varpi\)-measurable and
\(\sigma(\cG_T,\Lambda_X)\subseteq \cH_\varpi\), the tower property gives
\[
        \EE_\QQ[U_X^\QQ\mid \cG_T,\Lambda_X]
        =
        \EE_\QQ\bigl[\EE_\QQ[U_X^\QQ\mid \cH_\varpi]\mid \cG_T,\Lambda_X\bigr]
        =0.
\]
\end{proof}

\begin{corollary}[Affine variance-linked claims inherit aggregate-factor alignment]\label{cor:q-affine-factor}
Assume Assumption~\ref{ass:q-visible-aggregate}, and let
\[
        X=a+bQ_M
\]
for constants \(a,b\in\R\) such that \(X\in L^2(\QQ,\cF_T)\).  Then
\[
        L_X^\QQ
        =
        C_{\cG_T}^\QQ(L_X^\QQ)
        +
        b\,\alpha_{Q_M}^{\QQ,\star}\Lambda_M^\star,
\]
and \(R_X^{\QQ,\star}=0\).  Hence every affine variance-linked claim is aggregate-factor aligned.
\end{corollary}

\begin{proof}
Linearity of conditional expectation gives
\[
        L_X^\QQ
        =
        a+bL_{Q_M}^\QQ.
\]
Apply Proposition~\ref{prop:q-qm-factor} to \(L_{Q_M}^\QQ\):
\[
        L_X^\QQ
        =
        a+bV_{Q_M}^\QQ+b\,\alpha_{Q_M}^{\QQ,\star}\Lambda_M^\star.
\]
Since \(a+bV_{Q_M}^\QQ\in L^2(\QQ,\cG_T)\), this is exactly the decomposition from
Proposition~\ref{thm:factor-decomp} with zero orthogonal centered remainder.
\end{proof}

\begin{remark}[Why nonlinear variance-linked claims require a hierarchy]
Corollary~\ref{cor:q-affine-factor} is exact because affine transformations preserve the
one-factor structure of the centered market-variance loading.  A general nonlinear claim
\(X=\Psi(Q_M)\) need not do so. After conditioning onto \(\cH_\varpi\), the centered structural
loading of \(\Psi(Q_M)\) can acquire higher-order centered components beyond the span of
\(\Lambda_M^\star\).  Thus exact one-factor inheritance is the affine branch, while nonlinear
variance-linked claims are handled by the general compression and hedging theorems and, in the
one-factor experiment below, by the explicit centered-power hierarchy in
Proposition~\ref{thm:concrete-nonlinear-variance-compression}.
\end{remark}

\begin{lemma}[Nonlinearity plus compression of the priced structural loading]\label{lem:nonlinear-gap}
Let \(L\in L^2(\QQ,\cH_\varpi)\), let \(\cH\subseteq \cG_T\), and let \(\Phi\colon\R\to\R\) be
convex and twice continuously differentiable with
\[
        \inf_x \Phi''(x)\ge \underline\kappa>0.
\]
Assume \(\Phi(L),\Phi(C_{\cH}^\QQ(L))\in L^1(\QQ)\).  Then
\[
        \EE_\QQ[\Phi(L)]-\EE_\QQ[\Phi(C_{\cH}^\QQ(L))]
        \ge
        \frac{\underline\kappa}{2}
        \|L-C_{\cH}^\QQ(L)\|_{L^2(\QQ)}^2.
\]
If \(L=\alpha\Lambda\) with \(\alpha\ge0\) and \(\|\Lambda\|_{L^2(\QQ)}=1\), then
\[
        \EE_\QQ[\Phi(\alpha\Lambda)]-\EE_\QQ[\Phi(\alpha C_{\cH}^\QQ(\Lambda))]
        \ge
        \frac{\underline\kappa}{2}
        \|\alpha\Lambda-\alpha C_{\cH}^\QQ(\Lambda)\|_{L^2(\QQ)}^2.
\]
\end{lemma}

\begin{proof}
Since \(C_{\cH}^\QQ(L)=\EE_\QQ[L\mid \cH]\), conditional Jensen together with the strong convexity
bound gives
\[
        \EE_\QQ[\Phi(L)]-\EE_\QQ[\Phi(C_{\cH}^\QQ(L))]
        \ge
        \frac{\underline\kappa}{2}\EE_\QQ[\Var_\QQ(L\mid \cH)].
\]
The first claim follows from
\[
        \EE_\QQ[\Var_\QQ(L\mid \cH)]
        =
        \|L-C_{\cH}^\QQ(L)\|_{L^2(\QQ)}^2.
\]
The second is the specialization \(L=\alpha\Lambda\).
\end{proof}

The variance-compression theorem identifies the Jensen gap left after public conditioning of the
priced structural loading.

\begin{theorem}[Risk-neutral compression of priced structural claims]\label{thm:q-compression}
Let \(X\in L^2(\QQ,\cF_T)\), and let
\[
        X=L_X^\QQ+U_X^\QQ
\]
be the decomposition of Proposition~\ref{thm:claim-decomp}.  Let \(\Phi\colon\R\to\R\) be convex and
twice continuously differentiable with
\[
        \inf_x \Phi''(x)\ge \underline\kappa>0,
\]
and assume \(\Phi(X),\Phi(C_{\cH}^\QQ(X))\in L^1(\QQ)\) for every sub-\(\sigma\)-field
\(\cH\subseteq \cG_T\) considered below.  Then for every such \(\cH\),
\[
        \EE_\QQ[\Phi(X)]
        -
        \EE_\QQ[\Phi(C_{\cH}^\QQ(X))]
        \ge
        \frac{\underline\kappa}{2}
        \|X-C_{\cH}^\QQ(X)\|_{L^2(\QQ)}^2.
\]
Equivalently,
\[
        \EE_\QQ[\Phi(X)]
        -
        \EE_\QQ[\Phi(C_{\cH}^\QQ(X))]
        \ge
        \frac{\underline\kappa}{2}
        \left(
        \|L_X^\QQ-C_{\cH}^\QQ(L_X^\QQ)\|_{L^2(\QQ)}^2
        +
        \|U_X^\QQ\|_{L^2(\QQ)}^2
        \right).
\]
If \(L_X^\QQ=\alpha_X^\QQ\Lambda_X\), the same lower bound is
\[
        \frac{\underline\kappa}{2}
        \left(
        (\alpha_X^\QQ)^2
        \|\Lambda_X-C_{\cH}^\QQ(\Lambda_X)\|_{L^2(\QQ)}^2
        +
        \|U_X^\QQ\|_{L^2(\QQ)}^2
        \right).
\]
In particular, compression to any visible pricing summary
\(\sigma(\mathfrak S)\subseteq \cG_T\) leaves a strict positive gap whenever either
\[
        \|L_X^\QQ-C_{\sigma(\mathfrak S)}^\QQ(L_X^\QQ)\|_{L^2(\QQ)}^2>0
\]
or \(\|U_X^\QQ\|_{L^2(\QQ)}^2>0\).
\end{theorem}

\begin{proof}
For every \(x,y\in\R\), the curvature lower bound gives
\[
        \Phi(y)
        \ge
        \Phi(x)+\Phi'(x)(y-x)+\frac{\underline\kappa}{2}(y-x)^2.
\]
Set \(x=C_{\cH}^\QQ(X)\) and \(y=X\), and condition on \(\cH\).  The affine term vanishes because
\(\EE_\QQ[X-C_{\cH}^\QQ(X)\mid\cH]=0\).  Therefore
\[
        \EE_\QQ[\Phi(X)]-\EE_\QQ[\Phi(C_{\cH}^\QQ(X))]
        \ge
        \frac{\underline\kappa}{2}\EE_\QQ\bigl[\Var_\QQ(X\mid \cH)\bigr]
        =
        \frac{\underline\kappa}{2}\|X-C_{\cH}^\QQ(X)\|_{L^2(\QQ)}^2.
\]
Now \(L_X^\QQ\) is \(\cH_\varpi\)-measurable and Proposition~\ref{thm:claim-decomp} gives
\(\EE_\QQ[U_X^\QQ\mid \cH_\varpi]=0\).  Since \(\cH\subseteq\cG_T\subseteq\cH_\varpi\) and
\(\sigma(\cH,L_X^\QQ)\subseteq \cH_\varpi\),
\[
        \EE_\QQ[U_X^\QQ\mid \cH,L_X^\QQ]=0.
\]
Equivalently, \(L_X^\QQ-C_{\cH}^\QQ(L_X^\QQ)\in L^2(\QQ,\cH_\varpi)\) while
\(U_X^\QQ\perp L^2(\QQ,\cH_\varpi)\).  Since
\[
        X-C_{\cH}^\QQ(X)
        =
        L_X^\QQ-C_{\cH}^\QQ(L_X^\QQ)+U_X^\QQ,
\]
Pythagoras gives
\[
        \|X-C_{\cH}^\QQ(X)\|_{L^2(\QQ)}^2
        =
        \|L_X^\QQ-C_{\cH}^\QQ(L_X^\QQ)\|_{L^2(\QQ)}^2
        +
        \|U_X^\QQ\|_{L^2(\QQ)}^2.
\]
Combining the displays proves the second inequality.  The scalar-factor form follows from
\(L_X^\QQ=\alpha_X^\QQ\Lambda_X\).
\end{proof}

\begin{corollary}[Factor-level decomposition of the compression gap]\label{cor:q-compression-factor}
Under the assumptions of Proposition~\ref{thm:factor-decomp} and~\ref{thm:q-compression}, for every
public sub-\(\sigma\)-field \(\cH\subseteq\cG_T\),
\[
        \|L_X^\QQ-C_{\cH}^\QQ(L_X^\QQ)\|_{L^2(\QQ)}^2
        =
        \|V_X^\QQ-C_{\cH}^\QQ(V_X^\QQ)\|_{L^2(\QQ)}^2
        +
        (\alpha_X^{\QQ,\star})^2
        +
        \|R_X^{\QQ,\star}\|_{L^2(\QQ)}^2.
\]
If \(X\) is aggregate-factor aligned, so \(R_X^{\QQ,\star}=0\), this reduces to
\[
        \|L_X^\QQ-C_{\cH}^\QQ(L_X^\QQ)\|_{L^2(\QQ)}^2
        =
        \|V_X^\QQ-C_{\cH}^\QQ(V_X^\QQ)\|_{L^2(\QQ)}^2
        +
        (\alpha_X^{\QQ,\star})^2.
\]
Consequently,
\[
        \EE_\QQ[\Phi(X)]
        -
        \EE_\QQ[\Phi(C_{\cH}^\QQ(X))]
        \ge
        \frac{\underline\kappa}{2}
        \left(
                \|V_X^\QQ-C_{\cH}^\QQ(V_X^\QQ)\|_{L^2(\QQ)}^2
                +
                (\alpha_X^{\QQ,\star})^2
                +
                \|R_X^{\QQ,\star}\|_{L^2(\QQ)}^2
        \right).
\]
\end{corollary}

\begin{proof}
By Proposition~\ref{thm:factor-decomp},
\[
        L_X^\QQ
        =
        V_X^\QQ+\alpha_X^{\QQ,\star}\Lambda_M^\star+R_X^{\QQ,\star}.
\]
Because \(\cH\subseteq\cG_T\), Corollary~\ref{cor:factor-screening} yields
\[
        C_{\cH}^\QQ(L_X^\QQ)=C_{\cH}^\QQ(V_X^\QQ).
\]
Therefore
\[
        L_X^\QQ-C_{\cH}^\QQ(L_X^\QQ)
        =
        \bigl(V_X^\QQ-C_{\cH}^\QQ(V_X^\QQ)\bigr)
        +
        \alpha_X^{\QQ,\star}\Lambda_M^\star
        +
        R_X^{\QQ,\star}.
\]
The first term belongs to \(L^2(\QQ,\cG_T)\), while
\(\alpha_X^{\QQ,\star}\Lambda_M^\star+R_X^{\QQ,\star}\in\cM_\QQ^\circ\); hence it is orthogonal to
the first term.  Proposition~\ref{thm:factor-decomp} also gives
\(R_X^{\QQ,\star}\perp\Lambda_M^\star\).  The displayed norm identity follows, and the final lower
bound is Theorem~\ref{thm:q-compression} combined with that identity.
\end{proof}

\begin{corollary}[Market-variance application of the \(\QQ\)-side compression theorem]\label{cor:market-transfer}
Under Assumption~\ref{ass:q-visible-aggregate}, Theorem~\ref{thm:q-compression} applies in
particular to \(Q_M\) and to variance-linked claims \(X=\Psi(Q_M)\).  By
Proposition~\ref{prop:q-qm-factor} and Corollary~\ref{cor:q-affine-factor}, the centered
market-variance loading itself and every affine variance-linked claim are exactly aggregate-factor
aligned.  This is the risk-neutral analogue of the compression result in \cite{vidal2026market}.
The option block and the market-variance block therefore use the same projection structure.  Public
prices are visible \(\QQ\)-compressions of claims that may still carry priced structural loadings in
\(\cM_\QQ\).
\end{corollary}

\begin{assumption}[One-factor aggregate variance experiment]\label{ass:one-factor-variance}
Let \(\Xi\in L^4(\QQ)\) be independent of \(\cG_T\), with
\[
        \EE_\QQ[\Xi]=0,
        \qquad
        \EE_\QQ[\Xi^2]=1.
\]
Let \(f\in L^2([0,T])\) be deterministic and nonzero, and let \(m(t)\) be a square-integrable
\(\cG_T\)-measurable visible process.  Fix \(a\in\R\).  The aggregate volatility is
\[
        M(t)=m(t)+a f(t)\Xi,
        \qquad
        0\le t\le T.
\]
Set
\[
        Q_M:=\int_0^TM(t)^2\,\dd t,
        \qquad
        A_G:=\int_0^Tm(t)f(t)\,\dd t,
        \qquad
        B_f:=\int_0^Tf(t)^2\,\dd t,
\]
and
\[
        V_Q:=\int_0^Tm(t)^2\,\dd t+a^2B_f.
\]
\end{assumption}

\begin{proposition}[One-factor market-variance compression gap]\label{thm:concrete-variance-compression}
Under Assumption~\ref{ass:one-factor-variance},
\[
        \EE_\QQ[Q_M\mid\cG_T]=V_Q
\]
and
\[
        Q_M-V_Q
        =
        2aA_G\Xi+a^2B_f(\Xi^2-1).
\]
Consequently, for every public \(\cH\subseteq\cG_T\),
\[
        \|Q_M-\EE_\QQ[Q_M\mid\cH]\|_{L^2(\QQ)}^2
        =
        \|V_Q-\EE_\QQ[V_Q\mid\cH]\|_{L^2(\QQ)}^2
        +
        \Delta_Q^2,
\]
where
\[
        \Delta_Q^2
        :=
        4a^2\EE_\QQ[A_G^2]
        +
        4a^3B_f\EE_\QQ[A_G]\,\EE_\QQ[\Xi(\Xi^2-1)]
        +
        a^4B_f^2\Var_\QQ(\Xi^2).
\]
If \(\Xi\) is symmetric, this reduces to
\[
        \Delta_Q^2
        =
        4a^2\EE_\QQ[A_G^2]+a^4B_f^2\Var_\QQ(\Xi^2).
\]
Let \(\Phi:\R\to\R\) be twice continuously differentiable and strongly convex with
\(\inf_x\Phi''(x)\ge\underline\kappa>0\), and assume the displayed expectations are finite.  Then
\[
        \EE_\QQ[\Phi(Q_M)]-
        \EE_\QQ\!\bigl[\Phi(\EE_\QQ[Q_M\mid\cH])\bigr]
        \ge
        \frac{\underline\kappa}{2}
        \left(
                \|V_Q-\EE_\QQ[V_Q\mid\cH]\|_{L^2(\QQ)}^2+
                \Delta_Q^2
        \right).
\]
In particular, even when \(\cH=\cG_T\) and all visible variance information is used, the strongly
convex compression gap is bounded below by
\[
        \frac{\underline\kappa}{2}\Delta_Q^2.
\]
\end{proposition}

\begin{proof}
Expanding the square gives
\[
        Q_M
        =
        \int_0^Tm(t)^2\,\dd t
        +2a\Xi\int_0^Tm(t)f(t)\,\dd t
        +a^2\Xi^2\int_0^Tf(t)^2\,\dd t.
\]
Since \(\Xi\) is independent of \(\cG_T\), centered, and has unit variance,
\[
        \EE_\QQ[Q_M\mid\cG_T]=\int_0^Tm(t)^2\,\dd t+a^2B_f=V_Q,
\]
and the centered hidden part is
\[
        Q_M-V_Q=2aA_G\Xi+a^2B_f(\Xi^2-1).
\]
For \(\cH\subseteq\cG_T\), conditional independence gives
\[
        \EE_\QQ[Q_M\mid\cH]=\EE_\QQ[V_Q\mid\cH].
\]
Thus
\[
        Q_M-\EE_\QQ[Q_M\mid\cH]
        =
        \bigl(V_Q-\EE_\QQ[V_Q\mid\cH]\bigr)
        +2aA_G\Xi+a^2B_f(\Xi^2-1).
\]
The first term is \(\cG_T\)-measurable.  It is orthogonal to both hidden terms because
\(\EE_\QQ[\Xi]=0\), \(\EE_\QQ[\Xi^2-1]=0\), and \(\Xi\) is independent of \(\cG_T\).  Taking the
\(L^2\)-norm gives the norm identity; the cross term between \(2aA_G\Xi\) and
\(a^2B_f(\Xi^2-1)\) is exactly
\[
        4a^3B_f\EE_\QQ[A_G]\EE_\QQ[\Xi(\Xi^2-1)].
\]
The strong-convexity bound follows from conditional Jensen with the variance correction:
\[
        \EE_\QQ[\Phi(Q_M)]-
        \EE_\QQ[\Phi(\EE_\QQ[Q_M\mid\cH])]
        \ge
        \frac{\underline\kappa}{2}
        \EE_\QQ[\Var_\QQ(Q_M\mid\cH)]
        =
        \frac{\underline\kappa}{2}
        \|Q_M-\EE_\QQ[Q_M\mid\cH]\|_{L^2(\QQ)}^2.
\]
\end{proof}

\begin{corollary}[Pure hidden-factor variance floor]\label{cor:pure-hidden-variance-floor}
If, in Proposition~\ref{thm:concrete-variance-compression}, \(m\equiv0\), then
\[
        Q_M=a^2B_f\Xi^2,
        \qquad
        \EE_\QQ[Q_M\mid\cG_T]=a^2B_f,
\]
and for every \(\cH\subseteq\cG_T\),
\[
        \|Q_M-\EE_\QQ[Q_M\mid\cH]\|_{L^2(\QQ)}^2=a^4B_f^2\Var_\QQ(\Xi^2).
\]
If \(\Xi\sim\mathcal N(0,1)\), this is
\[
        \|Q_M-\EE_\QQ[Q_M\mid\cH]\|_{L^2(\QQ)}^2=2a^4B_f^2.
\]
\end{corollary}

\begin{proposition}[One-factor nonlinear variance-linked compression hierarchy]\label{thm:concrete-nonlinear-variance-compression}
Assume Assumption~\ref{ass:one-factor-variance}, and set
\[
        H_Q:=2aA_G\Xi+a^2B_f(\Xi^2-1).
\]
Then
\[
        Q_M=V_Q+H_Q,
        \qquad
        \EE_\QQ[H_Q\mid\cG_T]=0.
\]
Let \(\Psi:\R\to\R\) be Borel with \(\Psi(Q_M)\in L^2(\QQ)\), and define
\[
        V_\Psi:=\EE_\QQ[\Psi(V_Q+H_Q)\mid\cG_T],
        \qquad
        R_\Psi:=\Psi(V_Q+H_Q)-V_\Psi .
\]
Then
\[
        \EE_\QQ[\Psi(Q_M)\mid\cG_T]=V_\Psi,
        \qquad
        \EE_\QQ[R_\Psi\mid\cG_T]=0.
\]
Consequently, for every public \(\sigma\)-field \(\cH\subseteq\cG_T\),
\[
        \EE_\QQ[\Psi(Q_M)\mid\cH]=\EE_\QQ[V_\Psi\mid\cH],
\]
and for every closed visible hedge class \(\cS\subseteq L^2(\QQ,\cG_T)\),
\[
        \inf_{Y\in\cS}\EE_\QQ[(\Psi(Q_M)-Y)^2]
        =
        \dist_{L^2(\QQ)}(V_\Psi,\cS)^2+\|R_\Psi\|_{L^2(\QQ)}^2 .
\]
If, in addition,
\[
        \Psi(x)=\sum_{r=0}^p c_rx^r
\]
and the displayed polynomial terms are square-integrable, then
\[
        V_\Psi
        =
        \sum_{r=0}^p c_r
        \sum_{k=0}^r \binom{r}{k}
        V_Q^{r-k}\EE_\QQ[H_Q^k\mid\cG_T],
\]
while the full-visible hidden residual is
\[
        R_\Psi
        =
        \sum_{r=0}^p c_r
        \sum_{k=1}^r \binom{r}{k}
        V_Q^{r-k}
        \bigl(H_Q^k-\EE_\QQ[H_Q^k\mid\cG_T]\bigr).
\]
\end{proposition}

\begin{proof}
The identity \(Q_M=V_Q+H_Q\) and the centering
\(\EE_\QQ[H_Q\mid\cG_T]=0\) are exactly Proposition~\ref{thm:concrete-variance-compression}.  The
definitions of \(V_\Psi\) and \(R_\Psi\) give
\[
        \Psi(Q_M)=V_\Psi+R_\Psi,
        \qquad
        \EE_\QQ[R_\Psi\mid\cG_T]=0.
\]
The tower property gives, for \(\cH\subseteq\cG_T\),
\[
        \EE_\QQ[\Psi(Q_M)\mid\cH]
        =
        \EE_\QQ\bigl[\EE_\QQ[\Psi(Q_M)\mid\cG_T]\mid\cH\bigr]
        =
        \EE_\QQ[V_\Psi\mid\cH].
\]
If \(Y\in\cS\subseteq L^2(\QQ,\cG_T)\), then \(V_\Psi-Y\) is \(\cG_T\)-measurable and therefore
orthogonal to \(R_\Psi\).  Hence
\[
        \EE_\QQ[(\Psi(Q_M)-Y)^2]
        =
        \|V_\Psi-Y\|_{L^2(\QQ)}^2+\|R_\Psi\|_{L^2(\QQ)}^2.
\]
Taking the infimum over \(Y\in\cS\) gives the hedging identity.  The polynomial formulas follow by
expanding \((V_Q+H_Q)^r\), using that \(V_Q\) is \(\cG_T\)-measurable, and subtracting the
conditional expectation from the full polynomial.
\end{proof}

\begin{corollary}[Quadratic variance-linked claims expose the next hidden component]\label{cor:quadratic-variance-hidden-hierarchy}
Assume Assumption~\ref{ass:one-factor-variance}, and assume \(Q_M^2\in L^2(\QQ)\).  For
\[
        \Psi(x)=\alpha+\beta x+\gamma x^2
\]
one has
\[
        V_\Psi
        =
        \alpha+\beta V_Q+\gamma
        \bigl(V_Q^2+\EE_\QQ[H_Q^2\mid\cG_T]\bigr),
\]
and
\[
        R_\Psi
        =
        (\beta+2\gamma V_Q)H_Q+
        \gamma\bigl(H_Q^2-\EE_\QQ[H_Q^2\mid\cG_T]\bigr).
\]
Moreover,
\[
        \EE_\QQ[H_Q^2\mid\cG_T]
        =
        4a^2A_G^2+
        4a^3A_GB_f\,\EE_\QQ[\Xi(\Xi^2-1)]
        +
        a^4B_f^2\Var_\QQ(\Xi^2).
\]
If \(\Xi\) is symmetric, this conditional second moment reduces to
\[
        4a^2A_G^2+a^4B_f^2\Var_\QQ(\Xi^2).
\]
Thus the affine case \(\gamma=0\) is the single hidden-factor branch, while a genuinely quadratic
claim exposes the additional centered-square component
\[
        H_Q^2-\EE_\QQ[H_Q^2\mid\cG_T]
\]
with nonzero contribution whenever \(\gamma\neq0\) and this centered square is nonzero in
\(L^2(\QQ)\).
\end{corollary}

\begin{proof}
Apply Proposition~\ref{thm:concrete-nonlinear-variance-compression} with \(p=2\).  Since
\(\EE_\QQ[H_Q\mid\cG_T]=0\), the \(k=1\) term contributes only to the centered residual.  Finally,
write
\[
        H_Q=2aA_G\Xi+a^2B_f(\Xi^2-1)
\]
and use independence of \(\Xi\) from \(\cG_T\), together with
\(\EE_\QQ[\Xi^2]=1\), to compute \(\EE_\QQ[H_Q^2\mid\cG_T]\).  Symmetry of \(\Xi\) gives
\(\EE_\QQ[\Xi(\Xi^2-1)]=0\).
\end{proof}

\begin{assumption}[Conditionally Gaussian aggregate-variance reduction]\label{ass:cg-onefactor}
Let \(\cG\) be the visible market sigma-field.  Assume the integrated market-variance loading has
the one-factor form
\[
        Q_M=V+\alpha Z+S\xi,
\]
where \(V\) and \(S\ge0\) are \(\cG\)-measurable, \(\alpha\ge0\), and \((Z,\xi)\) are independent
standard normal variables independent of \(\cG\).  The canonical hidden aggregate factor is \(Z\).
For a Borel payoff \(\Psi\), define \(X=\Psi(Q_M)\), assume \(X\in L^2(\QQ)\), and write
\[
        P_s\Psi(x):=\EE[\Psi(x+s\xi)]
\]
whenever the expectation is finite.
\end{assumption}

\begin{proposition}[Exact one-factor Hermite hierarchy for nonlinear variance-linked claims]\label{thm:nonlinear-hermite}
Under Assumption~\ref{ass:cg-onefactor}, the priced structural loading of
\(X=\Psi(Q_M)\) relative to \(\cG\vee\sigma(Z)\) is
\[
        L_X^\QQ
        :=
        \EE_\QQ[X\mid \cG,Z]
        =
        P_S\Psi(V+\alpha Z).
\]
For each fixed \((V,S)\), it has the Hermite expansion
\[
        L_X^\QQ
        =
        c_0(V,S)+\sum_{m\ge1}c_m(V,S)\He_m(Z)
\]
in \(L^2\), where
\[
        c_m(V,S)
        =
        \frac{1}{m!}
        \EE\!\left[
                P_S\Psi(V+\alpha\zeta)\He_m(\zeta)
                \mid V,S
        \right],
        \qquad
        \zeta\sim N(0,1).
\]
If \(P_S\Psi\) has \(m\) weak derivatives with the required integrability, then
\[
        c_m(V,S)
        =
        \frac{\alpha^m}{m!}
        \EE\!\left[
                (P_S\Psi)^{(m)}(V+\alpha\zeta)
                \mid V,S
        \right].
\]
Thus arbitrary nonlinear \(\Psi\) contributes through the explicit Hermite coefficient sequence
\((c_m)_{m\ge0}\), as well as through the implicit conditional expectation.
\end{proposition}

\begin{proof}
Conditioning on \(\cG\vee\sigma(Z)\) freezes \(V,S,Z\) and integrates only the independent normal
residual \(\xi\), giving \(L_X^\QQ=P_S\Psi(V+\alpha Z)\).  For fixed \((V,S)\), the function
\(z\mapsto P_S\Psi(V+\alpha z)\) belongs to \(L^2\) under the standard normal law by
\(X\in L^2\) and Jensen's inequality.  Expanding it in the Hermite basis gives the coefficient
formula.  The derivative representation follows from Gaussian integration by parts:
\[
        \EE[f(\zeta)\He_m(\zeta)]=\EE[f^{(m)}(\zeta)]
\]
with \(f(z)=P_S\Psi(V+\alpha z)\).
\end{proof}

\begin{corollary}[Visible-beta first-factor law]\label{cor:visible-beta-law}
Assume the hypotheses of Proposition~\ref{thm:nonlinear-hermite}.  If \(P_S\Psi\) has two derivatives
and the corresponding second-derivative envelope is square-integrable for small \(\alpha\), then
\[
        L_X^\QQ
        =
        C_\cG^\QQ(L_X^\QQ)
        +
        \alpha (P_S\Psi)'(V)Z
        +
        R_{\Psi,\alpha},
\]
with
\[
        \|R_{\Psi,\alpha}\|_{L^2(\QQ)}=O(\alpha^2).
\]
Consequently the first hidden variance factor is the same aggregate factor \(Z\), with visible beta
\((P_S\Psi)'(V)\).  The nonlinear residual beyond that first factor is curvature-order.
\end{corollary}

\begin{proof}
The Hermite hierarchy gives
\[
        L_X^\QQ-C_\cG^\QQ(L_X^\QQ)=\sum_{m\ge1}c_m(V,S)\He_m(Z).
\]
The first coefficient is
\[
        c_1(V,S)
        =
        \alpha\EE[(P_S\Psi)'(V+\alpha\zeta)\mid V,S]
        =
        \alpha(P_S\Psi)'(V)+O_{L^2}(\alpha^2),
\]
where the last step is Taylor's formula under the assumed second-derivative envelope.  For
\(m\ge2\), the Hermite coefficient formula and the same envelope give the required
\(O_{L^2}(\alpha^m)\) bounds at this order.  Orthogonality of the Hermite polynomials then makes the
sum of all terms beyond the displayed first factor an \(L^2\)-residual of order \(\alpha^2\).
\end{proof}

\begin{corollary}[Simple nonlinear compression gap]\label{cor:nonlinear-gap-onefactor}
Under the assumptions of Corollary~\ref{cor:visible-beta-law}, if the visible hedge or price
summary is contained in \(L^2(\QQ,\cG)\), then
\[
        \|L_X^\QQ-C_\cG^\QQ(L_X^\QQ)\|_{L^2(\QQ)}^2
        =
        \alpha^2\EE_\QQ[((P_S\Psi)'(V))^2]
        +
        O(\alpha^3).
\]
If \((P_S\Psi)'(V)\) is nonzero with positive probability, the nonlinear variance-linked claim has
a positive first-order hidden aggregate-factor gap.
\end{corollary}

\begin{proof}
The factor \(Z\) is independent of \(\cG\) and centered, so
\[
        C_\cG^\QQ((P_S\Psi)'(V)Z)=0.
\]
The visible-beta law gives
\[
        L_X^\QQ-C_\cG^\QQ(L_X^\QQ)
        =
        \alpha(P_S\Psi)'(V)Z+O_{L^2}(\alpha^2).
\]
Taking squared norms and using \(\EE[Z^2]=1\) gives the expansion.
\end{proof}

\begin{corollary}[Affine claims are the only exact deterministic-beta scalar one-factor family]\label{cor:affine-only}
In the model of Assumption~\ref{ass:cg-onefactor} with \(S=0\) and nondegenerate \(V\), suppose
\[
        \Psi(V+\alpha Z)
        =
        A(V)+bZ
\]
for some \(\cG\)-measurable \(A(V)\) and deterministic scalar \(b\), for all sufficiently small
\(\alpha>0\).  Then \(\Psi\) is affine on the support of \(V+\alpha Z\).  Hence nonlinear claims
cannot generally have an exact deterministic-beta scalar one-factor law; the exact law is the
Hermite hierarchy, and the first-order law has visible beta.
\end{corollary}

\begin{proof}
For fixed \(V=v\), the map \(z\mapsto\Psi(v+\alpha z)\) is affine in \(z\) on the support of the
standard normal law.  Therefore \(x\mapsto\Psi(x)\) is affine on the translated line
\(x=v+\alpha z\).  Since the support is an interval and \(V\) is nondegenerate, these intervals
cover the support of \(V+\alpha Z\) under the stated nondegeneracy.  Thus a version of \(\Psi\) is affine on that support, up to null sets.  Non-affine
\(\Psi\) produces higher Hermite coefficients in Proposition~\ref{thm:nonlinear-hermite}.
\end{proof}

\begin{remark}[Nonlinear variance-linked law]\label{rem:nonlinear-hermite-headline}
For nonlinear variance-linked claims, the correct statement is not that every \(\Psi\) is scalar
one-factor.  That statement is false except in affine cases.  The correct formulation is: under a
conditionally Gaussian aggregate-variance reduction, every square-integrable nonlinear
variance-linked claim has an exact one-factor Hermite hierarchy, and its leading hidden component
is the same aggregate factor with beta \((P_S\Psi)'(V)\), up to curvature-order error.
\end{remark}

\section{Visible hedging and risk-neutral incompleteness}\label{sec:hedging}

\noindent
Visible hedging changes the space from public information to tradeable span after the variance compression of Section~\ref{sec:variance}.  By
Corollary~\ref{prop:compression-unifier}, a visible hedge acts on the same priced structural loading
\(L_X^\QQ\), but only through its projection onto a closed public hedge class.  Incompleteness is
therefore a projection residual.  A visible hedge class cannot eliminate structural loading outside
its span.  This is the Hilbert-space compression form of incomplete-market mean-square hedging
\cite{follmerSondermann1986,follmerSchweizer1991,schweizer1992meanVariance,
schweizer1995minimal,karatzasShreve1998}.

The hedging theorem identifies that residual exactly.

\begin{theorem}[Visible hedging as projection failure of priced structural loading]\label{thm:q-hedging}
Let \(X\in L^2(\QQ,\cF_T)\), and let
\[
        X=L_X^\QQ+U_X^\QQ
\]
be the decomposition of Proposition~\ref{thm:claim-decomp}.  If
\(\cS\subseteq L^2(\QQ,\cG_T)\) is a visible hedge class, then
\[
        \inf_{H\in \cS}\EE_\QQ[(X-H)^2]
        =
        \dist_{L^2(\QQ)}\!\bigl(L_X^\QQ,\cS\bigr)^2
        +
        \|U_X^\QQ\|_{L^2(\QQ)}^2.
\]
Equivalently,
\[
        \inf_{H\in \cS}\EE_\QQ[(X-H)^2]
        =
        \|L_X^\QQ-P_{\cS}^\QQ L_X^\QQ\|_{L^2(\QQ)}^2
        +
        \|U_X^\QQ\|_{L^2(\QQ)}^2.
\]
If \(L_X^\QQ=\alpha_X^\QQ\Lambda_X\), this becomes
\[
        \inf_{H\in \cS}\EE_\QQ[(X-H)^2]
        =
        (\alpha_X^\QQ)^2
        \|\Lambda_X-P_{\cS}^\QQ(\Lambda_X)\|_{L^2(\QQ)}^2
        +
        \|U_X^\QQ\|_{L^2(\QQ)}^2.
\]
Hence incompleteness is exactly the projection failure of the priced structural loading.
\end{theorem}

\begin{proof}
Fix \(H\in\cS\).  Since
\(
        H\in\cS\subseteq L^2(\QQ,\cG_T)\subseteq \cM_\QQ
\)
and \(L_X^\QQ\in\cM_\QQ\), one has \(L_X^\QQ-H\in\cM_\QQ\).  Because
\(U_X^\QQ\perp \cM_\QQ\), one has
\[
        \EE_\QQ[(L_X^\QQ-H)U_X^\QQ]=0.
\]
Therefore
\[
        \EE_\QQ[(X-H)^2]
        =
        \|L_X^\QQ-H\|_{L^2(\QQ)}^2
        +
        \|U_X^\QQ\|_{L^2(\QQ)}^2.
\]
The second term is independent of \(H\).  Taking the infimum over \(\cS\) gives the distance
identity, and closedness of \(\cS\) identifies the minimizer as \(P_{\cS}^\QQ L_X^\QQ\).  If
\(L_X^\QQ=\alpha_X^\QQ\Lambda_X\), linearity of orthogonal projection gives
\(P_{\cS}^\QQ L_X^\QQ=\alpha_X^\QQ P_{\cS}^\QQ\Lambda_X\), which yields the scalar-factor formula.
\end{proof}

\begin{proposition}[One-factor visible hedging residual for market variance]\label{thm:concrete-hedging-residual}
Assume Assumption~\ref{ass:one-factor-variance}.  Let
\(\cS\subseteq L^2(\QQ,\cG_T)\) be any closed visible hedge class.  Then
\[
        \inf_{H\in\cS}\EE_\QQ[(Q_M-H)^2]
        =
        \dist_{L^2(\QQ)}(V_Q,\cS)^2+
        \Delta_Q^2,
\]
where \(\Delta_Q^2\) is the hidden variance-residual constant from
Proposition~\ref{thm:concrete-variance-compression}.  In particular, if the visible hedge class is the
entire visible space \(\cS=L^2(\QQ,\cG_T)\), then
\[
        \inf_{H\in L^2(\QQ,\cG_T)}\EE_\QQ[(Q_M-H)^2]
        =
        \Delta_Q^2.
\]
If \(\Xi\) is symmetric, the full-visible hedge residual is
\[
        4a^2\EE_\QQ[A_G^2]+a^4B_f^2\Var_\QQ(\Xi^2).
\]
If moreover \(m\equiv0\) and \(\Xi\sim\mathcal N(0,1)\), the residual is the explicit floor
\[
        2a^4B_f^2.
\]
\end{proposition}

\begin{proof}
For any \(H\in\cS\subseteq L^2(\QQ,\cG_T)\), write
\[
        Q_M-H=(V_Q-H)+\bigl(2aA_G\Xi+a^2B_f(\Xi^2-1)\bigr).
\]
The visible term \(V_Q-H\) is orthogonal to the hidden centered term by the same independence and
centering argument used in Proposition~\ref{thm:concrete-variance-compression}.  Hence
\[
        \EE_\QQ[(Q_M-H)^2]
        =
        \|V_Q-H\|_{L^2(\QQ)}^2+
        \Delta_Q^2.
\]
Taking the infimum over \(H\in\cS\) gives the projection identity.  The special cases follow from
Proposition~\ref{thm:concrete-variance-compression} and Corollary~\ref{cor:pure-hidden-variance-floor}.
\end{proof}

\begin{remark}[Reading of the visible hedging residual]\label{rem:concrete-hedging-residual}
The abstract hedging theorem characterizes visible hedging failure as projection failure of the
priced structural loading.  Proposition~\ref{thm:concrete-hedging-residual} identifies that failure in
a one-factor aggregate-variance model.  The irreducible hidden piece is
\[
        2aA_G\Xi+a^2B_f(\Xi^2-1),
\]
which no \(\cG_T\)-measurable hedge can span.  The visible hedge class can improve only the visible
component \(V_Q\); it cannot reduce \(\Delta_Q^2\).
\end{remark}

\begin{corollary}[Factor-level decomposition of the visible-hedging residual]\label{cor:q-hedging-factor}
Under the assumptions of Proposition~\ref{thm:factor-decomp} and~\ref{thm:q-hedging}, if
\(\cS\subseteq L^2(\QQ,\cG_T)\) is a visible hedge class, then
\[
        \inf_{H\in \cS}\EE_\QQ[(X-H)^2]
        =
        \|V_X^\QQ-P_{\cS}^\QQ(V_X^\QQ)\|_{L^2(\QQ)}^2
        +
        (\alpha_X^{\QQ,\star})^2
        +
        \|R_X^{\QQ,\star}\|_{L^2(\QQ)}^2
        +
        \|U_X^\QQ\|_{L^2(\QQ)}^2.
\]
If \(X\) is aggregate-factor aligned, so \(R_X^{\QQ,\star}=0\), this becomes
\[
        \inf_{H\in \cS}\EE_\QQ[(X-H)^2]
        =
        \|V_X^\QQ-P_{\cS}^\QQ(V_X^\QQ)\|_{L^2(\QQ)}^2
        +
        (\alpha_X^{\QQ,\star})^2
        +
        \|U_X^\QQ\|_{L^2(\QQ)}^2.
\]
Thus the canonical centered aggregate factor contributes an exact irreducible term to the visible
hedging error whenever \(\alpha_X^{\QQ,\star}\neq0\).
\end{corollary}

\begin{proof}
By Proposition~\ref{thm:factor-decomp},
\[
        L_X^\QQ
        =
        V_X^\QQ+\alpha_X^{\QQ,\star}\Lambda_M^\star+R_X^{\QQ,\star}.
\]
Since \(\Lambda_M^\star\in\cM_\QQ^\circ=\cM_\QQ\ominus L^2(\QQ,\cG_T)\), the factor direction is
orthogonal to \(L^2(\QQ,\cG_T)\), hence to \(\cS\subseteq L^2(\QQ,\cG_T)\).  Therefore
\[
        P_{\cS}^\QQ(\Lambda_M^\star)=0.
\]
Also \(R_X^{\QQ,\star}\in \cM_\QQ^\circ\), so \(R_X^{\QQ,\star}\perp \cS\) and therefore
\[
        P_{\cS}^\QQ(R_X^{\QQ,\star})=0.
\]
Hence
\[
        L_X^\QQ-P_{\cS}^\QQ(L_X^\QQ)
        =
        \bigl(V_X^\QQ-P_{\cS}^\QQ(V_X^\QQ)\bigr)
        +
        \alpha_X^{\QQ,\star}\Lambda_M^\star
        +
        R_X^{\QQ,\star}.
\]
These three terms are mutually orthogonal for the same reason as in
Corollary~\ref{cor:q-compression-factor}.  Combining the resulting norm identity with
Theorem~\ref{thm:q-hedging} gives the claim.
\end{proof}

\begin{corollary}[Option first-variation claims inherit factor control of compression and hedging]\label{cor:option-factor-residuals}
Assume the hypotheses of Theorem~\ref{thm:option-factor}.  Fix a public
sub-\(\sigma\)-field \(\cH\subseteq\cG_T\), let \(\cS\subseteq L^2(\QQ,\cG_T)\) be a visible hedge
class, and let
\[
        V_{M,\tau}^\QQ:=C_{\cG_T}^\QQ(L_{M,\tau}^\QQ).
\]
Then
\[
        \|L_{M,\tau}^\QQ-C_{\cH}^\QQ(L_{M,\tau}^\QQ)\|_{L^2(\QQ)}^2
        \ge
        \|V_{M,\tau}^\QQ-C_{\cH}^\QQ(V_{M,\tau}^\QQ)\|_{L^2(\QQ)}^2
        +
        \bigl(|a_{M,\tau}^\star|-\delta_{M,\tau}^\star\bigr)_+^2,
\]
where \((x)_+:=\max\{x,0\}\).  Consequently, if \(\Phi\colon\R\to\R\) is convex and twice
continuously differentiable with \(\inf_x\Phi''(x)\ge \underline\kappa>0\), and if
\(\Phi(\Xi_{M,\tau}),\Phi(C_{\cH}^\QQ(\Xi_{M,\tau}))\in L^1(\QQ)\), then
\[
        \EE_\QQ[\Phi(\Xi_{M,\tau})]
        -
        \EE_\QQ[\Phi(C_{\cH}^\QQ(\Xi_{M,\tau}))]
        \ge
        \frac{\underline\kappa}{2}
        \left(
                \|V_{M,\tau}^\QQ-C_{\cH}^\QQ(V_{M,\tau}^\QQ)\|_{L^2(\QQ)}^2
                +
                \bigl(|a_{M,\tau}^\star|-\delta_{M,\tau}^\star\bigr)_+^2
        \right).
\]
Moreover,
\[
        \inf_{H\in\cS}\EE_\QQ[(\Xi_{M,\tau}-H)^2]
        \ge
        \|V_{M,\tau}^\QQ-P_{\cS}^\QQ(V_{M,\tau}^\QQ)\|_{L^2(\QQ)}^2
        +
        \bigl(|a_{M,\tau}^\star|-\delta_{M,\tau}^\star\bigr)_+^2
        +
        \|U_{M,\tau}^\QQ\|_{L^2(\QQ)}^2.
\]
If \(\delta_{M,\tau}^\star=0\), both bounds become exact identities with
\(\bigl(|a_{M,\tau}^\star|-\delta_{M,\tau}^\star\bigr)_+^2=(a_{M,\tau}^\star)^2\).
\end{corollary}

\begin{proof}
Under Theorem~\ref{thm:option-factor},
\[
        \widetilde L_{M,\tau}^\QQ
        =
        \alpha_{M,\tau}^{\QQ,\star}\Lambda_M^\star
        +
        R_{M,\tau}^{\QQ,\star},
\]
with
\[
        |\alpha_{M,\tau}^{\QQ,\star}-a_{M,\tau}^\star|
        \le
        \delta_{M,\tau}^\star,
        \qquad
        \|R_{M,\tau}^{\QQ,\star}\|_{L^2(\QQ)}
        \le
        \delta_{M,\tau}^\star.
\]
Hence
\[
        |\alpha_{M,\tau}^{\QQ,\star}|
        \ge
        \bigl(|a_{M,\tau}^\star|-\delta_{M,\tau}^\star\bigr)_+.
\]
Corollary~\ref{cor:q-compression-factor} gives
\[
        \|L_{M,\tau}^\QQ-C_{\cH}^\QQ(L_{M,\tau}^\QQ)\|_{L^2(\QQ)}^2
        =
        \|V_{M,\tau}^\QQ-C_{\cH}^\QQ(V_{M,\tau}^\QQ)\|_{L^2(\QQ)}^2
        +
        (\alpha_{M,\tau}^{\QQ,\star})^2
        +
        \|R_{M,\tau}^{\QQ,\star}\|_{L^2(\QQ)}^2,
\]
which implies the first lower bound.  The convex-compression estimate is then
Theorem~\ref{thm:q-compression} combined with that lower bound.

Likewise Corollary~\ref{cor:q-hedging-factor} yields
\[
        \inf_{H\in\cS}\EE_\QQ[(\Xi_{M,\tau}-H)^2]
        =
        \|V_{M,\tau}^\QQ-P_{\cS}^\QQ(V_{M,\tau}^\QQ)\|_{L^2(\QQ)}^2
        +
        (\alpha_{M,\tau}^{\QQ,\star})^2
        +
        \|R_{M,\tau}^{\QQ,\star}\|_{L^2(\QQ)}^2
        +
        \|U_{M,\tau}^\QQ\|_{L^2(\QQ)}^2,
\]
which implies the visible-hedging lower bound.  If \(\delta_{M,\tau}^\star=0\), then
Corollary~\ref{cor:option-factor-exact} gives
\[
        \alpha_{M,\tau}^{\QQ,\star}=a_{M,\tau}^\star,
        \qquad
        R_{M,\tau}^{\QQ,\star}=0,
\]
and the stated exact identities follow.
\end{proof}

\begin{corollary}[Visible-hedge-class residual]\label{cor:q-hedging-visible}
Under the assumptions of Theorem~\ref{thm:q-hedging},
\[
        \inf_{H\in \cS}\EE_\QQ[(X-H)^2]
        \ge
        \dist_{L^2(\QQ)}\!\bigl(L_X^\QQ,\cS\bigr)^2.
\]
Visible markets therefore do not complete structural claim content whenever the priced structural
loading does not belong to the visible hedge class.
\end{corollary}

\begin{proof}
Theorem~\ref{thm:q-hedging} gives the displayed lower bound after dropping the nonnegative term
\(\|U_X^\QQ\|_{L^2(\QQ)}^2\).  If \(L_X^\QQ\notin\cS\), the distance term is strictly positive by
closedness of \(\cS\); the visible hedge class leaves a residual fiber in the priced loading.
\end{proof}

\begin{corollary}[Structural compression across input geometries]\label{cor:program-closure}
Suppose the synchronized-perturbation contribution is screened in a visible summary family of the
type treated in Theorem~\ref{thm:market-visibility}, but \(L_X^\QQ\neq0\) for a variance-linked
claim \(X\).  Then:
\begin{enumerate}[label=(\roman*)]
\item the short-maturity pricing asymmetry of \cite{vidal2026latent} persists at market scale;
\item that short-maturity import is made either under a direct \(\QQ\)-Volterra specification with
the stated roughness exponents, or under an equivalent change of measure whose induced drift
correction is lower order at the short-maturity scale used in the skew expansion;
\item public screening under \(\QQ\) remains consistent with the observable-screening logic of
\cite{vidal2026dynamic};
\item the harmonized pairwise premium language of Section~\ref{sec:premium} states that the
surviving \(\PP\)-side screened object of \cite{vidal2026econometric} need not be recoverable from
the chosen visible \(\QQ\)-summary family, and that one-dimensional harmonized projected
information already determines premium support exactly once the pairwise normalizers and signs are
tracked;
\item visible prices and visible hedges remain compressions of structural claim content in the
sense of Corollary~\ref{prop:compression-unifier} and
Theorems~\ref{thm:q-compression} and~\ref{thm:q-hedging}; if the claim is aggregate-factor
aligned in the sense of Proposition~\ref{thm:factor-decomp}, then
Corollaries~\ref{cor:q-compression-factor} and~\ref{cor:q-hedging-factor} identify the canonical
centered aggregate factor as an exact additive source of the compression and hedging residuals.
\end{enumerate}
\end{corollary}

\begin{proof}
Part~(i) is the combination of Theorem~\ref{thm:market-visibility} with the threshold logic of
\cite{vidal2026latent}.  Part~(ii) is the short-maturity compatibility restriction needed when
the threshold is read under \(\QQ\).  Part~(iii) follows from the separation between structural
geometry and public visibility established in Sections~\ref{sec:geom} and~\ref{sec:compression}.  Part~(iv)
combines the existence of a nontrivial \(\PP\)-side screened object with the possibility of
risk-neutral screening at the public summary level and then expresses the resulting gap on the
harmonized projected basis of Section~\ref{sec:premium}.  Part~(v) is the common compression and
projection mechanism spelled out in Corollary~\ref{prop:compression-unifier},
Theorem~\ref{thm:q-compression}, and Theorem~\ref{thm:q-hedging}.
\end{proof}

\begin{corollary}[One structural loading, three public consequences]\label{cor:one-mechanism}
Fix a claim \(X\in L^2(\QQ,\cF_T)\) with priced structural loading \(L_X^\QQ=\alpha_X^\QQ\Lambda_X\),
fix a public sub-\(\sigma\)-field \(\cH\subseteq \cG_T\), and let
\(\cS\subseteq L^2(\QQ,\cG_T)\) be a visible hedge class.  Let \(\Phi\colon\R\to\R\) be convex and
twice continuously differentiable with
\[
        \inf_x \Phi''(x)\ge \underline\kappa>0,
\]
and assume \(\Phi(X),\Phi(C_{\cH}^\QQ(X))\in L^1(\QQ)\).  Then the three public mechanisms treated
here are generated by the same priced structural loading:
\begin{enumerate}[label=(\roman*)]
\item every visible summary functional \(\Psi_{\cH}\colon L^2(\QQ,\cH)\to\R^m\) acts on
\[
        \Psi_{\cH}(C_{\cH}^{\QQ}(X))
        =
        \Psi_{\cH}(\alpha_X^\QQ C_{\cH}^{\QQ}(\Lambda_X));
\]
\item every strongly convex pricing functional leaves a compression gap bounded below by
\[
        \frac{\underline\kappa}{2}
        (\alpha_X^\QQ)^2
        \|\Lambda_X-C_{\cH}^{\QQ}(\Lambda_X)\|_{L^2(\QQ)}^2;
\]
\item every visible hedge class leaves the exact residual
\[
        \inf_{H\in\cS}\EE_\QQ[(X-H)^2]
        =
        (\alpha_X^\QQ)^2
        \|\Lambda_X-P_{\cS}^{\QQ}(\Lambda_X)\|_{L^2(\QQ)}^2
        +
        \|U_X^\QQ\|_{L^2(\QQ)}^2.
\]
\end{enumerate}
Thus options, variance compression, and visible hedging are linear, nonlinear, and orthogonal
faces of the same public treatment of \(L_X^\QQ=\alpha_X^\QQ\Lambda_X\).  If one further resolves
\(L_X^\QQ\) through Proposition~\ref{thm:factor-decomp}, then
Corollaries~\ref{cor:q-compression-factor} and~\ref{cor:q-hedging-factor} show that the canonical
centered aggregate factor contributes an exact additive term \((\alpha_X^{\QQ,\star})^2\) to the
nonlinear compression gap and to the visible-hedging residual whenever \(X\) is
aggregate-factor aligned.
\end{corollary}

\begin{proof}
Part~(i) is Corollary~\ref{prop:compression-unifier}(iv).  Part~(ii) is the scalar-factor form of
Theorem~\ref{thm:q-compression}.  Part~(iii) is the scalar-factor form of
Theorem~\ref{thm:q-hedging}.
\end{proof}


\section{Scope and interpretation of the main results}
\label{sec:scope-gates}

The body uses one Hilbert-space structure.  A priced claim is first reduced to its structural
loading \(L_X^\QQ\).  Public option summaries, premium comparisons, variance claims, and visible
hedges contain only the corresponding \(\QQ\)-compression or projection of that loading.  The result
is narrower than a complete pricing model, a calibration theorem, or a market-completion statement.

Each endpoint has its own transfer condition.  The option branch requires a cumulant-class,
Hermite, jump, Gaussian/Bachelier, or local Black-window hypothesis connecting the quoted
coefficient to the compressed loading.  The premium branch requires signed coefficient vectors on a
common harmonized basis; unsigned projected information alone cannot give an unrestricted support
law.  The variance branch is exact for affine or explicitly reduced one-factor claims, while
nonlinear variance-linked claims expose a Hermite hierarchy rather than a universal scalar beta.
Black-implied-volatility transfer is local in maturity, moneyness, and vega.

Appendix~\ref{sec:detailed-scope-gates} records the endpoint obstructions and exact scope
conditions.  The structural relation remains fixed: option prices, premium comparisons,
variance-linked claims, and visible hedges are different public images of the same priced structural
loading.

\section{Conclusion}\label{sec:conclusion}

The \(\PP/\QQ\) compression link is generated by the Perron state.  Directed Volterra kernels,
roughness exponents, signed edge support, and the normalized Perron mode are primitive geometric
objects.  Equivalent pricing changes weighted projections of claims, not that geometry.
Finite-sample visibility weights and screened edge selections remain observational projections
unless they are imposed structurally.  The inherited object from \cite{vidal2026market} is the
Perron-generated sigma-field \(\cH_\varpi\); for market variance it also supplies the centered
increment
\[
        \Delta_{Q_M}
        =
        C_{\cH_\varpi}(Q_M)-C_{\cG_T}(Q_M).
\]
The market-scale construction of \cite{vidal2026market} derives the synchronized aggregate state.
This paper prices its public images under \(\QQ\).

Theorem~\ref{thm:main-risk-neutral-compression} is the organizing result.  Every treated public
pricing or hedging object is generated from the \(\QQ\)-side priced structural loading
\(L_X^\QQ\) by compression, projection, or measurable transformation.  Option summaries encode
conditional quote coefficients; premium results compare matched signed \(\PP/\QQ\) projected
loadings; variance claims expose nonlinear compression gaps; visible hedges leave projection
residuals.  Public prices and visible hedges do not recover the full latent state.  They select the
component that survives the specified risk-neutral compression.

\appendix

\section*{Technical appendices}
The appendices collect the option-transfer calculations and Gaussian benchmark reductions used
above.  They keep the option-surface branch checkable and leave the body focused on the compression
theorem, endpoint scope, and empirical diagnostic.



\section{Supporting non-Gaussian and jump option-transfer variants}
\label{sec:option-transfer-variants}

The option-transfer branch is stated through the compact cumulant-class theorem.  This appendix
records the Hermite, primitive non-Gaussian Volterra, and jump calculations that delimit the theorem's
scope.

\begin{assumption}[Hermite-admissible local non-Gaussian perturbation]\label{ass:hermite-admissible}
For every sufficiently small maturity \(\tau\), let
\[
        R_\tau^\varepsilon=X_\tau+\varepsilon Z_\tau+o_{L^2}(\varepsilon)
\]
under \(\QQ\), where
\[
        X_\tau\sim N(0,v_\tau),
        \qquad
        v_\tau=\vartheta\tau+o(\tau),
        \qquad
        \vartheta>0.
\]
Set \(Y_\tau:=X_\tau/\sqrt{v_\tau}\).  Assume \(Z_\tau\in L^2(\QQ)\),
\(\EE_\QQ[Z_\tau]=0\), and that its conditional projection onto \(\sigma(X_\tau)\) admits the
finite Hermite expansion
\[
        \EE_\QQ[Z_\tau\mid X_\tau]
        =
        \sum_{m=1}^{M_\tau} a_{m,\tau}\He_m(Y_\tau)+r_\tau(Y_\tau),
\]
where \(M_\tau\ge2\) and \(r_\tau\perp \He_m(Y_\tau)\) for \(0\le m\le M_\tau\).  For the
option-price statements below assume either \(r_\tau=0\), or that the displayed expansion holds in
\(L^2\) and the residual term is lower order on the moneyness window under consideration.
\end{assumption}

\begin{lemma}[Option-price derivative generated by the conditional projection]\label{lem:price-projection}
Under Assumption~\ref{ass:hermite-admissible}, for every strike coordinate \(k\) at which the
baseline distribution is continuous,
\[
        \partial_\varepsilon\EE_\QQ[(R_\tau^\varepsilon-k)_+]\big|_{\varepsilon=0}
        =
        \EE_\QQ[Z_\tau\one_{\{X_\tau>k\}}].
\]
Writing \(z=k/\sqrt{v_\tau}\), this equals
\[
        \phi(z)
        \sum_{m=1}^{M_\tau}a_{m,\tau}\He_{m-1}(z)
        +
        \EE_\QQ[r_\tau(Y_\tau)\one_{\{Y_\tau>z\}}].
\]
\end{lemma}

\begin{proof}
The directional derivative of \(x\mapsto(x-k)_+\) is \(\one_{\{x>k\}}\) at continuity points of
the baseline law.  The assumed \(L^2\)-differentiability gives the first identity by dominated
localization.  Taking conditional expectation onto \(\sigma(X_\tau)\), substituting the Hermite
expansion, and using
\[
        \EE[\He_m(Y)\one_{\{Y>z\}}]=\phi(z)\He_{m-1}(z),
        \qquad
        Y\sim N(0,1),
\]
gives the second identity.
\end{proof}

\begin{theorem}[Hermite option-price transfer and the cumulant obstruction]\label{thm:hermite-option-transfer}
Assume Assumption~\ref{ass:hermite-admissible} with \(r_\tau=0\).  Let
\(\sigma_N(\tau,k;\varepsilon)\) be the Bachelier implied normal volatility, and let
\(\sigma_{N,0}(\tau):=\sqrt{v_\tau/\tau}\).  Then, for fixed \(\tau\),
\[
        \partial_\varepsilon\sigma_N(\tau,k;\varepsilon)\big|_{\varepsilon=0}
        =
        \frac{1}{\sqrt\tau}
        \sum_{m=1}^{M_\tau}a_{m,\tau}\He_{m-1}\!\left(\frac{k}{\sqrt{v_\tau}}\right).
\]
Consequently the ATM normal-skew derivative is
\[
        \partial_k\partial_\varepsilon\sigma_N(\tau,k;\varepsilon)\big|_{k=0,\varepsilon=0}
        =
        \frac{1}{\sqrt\tau\sqrt{v_\tau}}
        \sum_{m=2}^{M_\tau}(m-1)a_{m,\tau}\He_{m-2}(0).
\]
The third-cumulant derivative is
\[
        \partial_\varepsilon \kappa_3^\QQ(R_\tau^\varepsilon)\big|_{\varepsilon=0}
        =
        6v_\tau a_{2,\tau}.
\]
Hence
\[
\begin{aligned}
        \partial_k\partial_\varepsilon\sigma_N(\tau,k;\varepsilon)\big|_{k=0,\varepsilon=0}
        &=
        \frac{\partial_\varepsilon \kappa_3^\QQ(R_\tau^\varepsilon)|_{\varepsilon=0}}
        {6v_\tau^{3/2}\sqrt\tau} \\
        &\quad+
        \frac{1}{\sqrt\tau\sqrt{v_\tau}}
        \sum_{\substack{m\ge4\\ m\le M_\tau}}
        (m-1)a_{m,\tau}\He_{m-2}(0),
\end{aligned}
\]
because the odd \(m\)-terms with \(m\ge3\) vanish at the money.  Therefore the cumulant-only
Bachelier transfer is exact if the conditional projection has no even Hermite components beyond
\(\He_2\), and it is asymptotically exact at scale \(b_\tau\) if
\[
        \frac{1}{\sqrt\tau\sqrt{v_\tau}}
        \sum_{\substack{m\ge4\\ m\le M_\tau}}
        (m-1)a_{m,\tau}\He_{m-2}(0)
        =
        o(b_\tau).
\]
\end{theorem}

\begin{proof}
At the baseline normal price, Bachelier vega is
\[
        \partial_\sigma C_N(\tau,k;\sigma_{N,0})
        =
        \sqrt\tau\,\phi(k/\sqrt{v_\tau}).
\]
Dividing the price derivative in Lemma~\ref{lem:price-projection} by this vega gives the
normal-volatility derivative.  Differentiating in \(k\) and using
\(\frac{\dd}{\dd z}\He_{m-1}(z)=(m-1)\He_{m-2}(z)\) gives the ATM formula.  Finally,
\[
        \partial_\varepsilon\kappa_3^\QQ(X_\tau+\varepsilon Z_\tau)|_{\varepsilon=0}
        =
        3\EE_\QQ[X_\tau^2Z_\tau]
        =
        3v_\tau\EE_\QQ[Y_\tau^2\EE_\QQ[Z_\tau\mid Y_\tau]].
\]
Since \(Y^2=\He_2(Y)+1\) and \(\EE_\QQ[Z_\tau]=0\), only the \(\He_2\)-coefficient survives, and
\(\EE[\He_2(Y)^2]=2\).  Thus the cumulant derivative is \(6v_\tau a_{2,\tau}\).  Substitution
proves the decomposition, and the parity statement follows from \(\He_n(0)=0\) for odd \(n\).
\end{proof}

\begin{corollary}[Second-chaos closure of the cumulant transfer]\label{cor:second-chaos-closure}
If
\[
        \EE_\QQ[Z_\tau\mid X_\tau]=a_{2,\tau}\He_2(X_\tau/\sqrt{v_\tau})
\]
for all sufficiently small \(\tau\), then the Bachelier ATM skew-cumulant transfer is exact:
\[
        \partial_k\partial_\varepsilon\sigma_N(\tau,k;\varepsilon)\big|_{k=0,\varepsilon=0}
        =
        \frac{\partial_\varepsilon \kappa_3^\QQ(R_\tau^\varepsilon)|_{\varepsilon=0}}
        {6v_\tau^{3/2}\sqrt\tau}.
\]
This is the continuous-return second-chaos branch generated by the Brownian-correlated Volterra
leverage terms.
\end{corollary}

\begin{proof}
Under the displayed assumption all Hermite coefficients in Theorem~\ref{thm:hermite-option-transfer} vanish
except the second one.  The skew derivative in that theorem therefore reduces to the
\(m=2\) term, while the cumulant identity in the same proof gives
\(\partial_\varepsilon\kappa_3^\QQ(R_\tau^\varepsilon)|_{\varepsilon=0}=6v_\tau a_{2,\tau}\).  Substitution
into the ATM Bachelier skew formula gives the displayed exact transfer.
\end{proof}

\begin{remark}[Sharp meaning of the non-Gaussian limitation]\label{rem:ng-hermite-obstruction}
A non-Gaussian option-price theorem cannot be universally reduced to the third cumulant: higher
even Hermite projections of the return perturbation also affect ATM skew.  The cumulant-class
theorem is therefore the correct theorem for the branch where the second Hermite projection
dominates, not a universal law for every non-Gaussian return family.
\end{remark}

\begin{assumption}[Non-Gaussian Volterra leverage primitives]\label{ass:ng-volterra-primitives}
Fix a finite active channel set \(\cE^{\mathrm{ng}}\).  On
\((\Omega,\cF,(\cF_t)_{t\ge0},\QQ)\), let \(B\) be the option-return Brownian motion.  For each
source index \(j\), let \(B^{\perp,j}\) be a Brownian motion independent of \(B\), and let
\(\widetilde J^j\) be a compensated compound-Poisson martingale,
\[
        \widetilde J_t^j
        =
        \sum_{m=1}^{N_t^j}\xi_m^j
        -\nu_jt\,\EE_\QQ[\xi_1^j],
\]
where \(N^j\) is a Poisson process of intensity \(\nu_j\), the jumps \(\xi_m^j\) have finite
fourth moment, and the families are mutually independent except through the common Brownian motion
stated below.  Let
\[
        L_t^j
        :=
        \rho_j B_t+\bar\rho_j B_t^{\perp,j}+\chi_j\widetilde J_t^j,
        \qquad
        \rho_j\in[-1,1],
        \quad
        \bar\rho_j\ge0,
        \quad
        \rho_j^2+\bar\rho_j^2\le1,
        \quad
        \chi_j\ge0 .
\]
Whenever \(\chi_j>0\) and the jump law is nondegenerate, \(L^j\) is non-Gaussian.

For each active channel \(e\in\cE^{\mathrm{ng}}\), write \(j(e)\) for its source index and assume
\[
        G_e(t,s)
        =
        d_e a_e(t,s)(t-s)^{H_e-1/2},
        \qquad
        0\le s<t\le T,
\]
where \(d_e\in\R\), \(H_e\in(0,1)\), and \(a_e\) is continuous with \(a_e(0,0)>0\) when
\(d_e\ne0\).  The baseline normal volatility \(\sigma_0\) is deterministic, strictly positive near
zero, and continuous.  There are constants \(C,t_0>0\) and \(\eta\in(0,1]\) such that, for
\(0\le r\le t\le t_0\),
\[
        |\sigma_0(t)-\sigma_0(0)|\le Ct^\eta,
        \qquad
        |\sigma_0(r)a_e(t,r)-\sigma_0(0)a_e(0,0)|\le Ct^\eta
        \quad(e\in\cE^{\mathrm{ng}}).
\]
Define the predictable Volterra leverage field by
\[
        Y_t
        :=
        \sum_{e\in\cE^{\mathrm{ng}}}\int_0^{t-}G_e(t,s)\,\dd L_s^{j(e)},
\]
and, for a local synchronized perturbation parameter \(\varepsilon\), define
\[
        R_\tau^\varepsilon
        :=
        \int_0^\tau (\sigma_0(t)+\varepsilon Y_t)\,\dd B_t,
        \qquad
        X_\tau:=R_\tau^0,
        \qquad
        Z_\tau:=\int_0^\tau Y_t\,\dd B_t .
\]
Assume \(X_\tau^2Z_\tau\in L^1(\QQ)\) for sufficiently small \(\tau\).
\end{assumption}

\begin{definition}[Primitive non-Gaussian Volterra coefficients]\label{def:ng-volterra-coefficients}
Under Assumption~\ref{ass:ng-volterra-primitives}, set
\[
        \vartheta:=\sigma_0(0)^2,
        \qquad
        \cE_\rho^{\mathrm{ng}}:=\{e\in\cE^{\mathrm{ng}}:d_e\rho_{j(e)}\ne0\}.
\]
For \(e\in\cE_\rho^{\mathrm{ng}}\), define
\[
        K_e
        :=
        \frac{6\sigma_0(0)^2d_e\rho_{j(e)}a_e(0,0)}{(H_e+1/2)(H_e+3/2)}.
\]
Let
\[
        \mathcal B_J:=\{H_e:e\in\cE_\rho^{\mathrm{ng}}\}.
\]
For \(\beta\in\mathcal B_J\), set
\[
        A_J(\beta)
        :=
        \sum_{e\in\cE_\rho^{\mathrm{ng}}:\,H_e=\beta}
        \frac{d_e\rho_{j(e)}a_e(0,0)}{\sigma_0(0)(H_e+1/2)(H_e+3/2)}
        =
        \frac{1}{6\vartheta^{3/2}}
        \sum_{e\in\cE_\rho^{\mathrm{ng}}:\,H_e=\beta}K_e .
\]
The primitive compressed support selected by the cumulant condition is
\[
        \operatorname{supp}_J:=\{\beta\in\mathcal B_J:A_J(\beta)\ne0\},
\]
with selected exponent
\[
        \beta_J^{\mathrm{comp}}:=\inf \operatorname{supp}_J
\]
when the support is nonempty.
\end{definition}

\begin{theorem}[Primitive non-Gaussian Volterra cumulant expansion]\label{thm:ng-volterra-cumulants}
Assume Assumption~\ref{ass:ng-volterra-primitives}.  Then
\[
        v(\tau):=\Var_\QQ(X_\tau)
        =
        \vartheta\tau+O(\tau^{1+\eta}).
\]
If \(\cE_\rho^{\mathrm{ng}}=\varnothing\), then
\[
        \partial_\varepsilon\kappa_3^\QQ(R_\tau^\varepsilon)\big|_{\varepsilon=0}=0
        \qquad
        \text{for all sufficiently small }\tau.
\]
If \(\cE_\rho^{\mathrm{ng}}\ne\varnothing\), then, as \(\tau\downarrow0\),
\[
        \partial_\varepsilon\kappa_3^\QQ(R_\tau^\varepsilon)\big|_{\varepsilon=0}
        =
        \sum_{e\in\cE_\rho^{\mathrm{ng}}}K_e\tau^{H_e+3/2}
        +
        O\!\left(
                \sum_{e\in\cE_\rho^{\mathrm{ng}}}|d_e\rho_{j(e)}|\tau^{H_e+3/2+\eta}
        \right).
\]
Equivalently,
\[
        \partial_\varepsilon\kappa_3^\QQ(R_\tau^\varepsilon)\big|_{\varepsilon=0}
        =
        6\vartheta^{3/2}
        \sum_{\beta\in\mathcal B_J}A_J(\beta)\tau^{\beta+3/2}
        +
        O(\tau^{\beta_0+3/2+\eta}),
\]
where \(\beta_0:=\min\mathcal B_J\).  In particular, if
\[
        \operatorname{supp}_J\ne\varnothing
        \quad\text{and}\quad
        \beta_0+\eta>\beta_J^{\mathrm{comp}},
\]
then the cumulant part of Assumption~\ref{ass:cumulant-class} holds with support
\(\operatorname{supp}_J\), selected exponent \(\beta_J^{\mathrm{comp}}\), and remainder
\[
        o(\tau^{\beta_J^{\mathrm{comp}}+3/2}).
\]
A sufficient simpler condition is \(A_J(\beta_0)\ne0\); in that case
\(\beta_J^{\mathrm{comp}}=\beta_0\), and the displayed remainder is automatically lower order.
\end{theorem}

\begin{proof}
The variance identity follows from It\^o isometry:
\[
        v(\tau)=\int_0^\tau \sigma_0(t)^2\,\dd t
        =\sigma_0(0)^2\tau+O(\tau^{1+\eta}).
\]
Since \(R_\tau^\varepsilon=X_\tau+\varepsilon Z_\tau\), \(X_\tau\) is centered Gaussian, and
\(\EE_\QQ[Z_\tau]=0\),
\[
        \partial_\varepsilon\kappa_3^\QQ(R_\tau^\varepsilon)\big|_{\varepsilon=0}
        =3\EE_\QQ[X_\tau^2Z_\tau].
\]
It\^o's product formula gives
\[
        \EE_\QQ[X_\tau^2Z_\tau]
        =
        2\int_0^\tau \sigma_0(t)\EE_\QQ[X_tY_t]\,\dd t .
\]
Only the \(\rho_jB\)-component of \(L^j\) has nonzero covariance with the return Brownian motion.
The orthogonal Brownian components and the compensated jump martingales have zero contribution to
\(\EE_\QQ[X_tY_t]\).  Hence
\[
        \EE_\QQ[X_tY_t]
        =
        \sum_{e\in\cE_\rho^{\mathrm{ng}}}
        d_e\rho_{j(e)}
        \int_0^t \sigma_0(r)a_e(t,r)(t-r)^{H_e-1/2}\,\dd r .
\]
The H\"older condition implies, uniformly over the finite active set,
\[
        \int_0^t \sigma_0(r)a_e(t,r)(t-r)^{H_e-1/2}\,\dd r
        =
        \frac{\sigma_0(0)a_e(0,0)}{H_e+1/2}t^{H_e+1/2}
        +O(t^{H_e+1/2+\eta}).
\]
Multiplying by the outer factor \(2\sigma_0(t)\), integrating in \(t\), and multiplying by the
third-cumulant factor \(3\) yields
\[
        3\EE_\QQ[X_\tau^2Z_\tau]
        =
        \sum_{e\in\cE_\rho^{\mathrm{ng}}}
        \frac{6\sigma_0(0)^2d_e\rho_{j(e)}a_e(0,0)}{(H_e+1/2)(H_e+3/2)}
        \tau^{H_e+3/2}
        +
        O\!\left(
                \sum_{e\in\cE_\rho^{\mathrm{ng}}}|d_e\rho_{j(e)}|\tau^{H_e+3/2+\eta}
        \right),
\]
which is the stated expansion.  Grouping equal powers gives the coefficient form.  If the grouped
support is nonempty and \(\beta_0+\eta>\beta_J^{\mathrm{comp}}\), then every lower-order
remainder, including remainders attached to powers whose leading coefficients cancel, is
\(o(\tau^{\beta_J^{\mathrm{comp}}+3/2})\).  The final statement follows.
\end{proof}

\begin{remark}[Role of the non-Gaussian jump component]\label{rem:ng-volterra-jump-role}
The jump sources make the latent Volterra driver non-Gaussian.  At first order in the synchronized
volatility perturbation, the third cumulant is selected by the predictable covariance between the
latent source and the return Brownian motion.  Thus the compound-Poisson component changes the law
and higher moments of the leverage field, while the leading first-order skew coefficient is carried
by the Brownian-correlated channel.  The pure rare-jump mechanism below gives a complementary class
in which the jump component itself carries the leading third cumulant.
\end{remark}

\begin{assumption}[Self-similar rare-jump option-facing return]\label{ass:rare-jump}
Let \(X_\tau=\int_0^\tau\sigma_0(t)\,\dd B_t\) be the same centered Gaussian baseline as above,
with \(\sigma_0(0)>0\).  For each channel \(e\in\cE^J\), fix \(\beta_e>0\), an intensity
coefficient \(q_e\ge0\), and a jump-size law \(F_e\) on \(\R\) with finite fourth moment and third
moment
\[
        m_{3,e}:=\int x^3F_e(\dd x).
\]
For a local intensity perturbation \(\varepsilon\ge0\), let \(N_e^\varepsilon\) be a Poisson random
measure on \((0,T]\times\R\), independent of \(B\), with compensator
\[
        \varepsilon q_e s^{\beta_e-1}\,\dd s\,F_e(\dd x).
\]
The rare-jump perturbation return is
\[
        J_\tau^\varepsilon
        :=
        \sum_{e\in\cE^J}\int_{(0,\tau]\times\R}s^{1/2}x\,
        \bigl(N_e^\varepsilon(\dd s,\dd x)-\varepsilon q_e s^{\beta_e-1}\dd sF_e(\dd x)\bigr),
\]
and
\[
        R_\tau^\varepsilon:=X_\tau+J_\tau^\varepsilon.
\]
\end{assumption}

\begin{theorem}[Exact rare-jump cumulant expansion]\label{thm:rare-jump-cumulants}
Under Assumption~\ref{ass:rare-jump},
\[
        \Var_\QQ(X_\tau)=\sigma_0(0)^2\tau+o(\tau).
\]
Moreover, for the right derivative at \(\varepsilon=0\) of the third cumulant of
\(R_\tau^\varepsilon\),
\[
        \partial_\varepsilon^+\kappa_3^\QQ(R_\tau^\varepsilon)\big|_{\varepsilon=0}
        =
        \sum_{e\in\cE^J}
        \frac{q_em_{3,e}}{\beta_e+3/2}\tau^{\beta_e+3/2}.
\]
Thus the cumulant support is
\[
        \operatorname{supp}_{J,\mathrm{rare}}
        =
        \left\{
        \beta:\sum_{e:\beta_e=\beta}\frac{q_em_{3,e}}{\beta+3/2}\ne0
        \right\},
\]
and, with \(\vartheta=\sigma_0(0)^2\), the corresponding coefficients are
\[
        A_{J,\mathrm{rare}}(\beta)
        =
        \frac{1}{6\vartheta^{3/2}}
        \sum_{e:\beta_e=\beta}\frac{q_em_{3,e}}{\beta+3/2}.
\]
There is no asymptotic cumulant remainder beyond grouping equal powers.
\end{theorem}

\begin{proof}
The baseline variance follows from It\^o's isometry and continuity of \(\sigma_0\) at zero.  Cumulants of
independent random variables add.  The centered Gaussian variable \(X_\tau\) has zero third
cumulant, and the compensated Poisson integral has third cumulant equal to the integral of the
third power of its jump amplitude against its compensator, the standard compound-Poisson cumulant
identity \cite{sato1999levy,contTankov2004jumps}.  Therefore
\[
        \kappa_3^\QQ(J_\tau^\varepsilon)
        =
        \varepsilon
        \sum_{e\in\cE^J}\int_0^\tau\int_\R (s^{1/2}x)^3
        q_es^{\beta_e-1}\,\dd sF_e(\dd x),
\]
which is exactly
\[
        \varepsilon
        \sum_{e\in\cE^J}\frac{q_em_{3,e}}{\beta_e+3/2}\tau^{\beta_e+3/2}.
\]
Taking the right derivative at \(\varepsilon=0\) gives the formula.
\end{proof}

\begin{remark}[Use of the two primitive non-Gaussian mechanisms]\label{rem:primitive-ng-lowerers}
The Volterra leverage class is the natural non-Gaussian extension of the option-facing dynamics in
the Gaussian block.  It supplies the cumulant power law from explicit primitive sources while
preserving the continuous-return second-chaos structure used in the Bachelier transfer below.  The
rare-jump class is a separate primitive mechanism for the cumulant condition.  It shows that the same
power structure can also arise from genuinely jump-carried third cumulants, with exact remainder
control; this is the local use of standard jump-process cumulant calculus
\cite{sato1999levy,contTankov2004jumps}.  For large rare jumps, option-price first variations need not be determined only by the
third cumulant, so the self-contained Bachelier transfer is stated for the continuous-return
second-chaos leverage class.
\end{remark}

\begin{remark}[Scope of the Poisson mechanisms]\label{rem:poisson-lowerer-scope}
The Poisson terms above are local primitive mechanisms, not a change in the baseline Gaussian or
risk-neutral compression model.  Setting the jump amplitudes to zero recovers the Brownian
Volterra leverage class.  The compensated compound-Poisson term in
Assumption~\ref{ass:ng-volterra-primitives} is used only to show that genuinely non-Gaussian
latent Volterra drivers can coexist with the same public Brownian leverage channel, while
Assumption~\ref{ass:rare-jump} is used only to exhibit an exact rare-jump cumulant power law.
Neither Poisson mechanism is invoked in the scalar Gaussian seam, the local maturity-anchor closure,
the premium transfer, or the variance and hedging compression theorems unless it is explicitly
referenced.
\end{remark}

\begin{assumption}[Compensated rare-jump perturbation of a normal baseline]\label{ass:jump-price}
Fix a maturity \(\tau\).  Let \(X_\tau\sim N(0,v_\tau)\), \(v_\tau>0\).  For a local intensity
parameter \(\varepsilon\ge0\), let the perturbation be a compensated compound-Poisson jump measure
on \(\R\) with finite first absolute moment and finite third absolute moment.  Its intensity
measure is \(\varepsilon\mu_\tau(\dd y)\), and
\[
        J_\tau^\varepsilon
        =
        \sum_{n=1}^{N_\tau^\varepsilon}Y_n
        -
        \varepsilon\int_\R y\,\mu_\tau(\dd y),
\]
where \(N_\tau^\varepsilon\) has intensity \(\varepsilon\mu_\tau(\R)\), the marks have law
\(\mu_\tau/\mu_\tau(\R)\) when \(\mu_\tau(\R)>0\), and \(J_\tau^\varepsilon\) is independent of
\(X_\tau\).  Set \(R_\tau^\varepsilon:=X_\tau+J_\tau^\varepsilon\).
\end{assumption}

\begin{theorem}[Exact rare-jump option-price first variation]\label{thm:exact-jump-price}
Under Assumption~\ref{ass:jump-price}, let
\[
        C_0(k):=\EE_\QQ[(X_\tau-k)_+].
\]
Then
\[
        \partial_\varepsilon\EE_\QQ[(R_\tau^\varepsilon-k)_+]\big|_{\varepsilon=0}
        =
        \int_\R
        \left[
                C_0(k-y)-C_0(k)+yC_0'(k)
        \right]
        \mu_\tau(\dd y).
\]
Consequently the Bachelier implied-normal-volatility derivative is
\[
        \partial_\varepsilon\sigma_N(\tau,k;\varepsilon)\big|_{\varepsilon=0}
        =
        \frac{1}{\sqrt\tau\,\phi(k/\sqrt{v_\tau})}
        \int_\R
        \left[
                C_0(k-y)-C_0(k)+yC_0'(k)
        \right]
        \mu_\tau(\dd y).
\]
The ATM normal-skew derivative is the exact tail-kernel expression
\[
        \partial_k\partial_\varepsilon\sigma_N(\tau,k;\varepsilon)\big|_{k=0,\varepsilon=0}
        =
        \frac{1}{\sqrt\tau\,\phi(0)}
        \int_\R
        \left[
                -\Phi\!\left(\frac{y}{\sqrt{v_\tau}}\right)
                +\frac12
                +\frac{y}{\sqrt{v_\tau}}\phi(0)
        \right]
        \mu_\tau(\dd y).
\]
\end{theorem}

\begin{proof}
To first order in intensity, either no jump occurs or one jump occurs.  The compensation
contributes the deterministic shift \(-\varepsilon\int y\mu_\tau(\dd y)\), whose first-order price
effect is \(-\int y\mu_\tau(\dd y)\QQ(X_\tau>k)\).  Since
\(C_0'(k)=-\QQ(X_\tau>k)\), the price derivative is
\[
        \int_\R
        \{\EE_\QQ[(X_\tau+y-k)_+]-\EE_\QQ[(X_\tau-k)_+]+yC_0'(k)\}
        \mu_\tau(\dd y).
\]
The identity \(\EE_\QQ[(X_\tau+y-k)_+]=C_0(k-y)\) gives the first formula.  Dividing by Bachelier
vega \(\sqrt\tau\phi(k/\sqrt{v_\tau})\) gives the implied-volatility derivative.  Differentiating
the price kernel in \(k\) and evaluating at zero gives
\[
        C_0'(-y)-C_0'(0)+yC_0''(0)
        =
        -\Phi(y/\sqrt{v_\tau})+\frac12+\frac{y}{\sqrt{v_\tau}}\phi(0),
\]
which proves the ATM formula.
\end{proof}

\begin{corollary}[Small-amplitude jumps recover the cumulant transfer]\label{cor:small-jumps-cumulant}
Assume the jump measure is concentrated on \(|y|\le \delta_\tau\sqrt{v_\tau}\), where
\(\delta_\tau\downarrow0\), and
\[
        \int |y|^5\,\mu_\tau(\dd y)<\infty.
\]
Then
\[
\begin{aligned}
        \partial_k\partial_\varepsilon\sigma_N(\tau,k;\varepsilon)\big|_{k=0,\varepsilon=0}
        &=
        \frac{1}{6v_\tau^{3/2}\sqrt\tau}
        \int y^3\mu_\tau(\dd y)  \\
        &\quad+
        O\!\left(
                \frac{1}{v_\tau^{5/2}\sqrt\tau}
                \int |y|^5\mu_\tau(\dd y)
        \right).
\end{aligned}
\]
Since
\[
        \partial_\varepsilon\kappa_3^\QQ(R_\tau^\varepsilon)\big|_{\varepsilon=0}
        =
        \int y^3\mu_\tau(\dd y),
\]
the leading term is exactly the cumulant transfer.
\end{corollary}

\begin{proof}
For \(u=y/\sqrt{v_\tau}\), Taylor expansion gives
\[
        -\Phi(u)+\frac12+u\phi(0)
        =
        \frac{\phi(0)}{6}u^3+O(|u|^5)
\]
uniformly for \(|u|\le\delta_\tau\).  Substitute this into the kernel of
Theorem~\ref{thm:exact-jump-price} and divide by \(\sqrt\tau\phi(0)\).  The cumulant identity is
the standard compound-Poisson cumulant formula for a compensated jump integral
\cite{sato1999levy,contTankov2004jumps}.
\end{proof}

\begin{proposition}[Large jumps are not determined by the third cumulant]\label{prop:large-jump-not-cumulant}
Fix \(v_\tau>0\).  There exist two finite jump intensity measures \(\mu_\tau^{(1)}\) and
\(\mu_\tau^{(2)}\) with
\[
        \int y^3\mu_\tau^{(1)}(\dd y)
        =
        \int y^3\mu_\tau^{(2)}(\dd y)
\]
but with different ATM jump-skew derivatives in Theorem~\ref{thm:exact-jump-price}.  Hence no
large-jump option-price transfer can depend only on the third cumulant.
\end{proposition}

\begin{proof}
The ATM jump-skew derivative is the integral of
\[
        K_v(y):=-\Phi(y/\sqrt v)+\frac12+\frac{y}{\sqrt v}\phi(0)
\]
against the jump measure, up to the fixed factor \(1/(\sqrt\tau\phi(0))\).  The function \(K_v\)
is not a constant multiple of \(y^3\) on \(\R\).  Hence there is a finite signed measure \(\nu\)
with \(\int y^3\nu(\dd y)=0\) but \(\int K_v(y)\nu(\dd y)\ne0\).  Adding a sufficiently small
multiple of \(\nu\) to any strictly positive finite measure with finite third moment gives two
positive finite measures with equal third moment and different kernel integrals.
\end{proof}

\begin{remark}[Jump-transfer scope]\label{rem:jump-transfer-scope}
Rare or small jumps are primitive mechanisms for cumulant powers and, under the small-amplitude
condition of Corollary~\ref{cor:small-jumps-cumulant}, also fall into the cumulant option-transfer
branch.  Large jumps have the exact transfer in Theorem~\ref{thm:exact-jump-price}; the governing
coefficient is the tail-kernel integral, not the third cumulant.
\end{remark}



\section{Premium implementation and support classification details}
\label{sec:premium-implementation-appendix}

This appendix collects the estimator, covariance-normalizer, finite-dimensional support, and
unsigned-information results that support the premium section without serving as main-text theorem
endpoints.

\begin{definition}[One-dimensional harmonized projected information]\label{def:premium-info}
Fix an ordered pair \((i,j)\) and suppose the harmonized reduced experiment under
\(\mu\in\{\PP,\QQ\}\) is one-dimensional on the common source-screened basis \(\varphi_{ij}\).  Let
\(\omega_{ij}^{\mu,\mathrm{harm}}>0\) denote the corresponding one-dimensional covariance
normalizer on that basis.  Define the harmonized one-dimensional projected information by
\[
        \mathcal I_{ij}^{\mu,\mathrm{harm}}
        :=
        \frac{(\gamma_{ij}^\mu)^2}{\omega_{ij}^{\mu,\mathrm{harm}}}.
\]
This is the one-dimensional harmonized counterpart of the projected information quantities used in
the operational \(\PP/\QQ\) classification template of \cite{vidal2026econometric}.
\end{definition}

\begin{assumption}[Estimated harmonized covariance normalizers]\label{ass:premium-info-est}
In the setup of Assumption~\ref{ass:premium}, suppose moreover that for each ordered pair
\((i,j)\) in the genuinely one-dimensional harmonized regime there exist positive estimators
\[
        \widehat\omega_{ij,n}^{\PP,\mathrm{harm}},
        \qquad
        \widehat\omega_{ij,n}^{\QQ,\mathrm{harm}},
\]
such that
\[
        \sqrt{n}
        \begin{pmatrix}
                \widehat\gamma_{ij,n}^\PP-\gamma_{ij}^\PP\\
                \widehat\gamma_{ij,n}^\QQ-\gamma_{ij}^\QQ\\
                \widehat\omega_{ij,n}^{\PP,\mathrm{harm}}-\omega_{ij}^{\PP,\mathrm{harm}}\\
                \widehat\omega_{ij,n}^{\QQ,\mathrm{harm}}-\omega_{ij}^{\QQ,\mathrm{harm}}
        \end{pmatrix}
        \Rightarrow
        \mathcal N(0,\Omega_{ij}^{\mathrm{info}})
\]
for some positive semidefinite \(4\times4\) matrix \(\Omega_{ij}^{\mathrm{info}}\).
\end{assumption}

\begin{proposition}[Empirical harmonized covariance normalizers from matched coefficient observations]\label{prop:premium-info-normalizers}
Fix an ordered pair \((i,j)\) in the genuinely one-dimensional harmonized regime and suppose the
matched coefficient-observation vectors
\[
        \left(
                \widehat g_{ij,d}^\PP,
                \widehat g_{ij,d}^\QQ
        \right),
        \qquad
        d=1,\dots,n,
\]
are i.i.d. across matched dates with finite fourth moments, with
\[
        \EE[\widehat g_{ij,d}^\mu]=\gamma_{ij}^\mu,
        \qquad
        \Var(\widehat g_{ij,d}^\mu)=\omega_{ij}^{\mu,\mathrm{harm}},
        \qquad
        \mu\in\{\PP,\QQ\}.
\]
Define the empirical one-dimensional covariance normalizers by
\[
        \widetilde\omega_{ij,n}^{\mu,\mathrm{harm}}
        :=
        \frac1n
        \sum_{d=1}^n
        \left(
                \widehat g_{ij,d}^\mu-\widehat\gamma_{ij,n}^\mu
        \right)^2,
        \qquad
        \mu\in\{\PP,\QQ\}.
\]
Then
\[
        \widetilde\omega_{ij,n}^{\mu,\mathrm{harm}}
        \to
        \omega_{ij}^{\mu,\mathrm{harm}}
        \qquad\text{in probability},
        \qquad
        \mu\in\{\PP,\QQ\},
\]
and
\[
        \sqrt{n}
        \begin{pmatrix}
                \widehat\gamma_{ij,n}^\PP-\gamma_{ij}^\PP\\
                \widehat\gamma_{ij,n}^\QQ-\gamma_{ij}^\QQ\\
                \widetilde\omega_{ij,n}^{\PP,\mathrm{harm}}-\omega_{ij}^{\PP,\mathrm{harm}}\\
                \widetilde\omega_{ij,n}^{\QQ,\mathrm{harm}}-\omega_{ij}^{\QQ,\mathrm{harm}}
        \end{pmatrix}
        \Rightarrow
        \mathcal N(0,\widetilde\Omega_{ij}^{\mathrm{info}}),
\]
where
\[
        \widetilde\Omega_{ij}^{\mathrm{info}}
        :=
        \Cov\!\left(
                \begin{pmatrix}
                        \widehat g_{ij,1}^\PP-\gamma_{ij}^\PP\\
                        \widehat g_{ij,1}^\QQ-\gamma_{ij}^\QQ\\
                        (\widehat g_{ij,1}^\PP-\gamma_{ij}^\PP)^2-\omega_{ij}^{\PP,\mathrm{harm}}\\
                        (\widehat g_{ij,1}^\QQ-\gamma_{ij}^\QQ)^2-\omega_{ij}^{\QQ,\mathrm{harm}}
                \end{pmatrix}
        \right).
\]
Consequently Assumption~\ref{ass:premium-info-est} holds with
\[
        \widehat\omega_{ij,n}^{\mu,\mathrm{harm}}
        =
        \widetilde\omega_{ij,n}^{\mu,\mathrm{harm}},
        \qquad
        \Omega_{ij}^{\mathrm{info}}
        =
        \widetilde\Omega_{ij}^{\mathrm{info}}.
\]
\end{proposition}

\begin{proof}
Write
\[
        \bar\omega_{ij,n}^{\mu,\mathrm{harm}}
        :=
        \frac1n
        \sum_{d=1}^n
        \left(
                \widehat g_{ij,d}^\mu-\gamma_{ij}^\mu
        \right)^2.
\]
Then the sample-variance identity gives
\[
        \widetilde\omega_{ij,n}^{\mu,\mathrm{harm}}
        =
        \bar\omega_{ij,n}^{\mu,\mathrm{harm}}
        -
        (\widehat\gamma_{ij,n}^\mu-\gamma_{ij}^\mu)^2.
\]
Since \(\widehat\gamma_{ij,n}^\mu-\gamma_{ij}^\mu=O_\PP(n^{-1/2})\), we have
\[
        \sqrt{n}
        \left(
                \widetilde\omega_{ij,n}^{\mu,\mathrm{harm}}
                -
                \bar\omega_{ij,n}^{\mu,\mathrm{harm}}
        \right)
        =
        -\sqrt{n}\,(\widehat\gamma_{ij,n}^\mu-\gamma_{ij}^\mu)^2
        =
        o_\PP(1).
\]
By the multivariate CLT applied to the i.i.d. vector
\[
        \begin{pmatrix}
                \widehat g_{ij,d}^\PP-\gamma_{ij}^\PP\\
                \widehat g_{ij,d}^\QQ-\gamma_{ij}^\QQ\\
                (\widehat g_{ij,d}^\PP-\gamma_{ij}^\PP)^2-\omega_{ij}^{\PP,\mathrm{harm}}\\
                (\widehat g_{ij,d}^\QQ-\gamma_{ij}^\QQ)^2-\omega_{ij}^{\QQ,\mathrm{harm}}
        \end{pmatrix},
\]
the sample averages built from \((\widehat\gamma_{ij,n}^\PP,\widehat\gamma_{ij,n}^\QQ,
\bar\omega_{ij,n}^{\PP,\mathrm{harm}},\bar\omega_{ij,n}^{\QQ,\mathrm{harm}})\) satisfy the stated
joint CLT.  Replacing \(\bar\omega\) by \(\widetilde\omega\) does not change the limit by the
displayed \(o_\PP(1)\) estimate.
\end{proof}

\begin{theorem}[Operational harmonized projected-information estimator and support classifier]\label{thm:premium-info-estimator}
Under Assumptions~\ref{ass:premium} and~\ref{ass:premium-info-est}, define
\[
        \widehat{\mathcal I}_{ij,n}^{\mu,\mathrm{harm}}
        :=
        \frac{(\widehat\gamma_{ij,n}^\mu)^2}{\widehat\omega_{ij,n}^{\mu,\mathrm{harm}}},
        \qquad
        \widehat\sigma_{ij,n}^\mu
        :=
        \operatorname{sgn}(\widehat\gamma_{ij,n}^\mu),
        \qquad
        \mu\in\{\PP,\QQ\}.
\]
We use the convention \(\operatorname{sgn}(0)=0\).
Then:
\begin{enumerate}[label=(\roman*)]
\item
\[
        \widehat{\mathcal I}_{ij,n}^{\mu,\mathrm{harm}}
        \to
        \mathcal I_{ij}^{\mu,\mathrm{harm}}
        \qquad\text{in probability},
        \qquad
        \mu\in\{\PP,\QQ\}.
\]
\item With
\[
        J_{ij}^{\mathrm{info}}
        :=
        \begin{pmatrix}
                \dfrac{2\gamma_{ij}^\PP}{\omega_{ij}^{\PP,\mathrm{harm}}}
                & 0
                & -\dfrac{(\gamma_{ij}^\PP)^2}{(\omega_{ij}^{\PP,\mathrm{harm}})^2}
                & 0\\[1.1em]
                0
                & \dfrac{2\gamma_{ij}^\QQ}{\omega_{ij}^{\QQ,\mathrm{harm}}}
                & 0
                & -\dfrac{(\gamma_{ij}^\QQ)^2}{(\omega_{ij}^{\QQ,\mathrm{harm}})^2}
        \end{pmatrix},
\]
we have
\[
        \sqrt{n}
        \begin{pmatrix}
                \widehat{\mathcal I}_{ij,n}^{\PP,\mathrm{harm}}-\mathcal I_{ij}^{\PP,\mathrm{harm}}\\
                \widehat{\mathcal I}_{ij,n}^{\QQ,\mathrm{harm}}-\mathcal I_{ij}^{\QQ,\mathrm{harm}}
        \end{pmatrix}
        \Rightarrow
        \mathcal N\!\left(
                0,\,
                J_{ij}^{\mathrm{info}}
                \Omega_{ij}^{\mathrm{info}}
                (J_{ij}^{\mathrm{info}})^\top
        \right).
\]
\item If
\[
        \widehat\Pi_{ij,n}^{\mathrm{info}}
        :=
        \widehat\sigma_{ij,n}^\QQ
        \sqrt{
                \widehat\omega_{ij,n}^{\QQ,\mathrm{harm}}\,
                \widehat{\mathcal I}_{ij,n}^{\QQ,\mathrm{harm}}
        }
        -
        \widehat\sigma_{ij,n}^\PP
        \sqrt{
                \widehat\omega_{ij,n}^{\PP,\mathrm{harm}}\,
                \widehat{\mathcal I}_{ij,n}^{\PP,\mathrm{harm}}
        },
\]
then
\[
        \widehat\Pi_{ij,n}^{\mathrm{info}}
        =
        \widehat\Pi_{ij,n}
        \qquad\text{identically}.
\]
\item If \(\gamma_{ij}^\mu\neq0\), then
\[
        \widehat\sigma_{ij,n}^\mu\to \sigma_{ij}^\mu
        \qquad\text{in probability}.
\]
\item For every deterministic sequence \(\eta_{\Pi,n}\downarrow0\) satisfying
\[
        \sqrt{n}\,\eta_{\Pi,n}\to\infty,
\]
we have
\[
        \Pi_{ij}=0
        \quad\Longrightarrow\quad
        \Pr\!\left(
                |\widehat\Pi_{ij,n}^{\mathrm{info}}|<\eta_{\Pi,n}
        \right)\to1,
\]
and
\[
        \Pi_{ij}\neq0
        \quad\Longrightarrow\quad
        \Pr\!\left(
                |\widehat\Pi_{ij,n}^{\mathrm{info}}|\ge \eta_{\Pi,n}
        \right)\to1.
\]
\end{enumerate}
Thus the exact information-side premium-support laws admit consistent plug-in empirical versions on
the matched cleaned \(\PP/\QQ\) surface.
\end{theorem}

\begin{proof}
Part~(i) follows from the continuous mapping theorem applied to the quotient map
\((x,u)\mapsto x^2/u\), with
\(\widehat\omega_{ij,n}^{\mu,\mathrm{harm}}>0\) by assumption.  Part~(ii) is the multivariate delta
method applied to the map
\[
        (x_P,x_Q,u_P,u_Q)
        \longmapsto
        \left(
                \frac{x_P^2}{u_P},
                \frac{x_Q^2}{u_Q}
        \right),
\]
whose Jacobian at the population point is \(J_{ij}^{\mathrm{info}}\).

For part~(iii), use the definition of
\(\widehat{\mathcal I}_{ij,n}^{\mu,\mathrm{harm}}\) to write
\[
        \widehat\sigma_{ij,n}^\mu
        \sqrt{
                \widehat\omega_{ij,n}^{\mu,\mathrm{harm}}\,
                \widehat{\mathcal I}_{ij,n}^{\mu,\mathrm{harm}}
        }
        =
        \widehat\sigma_{ij,n}^\mu\,
        \sqrt{(\widehat\gamma_{ij,n}^\mu)^2}
        =
        \widehat\gamma_{ij,n}^\mu,
        \qquad
        \mu\in\{\PP,\QQ\},
\]
by the sign convention.  Subtracting the \(\PP\)-term from the \(\QQ\)-term yields
\(\widehat\Pi_{ij,n}^{\mathrm{info}}=\widehat\Pi_{ij,n}\).

For part~(iv), if \(\gamma_{ij}^\mu\neq0\), consistency of \(\widehat\gamma_{ij,n}^\mu\) implies
sign stability, so \(\widehat\sigma_{ij,n}^\mu\to\sigma_{ij}^\mu\) in probability.  For part~(v),
if \(\Pi_{ij}=0\), Proposition~\ref{thm:premium-estimator} and part~(iii) give
\[
        \sqrt{n}\,\widehat\Pi_{ij,n}^{\mathrm{info}}
        =
        \sqrt{n}\,\widehat\Pi_{ij,n}
        =O_\PP(1),
\]
so \(|\widehat\Pi_{ij,n}^{\mathrm{info}}|<\eta_{\Pi,n}\) with probability tending to one because
\(\sqrt{n}\eta_{\Pi,n}\to\infty\).  If \(\Pi_{ij}\neq0\), part~(iii) and consistency of
\(\widehat\Pi_{ij,n}\) yield
\[
        \widehat\Pi_{ij,n}^{\mathrm{info}}
        \to
        \Pi_{ij},
\]
while \(\eta_{\Pi,n}\to0\), so eventually
\(|\widehat\Pi_{ij,n}^{\mathrm{info}}|\ge \eta_{\Pi,n}\) with probability tending to one.
\end{proof}

\begin{remark}[Operational exact information-side support classification]\label{rem:premium-info-estimator}
Theorem~\ref{thm:premium-info-estimator} is the estimator-side counterpart to
Corollaries~\ref{cor:premium-support-info-exact} and
\ref{cor:premium-support-info-exact-common}, while
Corollary~\ref{cor:premium-info-estimator-mixing} is the weak-dependence operational extension.
The population exact-support laws carry more than algebraic content.  Once the
one-dimensional covariance normalizers are estimated on the matched cleaned surface, the same
support conclusions can be implemented empirically either through the coefficient premium
\(\widehat\Pi_{ij,n}\) or through the equivalent projected-information representation
\(\widehat\Pi_{ij,n}^{\mathrm{info}}\), and they remain valid under the matched weak-dependence
extension once the long-run covariance is controlled.
\end{remark}

\begin{corollary}[Weakly dependent matched-date extension of the operational information-side premium theorem]\label{cor:premium-info-estimator-mixing}
Fix an ordered pair \((i,j)\) in the genuinely one-dimensional harmonized regime and define
\[
        Y_{ij,d}
        :=
        \begin{pmatrix}
                \widehat g_{ij,d}^\PP-\gamma_{ij}^\PP\\
                \widehat g_{ij,d}^\QQ-\gamma_{ij}^\QQ\\
                (\widehat g_{ij,d}^\PP-\gamma_{ij}^\PP)^2-\omega_{ij}^{\PP,\mathrm{harm}}\\
                (\widehat g_{ij,d}^\QQ-\gamma_{ij}^\QQ)^2-\omega_{ij}^{\QQ,\mathrm{harm}}
        \end{pmatrix}.
\]
Assume that \((Y_{ij,d})_{d\ge1}\) is strictly stationary, strongly mixing with summable mixing
coefficients, and has finite \((4+\delta)\)-moments for some \(\delta>0\).  Suppose moreover that
the long-run covariance matrix
\[
        \Omega_{ij}^{\mathrm{lr,info}}
        :=
        \sum_{h\in\mathbb Z}\Cov(Y_{ij,0},Y_{ij,h})
\]
is well defined and that the corresponding multivariate central limit theorem holds:
\[
        \sqrt{n}
        \begin{pmatrix}
                \widehat\gamma_{ij,n}^\PP-\gamma_{ij}^\PP\\
                \widehat\gamma_{ij,n}^\QQ-\gamma_{ij}^\QQ\\
                \widetilde\omega_{ij,n}^{\PP,\mathrm{harm}}-\omega_{ij}^{\PP,\mathrm{harm}}\\
                \widetilde\omega_{ij,n}^{\QQ,\mathrm{harm}}-\omega_{ij}^{\QQ,\mathrm{harm}}
        \end{pmatrix}
        \Rightarrow
        \mathcal N(0,\Omega_{ij}^{\mathrm{lr,info}}).
\]
Then the conclusions of Proposition~\ref{prop:premium-info-normalizers} and
Theorem~\ref{thm:premium-info-estimator} continue to hold with
\(\widetilde\Omega_{ij}^{\mathrm{info}}\) and \(\Omega_{ij}^{\mathrm{info}}\) replaced by
\(\Omega_{ij}^{\mathrm{lr,info}}\).  In particular,
\[
        \sqrt{n}
        \begin{pmatrix}
                \widehat{\mathcal I}_{ij,n}^{\PP,\mathrm{harm}}-\mathcal I_{ij}^{\PP,\mathrm{harm}}\\
                \widehat{\mathcal I}_{ij,n}^{\QQ,\mathrm{harm}}-\mathcal I_{ij}^{\QQ,\mathrm{harm}}
        \end{pmatrix}
        \Rightarrow
        \mathcal N\!\left(
                0,\,
                J_{ij}^{\mathrm{info}}
                \Omega_{ij}^{\mathrm{lr,info}}
                (J_{ij}^{\mathrm{info}})^\top
        \right),
\]
and the information-side premium-support classifier based on
\(\widehat\Pi_{ij,n}^{\mathrm{info}}=\widehat\Pi_{ij,n}\) remains consistent for any deterministic
threshold sequence \(\eta_{\Pi,n}\downarrow0\) with \(\sqrt{n}\eta_{\Pi,n}\to\infty\).
\end{corollary}

\begin{proof}
The proof is the same as for Proposition~\ref{prop:premium-info-normalizers} and
Theorem~\ref{thm:premium-info-estimator}, except that the i.i.d. multivariate CLT is replaced by
the assumed long-run covariance CLT for the strictly stationary mixing array.  The sample-variance
identity
\[
        \widetilde\omega_{ij,n}^{\mu,\mathrm{harm}}
        =
        \bar\omega_{ij,n}^{\mu,\mathrm{harm}}
        -
        (\widehat\gamma_{ij,n}^\mu-\gamma_{ij}^\mu)^2
\]
still gives the \(o_\PP(1)\) replacement of \(\bar\omega\) by \(\widetilde\omega\) after
multiplication by \(\sqrt{n}\), and the multivariate delta method applies unchanged.  The identity
\(\widehat\Pi_{ij,n}^{\mathrm{info}}=\widehat\Pi_{ij,n}\) is algebraic, so the support
classification conclusion again follows from the asymptotic law of \(\widehat\Pi_{ij,n}\) and the
threshold condition \(\sqrt{n}\eta_{\Pi,n}\to\infty\).
\end{proof}

\begin{remark}[Long-run covariance estimation for operational premium support]\label{rem:premium-info-hac}
Corollary~\ref{cor:premium-info-estimator-mixing} shows that the operational information-side
premium-support theorem is not tied to an i.i.d. matched-date design.  Under weak dependence, the
same plug-in classifier is valid once the long-run covariance matrix is estimated consistently, for
instance by a standard HAC procedure applied to the matched vector \(Y_{ij,d}\)
\cite{neweyWest1987,andrews1991}.
\end{remark}

\begin{remark}[Preferred operational reading and premium extension boundary]\label{rem:premium-info-default}
For implementation and empirical interpretation,
Corollary~\ref{cor:premium-info-estimator-mixing} is the preferred operational theorem.  It
contains Proposition~\ref{prop:premium-info-normalizers} and
Theorem~\ref{thm:premium-info-estimator} as the special case in which the long-run covariance
reduces to the contemporaneous covariance, while keeping the premium section inside the genuinely
one-dimensional harmonized regime solved here.  Extending the exact information-side support law
beyond that regime, for instance to low-rank or unrestricted matched local geometry under \(\QQ\),
would require new normalizer and sign tracking and is outside the present theorem.
\end{remark}

\begin{proposition}[One-dimensional harmonized projected information is coefficient-equivalent]\label{prop:premium-info-equiv}
Fix an ordered pair \((i,j)\) and suppose the harmonized reduced experiment under
\(\mu\in\{\PP,\QQ\}\) is one-dimensional on the common source-screened basis \(\varphi_{ij}\).
Then, with \(\mathcal I_{ij}^{\mu,\mathrm{harm}}\) defined by
Definition~\ref{def:premium-info},
\[
        \mathcal I_{ij}^{\mu,\mathrm{harm}}=0
        \qquad\Longleftrightarrow\qquad
        \gamma_{ij}^\mu=0
        \qquad\Longleftrightarrow\qquad
        D_{ij}^\mu=0.
\]
If, moreover,
\[
        0<\underline\omega_{ij}^\mu
        \le
        \omega_{ij}^{\mu,\mathrm{harm}}
        \le
        \overline\omega_{ij}^\mu
        <\infty,
\]
then for every \(\kappa_\mu>0\),
\[
        \mathcal I_{ij}^{\mu,\mathrm{harm}}\ge \kappa_\mu
        \quad\Longrightarrow\quad
        |\gamma_{ij}^\mu|
        \ge
        \sqrt{\kappa_\mu\,\underline\omega_{ij}^\mu},
\]
and
\[
        \mathcal I_{ij}^{\mu,\mathrm{harm}}< \kappa_\mu
        \quad\Longrightarrow\quad
        |\gamma_{ij}^\mu|
        <
        \sqrt{\kappa_\mu\,\overline\omega_{ij}^\mu}.
\]
\end{proposition}

\begin{proof}
Since \(\omega_{ij}^{\mu,\mathrm{harm}}>0\),
\[
        \mathcal I_{ij}^{\mu,\mathrm{harm}}=0
        \qquad\Longleftrightarrow\qquad
        \gamma_{ij}^\mu=0.
\]
The equivalence with \(D_{ij}^\mu=0\) is Proposition~\ref{prop:premium-coeff-nonzero}.  The lower
and upper bounds follow from
\[
        (\gamma_{ij}^\mu)^2
        =
        \mathcal I_{ij}^{\mu,\mathrm{harm}}
        \omega_{ij}^{\mu,\mathrm{harm}}
\]
and the assumed bounds on \(\omega_{ij}^{\mu,\mathrm{harm}}\).
\end{proof}

\begin{corollary}[Separated harmonized projected information forces premium support]\label{cor:premium-support-info}
Fix an ordered pair \((i,j)\) and suppose the harmonized reduced experiments under \(\PP\) and
\(\QQ\) are both one-dimensional on the common source-screened basis \(\varphi_{ij}\), with
positive covariance normalizers satisfying
\[
        0<\underline\omega_{ij}^\PP\le \omega_{ij}^{\PP,\mathrm{harm}}\le \overline\omega_{ij}^\PP,
        \qquad
        0<\underline\omega_{ij}^\QQ\le \omega_{ij}^{\QQ,\mathrm{harm}}\le \overline\omega_{ij}^\QQ.
\]
\begin{enumerate}[label=(\roman*)]
\item If
\[
        \mathcal I_{ij}^{\PP,\mathrm{harm}}\ge \kappa_P,
        \qquad
        \mathcal I_{ij}^{\QQ,\mathrm{harm}}<\kappa_Q,
\]
and
\[
        \sqrt{\kappa_P\,\underline\omega_{ij}^\PP}
        >
        \sqrt{\kappa_Q\,\overline\omega_{ij}^\QQ},
\]
then
\[
        |\Pi_{ij}|
        \ge
        \sqrt{\kappa_P\,\underline\omega_{ij}^\PP}
        -
        \sqrt{\kappa_Q\,\overline\omega_{ij}^\QQ}
        >0,
\]
so \((i,j)\in\cE_\Pi\).
\item If
\[
        \mathcal I_{ij}^{\QQ,\mathrm{harm}}\ge \kappa_Q,
        \qquad
        \mathcal I_{ij}^{\PP,\mathrm{harm}}<\kappa_P,
\]
and
\[
        \sqrt{\kappa_Q\,\underline\omega_{ij}^\QQ}
        >
        \sqrt{\kappa_P\,\overline\omega_{ij}^\PP},
\]
then
\[
        |\Pi_{ij}|
        \ge
        \sqrt{\kappa_Q\,\underline\omega_{ij}^\QQ}
        -
        \sqrt{\kappa_P\,\overline\omega_{ij}^\PP}
        >0,
\]
so \((i,j)\in\cE_\Pi\).
\end{enumerate}
In particular, separated projected-information regimes on the harmonized source-screened basis force
nontrivial premium support.
\end{corollary}

\begin{proof}
Proposition~\ref{prop:premium-info-equiv} converts the projected-information bounds into the
coefficient bounds
\[
        |\gamma_{ij}^\PP|
        \ge
        \sqrt{\kappa_P\,\underline\omega_{ij}^\PP},
        \qquad
        |\gamma_{ij}^\QQ|
        <
        \sqrt{\kappa_Q\,\overline\omega_{ij}^\QQ}
\]
in part~(i), and the analogous bounds with \(\PP\) and \(\QQ\) exchanged in part~(ii).  Apply the
reverse triangle inequality
\[
        |\Pi_{ij}|
        \ge
        \bigl||\gamma_{ij}^\QQ|-|\gamma_{ij}^\PP|\bigr|
\]
to conclude.
\end{proof}

\begin{remark}[Relation to the operational \(\PP/\QQ\) classification template]\label{rem:premium-support-info}
Corollary~\ref{cor:premium-support-info} is the premium-support analogue of the operational
\(\PP/\QQ\) classification template in \cite{vidal2026econometric}.  In the genuinely
one-dimensional harmonized regime, a path-detectable-only or pricing-visible-only information
pattern on the common source-screened basis is no longer only a qualitative label.  It forces
membership in the premium support once the one-dimensional covariance normalizers are controlled.
\end{remark}

\begin{proposition}[Exact premium reconstruction from harmonized projected information]\label{prop:premium-info-reconstruct}
Fix an ordered pair \((i,j)\) and suppose the harmonized reduced experiments under \(\PP\) and
\(\QQ\) are both one-dimensional on the common source-screened basis \(\varphi_{ij}\).
Define the sign indicators
\[
        \sigma_{ij}^\mu
        :=
        \begin{cases}
                1,& \gamma_{ij}^\mu>0,\\
                0,& \gamma_{ij}^\mu=0,\\
                -1,& \gamma_{ij}^\mu<0,
        \end{cases}
        \qquad
        \mu\in\{\PP,\QQ\}.
\]
Then
\[
        \gamma_{ij}^\mu
        =
        \sigma_{ij}^\mu\,
        \sqrt{
                \omega_{ij}^{\mu,\mathrm{harm}}\,
                \mathcal I_{ij}^{\mu,\mathrm{harm}}
        },
        \qquad
        \mu\in\{\PP,\QQ\},
\]
and therefore
\[
        \Pi_{ij}
        =
        \sigma_{ij}^\QQ
        \sqrt{
                \omega_{ij}^{\QQ,\mathrm{harm}}\,
                \mathcal I_{ij}^{\QQ,\mathrm{harm}}
        }
        -
        \sigma_{ij}^\PP
        \sqrt{
                \omega_{ij}^{\PP,\mathrm{harm}}\,
                \mathcal I_{ij}^{\PP,\mathrm{harm}}
        }.
\]
\end{proposition}

\begin{proof}
By Definition~\ref{def:premium-info},
\[
        (\gamma_{ij}^\mu)^2
        =
        \omega_{ij}^{\mu,\mathrm{harm}}
        \mathcal I_{ij}^{\mu,\mathrm{harm}}.
\]
Multiplying the positive square root by the sign indicator \(\sigma_{ij}^\mu\) recovers
\(\gamma_{ij}^\mu\).  Substituting that representation into
\(\Pi_{ij}=\gamma_{ij}^\QQ-\gamma_{ij}^\PP\) gives the displayed formula.
\end{proof}

\begin{corollary}[Exact premium-support law from harmonized projected information and normalizers]\label{cor:premium-support-info-exact}
Under the hypotheses of Proposition~\ref{prop:premium-info-reconstruct},
\[
        \Pi_{ij}=0
        \qquad\Longleftrightarrow\qquad
        \sigma_{ij}^\QQ
        \sqrt{
                \omega_{ij}^{\QQ,\mathrm{harm}}\,
                \mathcal I_{ij}^{\QQ,\mathrm{harm}}
        }
        =
        \sigma_{ij}^\PP
        \sqrt{
                \omega_{ij}^{\PP,\mathrm{harm}}\,
                \mathcal I_{ij}^{\PP,\mathrm{harm}}
        }.
\]
Equivalently,
\[
        (i,j)\in\cE_\Pi
        \qquad\Longleftrightarrow\qquad
        \sigma_{ij}^\QQ
        \sqrt{
                \omega_{ij}^{\QQ,\mathrm{harm}}\,
                \mathcal I_{ij}^{\QQ,\mathrm{harm}}
        }
        \neq
        \sigma_{ij}^\PP
        \sqrt{
                \omega_{ij}^{\PP,\mathrm{harm}}\,
                \mathcal I_{ij}^{\PP,\mathrm{harm}}
        }.
\]
In particular, if \(\sigma_{ij}^\QQ=-\sigma_{ij}^\PP\) and
\[
        \mathcal I_{ij}^{\QQ,\mathrm{harm}}
        +
        \mathcal I_{ij}^{\PP,\mathrm{harm}}
        >
        0,
\]
then \((i,j)\in\cE_\Pi\).
\end{corollary}

\begin{proof}
The equivalence is exactly the identity in Proposition~\ref{prop:premium-info-reconstruct}.  If
\(\sigma_{ij}^\QQ=-\sigma_{ij}^\PP\) and at least one information term is positive, the two signed
square-root terms cannot cancel.
\end{proof}

\begin{corollary}[Common-normalizer exact premium-support law from harmonized projected information]\label{cor:premium-support-info-exact-common}
Under the hypotheses of Proposition~\ref{prop:premium-info-reconstruct}, assume in addition that
\[
        \omega_{ij}^{\PP,\mathrm{harm}}
        =
        \omega_{ij}^{\QQ,\mathrm{harm}}
        =:
        \omega_{ij}^{\mathrm{harm}}
        >
        0.
\]
Then:
\begin{enumerate}[label=(\roman*)]
\item if \(\sigma_{ij}^\QQ=\sigma_{ij}^\PP\), then
\[
        |\Pi_{ij}|
        =
        \sqrt{\omega_{ij}^{\mathrm{harm}}}\,
        \left|
                \sqrt{\mathcal I_{ij}^{\QQ,\mathrm{harm}}}
                -
                \sqrt{\mathcal I_{ij}^{\PP,\mathrm{harm}}}
        \right|,
\]
so
\[
        \Pi_{ij}=0
        \qquad\Longleftrightarrow\qquad
        \mathcal I_{ij}^{\QQ,\mathrm{harm}}
        =
        \mathcal I_{ij}^{\PP,\mathrm{harm}}.
\]
Consequently,
\[
        (i,j)\in\cE_\Pi
        \qquad\Longleftrightarrow\qquad
        \mathcal I_{ij}^{\QQ,\mathrm{harm}}
        \neq
        \mathcal I_{ij}^{\PP,\mathrm{harm}}.
\]
\item if \(\sigma_{ij}^\QQ=-\sigma_{ij}^\PP\) and
\[
        \mathcal I_{ij}^{\QQ,\mathrm{harm}}
        +
        \mathcal I_{ij}^{\PP,\mathrm{harm}}
        >
        0,
\]
then
\[
        |\Pi_{ij}|
        =
        \sqrt{\omega_{ij}^{\mathrm{harm}}}
        \left(
                \sqrt{\mathcal I_{ij}^{\QQ,\mathrm{harm}}}
                +
                \sqrt{\mathcal I_{ij}^{\PP,\mathrm{harm}}}
        \right)
        >
        0,
\]
so \((i,j)\in\cE_\Pi\).
\end{enumerate}
\end{corollary}

\begin{proof}
Under the common-normalizer hypothesis,
Corollary~\ref{cor:premium-support-info-exact} simplifies to the displayed formulas.  If
\(\sigma_{ij}^\QQ=\sigma_{ij}^\PP\), factor out the common sign and take absolute values.  If
\(\sigma_{ij}^\QQ=-\sigma_{ij}^\PP\), the two signed square-root terms add in absolute value.
\end{proof}

\begin{remark}[From support inclusion to an exact information-side support law]\label{rem:premium-support-info-exact}
Corollary~\ref{cor:premium-support-info} is the threshold separation theorem: a strict
projected-information gap forces premium support without requiring exact coefficient matching.
Corollary~\ref{cor:premium-support-info-exact} is the exact law in the genuinely one-dimensional
harmonized regime once the covariance normalizers and coefficient signs are tracked explicitly, and
Corollary~\ref{cor:premium-support-info-exact-common} is the common-normalizer reduction: there,
premium support is characterized exactly by projected-information mismatch
on the common basis when the two projected coefficients have the same sign, and sign opposition
forces support whenever one side is nondegenerate.
\end{remark}

\begin{remark}[What the premium support theorem does and does not identify]\label{rem:premium-support-scope}
Section~\ref{sec:premium} is therefore no longer only a fixed-pair estimator theorem.  Definition~
\ref{def:premium-support}, Proposition~\ref{prop:premium-support-separated}, and
Corollaries~\ref{cor:premium-support-structural}, \ref{cor:premium-support-exact}, and
\ref{cor:premium-support-info} identify structural subsets of the pairwise premium support directly
from the \(\PP/\QQ\) visibility separation on the harmonized source-screened basis, while
Corollaries~\ref{cor:premium-support-info-exact} and
\ref{cor:premium-support-info-exact-common} give genuine if-and-only-if support laws in the
one-dimensional harmonized regime once the covariance normalizers and coefficient signs are
tracked, with a reduced expression in the common-normalizer case.  What is \emph{not} claimed
here is a full unrestricted if-and-only-if description of \(\cE_\Pi\) without those extra
one-dimensional normalization and sign primitives.  Such a characterization would require a
complete matched \(\QQ\)-side pairwise local geometry on the same harmonized basis, beyond the
support inclusion supplied by separated pricing/path visibility or the exact information-side law
available in the normalizer-controlled one-dimensional case.
\end{remark}

\begin{assumption}[Finite-dimensional matched harmonized premium experiment]\label{ass:fd-premium}
For a fixed ordered pair \((i,j)\), let
\[
        \Phi_{ij}=(\varphi_{ij,1},\dots,\varphi_{ij,r_{ij}})
\]
be a common source-screened basis used under both \(\PP\) and \(\QQ\).  The reduced projected
structural-loading representatives have coefficient vectors
\[
        \gamma_{ij}^\PP=(\gamma_{ij,1}^\PP,\dots,\gamma_{ij,r_{ij}}^\PP)^\top,
        \qquad
        \gamma_{ij}^\QQ=(\gamma_{ij,1}^\QQ,\dots,\gamma_{ij,r_{ij}}^\QQ)^\top.
\]
The finite-dimensional pairwise premium vector is
\[
        \Pi_{ij}^{\mathrm{vec}}
        :=
        \gamma_{ij}^\QQ-\gamma_{ij}^\PP.
\]
Let \(W_{ij}^\mu\) be the positive-definite covariance normalizer on the same basis under
\(\mu\in\{\PP,\QQ\}\).  The signed normalized coefficient vector is
\[
        \theta_{ij}^\mu:=(W_{ij}^\mu)^{-1/2}\gamma_{ij}^\mu.
\]
\end{assumption}

\begin{definition}[Finite-dimensional premium support]\label{def:fd-premium-support}
Under Assumption~\ref{ass:fd-premium}, define
\[
        (i,j)\in\cE_\Pi^{\mathrm{vec}}
        \qquad\Longleftrightarrow\qquad
        \Pi_{ij}^{\mathrm{vec}}\ne0.
\]
At level \(\eta\ge0\), define
\[
        (i,j)\in\cE_\Pi^{\mathrm{vec}}(\eta)
        \qquad\Longleftrightarrow\qquad
        \|\Pi_{ij}^{\mathrm{vec}}\|_2\ge\eta.
\]
\end{definition}

\begin{theorem}[Finite-dimensional if-and-only-if premium support law]\label{thm:fd-premium-iff}
Under Assumption~\ref{ass:fd-premium},
\[
        (i,j)\in\cE_\Pi^{\mathrm{vec}}
        \qquad\Longleftrightarrow\qquad
        \gamma_{ij}^\QQ\ne\gamma_{ij}^\PP.
\]
Equivalently, in signed normalized information coordinates,
\[
        (i,j)\in\cE_\Pi^{\mathrm{vec}}
        \qquad\Longleftrightarrow\qquad
        (W_{ij}^\QQ)^{1/2}\theta_{ij}^\QQ
        \ne
        (W_{ij}^\PP)^{1/2}\theta_{ij}^\PP.
\]
Moreover,
\[
        \|\Pi_{ij}^{\mathrm{vec}}\|_2^2
        =
        \|\gamma_{ij}^\QQ\|_2^2
        +
        \|\gamma_{ij}^\PP\|_2^2
        -
        2\langle \gamma_{ij}^\QQ,\gamma_{ij}^\PP\rangle.
\]
Thus a complete finite-dimensional support decision requires the relative orientation of the two
coefficient vectors as well as their separate squared lengths.
\end{theorem}

\begin{proof}
The first equivalence is the definition of \(\Pi_{ij}^{\mathrm{vec}}\).  The normalized-coordinate
form follows from \(\gamma_{ij}^\mu=(W_{ij}^\mu)^{1/2}\theta_{ij}^\mu\).  The norm identity is the
polarization identity in \(\R^{r_{ij}}\).
\end{proof}

\begin{corollary}[Operational finite-dimensional premium estimator]\label{cor:fd-premium-estimator}
Suppose estimators satisfy the joint central limit theorem
\[
        \sqrt n
        \begin{pmatrix}
                \widehat\gamma_{ij,n}^\PP-\gamma_{ij}^\PP\\
                \widehat\gamma_{ij,n}^\QQ-\gamma_{ij}^\QQ
        \end{pmatrix}
        \Rightarrow
        N(0,\Sigma_{ij}),
\]
where \(\Sigma_{ij}\) is a \(2r_{ij}\times2r_{ij}\) covariance matrix.  Then
\[
        \widehat\Pi_{ij,n}^{\mathrm{vec}}
        :=
        \widehat\gamma_{ij,n}^\QQ-
        \widehat\gamma_{ij,n}^\PP
\]
is consistent and
\[
        \sqrt n(\widehat\Pi_{ij,n}^{\mathrm{vec}}-\Pi_{ij}^{\mathrm{vec}})
        \Rightarrow
        N(0,D\Sigma_{ij}D^\top),
        \qquad
        D=(-I_{r_{ij}},I_{r_{ij}}).
\]
If \(\Pi_{ij}^{\mathrm{vec}}\ne0\), then
\[
        \sqrt n
        \left(
                \|\widehat\Pi_{ij,n}^{\mathrm{vec}}\|_2
                -
                \|\Pi_{ij}^{\mathrm{vec}}\|_2
        \right)
        \Rightarrow
        N\!\left(0,
                \frac{(\Pi_{ij}^{\mathrm{vec}})^\top D\Sigma_{ij}D^\top\Pi_{ij}^{\mathrm{vec}}}
                {\|\Pi_{ij}^{\mathrm{vec}}\|_2^2}
        \right).
\]
\end{corollary}

\begin{proof}
The vector premium estimator is a linear map of the joint estimator, so the first limit follows
from the multivariate delta method with derivative \(D\).  Away from zero, the Euclidean norm is
differentiable with gradient \(x/\|x\|_2\), giving the second limit.
\end{proof}

\begin{proposition}[Unsigned information cannot give unrestricted support]\label{prop:premium-nonidentification}
No unrestricted if-and-only-if premium-support law can be written using only the two unsigned
projected information levels
\[
        \|\theta_{ij}^\PP\|_2^2,
        \qquad
        \|\theta_{ij}^\QQ\|_2^2,
\]
even when the normalizers are known.
\end{proposition}

\begin{proof}
It is enough to take \(r_{ij}=1\) and common normalizer \(W=1\).  The two coefficient pairs
\[
        (\gamma^\PP,\gamma^\QQ)=(1,1),
        \qquad
        (\gamma^\PP,\gamma^\QQ)=(1,-1)
\]
have the same unsigned information levels under \(\PP\) and \(\QQ\), namely \(1\) and \(1\).  In
the first case \(\Pi=0\); in the second case \(\Pi=-2\).  Hence unsigned information alone cannot
distinguish no support from nontrivial support.  In higher dimension the same obstruction appears
as an unknown relative angle between \(\gamma^\PP\) and \(\gamma^\QQ\).
\end{proof}

\begin{remark}[Finite-dimensional premium scope]\label{rem:fd-premium-scope}
The one-dimensional exact law above is the first identifiable case using scalar projected
information plus normalizers and signs.  Theorem~\ref{thm:fd-premium-iff} is the natural
unrestricted statement once the full signed matched coefficient vector is part of the data
contract.  Proposition~\ref{prop:premium-nonidentification} explains why a purely
information-magnitude theorem cannot be unrestricted.
\end{remark}

\begin{remark}[Implementation data contract]
The theorem above is a compression theorem with an explicit data contract.  To instantiate it one
needs common observation dates, a maturity grid, a strike or moneyness grid, an
arbitrage-regularized option surface on each date, a local derivative extraction map on that cleaned
surface, and the same harmonized source-screened projection basis as on the realized-volatility
side.  Market microstructure is outside the theorem.  The matched \(\PP/\QQ\) comparison is therefore specified as a
data contract rather than only as an abstract possibility.
\end{remark}



\section{Detailed scope conditions and endpoint obstructions}
\label{sec:detailed-scope-gates}

The endpoint theorems are conditional.  Each endpoint requires its own regularity, model-class, or
data-alignment assumptions.  The material below records how those assumptions are used.  It adds no
primitive model.  Its role is to state when risk-neutral compression can be converted into a
third-cumulant skew law, a Hermite price law, a jump-kernel price law, a local Black surface law, a
signed premium comparison, or a nonlinear variance-linked Hermite hierarchy.

The assumptions have four roles.  The Perron reduction assumptions import the synchronized state
and the Volterra market geometry; in the empirical section that reduction is rebuilt from raw
Oxford--Man data.  The risk-neutral regularity assumptions make state-price-density changes,
\(L^2(\QQ)\) projections, and option first variations well-defined.  The option-transfer
assumptions specify which surface coordinate is read from the compressed loading: a cumulant-class
skew, a Hermite price coefficient, an exact jump-kernel coefficient, or a local
Black-implied-volatility coefficient.  The premium, variance, and hedging assumptions state how the
same loading is compared on a signed \(\PP/\QQ\) basis or passed through nonlinear and
projection-residual functionals.

The Gaussian Volterra branch has a separate role.  Assumptions~\ref{ass:gaussian-scalar-seam-package}
and~\ref{ass:gaussian-lambda-public-frontier} are not empirical calibration knobs; they connect the
body to Companion~A.  Companion~A proves the raw Volterra estimates and source-family verification
results.  The article uses only the scalar seam, the common-channel reduction, and the Gaussian
option-factor consequence.  This separation keeps the body focused on risk-neutral compression and
the option-surface diagnostic.

\subsection{Non-Gaussian option pricing: cumulant scope and obstructions}
\label{subsec:ng-option-obstructions}

The non-Gaussian option branch has two levels.  The abstract compression theorem is
model-agnostic.  The scalar third-cumulant formula is not.  The latter is the correct endpoint only
when the Hermite and jump conditions leave no leading term outside the third-cumulant channel.
This keeps the skew language close to the option skew-law literature \cite{bakshi2003} and makes
explicit the extra compression and quote-transfer conditions needed here.

\begin{definition}[Hermite ATM-skew defect]
\label{def:hermite-skew-defect}
Under Assumption~\ref{ass:hermite-admissible} with \(r_\tau=0\), define
\[
        \mathfrak h_\tau
        :=
        \frac{1}{\sqrt\tau\sqrt{v_\tau}}
        \sum_{\substack{m\ge4\\ m\le M_\tau}}
        (m-1)a_{m,\tau}\He_{m-2}(0).
\]
Then Theorem~\ref{thm:hermite-option-transfer} says that
\[
        \partial_k\partial_\varepsilon\sigma_N(\tau,k;\varepsilon)
        \big|_{k=0,\varepsilon=0}
        =
        \frac{\partial_\varepsilon\kappa_3^\QQ(R_\tau^\varepsilon)|_{\varepsilon=0}}
        {6v_\tau^{3/2}\sqrt\tau}
        +\mathfrak h_\tau.
\]
The cumulant transfer is exact when \(\mathfrak h_\tau=0\), and it is asymptotically exact at scale
\(b_\tau\) when \(\mathfrak h_\tau=o(b_\tau)\).
\end{definition}

\begin{proposition}[Same third cumulant, different ATM skew]
\label{prop:hermite-same-cumulant-different-skew}
Fix \(v_\tau>0\) and let \(Y_\tau=X_\tau/\sqrt{v_\tau}\sim N(0,1)\).  For any \(a_2\in\R\) and
any nonzero \(a_4\in\R\), the two conditional projections
\[
        \EE_\QQ[Z_\tau^{(1)}\mid X_\tau]=a_2\He_2(Y_\tau),
        \qquad
        \EE_\QQ[Z_\tau^{(2)}\mid X_\tau]=a_2\He_2(Y_\tau)+a_4\He_4(Y_\tau)
\]
have the same third-cumulant derivative but different Bachelier ATM-skew derivatives.
Consequently the third cumulant alone cannot determine non-Gaussian ATM skew even in the finite
Hermite-admissible class.
\end{proposition}

\begin{proof}
By Theorem~\ref{thm:hermite-option-transfer}, the third-cumulant derivative is \(6v_\tau a_2\) for
both perturbations, since it depends only on the \(\He_2\)-coefficient.  The ATM-skew derivative
differs by
\[
        \frac{3a_4\He_2(0)}{\sqrt\tau\sqrt{v_\tau}}
        =
        -\frac{3a_4}{\sqrt\tau\sqrt{v_\tau}},
\]
because \(\He_2(0)=-1\).  This is nonzero when \(a_4\ne0\).
\end{proof}

\begin{definition}[Cumulant-valid continuous perturbations]
\label{def:cumulant-pass-diagnostic}
Suppose the cumulant term in Definition~\ref{def:hermite-skew-defect} has selected scale
\(b_\tau\).  The continuous non-Gaussian perturbation is cumulant-valid at that scale if
\[
        \mathfrak h_\tau=o(b_\tau).
\]
If \(\mathfrak h_\tau\) is of the same or larger order, the cumulant reduction is not valid at the
selected scale.  In that
case the option-price endpoint is still explicit, but the selected coefficient is the Hermite sum,
not the third cumulant alone.
\end{definition}

\begin{definition}[Jump tail-kernel defect]
\label{def:jump-tail-kernel-defect}
For a jump intensity measure \(\mu_\tau\) in Theorem~\ref{thm:exact-jump-price}, set
\[
        K_{v_\tau}(y)
        :=
        -\Phi(y/\sqrt{v_\tau})+\frac12+\frac{y}{\sqrt{v_\tau}}\phi(0),
\]
and define the exact jump ATM-skew contribution, up to the common factor
\((\sqrt\tau\phi(0))^{-1}\), by
\[
        \mathfrak j_\tau
        :=
        \int K_{v_\tau}(y)\,\mu_\tau(\dd y).
\]
The small-jump cumulant approximation replaces \(K_{v_\tau}(y)\) by
\(\phi(0)y^3/(6v_\tau^{3/2})\).  Its residual is
\[
        \mathfrak r_\tau^J
        :=
        \int
        \left(
                K_{v_\tau}(y)-\frac{\phi(0)}{6v_\tau^{3/2}}y^3
        \right)
        \mu_\tau(\dd y).
\]
The jump branch is cumulant-valid when \(\mathfrak r_\tau^J\) is lower order at
the selected scale.  Otherwise the exact coefficient is \(\mathfrak j_\tau\), not the third
cumulant.
\end{definition}

\begin{remark}[Correct non-Gaussian formulation]
\label{rem:correct-ng-headline}
The formulation is therefore not ``non-Gaussian skew is always third-cumulant skew.''  Rather,
the third cumulant is the option-price coefficient on the cumulant-class branch, i.e.
when higher even Hermite projections and large-jump tail-kernel effects are absent or lower order.
Outside that branch there is still an option pricing statement, but the coefficient is the
visible Hermite or jump-kernel coefficient.
\end{remark}

\subsection{Premium recovery: what signed vectors identify and unsigned scalars do not}
\label{subsec:premium-recovery-scope}

The finite-dimensional premium theorem is a genuine strengthening because it moves from a scalar
support statement to a signed matched vector statement.  The price of that strengthening is a data
contract: without the signed relative orientation of the \(\PP\)- and \(\QQ\)-coefficient vectors,
there is no unrestricted support law.

\begin{proposition}[Sharp unsigned premium interval]
\label{prop:unsigned-premium-interval}
Let \(u,v\in\R^r\) be two possible signed coefficient vectors and suppose that only their lengths
\(p:=\|u\|_2\) and \(q:=\|v\|_2\) are known.  Then every possible premium magnitude
\(\|v-u\|_2\) lies in
\[
        [|p-q|,p+q].
\]
If \(r\ge2\), every value in this interval is attained by a suitable relative angle between \(u\)
and \(v\).  If \(r=1\), the only possible values are \(|p-q|\) and \(p+q\), corresponding to the
two relative signs.
\end{proposition}

\begin{proof}
The bounds are the reverse triangle inequality and the triangle inequality.  For \(r\ge2\), choose
unit vectors \(e_1,e_2\) and write
\[
        u=pe_1,
        \qquad
        v=q(\cos\theta\,e_1+\sin\theta\,e_2).
\]
Then
\[
        \|v-u\|_2^2=p^2+q^2-2pq\cos\theta,
\]
which ranges continuously from \((p-q)^2\) to \((p+q)^2\).  In dimension one only
\(\cos\theta\in\{-1,1\}\) is available.
\end{proof}

\begin{corollary}[Unsigned information is one-sided]
\label{cor:unsigned-one-sided-premium}
In the finite-dimensional premium experiment, unequal known signed-vector lengths can force
premium support, because \(\|\gamma_{ij}^\QQ-\gamma_{ij}^\PP\|_2\ge
\bigl|\|\gamma_{ij}^\QQ\|_2-\|\gamma_{ij}^\PP\|_2\bigr|\).  Equal lengths do not rule out premium
support.  Thus unsigned projected information can give sufficient separation tests, but not an
unrestricted if-and-only-if support law.
\end{corollary}

\begin{proof}
The lower bound is Proposition~\ref{prop:unsigned-premium-interval}.  When the two lengths are
equal and positive, the same proposition permits both zero premium, obtained by equal signed
vectors, and nonzero premium, obtained by rotating or changing sign.  This is exactly the
nonidentification mechanism of Proposition~\ref{prop:premium-nonidentification}.
\end{proof}

\begin{remark}[Premium data contract]
\label{rem:premium-data-contract-developed}
The unrestricted finite-dimensional support theorem therefore requires the signed matched vector
\(\Pi_{ij}^{\mathrm{vec}}=\gamma_{ij}^\QQ-\gamma_{ij}^\PP\).  Scalar or unsigned information may
be used for conservative screening, but the full support decision requires the relative
orientation term
\[
        \langle\gamma_{ij}^\QQ,\gamma_{ij}^\PP\rangle
\]
in the polarization identity of Theorem~\ref{thm:fd-premium-iff}.  This orientation term is the
exact obstruction; removability would require additional structure.
\end{remark}

\subsection{Nonlinear variance branch and scalar-factor scope}
\label{subsec:variance-hierarchy-scope}

The nonlinear variance-linked branch is a nonlinear theorem with model-class scope.  The
exact object under Assumption~\ref{ass:cg-onefactor} is the Hermite coefficient sequence of
Proposition~\ref{thm:nonlinear-hermite}.  A scalar one-factor representation is a special collapse
of that sequence rather than the generic law.

\begin{definition}[Exact visible-beta one-factor collapse]
\label{def:visible-beta-collapse}
Under Assumption~\ref{ass:cg-onefactor}, the nonlinear claim \(X=\Psi(Q_M)\) has an exact
visible-beta one-factor collapse if there exists a \(\cG\)-measurable random variable
\(b(V,S)\in L^2(\QQ)\) such that
\[
        L_X^\QQ-C_\cG^\QQ(L_X^\QQ)=b(V,S)Z.
\]
This is weaker than a deterministic-beta scalar law, because the coefficient may depend on visible
state variables.  It is still stronger than the general Hermite hierarchy.
\end{definition}

\begin{proposition}[Hermite criterion for exact one-factor collapse]
\label{prop:hermite-one-factor-criterion}
Under Assumption~\ref{ass:cg-onefactor}, the exact visible-beta one-factor collapse in
Definition~\ref{def:visible-beta-collapse} holds if and only if
\[
        c_m(V,S)=0
        \qquad
        \text{for every }m\ge2
\]
for the Hermite coefficients of Proposition~\ref{thm:nonlinear-hermite}.  Equivalently, for almost
every realized visible state \((V,S)\), the function
\[
        z\longmapsto P_S\Psi(V+\alpha z)
\]
is affine as an element of \(L^2\) under the standard normal law.  In that case
\(b(V,S)=c_1(V,S)\).
\end{proposition}

\begin{proof}
Proposition~\ref{thm:nonlinear-hermite} gives
\[
        L_X^\QQ-C_\cG^\QQ(L_X^\QQ)=\sum_{m\ge1}c_m(V,S)\He_m(Z).
\]
The Hermite polynomials are an orthogonal basis of \(L^2\) under the standard normal law.  Hence
the right-hand side belongs to the span of \(\He_1(Z)=Z\) if and only if all coefficients of
\(\He_m\), \(m\ge2\), vanish.  Completeness of the Hermite basis gives the equivalent affine
statement for the conditional function of \(z\).
\end{proof}

\begin{corollary}[Why the affine branch is the exact scalar branch]
\label{cor:affine-branch-exact-scalar}
If \(S=0\), \(\alpha>0\), and \(\Psi\) is non-affine on the support reached by \(V+\alpha Z\),
then the exact visible-beta one-factor collapse fails.  The leading first-factor approximation of
Corollary~\ref{cor:visible-beta-law} may still hold for small \(\alpha\), but the exact law is the
Hermite hierarchy.
\end{corollary}

\begin{proof}
When \(S=0\), the conditional function is \(z\mapsto \Psi(V+\alpha z)\).  If it were affine in
\(z\) for almost every visible state, \(\Psi\) would be affine on the corresponding translated
support intervals.  This is the obstruction isolated in Corollary~\ref{cor:affine-only}.  If
\(\Psi\) is not affine there, at least one higher Hermite coefficient is nonzero by
Proposition~\ref{prop:hermite-one-factor-criterion}.
\end{proof}

\begin{remark}[Variance-branch scope]
\label{rem:variance-branch-scope-developed}
The conditionally Gaussian theorem is therefore a hierarchy theorem with a first-factor
asymptotic, not a universal exact one-factor scalar theorem.  Outside
Assumption~\ref{ass:cg-onefactor}, the valid statement is the general compression and
hedging residual statement.  Inside that assumption, the exact nonlinear object is the full
sequence \((c_m(V,S))_{m\ge0}\).
\end{remark}

\subsection{Local Black surface transfer}
\label{subsec:black-local-scope}

The Black-implied-volatility conversion is local.  It preserves visibility on a
regular window where the price-to-volatility inverse is smooth and uniformly controlled.  It does
not claim a global implied-volatility surface theorem.  This is the conservative local analogue of
the Black and local-surface conventions \cite{black1976,dupire1994}.

\begin{definition}[Black-window regularity condition]
\label{def:black-window-certificate}
A Black-window regularity condition at scale \(a_\tau\) consists of
\[
        (\cD,m,a_\tau,\underline v_B,K_{m+1}),
\]
where \(\cD\subset(0,T]\times\R\) is compact, \(m\ge0\),
\[
        \inf_{(\tau,k)\in\cD}V_B(\tau,k)\ge \underline v_B>0,
\]
the inverse Black map has \(C^{m+1}\)-norm bounded by \(K_{m+1}\) on a neighborhood of
\(C^0(\cD)\), and the price expansion satisfies
\[
        C^\varepsilon=C^0+\varepsilon\dot C+\rho^\varepsilon,
        \qquad
        \|\rho^\varepsilon\|_{C^m(\cD)}=o(\varepsilon a_\tau).
\]
This condition is exactly the operational content of Assumption~\ref{ass:black-window}.
\end{definition}

\begin{proposition}[Why the Black theorem is not global]
\label{prop:black-not-global}
A scale-uniform price remainder does not imply a scale-uniform implied-volatility remainder without
a vega lower bound.  In particular, if along a sequence of points the Black vega satisfies
\(V_{B,n}\downarrow0\), then there are price remainders \(r_n=o(a_n)\) for which
\[
        \frac{r_n}{V_{B,n}}
\]
is not \(o(a_n)\).  Thus the local vega condition in Assumption~\ref{ass:black-window} is the
condition that makes the Taylor transfer scale-uniform, rather than a cosmetic regularity
assumption.
\end{proposition}

\begin{proof}
Take \(V_{B,n}=1/n\) and \(r_n=a_n/\sqrt n\).  Then \(r_n=o(a_n)\), but
\[
        \frac{r_n}{V_{B,n}}=\sqrt n\,a_n,
\]
which is not \(o(a_n)\).  Since first-order implied-volatility perturbations divide price
perturbations by vega, a lower bound on vega is necessary for a uniform price remainder to remain
lower order after inversion.
\end{proof}

\begin{remark}[Black-conversion scope]
\label{rem:black-conversion-scope-developed}
The correct Black statement is local: on a regular window satisfying Assumption~\ref{ass:black-window}, Black and normal visibility
agree for the surface jets whose linear conversion coefficient does not vanish.  The theorem does
not extend across maturity--strike regions where vega degenerates, the inverse map loses uniform
smoothness, or the Taylor remainder is not controlled at the selected scale.
\end{remark}

\subsection{Summary of scope}

\begin{remark}[Net effect of the scope analysis]
\label{rem:net-effect-weak-points}
The core structure remains risk-neutral compression of priced structural loading.  The surrounding
assumptions state which additional regularity, basis, tail, or model-class conditions are needed to
express that structure as a cumulant skew law, a local Black-implied-volatility law, a
finite-dimensional premium support law, or a nonlinear variance-linked Hermite law.  When one of
those conditions is unavailable, the compression theorem itself remains in force, but the
corresponding endpoint conversion should be replaced by the more general projection or price-kernel
statement proved earlier.
\end{remark}



\section{Market-scale short-maturity skew asymptotics}\label{sec:skew}

The structural option statement is Theorem~\ref{thm:option-abstract}; the explicit public skew law
is Theorem~\ref{thm:market-skew}.  This section verifies the non-Gaussian cumulant-class hypotheses
in the Gaussian-Volterra/Bachelier market-index model from Section~\ref{sec:qsetup} and computes the
aggregated coefficients in closed form.  At factor level, the Gaussian block reduces option
alignment to the centered structural loading of the perturbation return \(Z_\tau\) and the centered
quadratic-fluctuation remainder generated by \(X_\tau^2-v_\tau\).

The internal ladder is proof infrastructure.  The centered option-facing loading is first reduced
to a canonical Perron-generated proxy claim, then to a benchmark-profile claim built from the
leading Perron kernels, then to the benchmark-generated centered \(\varpi\)-profile generator, and
only at the last step to the aggregate factor benchmark \(\Lambda_M^\star\).  The consequence for
this article is the Companion~A source-family verification.  Its role is to prove the scalar seam,
the common-channel reduction, and the option-factor theorem, not to introduce additional endpoints.
The centered mixed visible/\(\varpi\) remainder and the visible-screening defect are fixed
structural obstructions.  The off-Perron and quadratic terms remain the maturity-scale defects.

\begin{definition}[Gaussian specialization of the option first-variation pricing object]\label{def:gaussian-option-object}
Under Assumption~\ref{ass:q-index}, define
\[
        \Xi_{M,\tau}
        :=
        \frac{1}{2(\sigma_0^M(0))^3\tau^2}\,X_\tau^2 Z_\tau.
\]
This is the Gaussian-Volterra specialization of the abstract option first-variation pricing object
from Assumption~\ref{ass:option-fv}; its priced structural loading is the corresponding
\(L_{M,\tau}^\QQ=C_{\cH_\varpi}^\QQ(\Xi_{M,\tau})\) from
Definition~\ref{def:option-loading}.
\end{definition}

\begin{lemma}[Market leverage covariance asymptotics]\label{lem:market-cov}
Under Assumption~\ref{ass:q-index}, if \(\cE_M^{\mathrm{ret}}=\varnothing\), then
\[
        \Cov_\QQ(Y_M(t),B_t^\QQ)=0
        \qquad
        \text{for every }t\in[0,T].
\]
If \(\cE_M^{\mathrm{ret}}\neq\varnothing\), then, as \(t\downarrow0\),
\[
        \Cov_\QQ(Y_M(t),B_t^\QQ)
        =
        \sum_{(i,j)\in\cE_M^{\mathrm{ret}}} \ell_{ij}^M\,t^{H_{ij}+1/2}
        +
        o\!\left(t^{\beta_M+1/2}\right).
\]
\end{lemma}

\begin{proof}
By the covariance formula for correlated Brownian integrals,
\[
        \Cov_\QQ(Y_M(t),B_t^\QQ)
        =
        \sum_{j=1}^N \rho_j \int_0^t G_j(t,s)\,\dd s.
\]
If \(\cE_M^{\mathrm{ret}}=\varnothing\), every coefficient
\(b_i\dot\beta_{ij}\rho_j\one_{\{(i,j)\in\cE_M\}}\) vanishes, so this sum is identically zero.
Assume now that \(\cE_M^{\mathrm{ret}}\neq\varnothing\).
Substituting the definition of \(G_j\) and using continuity of \(a_{ij}\) along the diagonal gives
\[
        \int_0^t a_{ij}(t,s)\,(t-s)^{H_{ij}-1/2}\,\dd s
        =
        \frac{a_{ij}(0,0)}{H_{ij}+1/2}\,t^{H_{ij}+1/2}
        +
        o\!\left(t^{H_{ij}+1/2}\right),
\]
uniformly over the finitely many active channels.  Summing over \((i,j)\in\cE_M\) yields the
stated expansion, and channels with \(b_i\dot\beta_{ij}\rho_j=0\) drop out by definition of
\(\cE_M^{\mathrm{ret}}\).
\end{proof}

\begin{proposition}[First variation of the third cumulant of the perturbed market return]\label{prop:market-kappa3}
Under Assumption~\ref{ass:q-index}, if \(\cE_M^{\mathrm{ret}}=\varnothing\), then
\[
        \partial_\varepsilon
        \kappa_3^\QQ(R_\tau^\varepsilon)\big|_{\varepsilon=0}
        =
        0
        \qquad
        \text{for every }\tau\in[0,T].
\]
If \(\cE_M^{\mathrm{ret}}\neq\varnothing\), then, as \(\tau\downarrow0\),
\[
        \partial_\varepsilon
        \kappa_3^\QQ(R_\tau^\varepsilon)\big|_{\varepsilon=0}
        =
        \sum_{(i,j)\in\cE_M^{\mathrm{ret}}}
        \kappa_{ij}^M\,\tau^{H_{ij}+3/2}
        +
        o\!\left(\tau^{\beta_M+3/2}\right).
\]
\end{proposition}

\begin{proof}
Because \(R_\tau^\varepsilon=X_\tau+\varepsilon Z_\tau\) and \(X_\tau\) is centered Gaussian,
\[
        \partial_\varepsilon
        \kappa_3^\QQ(R_\tau^\varepsilon)\big|_{\varepsilon=0}
        =
        3\,\EE_\QQ[X_\tau^2 Z_\tau].
\]
Apply It\^o's product formula to \(X_t^2 Z_t\).  Since
\[
        \dd X_t=\sigma_0^M(t)\,\dd B_t^\QQ,
        \qquad
        \dd Z_t=Y_M(t)\,\dd B_t^\QQ,
\]
taking expectations yields
\[
        \EE_\QQ[X_\tau^2 Z_\tau]
        =
        2\int_0^\tau \sigma_0^M(u)\,\EE_\QQ[X_uY_M(u)]\,\dd u.
\]
By correlated It\^o isometry,
\[
        \EE_\QQ[X_uY_M(u)]
        =
        \sum_{j=1}^N \rho_j\int_0^u \sigma_0^M(r)\,G_j(u,r)\,\dd r.
\]
If \(\cE_M^{\mathrm{ret}}=\varnothing\), the last sum is identically zero, so the first claim
follows.  Assume now that \(\cE_M^{\mathrm{ret}}\neq\varnothing\).
Using the continuity of \(\sigma_0^M\) at \(0\), the continuity of \(a_{ij}\) along the diagonal,
and the definition of \(G_j\), one obtains
\[
        \EE_\QQ[X_uY_M(u)]
        =
        \sigma_0^M(0)
        \sum_{(i,j)\in\cE_M^{\mathrm{ret}}}
        \frac{b_i\dot\beta_{ij}a_{ij}(0,0)\rho_j}{H_{ij}+1/2}\,
        u^{H_{ij}+1/2}
        +
        o\!\left(u^{\beta_M+1/2}\right).
\]
Integrating in \(u\), multiplying by the outer factor \(2\sigma_0^M(0)\), and then multiplying by
\(3\) gives the stated coefficient \(\kappa_{ij}^M\).
\end{proof}

\begin{remark}[Where the Gaussian structure enters]
The first-variation identity
\[
        \partial_\varepsilon
        \kappa_3^\QQ(R_\tau^\varepsilon)\big|_{\varepsilon=0}
        =
        3\,\EE_\QQ[X_\tau^2 Z_\tau]
\]
is a specialization to the present Gaussian-Volterra \(\QQ\)-model.  Gaussianity is absent from
the abstract compression statement of Theorem~\ref{thm:option-abstract}; it enters only in the
explicit closed-form cumulant law derived here.
\end{remark}

\begin{lemma}[Market variance expansion]\label{lem:market-var}
Under Assumption~\ref{ass:q-index},
\[
        \Var_\QQ(R_\tau^0)
        =
        (\sigma_0^M(0))^2\,\tau + o(\tau),
        \qquad
        \tau\downarrow0.
\]
\end{lemma}

\begin{proof}
By It\^o isometry,
\[
        \Var_\QQ(R_\tau^0)=\int_0^\tau (\sigma_0^M(u))^2\,\dd u.
\]
Continuity of \(\sigma_0^M\) at \(0\) gives the claim.
\end{proof}

\begin{lemma}[Projection of second chaos onto one Gaussian coordinate]\label{lem:second-chaos-projection}
Let \(X\) be centered Gaussian with variance \(v>0\), and let \(Z\in L^2(\QQ)\) be a centered
second-Wiener-chaos random variable for the same Gaussian system.  Then
\[
        \EE_\QQ[Z\mid X]
        =
        \frac{\EE_\QQ[X^2Z]}{2v^2}(X^2-v).
\]
Equivalently, if \(\EE_\QQ[Z\mid X]=a(X^2-v)\), then
\[
        a=\frac{\EE_\QQ[X^2Z]}{2v^2}
        =
        \frac{\partial_\varepsilon\kappa_3^\QQ(X+\varepsilon Z)|_{\varepsilon=0}}{6v^2}.
\]
Consequently, for every \(k\in\R\),
\[
        \EE_\QQ[Z\one_{\{X>k\}}]
        =
        \frac{\partial_\varepsilon\kappa_3^\QQ(X+\varepsilon Z)|_{\varepsilon=0}}{6v}
        \frac{k}{\sqrt v}\phi\!\left(\frac{k}{\sqrt v}\right),
\]
where \(\phi\) is the standard normal density.
\end{lemma}

\begin{proof}
The closed subspace \(L^2(\QQ,\sigma(X))\) has orthogonal basis
\(\{\He_m(X/\sqrt v):m\ge0\}\).  Since \(Z\) lies in the second Wiener chaos, it is orthogonal to
every Hermite direction except possibly \(m=2\).  Therefore
\[
        \EE_\QQ[Z\mid X]
        =
        a\,\He_2(X/\sqrt v)
        =
        a\left(\frac{X^2}{v}-1\right)
\]
for some \(a\in\R\).  Equivalently,
\[
        \EE_\QQ[Z\mid X]=\frac{a}{v}(X^2-v).
\]
For a centered normal variable \(X\) with variance \(v\),
\[
        \EE_\QQ[(X^2-v)^2]=2v^2.
\]
Taking inner products with \(X^2-v\) gives
\[
        \frac{a}{v}
        =
        \frac{\EE_\QQ[Z(X^2-v)]}{2v^2}
        =
        \frac{\EE_\QQ[X^2Z]}{2v^2},
\]
which is the displayed conditional projection formula after renaming \(a/v\) as the coefficient
of \(X^2-v\).
Since \(X\) is centered Gaussian and \(\EE_\QQ[Z]=0\),
\[
        \partial_\varepsilon\kappa_3^\QQ(X+\varepsilon Z)|_{\varepsilon=0}=3\EE_\QQ[X^2Z],
\]
which proves the coefficient identity.  Finally, if \(Y=X/\sqrt v\), then the standard normal
integration-by-parts identity gives
\[
        \EE_\QQ[(X^2-v)\one_{\{X>k\}}]
        =
        v\EE_\QQ\!\left[(Y^2-1)\one_{\{Y>k/\sqrt v\}}\right]
        =
        v\frac{k}{\sqrt v}\phi\!\left(\frac{k}{\sqrt v}\right).
\]
Combining the displays yields the claim.
\end{proof}

\begin{proposition}[Gaussian-Volterra leverage perturbations satisfy the projection condition]\label{prop:gaussian-volterra-projection-condition}
Under Assumption~\ref{ass:q-index}, for each sufficiently small \(\tau>0\), the pair
\((X_\tau,Z_\tau)\) satisfies the projection condition of
Lemma~\ref{lem:second-chaos-projection}.  If \(\cE_M^{\mathrm{ret}}=\varnothing\), then this
projection is zero.
\end{proposition}

\begin{proof}
Decompose each \(W^{\QQ,j}\) into its component correlated with \(B^\QQ\) and an orthogonal
Gaussian component.  The Brownian-correlated part of \(Y_M(t)\) is a deterministic Volterra
integral against \(B^\QQ\); integrating it again against \(B^\QQ\) places its contribution to
\(Z_\tau\) in the second Wiener chaos of the return Brownian motion.  Conditional expectation onto
\(\sigma(X_\tau)\), where \(X_\tau\) is a first-chaos variable, keeps only the second Hermite
direction generated by \(X_\tau/\sqrt{\Var_\QQ(X_\tau)}\).  The orthogonal Gaussian components are
independent of \(X_\tau\) and have zero mean, so their \(\sigma(X_\tau)\)-projection vanishes.  If
\(\cE_M^{\mathrm{ret}}=\varnothing\), all Brownian-correlated coefficients vanish, hence the whole
projection is zero.
\end{proof}

\begin{theorem}[Self-contained Bachelier skew--cumulant transfer]\label{thm:self-contained-bachelier}
Let \(X_\tau\) be centered Gaussian with variance \(v(\tau)>0\), let \(Z_\tau\in L^2(\QQ)\), and
set \(R_\tau^\varepsilon:=X_\tau+\varepsilon Z_\tau\).  Define
\[
        C(\tau,k;\varepsilon):=\EE_\QQ[(X_\tau+\varepsilon Z_\tau-k)_+],
\]
and let \(\sigma_N(\tau,k;\varepsilon)\) be the Bachelier implied normal volatility defined by
\[
        C(\tau,k;\varepsilon)
        =
        \sigma_N(\tau,k;\varepsilon)\sqrt\tau\,
        \phi\!\left(\frac{k}{\sigma_N(\tau,k;\varepsilon)\sqrt\tau}\right)
        -
        k\Phi\!\left(-\frac{k}{\sigma_N(\tau,k;\varepsilon)\sqrt\tau}\right),
\]
where \(\Phi\) is the standard normal distribution function.  Assume the projection condition of
Lemma~\ref{lem:second-chaos-projection} holds for \((X_\tau,Z_\tau)\).  Then, for fixed small
\(\tau\),
\[
        \partial_\varepsilon C(\tau,k;\varepsilon)\big|_{\varepsilon=0}
        =
        \frac{\partial_\varepsilon\kappa_3^\QQ(R_\tau^\varepsilon)|_{\varepsilon=0}}{6v(\tau)}
        \frac{k}{\sqrt{v(\tau)}}
        \phi\!\left(\frac{k}{\sqrt{v(\tau)}}\right).
\]
Consequently,
\[
        \partial_\varepsilon\sigma_N(\tau,k;\varepsilon)\big|_{\varepsilon=0}
        =
        \frac{\partial_\varepsilon\kappa_3^\QQ(R_\tau^\varepsilon)|_{\varepsilon=0}}
        {6v(\tau)^{3/2}\sqrt\tau}\,k,
\]
and the ATM normal skew transfer is exact:
\[
        \partial_\varepsilon\partial_k\sigma_N(\tau,k;\varepsilon)
        \big|_{k=0,\varepsilon=0}
        =
        \frac{\partial_\varepsilon\kappa_3^\QQ(R_\tau^\varepsilon)|_{\varepsilon=0}}
        {6v(\tau)^{3/2}\sqrt\tau}.
\]
\end{theorem}

\begin{proof}
Because the law of \(X_\tau\) has no atom at \(k\) and \(Z_\tau\in L^2(\QQ)\), differentiating the
call payoff at \(\varepsilon=0\) gives
\[
        \partial_\varepsilon C(\tau,k;\varepsilon)|_{\varepsilon=0}
        =
        \EE_\QQ[Z_\tau\one_{\{X_\tau>k\}}].
\]
Lemma~\ref{lem:second-chaos-projection} gives the first displayed formula.

At \(\varepsilon=0\), the Bachelier implied normal volatility is
\[
        \sigma_N(\tau,k;0)=\sqrt{v(\tau)/\tau},
\]
because \(X_\tau\) is Gaussian.  The Bachelier vega at this point is
\[
        \partial_\sigma C_N(\tau,k;\sigma)|_{\sigma=\sqrt{v(\tau)/\tau}}
        =
        \sqrt\tau\,\phi\!\left(\frac{k}{\sqrt{v(\tau)}}\right).
\]
Implicit differentiation of the defining equation for implied normal volatility therefore gives
\[
        \partial_\varepsilon\sigma_N(\tau,k;\varepsilon)|_{\varepsilon=0}
        =
        \frac{\partial_\varepsilon C(\tau,k;\varepsilon)|_{\varepsilon=0}}
        {\sqrt\tau\,\phi(k/\sqrt{v(\tau)})}
        =
        \frac{\partial_\varepsilon\kappa_3^\QQ(R_\tau^\varepsilon)|_{\varepsilon=0}}
        {6v(\tau)^{3/2}\sqrt\tau}k .
\]
Differentiating in \(k\) at zero proves the exact ATM skew transfer.
\end{proof}

\begin{proposition}[The non-Gaussian Volterra leverage perturbation satisfies the projection condition]\label{prop:ng-volterra-projection-condition}
Under Assumption~\ref{ass:ng-volterra-primitives}, for each sufficiently small \(\tau>0\), the
perturbation return \(Z_\tau\) satisfies the projection condition of
Lemma~\ref{lem:second-chaos-projection}.  The compound-Poisson and orthogonal Brownian components
have zero contribution to \(\EE_\QQ[Z_\tau\mid X_\tau]\); the Brownian-correlated Volterra
component contributes only through the second Hermite direction \(X_\tau^2-v(\tau)\).
\end{proposition}

\begin{proof}
The Brownian-correlated part of \(Y_t\) is a deterministic Volterra integral against \(B\).  Hence
its stochastic integral against \(B\) is a second-Wiener-chaos random variable.  Conditional
expectation onto \(\sigma(X_\tau)\), where \(X_\tau\) is a first-chaos variable, keeps only the
component in the second Hermite polynomial of \(X_\tau/\sqrt{v(\tau)}\).  The orthogonal Brownian
components are independent of \(X_\tau\) and have zero mean.  Conditional on the path of a
compensated jump-source Volterra integrand, the stochastic integral against \(B\) is jointly
Gaussian with \(X_\tau\), and its conditional covariance is linear in that jump-source integrand.
Taking expectation over the compensated jump source gives zero because the integrand has mean
zero.  Thus the jump and orthogonal Brownian parts have zero \(\sigma(X_\tau)\)-projection, and
the stated projection form follows.
\end{proof}

\begin{corollary}[Self-contained Bachelier expansion for the primitive non-Gaussian Volterra class]\label{cor:self-contained-bachelier-ng-volterra}
Under Assumption~\ref{ass:ng-volterra-primitives}, in the nonempty support case of
Theorem~\ref{thm:ng-volterra-cumulants},
\[
        \partial_\varepsilon\partial_k\sigma_N(\tau,k;\varepsilon)
        \big|_{k=0,\varepsilon=0}
        =
        \sum_{\beta\in\mathcal B_J}A_J(\beta)\tau^{\beta-1/2}
        +
        O(\tau^{\beta_0-1/2+\eta}),
\]
with the same support and cancellation qualifications as in
Theorem~\ref{thm:ng-volterra-cumulants}.
\end{corollary}

\begin{proof}
Theorem~\ref{thm:ng-volterra-cumulants} gives the cumulant expansion with the displayed support,
and Proposition~\ref{prop:ng-volterra-projection-condition} is the admissibility condition that
lets that cumulant coefficient enter the Bachelier skew transfer.  Theorem~\ref{thm:self-contained-bachelier}
then divides by the leading variance scale.  Since
\[
        v(\tau)=\sigma_0(0)^2\tau+O(\tau^{1+\eta}).
\]
it follows that the division shifts each cumulant exponent by \(-1/2\) and leaves the stated
remainder order.
\end{proof}

\begin{corollary}[Gaussian-Volterra/Bachelier verification of the cumulant-class transfer]\label{cor:bachelier-transfer}
Under Assumption~\ref{ass:q-index}, if \(\cE_M^{\mathrm{ret}}=\varnothing\),
then
\[
        \partial_\varepsilon
        \kappa_3^\QQ(R_\tau^\varepsilon)\big|_{\varepsilon=0}
        =
        \partial_\varepsilon
        \mathfrak S_{M,N}(\tau;\varepsilon)\big|_{\varepsilon=0}
        =
        0
        \qquad
        \text{for every sufficiently small }\tau>0.
\]
If instead \(\cE_M^{\mathrm{ret}}\neq\varnothing\), so that \(\beta_M<\infty\), then as
\(\tau\downarrow0\),
\[
        \partial_\varepsilon
        \mathfrak S_{M,N}(\tau;\varepsilon)\big|_{\varepsilon=0}
        =
        \frac{1}{6(\sigma_0^M(0))^3\tau^2}\,
        \partial_\varepsilon \kappa_3^\QQ(R_\tau^\varepsilon)\big|_{\varepsilon=0}
        +
        o\!\left(\tau^{\beta_M-1/2}\right).
\]
Equivalently,
\[
        \partial_\varepsilon
        \mathfrak S_{M,N}(\tau;\varepsilon)\big|_{\varepsilon=0}
        =
        \EE_\QQ[\Xi_{M,\tau}]
        +
        o\!\left(\tau^{\beta_M-1/2}\right),
\]
where \(\Xi_{M,\tau}\) is the Gaussian-Volterra first-variation pricing object from
Definition~\ref{def:gaussian-option-object}.  In particular, the Gaussian-Volterra/Bachelier
option block satisfies Assumption~\ref{ass:option-fv} with trivial public option sigma-field,
identity summary functional, and, in the nontrivial case \(\cE_M^{\mathrm{ret}}\neq\varnothing\),
scale
\[
        a_{M,\tau}:=\tau^{\beta_M-1/2}.
\]
Moreover, with
\[
        G_M(\tau)
        :=
        \frac{\Var_\QQ(R_\tau^0)^{3/2}}{(\sigma_0^M(0))^3\tau^{3/2}},
\]
one has \(G_M(\tau)\to1\).  Consequently, after the coefficient grouping used in
Corollary~\ref{cor:gaussian-market-skew}, the Gaussian-Volterra block verifies
Assumption~\ref{ass:cumulant-class} whenever the grouped cumulant and quote-transfer remainders are
lower order at the resulting selected exponent \(\beta_M^{\mathrm{comp}}\).
\end{corollary}

\begin{proof}
If \(\cE_M^{\mathrm{ret}}=\varnothing\), then
\[
        \EE_\QQ[X_uY_M(u)]=0
        \qquad
        \text{for every }u\in[0,T],
\]
by the correlated It\^o-isometry formula used in the proof of
Proposition~\ref{prop:market-kappa3}.  Hence
\[
        \partial_\varepsilon\kappa_3^\QQ(R_\tau^\varepsilon)\big|_{\varepsilon=0}=0
\]
for all small \(\tau\).  Moreover, the nonzero channels in \(Y_M\) are orthogonal to
\(B^\QQ\); joint Gaussianity gives independence from the return Brownian channel.  Conditional on
the volatility perturbation path, \(R_\tau^\varepsilon\) is therefore centered Gaussian, hence
symmetric.  Equivalently, Proposition~\ref{prop:gaussian-volterra-projection-condition} gives a
zero second-Hermite projection, and Theorem~\ref{thm:self-contained-bachelier} gives vanishing of
the ATM Bachelier skew derivative.  Assume therefore that
\(\cE_M^{\mathrm{ret}}\neq\varnothing\).  Apply
Theorem~\ref{thm:self-contained-bachelier} to the additive Gaussian-Volterra market-index return
\(R_\tau^\varepsilon\), and combine it with Lemma~\ref{lem:market-var}.  The first display is the
resulting transfer law.  Substituting the
identity from Proposition~\ref{prop:market-kappa3} gives
\[
        \frac{1}{6(\sigma_0^M(0))^3\tau^2}\,
        \partial_\varepsilon \kappa_3^\QQ(R_\tau^\varepsilon)\big|_{\varepsilon=0}
        =
        \frac{1}{2(\sigma_0^M(0))^3\tau^2}\EE_\QQ[X_\tau^2 Z_\tau]
        =
        \EE_\QQ[\Xi_{M,\tau}],
\]
which is the second display.
\end{proof}

\begin{proposition}[First variation of the public skew summary as a compression of priced structural loading]\label{prop:option-compression}
Under the assumptions of Corollary~\ref{cor:bachelier-transfer}, define
\[
        \Gamma_{M,\tau}^\QQ
        :=
        C_{\cH_{M,\tau}^{\mathrm{opt}}}^\QQ(L_{M,\tau}^\QQ).
\]
If \(\alpha_{M,\tau}^\QQ>0\), then
\[
        \Gamma_{M,\tau}^\QQ
        =
        \alpha_{M,\tau}^\QQ
        C_{\cH_{M,\tau}^{\mathrm{opt}}}^\QQ(\Lambda_{M,\tau}).
\]
If \(\cE_M^{\mathrm{ret}}=\varnothing\), then
\[
        \partial_\varepsilon
        \mathfrak S_{M,N}(\tau;\varepsilon)\big|_{\varepsilon=0}
        =
        \Gamma_{M,\tau}^\QQ
        =
        0
        \qquad
        \text{for every sufficiently small }\tau>0.
\]
If \(\cE_M^{\mathrm{ret}}\neq\varnothing\), then, as \(\tau\downarrow0\),
\[
        \partial_\varepsilon
        \mathfrak S_{M,N}(\tau;\varepsilon)\big|_{\varepsilon=0}
        =
        \Gamma_{M,\tau}^\QQ
        +
        r_{M,\tau},
        \qquad
        |r_{M,\tau}|=o(a_{M,\tau}),
\]
with \(a_{M,\tau}=\tau^{\beta_M-1/2}\).  Thus the public skew summary is the public compression of
the option-facing priced structural loading up to a higher-order remainder at the Gaussian market
scale.
\end{proposition}

\begin{proof}
Under Corollary~\ref{cor:bachelier-transfer}, the Gaussian-Volterra/Bachelier option family
satisfies Assumption~\ref{ass:option-fv} with trivial public option sigma-field,
identity summary functional, and option first-variation pricing object \(\Xi_{M,\tau}\) from
Definition~\ref{def:gaussian-option-object}.  The compression identity is therefore the abstract
statement of Theorem~\ref{thm:option-abstract} specialized to this model.  Since the public option
sigma-field is trivial, \(\Gamma_{M,\tau}^\QQ\) is the constant \(\EE_\QQ[\Xi_{M,\tau}]\).  The
empty-channel zero follows from Corollary~\ref{cor:bachelier-transfer}; the nonempty remainder
scale is the scale identified there.
\end{proof}

\begin{corollary}[Option visibility is norm visibility at the quote scale]\label{cor:option-norm}
Under the assumptions of Proposition~\ref{prop:option-compression},
\[
        \|C_{\cH_{M,\tau}^{\mathrm{opt}}}^{\QQ}(L_{M,\tau}^\QQ)\|_{L^2(\QQ)}
        =
        |\Gamma_{M,\tau}^\QQ|.
\]
Because \(\cH_{M,\tau}^{\mathrm{opt}}\) is trivial, the option-facing structural contribution
induced by the synchronized perturbation is \(\QQ\)-visible at maturity \(\tau\) exactly when
\(\Gamma_{M,\tau}^\QQ\neq0\), and it is screened at that quote scale exactly when
\(\Gamma_{M,\tau}^\QQ=0\) while
\(L_{M,\tau}^\QQ\neq0\).
\end{corollary}

\begin{proof}
The compressed loading \(C_{\cH_{M,\tau}^{\mathrm{opt}}}^{\QQ}(L_{M,\tau}^\QQ)\) is the constant
random variable \(\Gamma_{M,\tau}^\QQ\).  The \(L^2(\QQ)\)-norm of a constant equals its absolute
value, and Proposition~\ref{thm:q-trichotomy} then gives the stated visibility criterion.
\end{proof}

\begin{theorem}[Aggregate-factor approximation for the option-facing centered loading]\label{thm:option-factor}
Assume Assumption~\ref{ass:option-fv}, and suppose the aggregate factor benchmark is nontrivial, so
\[
        \Lambda_M^\star\in \cM_\QQ^\circ,
        \qquad
        \|\Lambda_M^\star\|_{L^2(\QQ)}=1.
\]
Fix a sufficiently small maturity \(\tau>0\), and suppose
that the option first-variation pricing object admits a decomposition
\[
        \Xi_{M,\tau}
        =
        X_{M,\tau}^{\mathrm{vis}}
        +
        a_{M,\tau}^\star \Lambda_M^\star
        +
        E_{M,\tau},
\]
where
\[
        X_{M,\tau}^{\mathrm{vis}}\in L^2(\QQ,\cG_T),
        \qquad
        a_{M,\tau}^\star\in\R,
        \qquad
        E_{M,\tau}\in L^2(\QQ,\cF_T).
\]
Define the centered projected remainder
\[
        \widetilde E_{M,\tau}^\QQ
        :=
        C_{\cH_\varpi}^\QQ(E_{M,\tau})
        -
        C_{\cG_T}^\QQ\bigl(C_{\cH_\varpi}^\QQ(E_{M,\tau})\bigr),
\]
and set
\[
        \delta_{M,\tau}^\star:=\|\widetilde E_{M,\tau}^\QQ\|_{L^2(\QQ)}.
\]
Then the centered option-facing priced structural loading satisfies
\[
        \widetilde L_{M,\tau}^\QQ
        :=
        L_{M,\tau}^\QQ-C_{\cG_T}^\QQ(L_{M,\tau}^\QQ)
        =
        a_{M,\tau}^\star\Lambda_M^\star+\widetilde E_{M,\tau}^\QQ.
\]
Consequently, in the decomposition of Proposition~\ref{thm:factor-decomp},
\[
        \widetilde L_{M,\tau}^\QQ
        =
        \alpha_{M,\tau}^{\QQ,\star}\Lambda_M^\star
        +
        R_{M,\tau}^{\QQ,\star},
\]
one has
\[
        |\alpha_{M,\tau}^{\QQ,\star}-a_{M,\tau}^\star|
        \le
        \delta_{M,\tau}^\star,
        \qquad
        \|R_{M,\tau}^{\QQ,\star}\|_{L^2(\QQ)}
        \le
        \delta_{M,\tau}^\star.
\]
If, in addition,
\[
        a_{M,\tau}^\star\neq0
        \quad\text{for all sufficiently small }\tau,
        \qquad
        \delta_{M,\tau}^\star=o(|a_{M,\tau}^\star|),
\]
then the centered option-facing loading is asymptotically aggregate-factor aligned:
\[
        \frac{\|R_{M,\tau}^{\QQ,\star}\|_{L^2(\QQ)}}{|\alpha_{M,\tau}^{\QQ,\star}|}
        \to 0
        \qquad
        \text{as }\tau\downarrow0,
\]
and
\[
        \left\|
                \frac{\widetilde L_{M,\tau}^\QQ}{\|\widetilde L_{M,\tau}^\QQ\|_{L^2(\QQ)}}
                -
                \frac{a_{M,\tau}^\star}{|a_{M,\tau}^\star|}\Lambda_M^\star
        \right\|_{L^2(\QQ)}
        \to 0.
\]
\end{theorem}

\begin{proof}
By linearity of conditional expectation and because
\(X_{M,\tau}^{\mathrm{vis}}\in L^2(\QQ,\cG_T)\subseteq L^2(\QQ,\cH_\varpi)\),
\[
        L_{M,\tau}^\QQ
        =
        X_{M,\tau}^{\mathrm{vis}}
        +
        a_{M,\tau}^\star\Lambda_M^\star
        +
        C_{\cH_\varpi}^\QQ(E_{M,\tau}).
\]
Apply \(C_{\cG_T}^\QQ\) and use \(C_{\cG_T}^\QQ(\Lambda_M^\star)=0\) from
Corollary~\ref{cor:factor-screening}:
\[
        C_{\cG_T}^\QQ(L_{M,\tau}^\QQ)
        =
        X_{M,\tau}^{\mathrm{vis}}
        +
        C_{\cG_T}^\QQ\bigl(C_{\cH_\varpi}^\QQ(E_{M,\tau})\bigr).
\]
Subtracting gives the first displayed identity for \(\widetilde L_{M,\tau}^\QQ\).  Now apply
Proposition~\ref{thm:factor-decomp} to \(\widetilde L_{M,\tau}^\QQ\).  Since
\[
        \alpha_{M,\tau}^{\QQ,\star}
        =
        \langle \widetilde L_{M,\tau}^\QQ,\Lambda_M^\star\rangle_{L^2(\QQ)}
        =
        a_{M,\tau}^\star
        +
        \langle \widetilde E_{M,\tau}^\QQ,\Lambda_M^\star\rangle_{L^2(\QQ)},
\]
Cauchy--Schwarz yields
\[
        |\alpha_{M,\tau}^{\QQ,\star}-a_{M,\tau}^\star|
        \le
        \delta_{M,\tau}^\star.
\]
The orthogonal residual is
\[
        R_{M,\tau}^{\QQ,\star}
        =
        \widetilde E_{M,\tau}^\QQ
        -
        \langle \widetilde E_{M,\tau}^\QQ,\Lambda_M^\star\rangle_{L^2(\QQ)}\Lambda_M^\star,
\]
so
\[
        \|R_{M,\tau}^{\QQ,\star}\|_{L^2(\QQ)}
        \le
        \delta_{M,\tau}^\star.
\]
If \(\delta_{M,\tau}^\star=o(|a_{M,\tau}^\star|)\), then the previous coefficient bound implies
\[
        |\alpha_{M,\tau}^{\QQ,\star}|
        \ge
        |a_{M,\tau}^\star|-\delta_{M,\tau}^\star
        \sim
        |a_{M,\tau}^\star|,
\]
hence
\[
        \frac{\|R_{M,\tau}^{\QQ,\star}\|_{L^2(\QQ)}}{|\alpha_{M,\tau}^{\QQ,\star}|}
        \to 0.
\]
In particular, \(\alpha_{M,\tau}^{\QQ,\star}\) has the same sign as \(a_{M,\tau}^\star\) for all
sufficiently small \(\tau\).  Finally,
\[
        \frac{\widetilde L_{M,\tau}^\QQ}{\|\widetilde L_{M,\tau}^\QQ\|_{L^2(\QQ)}}
        =
        \frac{\alpha_{M,\tau}^{\QQ,\star}}{\sqrt{(\alpha_{M,\tau}^{\QQ,\star})^2+\|R_{M,\tau}^{\QQ,\star}\|_{L^2(\QQ)}^2}}
        \Lambda_M^\star
        +
        \frac{R_{M,\tau}^{\QQ,\star}}{\|\widetilde L_{M,\tau}^\QQ\|_{L^2(\QQ)}},
\]
where the first scalar coefficient converges to
\(
        a_{M,\tau}^\star/|a_{M,\tau}^\star|
\)
and the second term converges to \(0\) in \(L^2(\QQ)\).  This proves the normalized-direction
claim.
\end{proof}

\begin{corollary}[Exact option-factor alignment criterion]\label{cor:option-factor-exact}
Under the hypotheses of Theorem~\ref{thm:option-factor}, if
\[
        C_{\cH_\varpi}^\QQ(E_{M,\tau})
        =
        C_{\cG_T}^\QQ\bigl(C_{\cH_\varpi}^\QQ(E_{M,\tau})\bigr),
\]
equivalently \(\delta_{M,\tau}^\star=0\), then
\[
        \widetilde L_{M,\tau}^\QQ
        =
        a_{M,\tau}^\star\Lambda_M^\star,
        \qquad
        R_{M,\tau}^{\QQ,\star}=0.
\]
In particular, the option-facing centered loading is exactly aggregate-factor aligned whenever the
non-factor remainder contributes no centered \(\cH_\varpi\)-content beyond the visible sigma-field.
\end{corollary}

\begin{proof}
If \(\delta_{M,\tau}^\star=0\), then \(\widetilde E_{M,\tau}^\QQ=0\) in
Theorem~\ref{thm:option-factor}.  The factor decomposition in that theorem therefore reduces to
the single span component
\(\widetilde L_{M,\tau}^\QQ=a_{M,\tau}^\star\Lambda_M^\star\), and the orthogonal residual
\(R_{M,\tau}^{\QQ,\star}\) is zero by uniqueness of the Hilbert-space decomposition.
\end{proof}

\begin{remark}[Meaning for public option summaries]
Theorem~\ref{thm:option-factor} is stated only in the nontrivial factor regime
\(\Lambda_M^\star\neq0\).  When the centered aggregate benchmark collapses, there is no nonzero
aggregate factor direction for the centered option-facing loading to align with, so the theorem's
factor statement becomes vacuous.

Theorem~\ref{thm:option-factor} does not claim direct public visibility of the aggregate factor.
Corollary~\ref{cor:factor-screening} still gives
\[
        C_{\cH_{M,\tau}^{\mathrm{opt}}}^\QQ(\Lambda_M^\star)=0
        \qquad
        \text{for }
        \cH_{M,\tau}^{\mathrm{opt}}\subseteq\cG_T.
\]
Its content is different.  When the option first-variation object is generated by a
visible component plus the aggregate factor benchmark up to a lower-order centered
remainder, the centered option-facing loading itself points asymptotically in the aggregate factor
direction.  The quoted option summary still acts on the visible representative of that loading, but
the nonlinear and hedging consequences of the option-facing claim inherit the same factor
dominance.
\end{remark}

\begin{proposition}[Gaussian reduction of option-factor alignment to the perturbation return]\label{prop:gaussian-option-factor}
Assume Assumption~\ref{ass:q-index}, and suppose the aggregate factor benchmark is nontrivial, so
\[
        \Lambda_M^\star\in \cM_\QQ^\circ,
        \qquad
        \|\Lambda_M^\star\|_{L^2(\QQ)}=1.
\]
Set
\[
        q_\tau:=\frac{1}{2(\sigma_0^M(0))^3\tau^2},
        \qquad
        v_\tau:=\Var_\QQ(R_\tau^0).
\]
Let
\[
        L_{Z_\tau}^\QQ:=C_{\cH_\varpi}^\QQ(Z_\tau),
        \qquad
        V_{Z_\tau}^\QQ:=C_{\cG_T}^\QQ(L_{Z_\tau}^\QQ),
\]
and write the canonical factor decomposition of the perturbation return loading as
\[
        L_{Z_\tau}^\QQ
        =
        V_{Z_\tau}^\QQ
        +
        b_{M,\tau}^{\QQ,\star}\Lambda_M^\star
        +
        R_{Z_\tau}^{\QQ,\star},
\]
where \(R_{Z_\tau}^{\QQ,\star}\in \cM_\QQ^\circ\) and
\[
        R_{Z_\tau}^{\QQ,\star}\perp_{L^2(\QQ)}\Lambda_M^\star.
\]
Define the quadratic-fluctuation remainder
\[
        J_{M,\tau}
        :=
        q_\tau\bigl(X_\tau^2-v_\tau\bigr)Z_\tau,
\]
and its centered structural projection
\[
        \widetilde J_{M,\tau}^\QQ
        :=
        C_{\cH_\varpi}^\QQ(J_{M,\tau})
        -
        C_{\cG_T}^\QQ\bigl(C_{\cH_\varpi}^\QQ(J_{M,\tau})\bigr).
\]
Then the centered option-facing priced structural loading satisfies the exact identity
\[
        \widetilde L_{M,\tau}^\QQ
        =
        q_\tau v_\tau\,b_{M,\tau}^{\QQ,\star}\Lambda_M^\star
        +
        q_\tau v_\tau\,R_{Z_\tau}^{\QQ,\star}
        +
        \widetilde J_{M,\tau}^\QQ.
\]
Consequently,
\[
        \left|
                \alpha_{M,\tau}^{\QQ,\star}
                -
                q_\tau v_\tau\,b_{M,\tau}^{\QQ,\star}
        \right|
        \le
        \|\widetilde J_{M,\tau}^\QQ\|_{L^2(\QQ)},
\]
and
\[
        \|R_{M,\tau}^{\QQ,\star}\|_{L^2(\QQ)}
        \le
        q_\tau v_\tau\,\|R_{Z_\tau}^{\QQ,\star}\|_{L^2(\QQ)}
        +
        \|\widetilde J_{M,\tau}^\QQ\|_{L^2(\QQ)}.
\]
If, in addition,
\[
        b_{M,\tau}^{\QQ,\star}\neq0
        \quad\text{for all sufficiently small }\tau,
\]
and
\[
        q_\tau v_\tau\,\|R_{Z_\tau}^{\QQ,\star}\|_{L^2(\QQ)}
        +
        \|\widetilde J_{M,\tau}^\QQ\|_{L^2(\QQ)}
        =
        o\!\left(q_\tau v_\tau\,|b_{M,\tau}^{\QQ,\star}|\right),
\]
then the centered option-facing loading is asymptotically aggregate-factor aligned.
\end{proposition}

\begin{proof}
By Definition~\ref{def:gaussian-option-object},
\[
        \Xi_{M,\tau}
        =
        q_\tau X_\tau^2 Z_\tau
        =
        q_\tau v_\tau Z_\tau
        +
        J_{M,\tau}.
\]
Applying \(C_{\cH_\varpi}^\QQ\) and using the definition of \(L_{Z_\tau}^\QQ\) gives
\[
        L_{M,\tau}^\QQ
        =
        q_\tau v_\tau L_{Z_\tau}^\QQ
        +
        C_{\cH_\varpi}^\QQ(J_{M,\tau}).
\]
Apply \(C_{\cG_T}^\QQ\) and subtract:
\[
        \widetilde L_{M,\tau}^\QQ
        :=
        L_{M,\tau}^\QQ-C_{\cG_T}^\QQ(L_{M,\tau}^\QQ)
        =
        q_\tau v_\tau
        \bigl(L_{Z_\tau}^\QQ-C_{\cG_T}^\QQ(L_{Z_\tau}^\QQ)\bigr)
        +
        \widetilde J_{M,\tau}^\QQ.
\]
Using the centered part of the factor decomposition of \(L_{Z_\tau}^\QQ\),
\[
        L_{Z_\tau}^\QQ-C_{\cG_T}^\QQ(L_{Z_\tau}^\QQ)
        =
        b_{M,\tau}^{\QQ,\star}\Lambda_M^\star
        +
        R_{Z_\tau}^{\QQ,\star},
\]
yields the exact identity for \(\widetilde L_{M,\tau}^\QQ\).  Now apply
Proposition~\ref{thm:factor-decomp} to \(\widetilde L_{M,\tau}^\QQ\).  Since
\[
        \alpha_{M,\tau}^{\QQ,\star}
        =
        q_\tau v_\tau\,b_{M,\tau}^{\QQ,\star}
        +
        q_\tau v_\tau
        \langle R_{Z_\tau}^{\QQ,\star},\Lambda_M^\star\rangle_{L^2(\QQ)}
        +
        \langle \widetilde J_{M,\tau}^\QQ,\Lambda_M^\star\rangle_{L^2(\QQ)},
\]
and \(R_{Z_\tau}^{\QQ,\star}\perp\Lambda_M^\star\), the coefficient error reduces to
\[
        \alpha_{M,\tau}^{\QQ,\star}
        -
        q_\tau v_\tau\,b_{M,\tau}^{\QQ,\star}
        =
        \langle \widetilde J_{M,\tau}^\QQ,\Lambda_M^\star\rangle_{L^2(\QQ)},
\]
so Cauchy--Schwarz gives the coefficient bound.
The residual bound follows by taking the orthogonal complement of
\(\mathrm{span}\{\Lambda_M^\star\}\) in the displayed identity for
\(\widetilde L_{M,\tau}^\QQ\).  If the last displayed smallness condition holds, then
\[
        \left|
                \alpha_{M,\tau}^{\QQ,\star}
                -
                q_\tau v_\tau\,b_{M,\tau}^{\QQ,\star}
        \right|
        =
        o\!\left(q_\tau v_\tau\,|b_{M,\tau}^{\QQ,\star}|\right),
\]
and similarly
\[
        \|R_{M,\tau}^{\QQ,\star}\|_{L^2(\QQ)}
        =
        o\!\left(q_\tau v_\tau\,|b_{M,\tau}^{\QQ,\star}|\right).
\]
Therefore
\[
        \frac{\|R_{M,\tau}^{\QQ,\star}\|_{L^2(\QQ)}}{|\alpha_{M,\tau}^{\QQ,\star}|}
        \to0,
\]
which is exactly aggregate-factor alignment in the sense of
Theorem~\ref{thm:option-factor}.
\end{proof}

\begin{proposition}[Moment control of the quadratic-fluctuation remainder]\label{prop:gaussian-J-bound}
Under Assumption~\ref{ass:q-index},
\[
        \|\widetilde J_{M,\tau}^\QQ\|_{L^2(\QQ)}
        \le
        \|J_{M,\tau}\|_{L^2(\QQ)}
        \le
        q_\tau\|X_\tau^2-v_\tau\|_{L^4(\QQ)}\|Z_\tau\|_{L^4(\QQ)}.
\]
Moreover, if \(G\sim N(0,1)\) and \(c_G:=\|G^2-1\|_{L^4}\), then
\[
        \|X_\tau^2-v_\tau\|_{L^4(\QQ)}=c_G v_\tau,
\]
so
\[
        \|\widetilde J_{M,\tau}^\QQ\|_{L^2(\QQ)}
        \le
        c_G q_\tau v_\tau \|Z_\tau\|_{L^4(\QQ)}.
\]
In particular, by Burkholder--Davis--Gundy there exists a universal constant \(C_{\mathrm{BDG}}>0\)
such that
\[
        \|\widetilde J_{M,\tau}^\QQ\|_{L^2(\QQ)}
        \le
        C_{\mathrm{BDG}} c_G q_\tau v_\tau
        \left(
                \EE_\QQ\!\left[
                        \left(
                                \int_0^\tau Y_M(u)^2\,\dd u
                        \right)^2
                \right]
        \right)^{1/4}.
\]
\end{proposition}

\begin{proof}
Because \(L^2(\QQ,\cG_T)\subseteq L^2(\QQ,\cH_\varpi)\), the term
\[
        \widetilde J_{M,\tau}^\QQ
        =
        C_{\cH_\varpi}^\QQ(J_{M,\tau})
        -
        C_{\cG_T}^\QQ\bigl(C_{\cH_\varpi}^\QQ(J_{M,\tau})\bigr)
\]
is the orthogonal remainder of \(C_{\cH_\varpi}^\QQ(J_{M,\tau})\) relative to
\(L^2(\QQ,\cG_T)\).  Therefore
\[
        \|\widetilde J_{M,\tau}^\QQ\|_{L^2(\QQ)}
        \le
        \|C_{\cH_\varpi}^\QQ(J_{M,\tau})\|_{L^2(\QQ)}
        \le
        \|J_{M,\tau}\|_{L^2(\QQ)}.
\]
Now
\[
        J_{M,\tau}=q_\tau(X_\tau^2-v_\tau)Z_\tau,
\]
so H\"older gives
\[
        \|J_{M,\tau}\|_{L^2(\QQ)}
        \le
        q_\tau\|X_\tau^2-v_\tau\|_{L^4(\QQ)}\|Z_\tau\|_{L^4(\QQ)}.
\]
Since \(X_\tau/\sqrt{v_\tau}\sim N(0,1)\),
\[
        X_\tau^2-v_\tau
        =
        v_\tau\left(\left(\frac{X_\tau}{\sqrt{v_\tau}}\right)^2-1\right),
\]
which yields
\[
        \|X_\tau^2-v_\tau\|_{L^4(\QQ)}=c_G v_\tau.
\]
Finally, the \(L^4\) Burkholder--Davis--Gundy inequality gives
\[
        \|Z_\tau\|_{L^4(\QQ)}
        \le
        C_{\mathrm{BDG}}
        \left(
                \EE_\QQ\!\left[
                        \left(
                                \int_0^\tau Y_M(u)^2\,\dd u
                        \right)^2
                \right]
        \right)^{1/4},
\]
and substitution proves the last display.
\end{proof}

\begin{proposition}[Kernel-energy control of the quadratic-fluctuation remainder]\label{prop:gaussian-J-kernel}
Under Assumption~\ref{ass:q-index}, define the deterministic synchronized-perturbation kernel
energy
\[
        \mathfrak D_{M,\tau}^{\mathrm{pert}}
        :=
        \sum_{j=1}^N \int_0^\tau \int_0^u
        |G_j(u,s)|^2\,\dd s\,\dd u.
\]
Then
\[
        \left(
                \EE_\QQ\!\left[
                        \left(
                                \int_0^\tau Y_M(u)^2\,\dd u
                        \right)^2
                \right]
        \right)^{1/4}
        \le
        3^{1/4}
        \bigl(\mathfrak D_{M,\tau}^{\mathrm{pert}}\bigr)^{1/2},
\]
and therefore
\[
        \|\widetilde J_{M,\tau}^\QQ\|_{L^2(\QQ)}
        \le
        3^{1/4} C_{\mathrm{BDG}} c_G q_\tau v_\tau
        \bigl(\mathfrak D_{M,\tau}^{\mathrm{pert}}\bigr)^{1/2}.
\]
If, in addition,
\[
        \mathfrak D_{M,\tau}^{\mathrm{pert}}=0,
\]
then
\[
        \widetilde J_{M,\tau}^\QQ=0.
\]
\end{proposition}

\begin{proof}
Set
\[
        \Sigma_{M,\tau}(u)
        :=
        \EE_\QQ[Y_M(u)^2]
        =
        \sum_{j=1}^N \int_0^u |G_j(u,s)|^2\,\dd s,
        \qquad
        0\le u\le \tau,
\]
where the second identity is It\^o isometry and independence of the \(\QQ\)-Brownian
coordinates.  Because \(Y_M(u)\) is centered Gaussian for every \(u\), Wick's formula gives
\[
        \EE_\QQ[Y_M(u)^2Y_M(v)^2]
        =
        \Sigma_{M,\tau}(u)\Sigma_{M,\tau}(v)
        +
        2\,\Cov_\QQ(Y_M(u),Y_M(v))^2.
\]
By Cauchy--Schwarz,
\[
        \Cov_\QQ(Y_M(u),Y_M(v))^2
        \le
        \Sigma_{M,\tau}(u)\Sigma_{M,\tau}(v),
\]
so
\[
        \EE_\QQ[Y_M(u)^2Y_M(v)^2]
        \le
        3\,\Sigma_{M,\tau}(u)\Sigma_{M,\tau}(v).
\]
Integrating over \(u,v\in[0,\tau]\) yields
\[
        \EE_\QQ\!\left[
                \left(
                        \int_0^\tau Y_M(u)^2\,\dd u
                \right)^2
        \right]
        \le
        3
        \left(
                \int_0^\tau \Sigma_{M,\tau}(u)\,\dd u
        \right)^2
        =
        3\bigl(\mathfrak D_{M,\tau}^{\mathrm{pert}}\bigr)^2,
\]
which proves the first display after taking fourth roots.  The second display then follows from
the last bound in Proposition~\ref{prop:gaussian-J-bound}.  If
\(\mathfrak D_{M,\tau}^{\mathrm{pert}}=0\), the second display gives
\(\|\widetilde J_{M,\tau}^\QQ\|_{L^2(\QQ)}=0\), hence \(\widetilde J_{M,\tau}^\QQ=0\).
\end{proof}

\begin{remark}[What the Gaussian reduction theorem localizes]
Proposition~\ref{prop:gaussian-option-factor} localizes the option-factor problem in the
Gaussian-Volterra block.  The two primitive ingredients are the factor decomposition of the
perturbation return loading \(L_{Z_\tau}^\QQ\) and the centered size of the quadratic-fluctuation
remainder \(J_{M,\tau}\).  Thus the localized task is
to prove that the centered loading of \(Z_\tau\) is itself aggregate-factor dominated.  The
Perron-projected kernel construction introduced below sharpens that statement one step further.  It
isolates a canonical Perron-generated proxy claim for the perturbation return and an explicit
off-Perron kernel defect.  Proposition~\ref{prop:gaussian-J-bound} shows that the quadratic
remainder is already controlled by a fourth-moment perturbation-energy bound, so the remaining
work is pushed into the proxy-compatibility error and the off-Perron kernel defect.
\end{remark}

\begin{proposition}[Claim-level Perron approximation controls the perturbation-return loading]\label{prop:gaussian-z-approx}
Assume Assumption~\ref{ass:q-index}, and suppose the aggregate factor benchmark is nontrivial, so
\[
        \Lambda_M^\star\in \cM_\QQ^\circ,
        \qquad
        \|\Lambda_M^\star\|_{L^2(\QQ)}=1.
\]
Fix a sufficiently small maturity \(\tau>0\), and suppose the perturbation return admits a
decomposition
\[
        Z_\tau
        =
        Z_{M,\tau}^{\varpi}
        +
        E_{Z,\tau},
\]
where \(Z_{M,\tau}^{\varpi}\in L^2(\QQ,\cH_\varpi)\) and
\[
        Z_{M,\tau}^{\varpi}
        =
        Z_{M,\tau}^{\mathrm{vis}}
        +
        b_{M,\tau}^{\varpi,\star}\Lambda_M^\star,
\]
with
\[
        Z_{M,\tau}^{\mathrm{vis}}\in L^2(\QQ,\cG_T),
        \qquad
        b_{M,\tau}^{\varpi,\star}\in\R,
        \qquad
        E_{Z,\tau}\in L^2(\QQ,\cF_T).
\]
Define
\[
        \widetilde E_{Z,\tau}^\QQ
        :=
        C_{\cH_\varpi}^\QQ(E_{Z,\tau})
        -
        C_{\cG_T}^\QQ\bigl(C_{\cH_\varpi}^\QQ(E_{Z,\tau})\bigr).
\]
Then
\[
        L_{Z_\tau}^\QQ
        =
        Z_{M,\tau}^{\mathrm{vis}}
        +
        b_{M,\tau}^{\varpi,\star}\Lambda_M^\star
        +
        C_{\cH_\varpi}^\QQ(E_{Z,\tau}),
\]
and
\[
        L_{Z_\tau}^\QQ-C_{\cG_T}^\QQ(L_{Z_\tau}^\QQ)
        =
        b_{M,\tau}^{\varpi,\star}\Lambda_M^\star
        +
        \widetilde E_{Z,\tau}^\QQ.
\]
Consequently, in the canonical factor decomposition
\[
        L_{Z_\tau}^\QQ
        =
        V_{Z_\tau}^\QQ
        +
        b_{M,\tau}^{\QQ,\star}\Lambda_M^\star
        +
        R_{Z_\tau}^{\QQ,\star},
\]
one has
\[
        \left|
                b_{M,\tau}^{\QQ,\star}
                -
                b_{M,\tau}^{\varpi,\star}
        \right|
        \le
        \|E_{Z,\tau}\|_{L^2(\QQ)},
        \qquad
        \|R_{Z_\tau}^{\QQ,\star}\|_{L^2(\QQ)}
        \le
        \|E_{Z,\tau}\|_{L^2(\QQ)}.
\]
If, in addition,
\[
        b_{M,\tau}^{\varpi,\star}\neq0
        \quad\text{for all sufficiently small }\tau,
        \qquad
        \|E_{Z,\tau}\|_{L^2(\QQ)}
        =
        o\!\left(|b_{M,\tau}^{\varpi,\star}|\right),
\]
then the centered perturbation-return loading is asymptotically aggregate-factor aligned.
\end{proposition}

\begin{proof}
Because \(Z_{M,\tau}^{\varpi}\) is \(\cH_\varpi\)-measurable,
\[
        L_{Z_\tau}^\QQ
        =
        C_{\cH_\varpi}^\QQ(Z_\tau)
        =
        Z_{M,\tau}^{\varpi}
        +
        C_{\cH_\varpi}^\QQ(E_{Z,\tau}),
\]
which gives the first displayed identity.  Applying \(C_{\cG_T}^\QQ\) and using
\[
        C_{\cG_T}^\QQ(\Lambda_M^\star)=0
\]
from Corollary~\ref{cor:factor-screening} yields
\[
        C_{\cG_T}^\QQ(L_{Z_\tau}^\QQ)
        =
        Z_{M,\tau}^{\mathrm{vis}}
        +
        C_{\cG_T}^\QQ\bigl(C_{\cH_\varpi}^\QQ(E_{Z,\tau})\bigr),
\]
so subtracting gives the centered identity with \(\widetilde E_{Z,\tau}^\QQ\).  Since
\[
        \|\widetilde E_{Z,\tau}^\QQ\|_{L^2(\QQ)}
        \le
        \|E_{Z,\tau}\|_{L^2(\QQ)}
\]
by the same orthogonal-remainder argument as in Proposition~\ref{prop:gaussian-J-bound},
Proposition~\ref{thm:factor-decomp} gives
\[
        \left|
                b_{M,\tau}^{\QQ,\star}
                -
                b_{M,\tau}^{\varpi,\star}
        \right|
        \le
        \|\widetilde E_{Z,\tau}^\QQ\|_{L^2(\QQ)},
        \qquad
        \|R_{Z_\tau}^{\QQ,\star}\|_{L^2(\QQ)}
        \le
        \|\widetilde E_{Z,\tau}^\QQ\|_{L^2(\QQ)},
\]
hence the displayed bounds.  If
\(\|E_{Z,\tau}\|_{L^2(\QQ)}=o(|b_{M,\tau}^{\varpi,\star}|)\), then
\[
        \|R_{Z_\tau}^{\QQ,\star}\|_{L^2(\QQ)}
        =
        o\!\left(|b_{M,\tau}^{\varpi,\star}|\right),
\]
and
\[
        |b_{M,\tau}^{\QQ,\star}|
        \ge
        |b_{M,\tau}^{\varpi,\star}|-\|E_{Z,\tau}\|_{L^2(\QQ)}
        \sim
        |b_{M,\tau}^{\varpi,\star}|,
\]
which proves asymptotic aggregate-factor alignment of the centered perturbation-return loading.
\end{proof}

\begin{proposition}[Off-Perron perturbation energy controls the perturbation-return factor error]\label{prop:gaussian-z-energy}
Assume the hypotheses of Proposition~\ref{prop:gaussian-z-approx}.  Suppose, moreover, that the
claim-level Perron approximation error is realized as a stochastic integral remainder:
\[
        E_{Z,\tau}
        =
        \int_0^\tau Y_{M,\tau}^{\perp}(u)\,\dd B_u^\QQ
\]
for some predictable process \(Y_{M,\tau}^{\perp}\) satisfying
\[
        \EE_\QQ\!\left[\int_0^\tau \bigl(Y_{M,\tau}^{\perp}(u)\bigr)^2\,\dd u\right]<\infty.
\]
Then
\[
        \|E_{Z,\tau}\|_{L^2(\QQ)}
        =
        \left(
                \EE_\QQ\!\left[
                        \int_0^\tau \bigl(Y_{M,\tau}^{\perp}(u)\bigr)^2\,\dd u
                \right]
        \right)^{1/2}.
\]
Consequently,
\[
        \left|
                b_{M,\tau}^{\QQ,\star}
                -
                b_{M,\tau}^{\varpi,\star}
        \right|
        \le
        \left(
                \EE_\QQ\!\left[
                        \int_0^\tau \bigl(Y_{M,\tau}^{\perp}(u)\bigr)^2\,\dd u
                \right]
        \right)^{1/2},
\]
and
\[
        \|R_{Z_\tau}^{\QQ,\star}\|_{L^2(\QQ)}
        \le
        \left(
                \EE_\QQ\!\left[
                        \int_0^\tau \bigl(Y_{M,\tau}^{\perp}(u)\bigr)^2\,\dd u
                \right]
        \right)^{1/2}.
\]
If, in addition,
\[
        b_{M,\tau}^{\varpi,\star}\neq0
        \quad\text{for all sufficiently small }\tau,
\]
and
\[
        \EE_\QQ\!\left[\int_0^\tau \bigl(Y_{M,\tau}^{\perp}(u)\bigr)^2\,\dd u\right]
        =
        o\!\left(|b_{M,\tau}^{\varpi,\star}|^2\right),
\]
then the centered perturbation-return loading is asymptotically aggregate-factor aligned.
\end{proposition}

\begin{proof}
The identity for \(\|E_{Z,\tau}\|_{L^2(\QQ)}\) is It\^o isometry.  Substituting that identity into
Proposition~\ref{prop:gaussian-z-approx} yields the two displayed bounds.  If the last displayed
smallness condition holds, then
\[
        \|E_{Z,\tau}\|_{L^2(\QQ)}
        =
        o\!\left(|b_{M,\tau}^{\varpi,\star}|\right),
\]
so the final implication of Proposition~\ref{prop:gaussian-z-approx} gives asymptotic
aggregate-factor alignment of the centered perturbation-return loading.
\end{proof}

\begin{proposition}[Channel-side decomposition of the synchronized perturbation implies perturbation-return factor control]\label{prop:gaussian-y-to-z}
Assume Assumption~\ref{ass:q-index}, and suppose the aggregate factor benchmark is nontrivial.  Fix
a sufficiently small maturity \(\tau>0\), and suppose there exist predictable processes
\[
        Y_{M,\tau}^{\varpi},
        \qquad
        Y_{M,\tau}^{\perp},
\]
such that
\[
        Y_M(u)=Y_{M,\tau}^{\varpi}(u)+Y_{M,\tau}^{\perp}(u),
        \qquad
        0\le u\le \tau.
\]
Define
\[
        Z_{M,\tau}^{\varpi}
        :=
        \int_0^\tau Y_{M,\tau}^{\varpi}(u)\,\dd B_u^\QQ,
        \qquad
        E_{Z,\tau}
        :=
        \int_0^\tau Y_{M,\tau}^{\perp}(u)\,\dd B_u^\QQ.
\]
Assume further that
\[
        Z_{M,\tau}^{\varpi}
        =
        Z_{M,\tau}^{\mathrm{vis}}
        +
        b_{M,\tau}^{\varpi,\star}\Lambda_M^\star
\]
for some
\[
        Z_{M,\tau}^{\mathrm{vis}}\in L^2(\QQ,\cG_T),
        \qquad
        b_{M,\tau}^{\varpi,\star}\in\R.
\]
Then the hypotheses of Proposition~\ref{prop:gaussian-z-energy} hold with this choice of
\(E_{Z,\tau}\), and
\[
        \left|
                b_{M,\tau}^{\QQ,\star}
                -
                b_{M,\tau}^{\varpi,\star}
        \right|
        \le
        \left(
                \EE_\QQ\!\left[
                        \int_0^\tau \bigl(Y_{M,\tau}^{\perp}(u)\bigr)^2\,\dd u
                \right]
        \right)^{1/2},
\]
and
\[
        \|R_{Z_\tau}^{\QQ,\star}\|_{L^2(\QQ)}
        \le
        \left(
                \EE_\QQ\!\left[
                        \int_0^\tau \bigl(Y_{M,\tau}^{\perp}(u)\bigr)^2\,\dd u
                \right]
        \right)^{1/2}.
\]
If, in addition,
\[
        b_{M,\tau}^{\varpi,\star}\neq0
        \quad\text{for all sufficiently small }\tau,
\]
and
\[
        \EE_\QQ\!\left[
                \int_0^\tau \bigl(Y_{M,\tau}^{\perp}(u)\bigr)^2\,\dd u
        \right]
        =
        o\!\left(|b_{M,\tau}^{\varpi,\star}|^2\right),
\]
then the centered perturbation-return loading is asymptotically aggregate-factor aligned.
\end{proposition}

\begin{proof}
The definitions of \(Z_{M,\tau}^{\varpi}\) and \(E_{Z,\tau}\) give
\[
        Z_\tau
        =
        \int_0^\tau Y_M(u)\,\dd B_u^\QQ
        =
        Z_{M,\tau}^{\varpi}
        +
        E_{Z,\tau}.
\]
The assumed representation of \(Z_{M,\tau}^{\varpi}\) is exactly the decomposition required in
Proposition~\ref{prop:gaussian-z-approx}.  Proposition~\ref{prop:gaussian-z-energy} then applies
with the present stochastic-integral remainder and gives the displayed bounds and the final
alignment implication.
\end{proof}


\begin{remark}[Descending Gaussian option chain]
The Gaussian option block now follows a descending chain of reductions.  First,
Proposition~\ref{prop:gaussian-option-factor} reduces option-factor alignment to the perturbation
return \(Z_\tau\) and the quadratic-fluctuation remainder \(\widetilde J_{M,\tau}^\QQ\).  Next,
Propositions~\ref{prop:gaussian-z-approx}, \ref{prop:gaussian-z-energy}, and
\ref{prop:gaussian-y-to-z} reduce the perturbation-return factor problem to the off-Perron energy
of the synchronized perturbation process.  The canonical Perron construction below then isolates a
Perron-generated proxy claim and pushes the remaining structural uncertainty into the off-Perron
kernel defect.  The Gaussian subsection can therefore be read as one ladder:
\[
        \text{option claim}
        \;\Longrightarrow\;
        \text{perturbation return}
        \;\Longrightarrow\;
        \text{Perron proxy claim}
\]
\[
        \text{Perron proxy claim}
        \;\Longrightarrow\;
        \text{perturbation process}
        \;\Longrightarrow\;
        \text{kernel defect}.
\]
\end{remark}

\begin{corollary}[Kernel-side off-Perron defect criterion for option-factor alignment]\label{cor:gaussian-kernel-to-option}
Assume Assumption~\ref{ass:q-index}, and suppose that for each sufficiently small maturity
\(\tau>0\) and each \(1\le j\le N\) the synchronized perturbation kernel admits a decomposition
\[
        G_j(u,s)=G_{j,\tau}^{\varpi}(u,s)+G_{j,\tau}^{\perp}(u,s),
        \qquad
        0\le s\le u\le \tau,
\]
with corresponding processes
\[
        Y_{M,\tau}^{\varpi}(u)
        :=
        \sum_{j=1}^N \int_0^u G_{j,\tau}^{\varpi}(u,s)\,\dd W_s^{\QQ,j},
        \qquad
        Y_{M,\tau}^{\perp}(u)
        :=
        \sum_{j=1}^N \int_0^u G_{j,\tau}^{\perp}(u,s)\,\dd W_s^{\QQ,j}.
\]
Assume further that
\[
        \int_0^\tau Y_{M,\tau}^{\varpi}(u)\,\dd B_u^\QQ
        =
        Z_{M,\tau}^{\mathrm{vis}}
        +
        b_{M,\tau}^{\varpi,\star}\Lambda_M^\star
\]
for some \(Z_{M,\tau}^{\mathrm{vis}}\in L^2(\QQ,\cG_T)\), with
\[
        b_{M,\tau}^{\varpi,\star}\neq0
        \quad\text{for all sufficiently small }\tau.
\]
If
\[
        \sum_{j=1}^N \int_0^\tau \int_0^u
        \bigl|G_{j,\tau}^{\perp}(u,s)\bigr|^2\,\dd s\,\dd u
        =
        o\!\left(|b_{M,\tau}^{\varpi,\star}|^2\right),
\]
and
\[
        \left(
                \EE_\QQ\!\left[
                        \left(
                                \int_0^\tau Y_M(u)^2\,\dd u
                        \right)^2
                \right]
        \right)^{1/4}
        =
        o\!\left(|b_{M,\tau}^{\varpi,\star}|\right),
\]
then the centered option-facing loading is asymptotically aggregate-factor aligned.
\end{corollary}

\begin{proof}
Set
\[
        Y_{M,\tau}^{\perp}(u)
        :=
        \sum_{j=1}^N \int_0^u G_{j,\tau}^{\perp}(u,s)\,\dd W_s^{\QQ,j}.
\]
Because the \(\QQ\)-Brownian coordinates are independent, It\^o isometry gives
\[
        \EE_\QQ\!\left[
                \int_0^\tau \bigl(Y_{M,\tau}^{\perp}(u)\bigr)^2\,\dd u
        \right]
        =
        \sum_{j=1}^N \int_0^\tau \int_0^u
        \bigl|G_{j,\tau}^{\perp}(u,s)\bigr|^2\,\dd s\,\dd u.
\]
Hence the kernel defect hypothesis implies the off-Perron perturbation-energy condition in
Proposition~\ref{prop:gaussian-y-to-z}.  The visible-plus-factor representation of
\(\int_0^\tau Y_{M,\tau}^{\varpi}(u)\,\dd B_u^\QQ\) is exactly the structural hypothesis of that
proposition.  Therefore Proposition~\ref{prop:gaussian-y-to-z} gives
\[
        \|R_{Z_\tau}^{\QQ,\star}\|_{L^2(\QQ)}
        =
        o\!\left(|b_{M,\tau}^{\varpi,\star}|\right),
\]
and
\[
        |b_{M,\tau}^{\QQ,\star}|
        \sim
        |b_{M,\tau}^{\varpi,\star}|.
\]
Proposition~\ref{prop:gaussian-J-bound} then yields
\[
        \|\widetilde J_{M,\tau}^\QQ\|_{L^2(\QQ)}
        =
        o\!\left(q_\tau v_\tau\,|b_{M,\tau}^{\QQ,\star}|\right).
\]
Thus
\[
        q_\tau v_\tau\,\|R_{Z_\tau}^{\QQ,\star}\|_{L^2(\QQ)}
        +
        \|\widetilde J_{M,\tau}^\QQ\|_{L^2(\QQ)}
        =
        o\!\left(q_\tau v_\tau\,|b_{M,\tau}^{\QQ,\star}|\right),
\]
and Proposition~\ref{prop:gaussian-option-factor} gives asymptotic aggregate-factor alignment of
the centered option-facing loading.
\end{proof}

\begin{corollary}[Kernel-side off-Perron defect and deterministic synchronized-perturbation kernel energy criterion for option-factor alignment]\label{cor:gaussian-kernel-to-option-kernel}
Assume the hypotheses of Corollary~\ref{cor:gaussian-kernel-to-option}, except that instead of
assuming
\[
        \left(
                \EE_\QQ\!\left[
                        \left(
                                \int_0^\tau Y_M(u)^2\,\dd u
                        \right)^2
                \right]
        \right)^{1/4}
        =
        o\!\left(|b_{M,\tau}^{\varpi,\star}|\right),
\]
assume
\[
        \bigl(\mathfrak D_{M,\tau}^{\mathrm{pert}}\bigr)^{1/2}
        =
        o\!\left(|b_{M,\tau}^{\varpi,\star}|\right).
\]
Then the centered option-facing loading is asymptotically aggregate-factor aligned.
\end{corollary}

\begin{proof}
By Proposition~\ref{prop:gaussian-J-kernel},
\[
        \left(
                \EE_\QQ\!\left[
                        \left(
                                \int_0^\tau Y_M(u)^2\,\dd u
                        \right)^2
                \right]
        \right)^{1/4}
        =
        O\!\left(
                \bigl(\mathfrak D_{M,\tau}^{\mathrm{pert}}\bigr)^{1/2}
        \right),
\]
so the new hypothesis implies the moment condition in
Corollary~\ref{cor:gaussian-kernel-to-option}.  The conclusion follows.
\end{proof}

\begin{definition}[Canonical Perron-projected perturbation kernels and proxy claim]\label{def:gaussian-perron-proxy}
Under Definition~\ref{def:lambda-structure} and Assumption~\ref{ass:q-index}, let
\[
        P_\varpi:=vu^\top,
        \qquad
        P_\perp:=I_N-P_\varpi,
\]
where \(I_N\) is the identity on \(\R^N\).  For \(1\le j\le N\) and \(0\le s\le u\le T\), define
the channel vector
\[
        K_j(u,s)
        :=
        \bigl(
                \dot\beta_{ij}\,a_{ij}(u,s)\,(u-s)^{H_{ij}-1/2}\,
                \one_{\{(i,j)\in\cE_M\}}
        \bigr)_{i=1}^N
        \in\R^N.
\]
Then
\[
        G_j(u,s)=b^\top K_j(u,s).
\]
Define the canonical Perron-projected channel kernels
\[
        K_j^\varpi(u,s):=P_\varpi K_j(u,s),
        \qquad
        K_j^\perp(u,s):=P_\perp K_j(u,s),
\]
and the induced market kernels
\[
        G_j^\varpi(u,s):=b^\top K_j^\varpi(u,s),
        \qquad
        G_j^\perp(u,s):=b^\top K_j^\perp(u,s).
\]
For each maturity \(\tau>0\), define
\[
        Y_M^\varpi(u)
        :=
        \sum_{j=1}^N \int_0^u G_j^\varpi(u,s)\,\dd W_s^{\QQ,j},
        \qquad
        Y_M^\perp(u)
        :=
        \sum_{j=1}^N \int_0^u G_j^\perp(u,s)\,\dd W_s^{\QQ,j},
\]
\[
        Z_{M,\tau}^\varpi
        :=
        \int_0^\tau Y_M^\varpi(u)\,\dd B_u^\QQ,
        \qquad
        E_{Z,\tau}^\perp
        :=
        \int_0^\tau Y_M^\perp(u)\,\dd B_u^\QQ.
\]
The canonical off-Perron kernel defect is
\[
        \mathfrak D_{M,\tau}^\perp
        :=
        \sum_{j=1}^N \int_0^\tau \int_0^u
        \bigl|G_j^\perp(u,s)\bigr|^2\,\dd s\,\dd u.
\]
Finally, define the centered Perron-generated proxy claim
\[
        \Upsilon_{M,\tau}^\varpi
        :=
        C_{\cH_\varpi}^\QQ(Z_{M,\tau}^\varpi)
        -
        C_{\cG_T}^\QQ\bigl(C_{\cH_\varpi}^\QQ(Z_{M,\tau}^\varpi)\bigr)
        \in\cM_\QQ^\circ.
\]
\end{definition}

\begin{remark}[The Perron projection is algebraic]
The projection \(P_\varpi=vu^\top\) is the algebraic Perron projection associated with the
left-right normalization \(u^\top v=1\).  In general the projection is oblique.  Thus
\(P_\perp=I_N-P_\varpi\) denotes the complementary Perron residual map; norms of
\(P_\perp\)-terms are residual norms, with Pythagorean orthogonal-energy interpretations requiring
additional self-adjointness or a compatible-inner-product condition.
\end{remark}

\begin{proposition}[Rank-one Perron reduction of synchronized perturbation kernels]\label{prop:gaussian-rank-one}
Under Definition~\ref{def:gaussian-perron-proxy}, define the scalar Perron profiles
\[
        \psi_j(u,s):=u^\top K_j(u,s),
        \qquad
        1\le j\le N,
        \qquad
        0\le s\le u\le T.
\]
Then
\[
        K_j^\varpi(u,s)=v\,\psi_j(u,s),
        \qquad
        K_j^\perp(u,s)=K_j(u,s)-v\,\psi_j(u,s),
\]
and therefore
\[
        G_j^\varpi(u,s)
        =
        (b^\top v)\,\psi_j(u,s).
\]
Thus, after Perron projection, the synchronized perturbation kernels are rank one in the asset
index and the remaining structure is carried entirely by the scalar profiles \(\psi_j\).
\end{proposition}

\begin{proof}
By definition,
\[
        K_j^\varpi(u,s)
        =
        P_\varpi K_j(u,s)
        =
        vu^\top K_j(u,s)
        =
        v\,\psi_j(u,s).
\]
The identity for \(K_j^\perp\) follows from \(P_\perp=I_N-P_\varpi\), and then
\[
        G_j^\varpi(u,s)=b^\top K_j^\varpi(u,s)=(b^\top v)\psi_j(u,s).
\]
\end{proof}

\begin{proposition}[Canonical Perron proxy claim and kernel-defect control]\label{prop:gaussian-perron-proxy}
Under Definition~\ref{def:gaussian-perron-proxy} and Assumption~\ref{ass:q-index}, for every
sufficiently small maturity \(\tau>0\),
\[
        Y_M(u)=Y_M^\varpi(u)+Y_M^\perp(u),
        \qquad
        0\le u\le \tau,
\]
and
\[
        Z_\tau
        =
        Z_{M,\tau}^\varpi
        +
        E_{Z,\tau}^\perp.
\]
Moreover,
\[
        \|E_{Z,\tau}^\perp\|_{L^2(\QQ)}^2
        =
        \mathfrak D_{M,\tau}^\perp,
\]
\[
        \|\widetilde L_{Z_\tau}^\QQ-\Upsilon_{M,\tau}^\varpi\|_{L^2(\QQ)}
        \le
        \bigl(\mathfrak D_{M,\tau}^\perp\bigr)^{1/2},
\]
and, with \(q_\tau\) and \(v_\tau\) as in Proposition~\ref{prop:gaussian-option-factor},
\[
        \|\widetilde L_{M,\tau}^\QQ-q_\tau v_\tau\,\Upsilon_{M,\tau}^\varpi\|_{L^2(\QQ)}
        \le
        q_\tau v_\tau\,\bigl(\mathfrak D_{M,\tau}^\perp\bigr)^{1/2}
        +
        \|\widetilde J_{M,\tau}^\QQ\|_{L^2(\QQ)}.
\]
\end{proposition}

\begin{proof}
Because \(P_\varpi+P_\perp=I_N\), one has
\[
        K_j(u,s)=K_j^\varpi(u,s)+K_j^\perp(u,s),
        \qquad
        G_j(u,s)=G_j^\varpi(u,s)+G_j^\perp(u,s).
\]
Substituting these identities into the stochastic-integral definitions yields
\[
        Y_M(u)=Y_M^\varpi(u)+Y_M^\perp(u),
        \qquad
        Z_\tau=Z_{M,\tau}^\varpi+E_{Z,\tau}^\perp.
\]
Since \(E_{Z,\tau}^\perp\) is a stochastic integral against \(B^\QQ\), It\^o isometry gives
\[
        \|E_{Z,\tau}^\perp\|_{L^2(\QQ)}^2
        =
        \EE_\QQ\!\left[
                \int_0^\tau \bigl(Y_M^\perp(u)\bigr)^2\,\dd u
        \right].
\]
Using independence of the \(\QQ\)-Brownian coordinates,
\[
        \EE_\QQ\!\left[
                \int_0^\tau \bigl(Y_M^\perp(u)\bigr)^2\,\dd u
        \right]
        =
        \sum_{j=1}^N \int_0^\tau \int_0^u
        \bigl|G_j^\perp(u,s)\bigr|^2\,\dd s\,\dd u
        =
        \mathfrak D_{M,\tau}^\perp.
\]
Applying \(C_{\cH_\varpi}^\QQ\) to \(Z_\tau=Z_{M,\tau}^\varpi+E_{Z,\tau}^\perp\) gives
\[
        L_{Z_\tau}^\QQ
        =
        C_{\cH_\varpi}^\QQ(Z_{M,\tau}^\varpi)
        +
        C_{\cH_\varpi}^\QQ(E_{Z,\tau}^\perp).
\]
Subtracting the \(\cG_T\)-projection yields
\[
        \widetilde L_{Z_\tau}^\QQ
        =
        \Upsilon_{M,\tau}^\varpi
        +
        C_{\cH_\varpi}^\QQ(E_{Z,\tau}^\perp)
        -
        C_{\cG_T}^\QQ\bigl(C_{\cH_\varpi}^\QQ(E_{Z,\tau}^\perp)\bigr),
\]
so
\[
        \|\widetilde L_{Z_\tau}^\QQ-\Upsilon_{M,\tau}^\varpi\|_{L^2(\QQ)}
        \le
        \|E_{Z,\tau}^\perp\|_{L^2(\QQ)}
        =
        \bigl(\mathfrak D_{M,\tau}^\perp\bigr)^{1/2}.
\]
Finally, the proof of Proposition~\ref{prop:gaussian-option-factor} gives the exact identity
\[
        \widetilde L_{M,\tau}^\QQ
        =
        q_\tau v_\tau\,\widetilde L_{Z_\tau}^\QQ
        +
        \widetilde J_{M,\tau}^\QQ.
\]
Subtract \(q_\tau v_\tau\,\Upsilon_{M,\tau}^\varpi\) and apply the previous bound.
\end{proof}

\begin{theorem}[Canonical Perron proxy compatibility implies option-factor alignment]\label{thm:gaussian-proxy-factor}
Assume the hypotheses of Proposition~\ref{prop:gaussian-perron-proxy}, and suppose the aggregate
factor benchmark is nontrivial.  Fix a sufficiently small maturity \(\tau>0\), and suppose the
centered Perron-generated proxy claim admits the factor decomposition
\[
        \Upsilon_{M,\tau}^\varpi
        =
        c_{M,\tau}^\varpi\Lambda_M^\star
        +
        D_{M,\tau}^\varpi,
\]
where
\[
        c_{M,\tau}^\varpi\in\R,
        \qquad
        D_{M,\tau}^\varpi\in\cM_\QQ^\circ,
        \qquad
        D_{M,\tau}^\varpi\perp_{L^2(\QQ)}\Lambda_M^\star.
\]
Then
\[
        \left|
                \alpha_{M,\tau}^{\QQ,\star}
                -
                q_\tau v_\tau\,c_{M,\tau}^\varpi
        \right|
        \le
        q_\tau v_\tau\,\|D_{M,\tau}^\varpi\|_{L^2(\QQ)}
        +
        q_\tau v_\tau\,\bigl(\mathfrak D_{M,\tau}^\perp\bigr)^{1/2}
        +
        \|\widetilde J_{M,\tau}^\QQ\|_{L^2(\QQ)},
\]
and
\[
        \|R_{M,\tau}^{\QQ,\star}\|_{L^2(\QQ)}
        \le
        q_\tau v_\tau\,\|D_{M,\tau}^\varpi\|_{L^2(\QQ)}
        +
        q_\tau v_\tau\,\bigl(\mathfrak D_{M,\tau}^\perp\bigr)^{1/2}
        +
        \|\widetilde J_{M,\tau}^\QQ\|_{L^2(\QQ)}.
\]
If, moreover,
\[
        c_{M,\tau}^\varpi\neq0
        \quad\text{for all sufficiently small }\tau,
\]
and
\[
        \|D_{M,\tau}^\varpi\|_{L^2(\QQ)}
        +
        \bigl(\mathfrak D_{M,\tau}^\perp\bigr)^{1/2}
        +
        \frac{\|\widetilde J_{M,\tau}^\QQ\|_{L^2(\QQ)}}{q_\tau v_\tau}
        =
        o\!\left(|c_{M,\tau}^\varpi|\right),
\]
then the centered option-facing loading is asymptotically aggregate-factor aligned.  If in
addition
\[
        D_{M,\tau}^\varpi=0,
        \qquad
        \mathfrak D_{M,\tau}^\perp=0,
        \qquad
        \widetilde J_{M,\tau}^\QQ=0,
\]
then
\[
        \widetilde L_{M,\tau}^\QQ
        =
        q_\tau v_\tau\,c_{M,\tau}^\varpi\Lambda_M^\star,
        \qquad
        R_{M,\tau}^{\QQ,\star}=0,
\]
so the centered option-facing loading is exactly aggregate-factor aligned.
\end{theorem}

\begin{proof}
Set
\[
        \mathcal E_{M,\tau}
        :=
        \widetilde L_{M,\tau}^\QQ-q_\tau v_\tau\,\Upsilon_{M,\tau}^\varpi.
\]
By Proposition~\ref{prop:gaussian-perron-proxy},
\[
        \|\mathcal E_{M,\tau}\|_{L^2(\QQ)}
        \le
        q_\tau v_\tau\,\bigl(\mathfrak D_{M,\tau}^\perp\bigr)^{1/2}
        +
        \|\widetilde J_{M,\tau}^\QQ\|_{L^2(\QQ)}.
\]
Substituting the factor decomposition of \(\Upsilon_{M,\tau}^\varpi\) gives
\[
        \widetilde L_{M,\tau}^\QQ
        =
        q_\tau v_\tau\,c_{M,\tau}^\varpi\Lambda_M^\star
        +
        q_\tau v_\tau\,D_{M,\tau}^\varpi
        +
        \mathcal E_{M,\tau}.
\]
Apply Proposition~\ref{thm:factor-decomp} to the centered claim \(\widetilde L_{M,\tau}^\QQ\).  Since
\[
        q_\tau v_\tau\,D_{M,\tau}^\varpi+\mathcal E_{M,\tau}
\]
is the full perturbation away from the benchmark coefficient \(q_\tau v_\tau c_{M,\tau}^\varpi\),
the orthogonal projection onto \(\mathrm{span}\{\Lambda_M^\star\}\) yields
\[
        \left|
                \alpha_{M,\tau}^{\QQ,\star}
                -
                q_\tau v_\tau\,c_{M,\tau}^\varpi
        \right|
        \le
        \|q_\tau v_\tau\,D_{M,\tau}^\varpi+\mathcal E_{M,\tau}\|_{L^2(\QQ)},
\]
and similarly
\[
        \|R_{M,\tau}^{\QQ,\star}\|_{L^2(\QQ)}
        \le
        \|q_\tau v_\tau\,D_{M,\tau}^\varpi+\mathcal E_{M,\tau}\|_{L^2(\QQ)}.
\]
The displayed bounds follow by the triangle inequality and the bound on \(\mathcal E_{M,\tau}\).
If the smallness condition holds, then
\[
        \|q_\tau v_\tau\,D_{M,\tau}^\varpi+\mathcal E_{M,\tau}\|_{L^2(\QQ)}
        =
        o\!\left(q_\tau v_\tau\,|c_{M,\tau}^\varpi|\right),
\]
so
\[
        |\alpha_{M,\tau}^{\QQ,\star}|
        \ge
        q_\tau v_\tau\,|c_{M,\tau}^\varpi|
        -
        \|q_\tau v_\tau\,D_{M,\tau}^\varpi+\mathcal E_{M,\tau}\|_{L^2(\QQ)}
        \sim
        q_\tau v_\tau\,|c_{M,\tau}^\varpi|.
\]
Hence
\[
        \frac{\|R_{M,\tau}^{\QQ,\star}\|_{L^2(\QQ)}}{|\alpha_{M,\tau}^{\QQ,\star}|}\to0,
\]
which is exactly asymptotic aggregate-factor alignment.  If
\[
        D_{M,\tau}^\varpi=0,
        \qquad
        \mathfrak D_{M,\tau}^\perp=0,
        \qquad
        \widetilde J_{M,\tau}^\QQ=0,
\]
then \(\mathcal E_{M,\tau}=0\).  The preceding error inequality has zero right-hand side, so the
centered loading equals its aggregate-factor component.
\end{proof}

\begin{theorem}[Common-profile Perron benchmark kernels imply proxy compatibility]\label{thm:gaussian-common-profile}
Assume the hypotheses of Definition~\ref{def:gaussian-perron-proxy}, and suppose the aggregate
factor benchmark is nontrivial.  Let \(\bar G_1,\dots,\bar G_N\) be deterministic benchmark
kernels on \(\{0\le s\le u\le T\}\), and define
\[
        \bar Y_M(u)
        :=
        \sum_{j=1}^N \int_0^u \bar G_j(u,s)\,\dd W_s^{\QQ,j},
        \qquad
        \bar Z_{M,\tau}
        :=
        \int_0^\tau \bar Y_M(u)\,\dd B_u^\QQ.
\]
Write the centered benchmark-profile loading as
\[
        \widetilde L_{\bar Z_{M,\tau}}^\QQ
        :=
        C_{\cH_\varpi}^\QQ(\bar Z_{M,\tau})
        -
        C_{\cG_T}^\QQ\bigl(C_{\cH_\varpi}^\QQ(\bar Z_{M,\tau})\bigr)
        =
        \bar c_{M,\tau}^{\QQ,\star}\Lambda_M^\star
        +
        \bar R_{M,\tau}^{\QQ,\star},
\]
where
\[
        \bar R_{M,\tau}^{\QQ,\star}\perp_{L^2(\QQ)}\Lambda_M^\star.
\]
Suppose further that the canonical Perron-projected market kernels admit the common-profile decomposition
\[
        G_j^\varpi(u,s)
        =
        a_{M,\tau}^\varpi \bar G_j(u,s)
        +
        G_{j,\tau}^{\mathrm{comp}}(u,s),
        \qquad
        0\le s\le u\le \tau,
        \qquad
        1\le j\le N,
\]
for some scalar \(a_{M,\tau}^\varpi\in\R\).  Define the benchmark-compatibility defect
\[
        \mathfrak D_{M,\tau}^{\mathrm{comp}}
        :=
        \sum_{j=1}^N \int_0^\tau \int_0^u
        \bigl|G_{j,\tau}^{\mathrm{comp}}(u,s)\bigr|^2\,\dd s\,\dd u.
\]
Then, in the factor decomposition
\[
        \Upsilon_{M,\tau}^\varpi
        =
        c_{M,\tau}^\varpi\Lambda_M^\star
        +
        D_{M,\tau}^\varpi
\]
with
\[
        D_{M,\tau}^\varpi\perp_{L^2(\QQ)}\Lambda_M^\star,
\]
one has
\[
        \left|
                c_{M,\tau}^\varpi
                -
                a_{M,\tau}^\varpi \bar c_{M,\tau}^{\QQ,\star}
        \right|
        \le
        \bigl(\mathfrak D_{M,\tau}^{\mathrm{comp}}\bigr)^{1/2},
\]
and
\[
        \|D_{M,\tau}^\varpi\|_{L^2(\QQ)}
        \le
        |a_{M,\tau}^\varpi|\,\|\bar R_{M,\tau}^{\QQ,\star}\|_{L^2(\QQ)}
        +
        \bigl(\mathfrak D_{M,\tau}^{\mathrm{comp}}\bigr)^{1/2}.
\]
If, moreover,
\[
        a_{M,\tau}^\varpi \bar c_{M,\tau}^{\QQ,\star}\neq0
        \quad\text{for all sufficiently small }\tau,
\]
and
\[
        |a_{M,\tau}^\varpi|\,\|\bar R_{M,\tau}^{\QQ,\star}\|_{L^2(\QQ)}
        +
        \bigl(\mathfrak D_{M,\tau}^{\mathrm{comp}}\bigr)^{1/2}
        =
        o\!\left(|a_{M,\tau}^\varpi \bar c_{M,\tau}^{\QQ,\star}|\right),
\]
then the proxy-compatibility error is asymptotically negligible:
\[
        \|D_{M,\tau}^\varpi\|_{L^2(\QQ)}
        =
        o\!\left(|a_{M,\tau}^\varpi \bar c_{M,\tau}^{\QQ,\star}|\right),
\]
and
\[
        c_{M,\tau}^\varpi
        =
        a_{M,\tau}^\varpi \bar c_{M,\tau}^{\QQ,\star}
        +
        o\!\left(|a_{M,\tau}^\varpi \bar c_{M,\tau}^{\QQ,\star}|\right).
\]
If in addition
\[
        \bar R_{M,\tau}^{\QQ,\star}=0,
        \qquad
        \mathfrak D_{M,\tau}^{\mathrm{comp}}=0,
\]
then
\[
        \Upsilon_{M,\tau}^\varpi
        =
        a_{M,\tau}^\varpi \bar c_{M,\tau}^{\QQ,\star}\Lambda_M^\star
        \qquad
        \text{exactly.}
\]
\end{theorem}

\begin{proof}
Define
\[
        Y_{M,\tau}^{\mathrm{comp}}(u)
        :=
        \sum_{j=1}^N \int_0^u G_{j,\tau}^{\mathrm{comp}}(u,s)\,\dd W_s^{\QQ,j},
        \qquad
        E_{M,\tau}^{\mathrm{comp}}
        :=
        \int_0^\tau Y_{M,\tau}^{\mathrm{comp}}(u)\,\dd B_u^\QQ.
\]
By linearity of stochastic integration,
\[
        Z_{M,\tau}^\varpi
        =
        a_{M,\tau}^\varpi \bar Z_{M,\tau}
        +
        E_{M,\tau}^{\mathrm{comp}}.
\]
It\^o isometry and independence of the \(\QQ\)-Brownian coordinates give
\[
        \|E_{M,\tau}^{\mathrm{comp}}\|_{L^2(\QQ)}^2
        =
        \mathfrak D_{M,\tau}^{\mathrm{comp}}.
\]
Applying \(C_{\cH_\varpi}^\QQ\), subtracting the \(\cG_T\)-projection, and using the displayed
definition of \(\widetilde L_{\bar Z_{M,\tau}}^\QQ\), one obtains
\[
        \Upsilon_{M,\tau}^\varpi
        =
        a_{M,\tau}^\varpi \widetilde L_{\bar Z_{M,\tau}}^\QQ
        +
        \widetilde E_{M,\tau}^{\mathrm{comp},\QQ}
\]
and therefore
\[
        \Upsilon_{M,\tau}^\varpi
        =
        a_{M,\tau}^\varpi \bar c_{M,\tau}^{\QQ,\star}\Lambda_M^\star
        +
        a_{M,\tau}^\varpi \bar R_{M,\tau}^{\QQ,\star}
        +
        \widetilde E_{M,\tau}^{\mathrm{comp},\QQ},
\]
where
\[
        \widetilde E_{M,\tau}^{\mathrm{comp},\QQ}
        :=
        C_{\cH_\varpi}^\QQ(E_{M,\tau}^{\mathrm{comp}})
        -
        C_{\cG_T}^\QQ\bigl(C_{\cH_\varpi}^\QQ(E_{M,\tau}^{\mathrm{comp}})\bigr).
\]
As in Proposition~\ref{prop:gaussian-z-approx},
\[
        \|\widetilde E_{M,\tau}^{\mathrm{comp},\QQ}\|_{L^2(\QQ)}
        \le
        \|E_{M,\tau}^{\mathrm{comp}}\|_{L^2(\QQ)}
        =
        \bigl(\mathfrak D_{M,\tau}^{\mathrm{comp}}\bigr)^{1/2}.
\]
Taking the \(L^2(\QQ)\)-inner product against \(\Lambda_M^\star\), and using
\(
        \bar R_{M,\tau}^{\QQ,\star}\perp\Lambda_M^\star,
\)
gives
\[
        \left|
                c_{M,\tau}^\varpi
                -
                a_{M,\tau}^\varpi \bar c_{M,\tau}^{\QQ,\star}
        \right|
        \le
        \|\widetilde E_{M,\tau}^{\mathrm{comp},\QQ}\|_{L^2(\QQ)}.
\]
Subtracting the \(\Lambda_M^\star\)-component from the displayed decomposition of
\(\Upsilon_{M,\tau}^\varpi\) gives
\[
        D_{M,\tau}^\varpi
        =
        a_{M,\tau}^\varpi \bar R_{M,\tau}^{\QQ,\star}
        +
        \widetilde E_{M,\tau}^{\mathrm{comp},\QQ}
        -
        \langle \widetilde E_{M,\tau}^{\mathrm{comp},\QQ},\Lambda_M^\star\rangle_{L^2(\QQ)}
        \Lambda_M^\star,
\]
hence
\[
        \|D_{M,\tau}^\varpi\|_{L^2(\QQ)}
        \le
        |a_{M,\tau}^\varpi|\,\|\bar R_{M,\tau}^{\QQ,\star}\|_{L^2(\QQ)}
        +
        \|\widetilde E_{M,\tau}^{\mathrm{comp},\QQ}\|_{L^2(\QQ)}.
\]
Dividing the displayed bound by \(|a_{M,\tau}^\varpi|\) gives the stated asymptotic estimates under
the two smallness assumptions.  If
\[
        \bar R_{M,\tau}^{\QQ,\star}=0,
        \qquad
        \mathfrak D_{M,\tau}^{\mathrm{comp}}=0,
\]
then \(E_{M,\tau}^{\mathrm{comp}}=0\) in \(L^2(\QQ)\), so the exact identity follows.
\end{proof}

\begin{proposition}[Common-profile proxy compatibility with a nondegenerate scalar amplitude forces benchmark-loading coefficient control]\label{prop:gaussian-common-profile-loading}
Assume the hypotheses of Theorem~\ref{thm:gaussian-common-profile}.  Suppose there exist constants
\(\underline a_M,\overline a_M>0\) such that
\[
        \underline a_M
        \le
        |a_{M,\tau}^\varpi|
        \le
        \overline a_M
        \qquad
        \text{for all sufficiently small }\tau,
\]
and that
\[
        c_{M,\tau}^\varpi\neq0
        \qquad
        \text{for all sufficiently small }\tau.
\]
If
\[
        \|D_{M,\tau}^\varpi\|_{L^2(\QQ)}
        +
        \bigl(\mathfrak D_{M,\tau}^{\mathrm{comp}}\bigr)^{1/2}
        =
        o\!\left(
                |c_{M,\tau}^\varpi|
        \right),
\]
then
\[
        \left|
                \bar c_{M,\tau}^{\QQ,\star}
        \right|
        \asymp
        |c_{M,\tau}^\varpi|,
\]
and
\[
        \|\bar R_{M,\tau}^{\QQ,\star}\|_{L^2(\QQ)}
        =
        o\!\left(
                |c_{M,\tau}^\varpi|
        \right)
        =
        o\!\left(
                \left|
                        \bar c_{M,\tau}^{\QQ,\star}
                \right|
        \right).
\]
If, in addition,
\[
        \mathfrak d_{M,\tau}^{\varpi,\mathrm{bench}}
        =
        o\!\left(
                |c_{M,\tau}^\varpi|
        \right),
\]
then Proposition~\ref{prop:lambda-generator-separation-loading} applies with
\[
        a_\tau:=|c_{M,\tau}^\varpi|,
\]
so
\[
        d_{M,\tau}^{\mathrm{bench}}\neq0
        \qquad
        \text{for all sufficiently small }\tau,
\]
and
\[
        \left|
                d_{M,\tau}^{\mathrm{bench}}
        \right|
        \sim
        \left|
                \bar c_{M,\tau}^{\QQ,\star}
        \right|
        \sim
        |c_{M,\tau}^\varpi|.
\]
If, moreover,
\[
        D_{M,\tau}^\varpi=0,
        \qquad
        \mathfrak D_{M,\tau}^{\mathrm{comp}}=0,
\]
then
\[
        \bar R_{M,\tau}^{\QQ,\star}=0,
        \qquad
        c_{M,\tau}^\varpi
        =
        a_{M,\tau}^\varpi\bar c_{M,\tau}^{\QQ,\star}
        \qquad
        \text{for all sufficiently small }\tau.
\]
\end{proposition}

\begin{proof}
Set
\[
        e_{M,\tau}^{\mathrm{comp}}
        :=
        \bigl(\mathfrak D_{M,\tau}^{\mathrm{comp}}\bigr)^{1/2}.
\]
Theorem~\ref{thm:gaussian-common-profile} gives
\[
        \left|
                c_{M,\tau}^\varpi
                -
                a_{M,\tau}^\varpi \bar c_{M,\tau}^{\QQ,\star}
        \right|
        \le
        e_{M,\tau}^{\mathrm{comp}}.
\]
Hence
\[
        |a_{M,\tau}^\varpi|
        \left|
                \bar c_{M,\tau}^{\QQ,\star}
        \right|
        \ge
        |c_{M,\tau}^\varpi|
        -
        e_{M,\tau}^{\mathrm{comp}},
\]
and
\[
        |a_{M,\tau}^\varpi|
        \left|
                \bar c_{M,\tau}^{\QQ,\star}
        \right|
        \le
        |c_{M,\tau}^\varpi|
        +
        e_{M,\tau}^{\mathrm{comp}}.
\]
Because
\[
        e_{M,\tau}^{\mathrm{comp}}
        =
        o\!\left(
                |c_{M,\tau}^\varpi|
        \right),
\]
and \(|a_{M,\tau}^\varpi|\asymp1\), it follows that
\[
        \left|
                \bar c_{M,\tau}^{\QQ,\star}
        \right|
        \asymp
        |c_{M,\tau}^\varpi|.
\]
In the proof of Theorem~\ref{thm:gaussian-common-profile}, one has
\[
        D_{M,\tau}^\varpi
        =
        a_{M,\tau}^\varpi \bar R_{M,\tau}^{\QQ,\star}
        +
        \widetilde E_{M,\tau}^{\mathrm{comp},\QQ}
        -
        \langle \widetilde E_{M,\tau}^{\mathrm{comp},\QQ},\Lambda_M^\star\rangle_{L^2(\QQ)}
        \Lambda_M^\star,
\]
with
\[
        \|\widetilde E_{M,\tau}^{\mathrm{comp},\QQ}\|_{L^2(\QQ)}
        \le
        e_{M,\tau}^{\mathrm{comp}}.
\]
Therefore
\[
        |a_{M,\tau}^\varpi|
        \|\bar R_{M,\tau}^{\QQ,\star}\|_{L^2(\QQ)}
        \le
        \|D_{M,\tau}^\varpi\|_{L^2(\QQ)}
        +
        e_{M,\tau}^{\mathrm{comp}}
        =
        o\!\left(
                |c_{M,\tau}^\varpi|
        \right),
\]
so
\[
        \|\bar R_{M,\tau}^{\QQ,\star}\|_{L^2(\QQ)}
        =
        o\!\left(
                |c_{M,\tau}^\varpi|
        \right)
        =
        o\!\left(
                \left|
                        \bar c_{M,\tau}^{\QQ,\star}
                \right|
        \right).
\]
If, in addition,
\[
        \mathfrak d_{M,\tau}^{\varpi,\mathrm{bench}}
        =
        o\!\left(
                |c_{M,\tau}^\varpi|
        \right),
\]
then
\[
        \|\bar R_{M,\tau}^{\QQ,\star}\|_{L^2(\QQ)}
        =
        o(a_\tau),
        \qquad
        \mathfrak d_{M,\tau}^{\varpi,\mathrm{bench}}
        =
        o(a_\tau),
\]
with \(a_\tau=|c_{M,\tau}^\varpi|\), while the asymptotic equivalence
\[
        \left|
                \bar c_{M,\tau}^{\QQ,\star}
        \right|
        \asymp
        a_\tau
\]
provides the required lower bound on \(|\bar c_{M,\tau}^{\QQ,\star}|\).  Proposition~\ref{prop:lambda-generator-separation-loading}
therefore gives the stated nonvanishing and asymptotic equivalences for
\(d_{M,\tau}^{\mathrm{bench}}\).  If, moreover,
\[
        D_{M,\tau}^\varpi=0,
        \qquad
        \mathfrak D_{M,\tau}^{\mathrm{comp}}=0,
\]
then
\[
        \widetilde E_{M,\tau}^{\mathrm{comp},\QQ}=0,
\]
so the displayed decomposition of \(D_{M,\tau}^\varpi\) reduces to
\[
        0
        =
        a_{M,\tau}^\varpi \bar R_{M,\tau}^{\QQ,\star}.
\]
Since \(a_{M,\tau}^\varpi\neq0\), this gives
\[
        \bar R_{M,\tau}^{\QQ,\star}=0.
\]
The exact identity
\[
        c_{M,\tau}^\varpi
        =
        a_{M,\tau}^\varpi\bar c_{M,\tau}^{\QQ,\star}
\]
is the exact branch of Theorem~\ref{thm:gaussian-common-profile}.
\end{proof}

\begin{corollary}[Asymptotic benchmark-profile compatibility with the aggregate factor]\label{cor:gaussian-benchmark-direction}
Under the notation of Theorem~\ref{thm:gaussian-common-profile}, define the centered
benchmark-profile direction by
\[
        \Lambda_{M,\tau}^{\mathrm{bench}}
        :=
        \begin{cases}
                \widetilde L_{\bar Z_{M,\tau}}^\QQ
                \big/
                \|\widetilde L_{\bar Z_{M,\tau}}^\QQ\|_{L^2(\QQ)},
                &\text{if }\widetilde L_{\bar Z_{M,\tau}}^\QQ\neq0,\\[0.75em]
                0,&\text{if }\widetilde L_{\bar Z_{M,\tau}}^\QQ=0.
        \end{cases}
\]
If
\[
        \bar c_{M,\tau}^{\QQ,\star}\neq0
        \quad\text{for all sufficiently small }\tau,
        \qquad
        \|\bar R_{M,\tau}^{\QQ,\star}\|_{L^2(\QQ)}
        =
        o\!\left(|\bar c_{M,\tau}^{\QQ,\star}|\right),
\]
then
\[
        \left\|
                \Lambda_{M,\tau}^{\mathrm{bench}}
                -
                \frac{\bar c_{M,\tau}^{\QQ,\star}}{|\bar c_{M,\tau}^{\QQ,\star}|}\,
                \Lambda_M^\star
        \right\|_{L^2(\QQ)}
        \to 0.
\]
If, in addition, \(\bar R_{M,\tau}^{\QQ,\star}=0\), then
\[
        \Lambda_{M,\tau}^{\mathrm{bench}}
        =
        \frac{\bar c_{M,\tau}^{\QQ,\star}}{|\bar c_{M,\tau}^{\QQ,\star}|}\,
        \Lambda_M^\star
        \qquad
        \text{exactly}
\]
whenever \(\bar c_{M,\tau}^{\QQ,\star}\neq0\).
\end{corollary}

\begin{proof}
The factor decomposition of \(\widetilde L_{\bar Z_{M,\tau}}^\QQ\) is
\[
        \widetilde L_{\bar Z_{M,\tau}}^\QQ
        =
        \bar c_{M,\tau}^{\QQ,\star}\Lambda_M^\star
        +
        \bar R_{M,\tau}^{\QQ,\star}.
\]
If
\(
        \|\bar R_{M,\tau}^{\QQ,\star}\|_{L^2(\QQ)}
        =
        o(|\bar c_{M,\tau}^{\QQ,\star}|)
\),
then
\[
        \|\widetilde L_{\bar Z_{M,\tau}}^\QQ\|_{L^2(\QQ)}
        \sim
        |\bar c_{M,\tau}^{\QQ,\star}|,
\]
and the normalized-direction claim follows exactly as in the proof of
Theorem~\ref{thm:option-factor}.  If \(\bar R_{M,\tau}^{\QQ,\star}=0\), the orthogonal residual in
the benchmark-loading decomposition vanishes, so the exact identity follows from the displayed
projection formula.
\end{proof}

\begin{proposition}[Benchmark-loading coefficient lower bounds and residual smallness force benchmark/generator separation]\label{prop:lambda-generator-separation-loading}
Under the notation of Theorem~\ref{thm:gaussian-common-profile}, write
\[
        \widetilde L_{\bar Z_{M,\tau}}^\QQ
        =
        \bar c_{M,\tau}^{\QQ,\star}\Lambda_M^\star
        +
        \bar R_{M,\tau}^{\QQ,\star},
        \qquad
        \bar R_{M,\tau}^{\QQ,\star}\perp_{L^2(\QQ)}\Lambda_M^\star.
\]
Suppose there exists a deterministic positive scale \(a_\tau\) and a constant \(c_M>0\) such that
\[
        \left|
                \bar c_{M,\tau}^{\QQ,\star}
        \right|
        \ge
        c_M a_\tau
        \qquad
        \text{for all sufficiently small }\tau,
\]
and
\[
        \|\bar R_{M,\tau}^{\QQ,\star}\|_{L^2(\QQ)}
        =
        o(a_\tau).
\]
Then
\[
        \left\|
                \widetilde L_{\bar Z_{M,\tau}}^\QQ
        \right\|_{L^2(\QQ)}
        \ge
        c_M a_\tau
        \qquad
        \text{for all sufficiently small }\tau,
\]
and
\[
        \left\|
                \widetilde L_{\bar Z_{M,\tau}}^\QQ
        \right\|_{L^2(\QQ)}
        \sim
        \left|
                \bar c_{M,\tau}^{\QQ,\star}
        \right|.
\]
If, in addition,
\[
        \mathfrak d_{M,\tau}^{\varpi,\mathrm{bench}}
        =
        o(a_\tau),
\]
then the hypotheses of Proposition~\ref{prop:lambda-generator-separation} hold, and therefore
\[
        d_{M,\tau}^{\mathrm{bench}}\neq0
        \qquad
        \text{for all sufficiently small }\tau,
\]
with
\[
        \left|
                d_{M,\tau}^{\mathrm{bench}}
        \right|
        \sim
        \left\|
                \widetilde L_{\bar Z_{M,\tau}}^\QQ
        \right\|_{L^2(\QQ)}
        \sim
        \left|
                \bar c_{M,\tau}^{\QQ,\star}
        \right|.
\]
If, moreover,
\[
        \bar R_{M,\tau}^{\QQ,\star}=0,
        \qquad
        \mathfrak d_{M,\tau}^{\varpi,\mathrm{bench}}=0,
        \qquad
        \bar c_{M,\tau}^{\QQ,\star}\neq0
        \qquad
        \text{for all sufficiently small }\tau,
\]
then
\[
        \left|
                d_{M,\tau}^{\mathrm{bench}}
        \right|
        =
        \left|
                \bar c_{M,\tau}^{\QQ,\star}
        \right|
        \qquad
        \text{for all sufficiently small }\tau.
\]
\end{proposition}

\begin{proof}
By orthogonality of \(\bar R_{M,\tau}^{\QQ,\star}\) and \(\Lambda_M^\star\),
\[
        \left\|
                \widetilde L_{\bar Z_{M,\tau}}^\QQ
        \right\|_{L^2(\QQ)}^2
        =
        \left|
                \bar c_{M,\tau}^{\QQ,\star}
        \right|^2
        +
        \|\bar R_{M,\tau}^{\QQ,\star}\|_{L^2(\QQ)}^2.
\]
Therefore
\[
        \left\|
                \widetilde L_{\bar Z_{M,\tau}}^\QQ
        \right\|_{L^2(\QQ)}
        \ge
        \left|
                \bar c_{M,\tau}^{\QQ,\star}
        \right|
        \ge
        c_M a_\tau,
\]
which proves the displayed lower bound and, together with
\(
        \|\bar R_{M,\tau}^{\QQ,\star}\|_{L^2(\QQ)}=o(a_\tau)
\),
also gives
\[
        \left\|
                \widetilde L_{\bar Z_{M,\tau}}^\QQ
        \right\|_{L^2(\QQ)}
        \sim
        \left|
                \bar c_{M,\tau}^{\QQ,\star}
        \right|.
\]
If, in addition,
\[
        \mathfrak d_{M,\tau}^{\varpi,\mathrm{bench}}
        =
        o(a_\tau),
\]
then
\[
        \mathfrak d_{M,\tau}^{\varpi,\mathrm{bench}}
        =
        o\!\left(
                \left\|
                        \widetilde L_{\bar Z_{M,\tau}}^\QQ
                \right\|_{L^2(\QQ)}
        \right),
\]
so Proposition~\ref{prop:lambda-generator-separation} applies and gives the displayed
nonvanishing and asymptotic equivalences for \(d_{M,\tau}^{\mathrm{bench}}\).
If, moreover,
\[
        \bar R_{M,\tau}^{\QQ,\star}=0,
        \qquad
        \mathfrak d_{M,\tau}^{\varpi,\mathrm{bench}}=0,
        \qquad
        \bar c_{M,\tau}^{\QQ,\star}\neq0,
\]
then
\[
        \left\|
                \widetilde L_{\bar Z_{M,\tau}}^\QQ
        \right\|_{L^2(\QQ)}
        =
        \left|
                \bar c_{M,\tau}^{\QQ,\star}
        \right|
\]
and Proposition~\ref{prop:lambda-generator-separation} gives
\[
        \left|
                d_{M,\tau}^{\mathrm{bench}}
        \right|
        =
        \left\|
                \widetilde L_{\bar Z_{M,\tau}}^\QQ
        \right\|_{L^2(\QQ)},
\]
which proves the exact identity.
\end{proof}

\begin{definition}[Benchmark-generated centered \(\varpi\)-profile generator]\label{def:lambda-generator}
Under the notation of Theorem~\ref{thm:gaussian-common-profile}, define the centered
\(\varpi\)-shadow of the centered benchmark-profile loading by
\[
        \Theta_{M,\tau}^{\varpi,\mathrm{bench}}
        :=
        C_{\cL_T}^\QQ(\widetilde L_{\bar Z_{M,\tau}}^\QQ)
        -
        C_{\cG_T}^\QQ\!\left(
                C_{\cL_T}^\QQ(\widetilde L_{\bar Z_{M,\tau}}^\QQ)
        \right)
        \in \cM_\QQ^\circ.
\]
Define the benchmark-generated centered \(\varpi\)-profile generator by
\[
        \Gamma_{M,\tau}^\varpi
        :=
        \begin{cases}
                \Theta_{M,\tau}^{\varpi,\mathrm{bench}}
                \big/
                \|\Theta_{M,\tau}^{\varpi,\mathrm{bench}}\|_{L^2(\QQ)},
                &\text{if }\Theta_{M,\tau}^{\varpi,\mathrm{bench}}\neq0,\\[0.75em]
                0,&\text{if }\Theta_{M,\tau}^{\varpi,\mathrm{bench}}=0.
        \end{cases}
\]
and the associated canonical benchmark decomposition
\[
        \widetilde L_{\bar Z_{M,\tau}}^\QQ
        =
        d_{M,\tau}^{\mathrm{bench}}\Gamma_{M,\tau}^\varpi
        +
        E_{M,\tau}^{\mathrm{bench},\varpi},
\]
where
\[
        d_{M,\tau}^{\mathrm{bench}}
        :=
        \left\langle
                \widetilde L_{\bar Z_{M,\tau}}^\QQ,
                \Gamma_{M,\tau}^\varpi
        \right\rangle_{L^2(\QQ)},
        \qquad
        E_{M,\tau}^{\mathrm{bench},\varpi}
        :=
        \widetilde L_{\bar Z_{M,\tau}}^\QQ
        -
        d_{M,\tau}^{\mathrm{bench}}\Gamma_{M,\tau}^\varpi.
 \]
Then
\[
        E_{M,\tau}^{\mathrm{bench},\varpi}\in \cM_\QQ^\circ,
        \qquad
        E_{M,\tau}^{\mathrm{bench},\varpi}
        \perp_{L^2(\QQ)}
        \Gamma_{M,\tau}^\varpi.
\]
\end{definition}

\begin{proposition}[Canonical aggregate-factor benchmark decomposition relative to the benchmark-generated \(\varpi\)-profile generator]\label{prop:lambda-generator-qm}
Under the notation of Definition~\ref{def:lambda-generator}, define the benchmark
\(\varpi\)-generation defect by
\[
        \mathfrak d_{M,\tau}^{\varpi,\mathrm{bench}}
        :=
        \left\|
                \widetilde L_{\bar Z_{M,\tau}}^\QQ
                -
                \Theta_{M,\tau}^{\varpi,\mathrm{bench}}
        \right\|_{L^2(\QQ)},
\]
the canonical aggregate-factor coefficient relative to the benchmark-generated centered
\(\varpi\)-profile generator by
\[
        g_{M,\tau}^{\varpi,\star}
        :=
        \left\langle
                \Lambda_M^\star,
                \Gamma_{M,\tau}^\varpi
        \right\rangle_{L^2(\QQ)},
\]
and the canonical aggregate-factor residual by
\[
        H_{M,\tau}^{\varpi,\star}
        :=
        \Lambda_M^\star
        -
        g_{M,\tau}^{\varpi,\star}\Gamma_{M,\tau}^\varpi.
\]
Then
\[
        H_{M,\tau}^{\varpi,\star}\in \cM_\QQ^\circ,
        \qquad
        H_{M,\tau}^{\varpi,\star}
        \perp_{L^2(\QQ)}
        \Gamma_{M,\tau}^\varpi.
\]
Define the aggregate-factor \(\varpi\)-generation defect by
\[
        \mathfrak d_{M,\tau}^{\varpi,\star}
        :=
        \|H_{M,\tau}^{\varpi,\star}\|_{L^2(\QQ)}.
\]
Then
\[
        \|E_{M,\tau}^{\mathrm{bench},\varpi}\|_{L^2(\QQ)}
        \le
        \mathfrak d_{M,\tau}^{\varpi,\mathrm{bench}}.
\]
Moreover, under Proposition~\ref{prop:q-qm-factor},
\[
        \widetilde L_{Q_M}^\QQ
        =
        \alpha_{Q_M}^{\QQ,\star}\Lambda_M^\star,
\]
so the canonical market-variance coefficient and residual relative to
\(\Gamma_{M,\tau}^\varpi\) are
\[
        d_{M,\tau}^{\mathrm{var}}
        :=
        \left\langle
                \widetilde L_{Q_M}^\QQ,
                \Gamma_{M,\tau}^\varpi
        \right\rangle_{L^2(\QQ)}
        =
        \alpha_{Q_M}^{\QQ,\star}g_{M,\tau}^{\varpi,\star},
\]
\[
        E_{M,\tau}^{\mathrm{var},\varpi}
        :=
        \widetilde L_{Q_M}^\QQ
        -
        d_{M,\tau}^{\mathrm{var}}\Gamma_{M,\tau}^\varpi
        =
        \alpha_{Q_M}^{\QQ,\star}H_{M,\tau}^{\varpi,\star},
\]
and therefore
\[
        \|E_{M,\tau}^{\mathrm{var},\varpi}\|_{L^2(\QQ)}
        =
        \alpha_{Q_M}^{\QQ,\star}\mathfrak d_{M,\tau}^{\varpi,\star}.
\]
If \(g_{M,\tau}^{\varpi,\star}\neq0\) and \(\alpha_{Q_M}^{\QQ,\star}\neq0\), then
\[
        \frac{\|E_{M,\tau}^{\mathrm{var},\varpi}\|_{L^2(\QQ)}}{|d_{M,\tau}^{\mathrm{var}}|}
        =
        \frac{\mathfrak d_{M,\tau}^{\varpi,\star}}{|g_{M,\tau}^{\varpi,\star}|}.
\]
\end{proposition}

\begin{proof}
Because \(\Theta_{M,\tau}^{\varpi,\mathrm{bench}}\in \mathrm{span}\{\Gamma_{M,\tau}^\varpi\}\)
by Definition~\ref{def:lambda-generator}, the orthogonal projection property gives
\[
        \|E_{M,\tau}^{\mathrm{bench},\varpi}\|_{L^2(\QQ)}
        =
        \dist\!\left(
                \widetilde L_{\bar Z_{M,\tau}}^\QQ,
                \mathrm{span}\{\Gamma_{M,\tau}^\varpi\}
        \right)
        \le
        \mathfrak d_{M,\tau}^{\varpi,\mathrm{bench}}.
\]
The definitions of \(g_{M,\tau}^{\varpi,\star}\) and \(H_{M,\tau}^{\varpi,\star}\) make
\(H_{M,\tau}^{\varpi,\star}\) the orthogonal residual of \(\Lambda_M^\star\) relative to
\(\mathrm{span}\{\Gamma_{M,\tau}^\varpi\}\), which yields the displayed orthogonality.  The
formulas for \(d_{M,\tau}^{\mathrm{var}}\) and \(E_{M,\tau}^{\mathrm{var},\varpi}\) follow by
substituting \(\widetilde L_{Q_M}^\QQ=\alpha_{Q_M}^{\QQ,\star}\Lambda_M^\star\) from
Proposition~\ref{prop:q-qm-factor} into the definitions of the coefficient and residual relative
to \(\Gamma_{M,\tau}^\varpi\).  Orthogonality gives the Pythagorean norm identity, and the ratio
identity follows by substituting
\(d_{M,\tau}^{\mathrm{var}}=\alpha_{Q_M}^{\QQ,\star}g_{M,\tau}^{\varpi,\star}\) when both denominators are nonzero.
\end{proof}

\begin{proposition}[Benchmark-profile / generator norm separation forces a nontrivial benchmark coefficient]\label{prop:lambda-generator-separation}
Assume Definition~\ref{def:lambda-generator} and
Proposition~\ref{prop:lambda-generator-qm}.  If
\[
        \mathfrak d_{M,\tau}^{\varpi,\mathrm{bench}}
        <
        \left\|
                \widetilde L_{\bar Z_{M,\tau}}^\QQ
        \right\|_{L^2(\QQ)}
        \qquad
        \text{for all sufficiently small }\tau,
\]
then
\[
        d_{M,\tau}^{\mathrm{bench}}\neq0
        \qquad
        \text{for all sufficiently small }\tau.
\]
If, moreover,
\[
        \mathfrak d_{M,\tau}^{\varpi,\mathrm{bench}}
        =
        o\!\left(
                \left\|
                        \widetilde L_{\bar Z_{M,\tau}}^\QQ
                \right\|_{L^2(\QQ)}
        \right),
\]
then
\[
        \left|
                d_{M,\tau}^{\mathrm{bench}}
        \right|
        \sim
        \left\|
                \widetilde L_{\bar Z_{M,\tau}}^\QQ
        \right\|_{L^2(\QQ)}.
\]
\end{proposition}

\begin{proof}
If \(\Gamma_{M,\tau}^{\varpi}=0\), then Definition~\ref{def:lambda-generator} gives
\[
        d_{M,\tau}^{\mathrm{bench}}=0,
        \qquad
        E_{M,\tau}^{\mathrm{bench},\varpi}
        =
        \widetilde L_{\bar Z_{M,\tau}}^\QQ.
\]
Proposition~\ref{prop:lambda-generator-qm} would then imply
\[
        \left\|
                \widetilde L_{\bar Z_{M,\tau}}^\QQ
        \right\|_{L^2(\QQ)}
        =
        \|E_{M,\tau}^{\mathrm{bench},\varpi}\|_{L^2(\QQ)}
        \le
        \mathfrak d_{M,\tau}^{\varpi,\mathrm{bench}},
\]
contradicting the displayed strict inequality.  Therefore
\[
        \Gamma_{M,\tau}^{\varpi}\neq0
        \qquad
        \text{for all sufficiently small }\tau.
\]
Since \(\Gamma_{M,\tau}^{\varpi}\) is then a unit vector, the orthogonal decomposition in
Definition~\ref{def:lambda-generator} yields
\[
        \left\|
                \widetilde L_{\bar Z_{M,\tau}}^\QQ
        \right\|_{L^2(\QQ)}^2
        =
        \left|
                d_{M,\tau}^{\mathrm{bench}}
        \right|^2
        +
        \|E_{M,\tau}^{\mathrm{bench},\varpi}\|_{L^2(\QQ)}^2.
\]
Together with Proposition~\ref{prop:lambda-generator-qm}, this gives
\[
        \left|
                d_{M,\tau}^{\mathrm{bench}}
        \right|^2
        \ge
        \left\|
                \widetilde L_{\bar Z_{M,\tau}}^\QQ
        \right\|_{L^2(\QQ)}^2
        -
        \bigl(
                \mathfrak d_{M,\tau}^{\varpi,\mathrm{bench}}
        \bigr)^2.
\]
The displayed strict inequality makes the lower bound positive, which gives the first claim.  If,
moreover,
\[
        \mathfrak d_{M,\tau}^{\varpi,\mathrm{bench}}
        =
        o\!\left(
                \left\|
                        \widetilde L_{\bar Z_{M,\tau}}^\QQ
                \right\|_{L^2(\QQ)}
        \right),
\]
then
\[
        \frac{
                \|E_{M,\tau}^{\mathrm{bench},\varpi}\|_{L^2(\QQ)}
        }{
                \|\widetilde L_{\bar Z_{M,\tau}}^\QQ\|_{L^2(\QQ)}
        }
        \le
        \frac{
                \mathfrak d_{M,\tau}^{\varpi,\mathrm{bench}}
        }{
                \|\widetilde L_{\bar Z_{M,\tau}}^\QQ\|_{L^2(\QQ)}
        }
        \to
        0,
\]
and the orthogonal norm identity yields the asymptotic equivalence
\[
        \left|
                d_{M,\tau}^{\mathrm{bench}}
        \right|
        \sim
        \left\|
                \widetilde L_{\bar Z_{M,\tau}}^\QQ
        \right\|_{L^2(\QQ)}.
\]
\end{proof}

\begin{proposition}[\(Q_M\)-side \(\varpi\)-shadow control of the aggregate-factor defect]\label{prop:lambda-shadow}
Assume Proposition~\ref{prop:q-qm-factor} and Definition~\ref{def:lambda-generator}.  If
\[
        \alpha_{Q_M}^{\QQ,\star}\neq0,
\]
define the centered \(Q_M\)-side \(\varpi\)-shadow by
\[
        \Theta_{M,\tau}^{\varpi,Q}
        :=
        C_{\cL_T}^\QQ(\widetilde L_{Q_M}^\QQ)
        -
        C_{\cG_T}^\QQ\!\left(
                C_{\cL_T}^\QQ(\widetilde L_{Q_M}^\QQ)
        \right)
        \in
        \cM_\QQ^\circ
\]
and the normalized \(Q_M\)-side \(\varpi\)-shadow defect by
\[
        \mathfrak s_{M,\tau}^{\varpi,Q}
        :=
        \frac{
                \|\widetilde L_{Q_M}^\QQ-\Theta_{M,\tau}^{\varpi,Q}\|_{L^2(\QQ)}
                +
                \dist\!\left(
                        \Theta_{M,\tau}^{\varpi,Q},
                        \mathrm{span}\{\Gamma_{M,\tau}^\varpi\}
                \right)
        }{\alpha_{Q_M}^{\QQ,\star}}.
\]
Then
\[
        \mathfrak d_{M,\tau}^{\varpi,\star}
        \le
        \mathfrak s_{M,\tau}^{\varpi,Q}.
\]
Consequently, if
\[
        \|\widetilde L_{Q_M}^\QQ-\Theta_{M,\tau}^{\varpi,Q}\|_{L^2(\QQ)}
        =
        o\!\left(\alpha_{Q_M}^{\QQ,\star}\right),
        \qquad
        \dist\!\left(
                \Theta_{M,\tau}^{\varpi,Q},
                \mathrm{span}\{\Gamma_{M,\tau}^\varpi\}
        \right)
        =
        o\!\left(\alpha_{Q_M}^{\QQ,\star}\right),
\]
then
\[
        \mathfrak d_{M,\tau}^{\varpi,\star}\to 0.
\]
If both displayed terms vanish identically, then
\[
        \mathfrak d_{M,\tau}^{\varpi,\star}=0.
\]
\end{proposition}

\begin{proof}
Because
\[
        \widetilde L_{Q_M}^\QQ
        =
        \alpha_{Q_M}^{\QQ,\star}\Lambda_M^\star
\]
by Proposition~\ref{prop:q-qm-factor}, one has
\[
        \alpha_{Q_M}^{\QQ,\star}\mathfrak d_{M,\tau}^{\varpi,\star}
        =
        \dist\!\left(
                \widetilde L_{Q_M}^\QQ,
                \mathrm{span}\{\Gamma_{M,\tau}^\varpi\}
        \right).
\]
Applying the triangle inequality through \(\Theta_{M,\tau}^{\varpi,Q}\) yields
\[
        \dist\!\left(
                \widetilde L_{Q_M}^\QQ,
                \mathrm{span}\{\Gamma_{M,\tau}^\varpi\}
        \right)
        \le
        \|\widetilde L_{Q_M}^\QQ-\Theta_{M,\tau}^{\varpi,Q}\|_{L^2(\QQ)}
        +
        \dist\!\left(
                \Theta_{M,\tau}^{\varpi,Q},
                \mathrm{span}\{\Gamma_{M,\tau}^\varpi\}
        \right),
\]
which gives the first claim after dividing by \(\alpha_{Q_M}^{\QQ,\star}\).  If the two normalized
terms on the right vanish, the quotient defect also vanishes.  If both terms are zero, the
distance to \(\mathrm{span}\{\Gamma_{M,\tau}^\varpi\}\) is zero, which is the exact statement.
\end{proof}

\begin{proposition}[Explicit decomposition of the \(Q_M\)-side \(\varpi\)-shadow defect]\label{prop:lambda-shadow-explicit}
Assume Assumption~\ref{ass:q-visible-aggregate}, Proposition~\ref{prop:q-qm-explicit}, and
Definition~\ref{def:lambda-generator}.  Write
\[
        \zeta(t):=(w^\top v)\varpi(t),
        \qquad
        R_{\eta,M}^\QQ
        :=
        \int_0^T \EE_\QQ[\eta^\QQ(t)^2\mid \cH_\varpi]\,\dd t.
\]
Define the pure \(\varpi\)-benchmark
\[
        B_{M}^{\varpi,Q}
        :=
        \int_0^T \zeta(t)^2\,\dd t
        +
        C_{\cL_T}^\QQ(R_{\eta,M}^\QQ)
        \in L^2(\QQ,\cL_T),
\]
its centered image
\[
        \widetilde B_{M}^{\varpi,Q}
        :=
        B_{M}^{\varpi,Q}
        -
        C_{\cG_T}^\QQ(B_{M}^{\varpi,Q})
        \in \cM_\QQ^\circ,
\]
the mixed \(\varpi\)-visibility remainder
\[
        S_{M}^{\varpi,Q}
        :=
        2\int_0^T M^{\mathrm{vis},\QQ}(t)\zeta(t)\,\dd t
        +
        R_{\eta,M}^\QQ
        -
        C_{\cL_T}^\QQ(R_{\eta,M}^\QQ),
\]
and the visible-shadow defect
\[
        \Delta_{M}^{\varpi,\mathrm{vis}}
        :=
        \left\|
                C_{\cL_T}^\QQ\!\left(
                        C_{\cG_T}^\QQ(B_{M}^{\varpi,Q})
                \right)
                -
                C_{\cG_T}^\QQ\!\left(
                        C_{\cL_T}^\QQ\!\left(
                                C_{\cG_T}^\QQ(B_{M}^{\varpi,Q})
                        \right)
                \right)
        \right\|_{L^2(\QQ)}.
\]
Then
\[
        L_{Q_M}^\QQ
        =
        \int_0^T \bigl(M^{\mathrm{vis},\QQ}(t)\bigr)^2\,\dd t
        +
        B_{M}^{\varpi,Q}
        +
        S_{M}^{\varpi,Q},
\]
and therefore
\[
        \widetilde L_{Q_M}^\QQ
        =
        \widetilde B_{M}^{\varpi,Q}
        +
        \widetilde S_{M}^{\varpi,Q},
        \qquad
        \widetilde S_{M}^{\varpi,Q}
        :=
        S_{M}^{\varpi,Q}
        -
        C_{\cG_T}^\QQ(S_{M}^{\varpi,Q}).
\]
If \(\alpha_{Q_M}^{\QQ,\star}\neq0\), then
\[
        \mathfrak s_{M,\tau}^{\varpi,Q}
        \le
        \frac{
                5\,\|\widetilde S_{M}^{\varpi,Q}\|_{L^2(\QQ)}
                +
                2\,\Delta_{M}^{\varpi,\mathrm{vis}}
                +
                \dist\!\left(
                        \widetilde B_{M}^{\varpi,Q},
                        \mathrm{span}\{\Gamma_{M,\tau}^\varpi\}
                \right)
        }{\alpha_{Q_M}^{\QQ,\star}}.
\]
Consequently, if
\[
        \|\widetilde S_{M}^{\varpi,Q}\|_{L^2(\QQ)}
        =
        o\!\left(\alpha_{Q_M}^{\QQ,\star}\right),
        \qquad
        \Delta_{M}^{\varpi,\mathrm{vis}}
        =
        o\!\left(\alpha_{Q_M}^{\QQ,\star}\right),
\]
and
\[
        \dist\!\left(
                \widetilde B_{M}^{\varpi,Q},
                \mathrm{span}\{\Gamma_{M,\tau}^\varpi\}
        \right)
        =
        o\!\left(\alpha_{Q_M}^{\QQ,\star}\right),
\]
then
\[
        \mathfrak s_{M,\tau}^{\varpi,Q}=o(1).
\]
If all three displayed defects vanish identically, then
\[
        \mathfrak s_{M,\tau}^{\varpi,Q}=0.
\]
\end{proposition}

\begin{proof}
By Proposition~\ref{prop:q-qm-explicit},
\[
        L_{Q_M}^\QQ
        =
        \int_0^T \bigl(M^{\mathrm{vis},\QQ}(t)+\zeta(t)\bigr)^2\,\dd t
        +
        R_{\eta,M}^\QQ.
\]
Expanding the square and regrouping the \(\cL_T\)-measurable part gives the first displayed
decomposition of \(L_{Q_M}^\QQ\).  Since
\(\int_0^T(M^{\mathrm{vis},\QQ}(t))^2\,\dd t\in L^2(\QQ,\cG_T)\), subtracting
\(C_{\cG_T}^\QQ(L_{Q_M}^\QQ)\) yields
\[
        \widetilde L_{Q_M}^\QQ
        =
        \widetilde B_{M}^{\varpi,Q}
        +
        \widetilde S_{M}^{\varpi,Q}.
\]
Write
\[
        \Pi_\varpi^\QQ(Y)
        :=
        C_{\cL_T}^\QQ(Y)
        -
        C_{\cG_T}^\QQ\!\left(C_{\cL_T}^\QQ(Y)\right).
\]
Then \(\Theta_{M,\tau}^{\varpi,Q}=\Pi_\varpi^\QQ(\widetilde L_{Q_M}^\QQ)\).  Because
\(B_{M}^{\varpi,Q}\in L^2(\QQ,\cL_T)\),
\[
        \Pi_\varpi^\QQ(B_{M}^{\varpi,Q})
        =
        \widetilde B_{M}^{\varpi,Q},
\]
and hence
\[
        \Pi_\varpi^\QQ(\widetilde B_{M}^{\varpi,Q})
        =
        \widetilde B_{M}^{\varpi,Q}
        -
        \Pi_\varpi^\QQ\!\left(
                C_{\cG_T}^\QQ(B_{M}^{\varpi,Q})
        \right).
\]
Therefore
\[
        \Theta_{M,\tau}^{\varpi,Q}
        =
        \widetilde B_{M}^{\varpi,Q}
        -
        \Pi_\varpi^\QQ\!\left(
                C_{\cG_T}^\QQ(B_{M}^{\varpi,Q})
        \right)
        +
        \Pi_\varpi^\QQ(\widetilde S_{M}^{\varpi,Q}).
\]
Since conditional expectation is an \(L^2\)-contraction,
\[
        \|\Pi_\varpi^\QQ(\widetilde S_{M}^{\varpi,Q})\|_{L^2(\QQ)}
        \le
        2\|\widetilde S_{M}^{\varpi,Q}\|_{L^2(\QQ)}
\]
Using the displayed decomposition of \(\Theta_{M,\tau}^{\varpi,Q}\), one gets
\[
        \|\widetilde L_{Q_M}^\QQ-\Theta_{M,\tau}^{\varpi,Q}\|_{L^2(\QQ)}
        \le
        \|\widetilde S_{M}^{\varpi,Q}\|_{L^2(\QQ)}
        +
        \Delta_{M}^{\varpi,\mathrm{vis}}
        +
        \|\Pi_\varpi^\QQ(\widetilde S_{M}^{\varpi,Q})\|_{L^2(\QQ)}
        \le
        3\|\widetilde S_{M}^{\varpi,Q}\|_{L^2(\QQ)}
        +
        \Delta_{M}^{\varpi,\mathrm{vis}}.
\]
Likewise,
\[
\begin{aligned}
        \dist\!\left(
                \Theta_{M,\tau}^{\varpi,Q},
                \mathrm{span}\{\Gamma_{M,\tau}^\varpi\}
        \right)
        &\le
        \dist\!\left(
                \widetilde B_{M}^{\varpi,Q},
                \mathrm{span}\{\Gamma_{M,\tau}^\varpi\}
        \right)
        +
        \Delta_{M}^{\varpi,\mathrm{vis}}
        +
        \|\Pi_\varpi^\QQ(\widetilde S_{M}^{\varpi,Q})\|_{L^2(\QQ)} \\
        &\le
        \dist\!\left(
                \widetilde B_{M}^{\varpi,Q},
                \mathrm{span}\{\Gamma_{M,\tau}^\varpi\}
        \right)
        +
        \Delta_{M}^{\varpi,\mathrm{vis}}
        +
        2\|\widetilde S_{M}^{\varpi,Q}\|_{L^2(\QQ)}.
\end{aligned}
\]
Adding the two bounds and dividing by \(\alpha_{Q_M}^{\QQ,\star}\) gives the stated estimate for
\(\mathfrak s_{M,\tau}^{\varpi,Q}\).  Vanishing of each normalized term on the right gives the
asymptotic claim; simultaneous zero defects give the exact claim.
\end{proof}

\begin{theorem}[Common centered \(\varpi\)-profile generator implies benchmark-factor compatibility]\label{thm:lambda-generator}
Under Definition~\ref{def:lambda-generator} and
Proposition~\ref{prop:lambda-generator-qm}, let
\[
        \widetilde L_{\bar Z_{M,\tau}}^\QQ
        =
        \bar c_{M,\tau}^{\QQ,\star}\Lambda_M^\star
        +
        \bar R_{M,\tau}^{\QQ,\star}
\]
be the canonical aggregate-factor decomposition of the centered benchmark-profile loading, with
\[
        \bar R_{M,\tau}^{\QQ,\star}\perp_{L^2(\QQ)}\Lambda_M^\star.
\]
If
\[
        d_{M,\tau}^{\mathrm{bench}}\neq0,
        \qquad
        g_{M,\tau}^{\varpi,\star}\neq0
        \qquad
        \text{for all sufficiently small }\tau,
\]
define
\[
        s_{M,\tau}^{\varpi,\star}
        :=
        \frac{g_{M,\tau}^{\varpi,\star}}{|g_{M,\tau}^{\varpi,\star}|},
        \qquad
        \delta_{M,\tau}^{\varpi,\mathrm{gen}}
        :=
        \mathfrak d_{M,\tau}^{\varpi,\mathrm{bench}}
        +
        |d_{M,\tau}^{\mathrm{bench}}|
        \frac{\mathfrak d_{M,\tau}^{\varpi,\star}}{|g_{M,\tau}^{\varpi,\star}|}.
\]
Then
\[
        \left|
                \bar c_{M,\tau}^{\QQ,\star}
                -
                s_{M,\tau}^{\varpi,\star}\,d_{M,\tau}^{\mathrm{bench}}
        \right|
        \le
        \delta_{M,\tau}^{\varpi,\mathrm{gen}},
\]
and
\[
        \|\bar R_{M,\tau}^{\QQ,\star}\|_{L^2(\QQ)}
        \le
        \delta_{M,\tau}^{\varpi,\mathrm{gen}}.
\]
If, moreover,
\[
        \delta_{M,\tau}^{\varpi,\mathrm{gen}}
        =
        o\!\left(|d_{M,\tau}^{\mathrm{bench}}|\right),
\]
then the benchmark-profile direction is asymptotically aggregate-factor compatible:
\[
        \left\|
                \Lambda_{M,\tau}^{\mathrm{bench}}
                -
                \frac{d_{M,\tau}^{\mathrm{bench}}g_{M,\tau}^{\varpi,\star}}
                {|d_{M,\tau}^{\mathrm{bench}}g_{M,\tau}^{\varpi,\star}|}\,
                \Lambda_M^\star
        \right\|_{L^2(\QQ)}
        \to0.
\]
If, in addition,
\[
        \mathfrak d_{M,\tau}^{\varpi,\mathrm{bench}}=0,
        \qquad
        \mathfrak d_{M,\tau}^{\varpi,\star}=0,
\]
then
\[
        \bar c_{M,\tau}^{\QQ,\star}
        =
        s_{M,\tau}^{\varpi,\star}\,d_{M,\tau}^{\mathrm{bench}},
        \qquad
        \bar R_{M,\tau}^{\QQ,\star}=0,
\]
and therefore
\[
        \Lambda_{M,\tau}^{\mathrm{bench}}
        =
        \frac{d_{M,\tau}^{\mathrm{bench}}g_{M,\tau}^{\varpi,\star}}
        {|d_{M,\tau}^{\mathrm{bench}}g_{M,\tau}^{\varpi,\star}|}\,
        \Lambda_M^\star
        \qquad
        \text{exactly.}
\]
\end{theorem}

\begin{proof}
Set
\[
        \eta_{M,\tau}^{\varpi,\star}
        :=
        \left\|
                \Gamma_{M,\tau}^\varpi
                -
                s_{M,\tau}^{\varpi,\star}\Lambda_M^\star
        \right\|_{L^2(\QQ)}.
\]
Because \(\Gamma_{M,\tau}^\varpi\) and \(\Lambda_M^\star\) are unit vectors under the present
nondegeneracy assumptions and
\[
        \Lambda_M^\star
        =
        g_{M,\tau}^{\varpi,\star}\Gamma_{M,\tau}^\varpi
        +
        H_{M,\tau}^{\varpi,\star},
        \qquad
        H_{M,\tau}^{\varpi,\star}\perp_{L^2(\QQ)}\Gamma_{M,\tau}^\varpi,
\]
one has
\[
        1
        =
        \bigl(g_{M,\tau}^{\varpi,\star}\bigr)^2
        +
        \|H_{M,\tau}^{\varpi,\star}\|_{L^2(\QQ)}^2.
\]
Writing
\[
        s_{M,\tau}^{\varpi,\star}\Lambda_M^\star
        =
        |g_{M,\tau}^{\varpi,\star}|\,\Gamma_{M,\tau}^\varpi
        +
        s_{M,\tau}^{\varpi,\star}H_{M,\tau}^{\varpi,\star},
\]
and using again the orthogonality of \(H_{M,\tau}^{\varpi,\star}\) and
\(\Gamma_{M,\tau}^\varpi\), one gets
\[
        \bigl(\eta_{M,\tau}^{\varpi,\star}\bigr)^2
        =
        \left(
                1-|g_{M,\tau}^{\varpi,\star}|
        \right)^2
        +
        \|H_{M,\tau}^{\varpi,\star}\|_{L^2(\QQ)}^2
        =
        2\left(
                1-|g_{M,\tau}^{\varpi,\star}|
        \right)
        \le
        \frac{\|H_{M,\tau}^{\varpi,\star}\|_{L^2(\QQ)}^2}
        {\bigl(g_{M,\tau}^{\varpi,\star}\bigr)^2},
\]
so
\[
        \eta_{M,\tau}^{\varpi,\star}
        \le
        \frac{\mathfrak d_{M,\tau}^{\varpi,\star}}{|g_{M,\tau}^{\varpi,\star}|}.
\]
By Definition~\ref{def:lambda-generator},
\[
        \widetilde L_{\bar Z_{M,\tau}}^\QQ
        =
        s_{M,\tau}^{\varpi,\star}\,d_{M,\tau}^{\mathrm{bench}}\Lambda_M^\star
        +
        d_{M,\tau}^{\mathrm{bench}}
        \bigl(
                \Gamma_{M,\tau}^\varpi-s_{M,\tau}^{\varpi,\star}\Lambda_M^\star
        \bigr)
        +
        E_{M,\tau}^{\mathrm{bench},\varpi}.
\]
Taking the \(L^2(\QQ)\)-inner product against \(\Lambda_M^\star\) gives
\[
        \left|
                \bar c_{M,\tau}^{\QQ,\star}
                -
                s_{M,\tau}^{\varpi,\star}\,d_{M,\tau}^{\mathrm{bench}}
        \right|
        \le
        |d_{M,\tau}^{\mathrm{bench}}|\,\eta_{M,\tau}^{\varpi,\star}
        +
        \|E_{M,\tau}^{\mathrm{bench},\varpi}\|_{L^2(\QQ)}
        \le
        \delta_{M,\tau}^{\varpi,\mathrm{gen}}.
\]
Subtracting the \(\Lambda_M^\star\)-component from the same decomposition yields
\[
        \|\bar R_{M,\tau}^{\QQ,\star}\|_{L^2(\QQ)}
        \le
        |d_{M,\tau}^{\mathrm{bench}}|\,\eta_{M,\tau}^{\varpi,\star}
        +
        \|E_{M,\tau}^{\mathrm{bench},\varpi}\|_{L^2(\QQ)}
        \le
        \delta_{M,\tau}^{\varpi,\mathrm{gen}}.
\]
If \(\delta_{M,\tau}^{\varpi,\mathrm{gen}}=o(|d_{M,\tau}^{\mathrm{bench}}|)\), then
\[
        \bar c_{M,\tau}^{\QQ,\star}
        =
        s_{M,\tau}^{\varpi,\star}\,d_{M,\tau}^{\mathrm{bench}}
        +
        o\!\left(|d_{M,\tau}^{\mathrm{bench}}|\right),
\]
and
\[
        \|\bar R_{M,\tau}^{\QQ,\star}\|_{L^2(\QQ)}
        =
        o\!\left(|d_{M,\tau}^{\mathrm{bench}}|\right).
\]
Since \(|\bar c_{M,\tau}^{\QQ,\star}|\sim |d_{M,\tau}^{\mathrm{bench}}|\), the displayed
benchmark-direction compatibility is exactly Corollary~\ref{cor:gaussian-benchmark-direction}.  If
\[
        \mathfrak d_{M,\tau}^{\varpi,\mathrm{bench}}=0,
        \qquad
        \mathfrak d_{M,\tau}^{\varpi,\star}=0,
\]
then \(\eta_{M,\tau}^{\varpi,\star}=0\), so the exact identities follow.
\end{proof}

\begin{corollary}[Benchmark-factor compatibility from the \(Q_M\)-side \(\varpi\)-shadow]\label{cor:lambda-shadow}
Assume the hypotheses of Theorem~\ref{thm:lambda-generator} and
Propositions~\ref{prop:lambda-shadow} and \ref{prop:lambda-shadow-explicit}.  If
\[
        \alpha_{Q_M}^{\QQ,\star}\neq0,
\]
and if
\[
        d_{M,\tau}^{\mathrm{bench}}\neq0
        \qquad
        \text{for all sufficiently small }\tau,
\]
and
\[
        \mathfrak d_{M,\tau}^{\varpi,\mathrm{bench}}
        =
        o\!\left(|d_{M,\tau}^{\mathrm{bench}}|\right),
        \qquad
        \mathfrak s_{M,\tau}^{\varpi,Q}=o(1),
\]
then the benchmark-profile direction is asymptotically aggregate-factor compatible:
\[
        \left\|
                \Lambda_{M,\tau}^{\mathrm{bench}}
                -
                \frac{d_{M,\tau}^{\mathrm{bench}}g_{M,\tau}^{\varpi,\star}}
                {|d_{M,\tau}^{\mathrm{bench}}g_{M,\tau}^{\varpi,\star}|}\,
                \Lambda_M^\star
        \right\|_{L^2(\QQ)}
        \to0.
\]
In particular, it is enough that
\[
        \mathfrak d_{M,\tau}^{\varpi,\mathrm{bench}}
        =
        o\!\left(|d_{M,\tau}^{\mathrm{bench}}|\right)
\]
and
\[
        \|S_{M}^{\varpi,Q}\|_{L^2(\QQ)}
        +
        \Delta_{M}^{\varpi,\mathrm{vis}}
        +
        \dist\!\left(
                \widetilde B_{M}^{\varpi,Q},
                \mathrm{span}\{\Gamma_{M,\tau}^\varpi\}
        \right)
        =
        o\!\left(\alpha_{Q_M}^{\QQ,\star}\right).
\]
If, in addition,
\[
        \mathfrak d_{M,\tau}^{\varpi,\mathrm{bench}}=0,
        \qquad
        \mathfrak s_{M,\tau}^{\varpi,Q}=0,
\]
then the benchmark-profile direction is exactly aggregate-factor compatible.
In particular, exact compatibility follows if
\[
        \mathfrak d_{M,\tau}^{\varpi,\mathrm{bench}}=0,
        \qquad
        S_{M}^{\varpi,Q}=0,
        \qquad
        \Delta_{M}^{\varpi,\mathrm{vis}}=0,
        \qquad
        \widetilde B_{M}^{\varpi,Q}\in \mathrm{span}\{\Gamma_{M,\tau}^\varpi\}.
\]
\end{corollary}

\begin{proof}
Proposition~\ref{prop:lambda-shadow} gives
\[
        \mathfrak d_{M,\tau}^{\varpi,\star}
        \le
        \mathfrak s_{M,\tau}^{\varpi,Q}
        =
        o(1),
\]
so
\[
        1
        =
        \bigl(g_{M,\tau}^{\varpi,\star}\bigr)^2
        +
        \bigl(\mathfrak d_{M,\tau}^{\varpi,\star}\bigr)^2
\]
implies \(|g_{M,\tau}^{\varpi,\star}|\to1\), hence \(g_{M,\tau}^{\varpi,\star}\neq0\) for all
sufficiently small \(\tau\).  Therefore
\[
        \delta_{M,\tau}^{\varpi,\mathrm{gen}}
        =
        \mathfrak d_{M,\tau}^{\varpi,\mathrm{bench}}
        +
        |d_{M,\tau}^{\mathrm{bench}}|
        \frac{\mathfrak d_{M,\tau}^{\varpi,\star}}{|g_{M,\tau}^{\varpi,\star}|}
        =
        o\!\left(|d_{M,\tau}^{\mathrm{bench}}|\right),
\]
and Theorem~\ref{thm:lambda-generator} yields the asymptotic claim.  The explicit sufficient
conditions follow from Proposition~\ref{prop:lambda-shadow-explicit}.  If both displayed defects
vanish, then Proposition~\ref{prop:lambda-shadow} gives
\(\mathfrak d_{M,\tau}^{\varpi,\star}=0\), and the exact branch of
Theorem~\ref{thm:lambda-generator} applies; the exact sufficient conditions are again those of
Proposition~\ref{prop:lambda-shadow-explicit}.
\end{proof}

\begin{proposition}[Mixed \(\varpi\)-shadow remainder split]\label{prop:lambda-shadow-mixed}
Under the notation of Proposition~\ref{prop:lambda-shadow-explicit}, define
\[
        S_{M,\times}^{\varpi,Q}
        :=
        2\int_0^T M^{\mathrm{vis},\QQ}(t)\zeta(t)\,\dd t,
        \qquad
        S_{M,\eta}^{\varpi,Q}
        :=
        R_{\eta,M}^\QQ-C_{\cL_T}^\QQ(R_{\eta,M}^\QQ).
\]
Then
\[
        S_{M}^{\varpi,Q}
        =
        S_{M,\times}^{\varpi,Q}
        +
        S_{M,\eta}^{\varpi,Q}.
\]
Consequently,
\[
        \|S_{M}^{\varpi,Q}\|_{L^2(\QQ)}
        \le
        \|S_{M,\times}^{\varpi,Q}\|_{L^2(\QQ)}
        +
        \|S_{M,\eta}^{\varpi,Q}\|_{L^2(\QQ)}.
\]
If
\[
        \|S_{M,\times}^{\varpi,Q}\|_{L^2(\QQ)}
        +
        \|S_{M,\eta}^{\varpi,Q}\|_{L^2(\QQ)}
        =
        o\!\left(\alpha_{Q_M}^{\QQ,\star}\right),
\]
then
\[
        \|S_{M}^{\varpi,Q}\|_{L^2(\QQ)}
        =
        o\!\left(\alpha_{Q_M}^{\QQ,\star}\right).
\]
If
\[
        S_{M,\times}^{\varpi,Q}=0,
        \qquad
        S_{M,\eta}^{\varpi,Q}=0,
\]
then
\[
        S_{M}^{\varpi,Q}=0.
\]
\end{proposition}

\begin{proof}
Substituting the definition of \(S_{M}^{\varpi,Q}\) from
Proposition~\ref{prop:lambda-shadow-explicit} gives the identity.  Applying the triangle inequality
to the three summands gives the displayed bounds; if the summands vanish, the sum is zero.
\end{proof}

\begin{proposition}[Centered split of the mixed \(\varpi\)-shadow remainder]\label{prop:lambda-shadow-mixed-centered}
Under the notation of Proposition~\ref{prop:lambda-shadow-mixed}, define
\[
        \widetilde S_{M,\times}^{\varpi,Q}
        :=
        S_{M,\times}^{\varpi,Q}
        -
        C_{\cG_T}^\QQ(S_{M,\times}^{\varpi,Q}),
        \qquad
        \widetilde S_{M,\eta}^{\varpi,Q}
        :=
        S_{M,\eta}^{\varpi,Q}
        -
        C_{\cG_T}^\QQ(S_{M,\eta}^{\varpi,Q}).
\]
Then
\[
        \widetilde S_{M}^{\varpi,Q}
        =
        \widetilde S_{M,\times}^{\varpi,Q}
        +
        \widetilde S_{M,\eta}^{\varpi,Q},
\]
and therefore
\[
        \|\widetilde S_{M}^{\varpi,Q}\|_{L^2(\QQ)}
        \le
        \|\widetilde S_{M,\times}^{\varpi,Q}\|_{L^2(\QQ)}
        +
        \|\widetilde S_{M,\eta}^{\varpi,Q}\|_{L^2(\QQ)}.
\]
Moreover,
\[
        \|\widetilde S_{M,\times}^{\varpi,Q}\|_{L^2(\QQ)}
        \le
        2\|S_{M,\times}^{\varpi,Q}\|_{L^2(\QQ)},
        \qquad
        \|\widetilde S_{M,\eta}^{\varpi,Q}\|_{L^2(\QQ)}
        \le
        2\|S_{M,\eta}^{\varpi,Q}\|_{L^2(\QQ)}.
\]
If
\[
        S_{M,\times}^{\varpi,Q}\in L^2(\QQ,\cG_T),
        \qquad
        S_{M,\eta}^{\varpi,Q}\in L^2(\QQ,\cG_T),
\]
then
\[
        \widetilde S_{M,\times}^{\varpi,Q}=0,
        \qquad
        \widetilde S_{M,\eta}^{\varpi,Q}=0.
\]
\end{proposition}

\begin{proof}
The identity follows from
\[
        \widetilde S_{M}^{\varpi,Q}
        =
        S_{M}^{\varpi,Q}
        -
        C_{\cG_T}^\QQ(S_{M}^{\varpi,Q})
\]
and the decomposition \(S_{M}^{\varpi,Q}=S_{M,\times}^{\varpi,Q}+S_{M,\eta}^{\varpi,Q}\) from
Proposition~\ref{prop:lambda-shadow-mixed}.  The norm bound is the triangle inequality.  Since
\(C_{\cG_T}^\QQ\) is an \(L^2\)-contraction,
\[
        \|Y-C_{\cG_T}^\QQ(Y)\|_{L^2(\QQ)}
        \le
        \|Y\|_{L^2(\QQ)}+\|C_{\cG_T}^\QQ(Y)\|_{L^2(\QQ)}
        \le
        2\|Y\|_{L^2(\QQ)}
\]
for every \(Y\in L^2(\QQ,\cF_T)\), which yields the displayed bounds.  If
\(Y\in L^2(\QQ,\cG_T)\), then \(C_{\cG_T}^\QQ(Y)=Y\), so the centered residual is zero and the
exact claim follows.
\end{proof}

\begin{proposition}[Centered mixed term as visible--public projection interference]\label{prop:lambda-shadow-cross-public}
Assume Assumption~\ref{ass:q-visible-aggregate} and the notation of
Propositions~\ref{prop:q-qm-explicit}, \ref{prop:lambda-shadow-mixed}, and
\ref{prop:lambda-shadow-mixed-centered}.  Define the public projection of the synchronized aggregate
channel by
\[
        \widetilde\zeta^\QQ(t)
        :=
        \zeta(t)-C_{\cG_T}^\QQ(\zeta(t)),
        \qquad
        0\le t\le T.
\]
Then
\[
        \widetilde S_{M,\times}^{\varpi,Q}
        =
        2\int_0^T M^{\mathrm{vis},\QQ}(t)\widetilde\zeta^\QQ(t)\,\dd t.
\]
Consequently,
\[
        \|\widetilde S_{M,\times}^{\varpi,Q}\|_{L^2(\QQ)}
        \le
        2
        \left\|
                \int_0^T \bigl(M^{\mathrm{vis},\QQ}(t)\bigr)^2\,\dd t
        \right\|_{L^2(\QQ)}^{1/2}
        \left\|
                \int_0^T \bigl(\widetilde\zeta^\QQ(t)\bigr)^2\,\dd t
        \right\|_{L^2(\QQ)}^{1/2}.
\]
If
\[
        \zeta(t)\in L^2(\QQ,\cG_T)
        \qquad
        \text{for almost every }t\in[0,T],
\]
then
\[
        S_{M,\times}^{\varpi,Q}\in L^2(\QQ,\cG_T),
        \qquad
        \widetilde S_{M,\times}^{\varpi,Q}=0.
\]
\end{proposition}

\begin{proof}
Because \(M^{\mathrm{vis},\QQ}(t)\) is \(\cG_t\)-measurable for every \(t\), it is also
\(\cG_T\)-measurable.  Therefore
\[
        C_{\cG_T}^\QQ\!\left(
                M^{\mathrm{vis},\QQ}(t)\zeta(t)
        \right)
        =
        M^{\mathrm{vis},\QQ}(t)\,C_{\cG_T}^\QQ(\zeta(t))
\]
for almost every \(t\).  Applying conditional Fubini gives
\[
        C_{\cG_T}^\QQ\!\left(
                S_{M,\times}^{\varpi,Q}
        \right)
        =
        2\int_0^T
                M^{\mathrm{vis},\QQ}(t)\,
                C_{\cG_T}^\QQ(\zeta(t))\,\dd t.
\]
Subtracting this from
\[
        S_{M,\times}^{\varpi,Q}
        =
        2\int_0^T M^{\mathrm{vis},\QQ}(t)\zeta(t)\,\dd t
\]
yields the displayed identity for \(\widetilde S_{M,\times}^{\varpi,Q}\).  The \(L^2\) bound
then follows from pathwise Cauchy--Schwarz:
\[
        \left|
                \int_0^T M^{\mathrm{vis},\QQ}(t)\widetilde\zeta^\QQ(t)\,\dd t
        \right|^2
        \le
        \left(
                \int_0^T \bigl(M^{\mathrm{vis},\QQ}(t)\bigr)^2\,\dd t
        \right)
        \left(
                \int_0^T \bigl(\widetilde\zeta^\QQ(t)\bigr)^2\,\dd t
        \right),
\]
followed by Cauchy--Schwarz in expectation.  If \(\zeta(t)\) is \(\cG_T\)-measurable for almost
every \(t\), then \(M^{\mathrm{vis},\QQ}(t)\zeta(t)\) is \(\cG_T\)-measurable for almost every
\(t\), so \(S_{M,\times}^{\varpi,Q}\in L^2(\QQ,\cG_T)\), and the centered term vanishes by
definition.
\end{proof}

\begin{proposition}[Visible-screening split of the pure \(\varpi\)-benchmark defect]\label{prop:lambda-shadow-visible}
Under the notation of Proposition~\ref{prop:lambda-shadow-explicit}, define
\[
        B_{M,\zeta}^{\varpi,Q}
        :=
        \int_0^T \zeta(t)^2\,\dd t,
        \qquad
        B_{M,\eta}^{\varpi,Q}
        :=
        C_{\cL_T}^\QQ(R_{\eta,M}^\QQ),
\]
so that
\[
        B_{M}^{\varpi,Q}
        =
        B_{M,\zeta}^{\varpi,Q}
        +
        B_{M,\eta}^{\varpi,Q}.
\]
Write
\[
        \Pi_\varpi^\QQ(Y)
        :=
        C_{\cL_T}^\QQ(Y)
        -
        C_{\cG_T}^\QQ\!\left(C_{\cL_T}^\QQ(Y)\right).
\]
Define
\[
        \Delta_{M,\zeta}^{\varpi,\mathrm{vis}}
        :=
        \left\|
                \Pi_\varpi^\QQ\!\left(
                        C_{\cG_T}^\QQ(B_{M,\zeta}^{\varpi,Q})
                \right)
        \right\|_{L^2(\QQ)},
        \qquad
        \Delta_{M,\eta}^{\varpi,\mathrm{vis}}
        :=
        \left\|
                \Pi_\varpi^\QQ\!\left(
                        C_{\cG_T}^\QQ(B_{M,\eta}^{\varpi,Q})
                \right)
        \right\|_{L^2(\QQ)}.
\]
Then
\[
        \Delta_{M}^{\varpi,\mathrm{vis}}
        \le
        \Delta_{M,\zeta}^{\varpi,\mathrm{vis}}
        +
        \Delta_{M,\eta}^{\varpi,\mathrm{vis}}.
\]
Consequently, if
\[
        \Delta_{M,\zeta}^{\varpi,\mathrm{vis}}
        +
        \Delta_{M,\eta}^{\varpi,\mathrm{vis}}
        =
        o\!\left(\alpha_{Q_M}^{\QQ,\star}\right),
\]
then
\[
        \Delta_{M}^{\varpi,\mathrm{vis}}
        =
        o\!\left(\alpha_{Q_M}^{\QQ,\star}\right).
\]
If
\[
        \Delta_{M,\zeta}^{\varpi,\mathrm{vis}}=0,
        \qquad
        \Delta_{M,\eta}^{\varpi,\mathrm{vis}}=0,
\]
then
\[
        \Delta_{M}^{\varpi,\mathrm{vis}}=0.
\]
\end{proposition}

\begin{proof}
Because \(B_{M}^{\varpi,Q}=B_{M,\zeta}^{\varpi,Q}+B_{M,\eta}^{\varpi,Q}\), linearity of
\(\Pi_\varpi^\QQ\) gives
\[
        \Pi_\varpi^\QQ\!\left(
                C_{\cG_T}^\QQ(B_{M}^{\varpi,Q})
        \right)
        =
        \Pi_\varpi^\QQ\!\left(
                C_{\cG_T}^\QQ(B_{M,\zeta}^{\varpi,Q})
        \right)
        +
        \Pi_\varpi^\QQ\!\left(
                C_{\cG_T}^\QQ(B_{M,\eta}^{\varpi,Q})
        \right).
\]
Taking \(L^2(\QQ)\)-norms and applying the triangle inequality yields the displayed bound.  If the
normalized right-hand terms vanish, the left-hand defect vanishes asymptotically; if the component
terms are zero, the centered pure benchmark defect is zero.
\end{proof}

\begin{proposition}[Centered \(\varpi\)-benchmark split relative to the benchmark-generated generator]\label{prop:lambda-shadow-benchmark}
Under the notation of Proposition~\ref{prop:lambda-shadow-visible}, define the centered pieces
\[
        \widetilde B_{M,\zeta}^{\varpi,Q}
        :=
        B_{M,\zeta}^{\varpi,Q}
        -
        C_{\cG_T}^\QQ(B_{M,\zeta}^{\varpi,Q}),
        \qquad
        \widetilde B_{M,\eta}^{\varpi,Q}
        :=
        B_{M,\eta}^{\varpi,Q}
        -
        C_{\cG_T}^\QQ(B_{M,\eta}^{\varpi,Q}),
\]
so that
\[
        \widetilde B_{M}^{\varpi,Q}
        =
        \widetilde B_{M,\zeta}^{\varpi,Q}
        +
        \widetilde B_{M,\eta}^{\varpi,Q}.
\]
Define the canonical centered \(\varpi\)-benchmark coefficient and residual by
\[
        b_{M,\tau}^{\varpi,Q}
        :=
        \left\langle
                \widetilde B_{M}^{\varpi,Q},
                \Gamma_{M,\tau}^\varpi
        \right\rangle_{L^2(\QQ)},
        \qquad
        E_{M,\tau}^{\varpi,B}
        :=
        \widetilde B_{M}^{\varpi,Q}
        -
        b_{M,\tau}^{\varpi,Q}\Gamma_{M,\tau}^\varpi.
\]
Define the centered \(\varpi\)-benchmark generator defect by
\[
        \mathfrak d_{M,\tau}^{\varpi,B}
        :=
        \dist\!\left(
                \widetilde B_{M,\zeta}^{\varpi,Q},
                \mathrm{span}\{\Gamma_{M,\tau}^\varpi\}
        \right)
        +
        \dist\!\left(
                \widetilde B_{M,\eta}^{\varpi,Q},
                \mathrm{span}\{\Gamma_{M,\tau}^\varpi\}
        \right).
\]
Then
\[
        \dist\!\left(
                \widetilde B_{M}^{\varpi,Q},
                \mathrm{span}\{\Gamma_{M,\tau}^\varpi\}
        \right)
        \le
        \mathfrak d_{M,\tau}^{\varpi,B}.
\]
Moreover,
\[
        E_{M,\tau}^{\varpi,B}\in \cM_\QQ^\circ,
        \qquad
        E_{M,\tau}^{\varpi,B}
        \perp_{L^2(\QQ)}
        \Gamma_{M,\tau}^\varpi,
        \qquad
        \|E_{M,\tau}^{\varpi,B}\|_{L^2(\QQ)}
        =
        \dist\!\left(
                \widetilde B_{M}^{\varpi,Q},
                \mathrm{span}\{\Gamma_{M,\tau}^\varpi\}
        \right).
\]
In particular, if
\[
        \mathfrak d_{M,\tau}^{\varpi,B}
        =
        o\!\left(\alpha_{Q_M}^{\QQ,\star}\right),
\]
then
\[
        \dist\!\left(
                \widetilde B_{M}^{\varpi,Q},
                \mathrm{span}\{\Gamma_{M,\tau}^\varpi\}
        \right)
        =
        o\!\left(\alpha_{Q_M}^{\QQ,\star}\right).
\]
If
\[
        \mathfrak d_{M,\tau}^{\varpi,B}=0,
\]
then
\[
        \widetilde B_{M}^{\varpi,Q}\in \mathrm{span}\{\Gamma_{M,\tau}^\varpi\}.
\]
\end{proposition}

\begin{proof}
Centering the decomposition of \(B_{M}^{\varpi,Q}\) gives the displayed split for
\(\widetilde B_{M}^{\varpi,Q}\).  Applying the triangle inequality for the distance to the subspace
\(\mathrm{span}\{\Gamma_{M,\tau}^\varpi\}\) gives the displayed bound.  If the two distance terms
vanish asymptotically, the total distance does also; if both centered components lie in the span,
so does their sum.
\end{proof}

\begin{proposition}[Reduction of the centered \(\varpi\)-benchmark defect to the centered \(\varpi^2\) benchmark]\label{prop:lambda-shadow-lambda-square}
Under the notation of Propositions~\ref{prop:lambda-shadow-visible} and
\ref{prop:lambda-shadow-benchmark}, define the centered synchronized-square benchmark by
\[
        Q_{\varpi}^\QQ
        :=
        \int_0^T \varpi(t)^2\,\dd t,
        \qquad
        \widetilde Q_{\varpi}^\QQ
        :=
        Q_{\varpi}^\QQ-C_{\cG_T}^\QQ(Q_{\varpi}^\QQ).
\]
Then
\[
        \widetilde B_{M,\zeta}^{\varpi,Q}
        =
        (w^\top v)^2\widetilde Q_{\varpi}^\QQ,
\]
and therefore
\[
        \mathfrak d_{M,\tau}^{\varpi,B}
        =
        |w^\top v|^2
        \dist\!\left(
                \widetilde Q_{\varpi}^\QQ,
                \mathrm{span}\{\Gamma_{M,\tau}^\varpi\}
        \right)
        +
        \dist\!\left(
                \widetilde B_{M,\eta}^{\varpi,Q},
                \mathrm{span}\{\Gamma_{M,\tau}^\varpi\}
        \right).
\]
In particular, if
\[
        \widetilde B_{M,\eta}^{\varpi,Q}\in \mathrm{span}\{\Gamma_{M,\tau}^\varpi\},
\]
then
\[
        \mathfrak d_{M,\tau}^{\varpi,B}
        =
        |w^\top v|^2
        \dist\!\left(
                \widetilde Q_{\varpi}^\QQ,
                \mathrm{span}\{\Gamma_{M,\tau}^\varpi\}
        \right).
\]
If, more strongly,
\[
        B_{M,\eta}^{\varpi,Q}\in L^2(\QQ,\cG_T),
\]
then
\[
        \widetilde B_{M,\eta}^{\varpi,Q}=0
        \qquad
        \text{and hence}\qquad
        \mathfrak d_{M,\tau}^{\varpi,B}
        =
        |w^\top v|^2
        \dist\!\left(
                \widetilde Q_{\varpi}^\QQ,
                \mathrm{span}\{\Gamma_{M,\tau}^\varpi\}
        \right).
\]
\end{proposition}

\begin{proof}
Because \(\zeta(t)=(w^\top v)\varpi(t)\),
\[
        B_{M,\zeta}^{\varpi,Q}
        =
        \int_0^T \zeta(t)^2\,\dd t
        =
        (w^\top v)^2\int_0^T \varpi(t)^2\,\dd t
        =
        (w^\top v)^2 Q_{\varpi}^\QQ.
\]
Subtracting \(C_{\cG_T}^\QQ(\cdot)\) gives
\[
        \widetilde B_{M,\zeta}^{\varpi,Q}
        =
        (w^\top v)^2\widetilde Q_{\varpi}^\QQ.
\]
Since distance to a linear subspace is homogeneous,
\[
        \dist\!\left(
                \widetilde B_{M,\zeta}^{\varpi,Q},
                \mathrm{span}\{\Gamma_{M,\tau}^\varpi\}
        \right)
        =
        |w^\top v|^2
        \dist\!\left(
                \widetilde Q_{\varpi}^\QQ,
                \mathrm{span}\{\Gamma_{M,\tau}^\varpi\}
        \right).
\]
The displayed identity for \(\mathfrak d_{M,\tau}^{\varpi,B}\) then follows from the definition
in Proposition~\ref{prop:lambda-shadow-benchmark}.  If
\(\widetilde B_{M,\eta}^{\varpi,Q}\in \mathrm{span}\{\Gamma_{M,\tau}^\varpi\}\), the second
distance term vanishes, giving the first refinement.  If
\(B_{M,\eta}^{\varpi,Q}\in L^2(\QQ,\cG_T)\), then
\[
        \widetilde B_{M,\eta}^{\varpi,Q}
        =
        B_{M,\eta}^{\varpi,Q}-C_{\cG_T}^\QQ(B_{M,\eta}^{\varpi,Q})
        =
        0,
\]
which gives the exact final identity.
\end{proof}

\begin{proposition}[Centered \(\varpi^2\) benchmark coefficient-residual decomposition]\label{prop:lambda-shadow-lambda-square-coeff}
Under the notation of Proposition~\ref{prop:lambda-shadow-lambda-square}, define the centered
\(\varpi^2\)-benchmark coefficient and residual by
\[
        b_{M,\tau}^{\varpi,2}
        :=
        \left\langle
                \widetilde Q_{\varpi}^\QQ,
                \Gamma_{M,\tau}^\varpi
        \right\rangle_{L^2(\QQ)},
        \qquad
        E_{M,\tau}^{\varpi,2}
        :=
        \widetilde Q_{\varpi}^\QQ
        -
        b_{M,\tau}^{\varpi,2}\Gamma_{M,\tau}^\varpi.
\]
Then
\[
        E_{M,\tau}^{\varpi,2}\in \cM_\QQ^\circ,
        \qquad
        E_{M,\tau}^{\varpi,2}
        \perp_{L^2(\QQ)}
        \Gamma_{M,\tau}^\varpi,
        \qquad
        \|E_{M,\tau}^{\varpi,2}\|_{L^2(\QQ)}
        =
        \dist\!\left(
                \widetilde Q_{\varpi}^\QQ,
                \mathrm{span}\{\Gamma_{M,\tau}^\varpi\}
        \right).
\]
If, in addition,
\[
        B_{M,\eta}^{\varpi,Q}\in L^2(\QQ,\cG_T),
\]
then
\[
        \widetilde B_{M}^{\varpi,Q}
        =
        (w^\top v)^2\widetilde Q_{\varpi}^\QQ,
        \qquad
        b_{M,\tau}^{\varpi,Q}
        =
        (w^\top v)^2 b_{M,\tau}^{\varpi,2},
        \qquad
        E_{M,\tau}^{\varpi,B}
        =
        (w^\top v)^2 E_{M,\tau}^{\varpi,2},
\]
and
\[
        \mathfrak d_{M,\tau}^{\varpi,B}
        =
        |w^\top v|^2\|E_{M,\tau}^{\varpi,2}\|_{L^2(\QQ)}.
\]
\end{proposition}

\begin{proof}
Because \(\Gamma_{M,\tau}^\varpi\in \cM_\QQ^\circ\), the centered residual
\(E_{M,\tau}^{\varpi,2}\) belongs to \(\cM_\QQ^\circ\).  The coefficient definition gives
\[
        \left\langle
                E_{M,\tau}^{\varpi,2},
                \Gamma_{M,\tau}^\varpi
        \right\rangle_{L^2(\QQ)}
        =
        \left\langle
                \widetilde Q_{\varpi}^\QQ,
                \Gamma_{M,\tau}^\varpi
        \right\rangle_{L^2(\QQ)}
        -
        b_{M,\tau}^{\varpi,2}
        \|\Gamma_{M,\tau}^\varpi\|_{L^2(\QQ)}^2
        =
        0,
\]
where the last identity is trivial if \(\Gamma_{M,\tau}^\varpi=0\).  Hence
\(\|E_{M,\tau}^{\varpi,2}\|_{L^2(\QQ)}\) is the distance from
\(\widetilde Q_{\varpi}^\QQ\) to \(\mathrm{span}\{\Gamma_{M,\tau}^\varpi\}\).
If \(B_{M,\eta}^{\varpi,Q}\in L^2(\QQ,\cG_T)\), Proposition~\ref{prop:lambda-shadow-lambda-square}
gives
\[
        \widetilde B_{M}^{\varpi,Q}
        =
        \widetilde B_{M,\zeta}^{\varpi,Q}
        =
        (w^\top v)^2\widetilde Q_{\varpi}^\QQ.
\]
Taking inner products against \(\Gamma_{M,\tau}^\varpi\) yields
\[
        b_{M,\tau}^{\varpi,Q}
        =
        (w^\top v)^2 b_{M,\tau}^{\varpi,2},
\]
and subtracting the aligned component gives
\[
        E_{M,\tau}^{\varpi,B}
        =
        \widetilde B_{M}^{\varpi,Q}
        -
        b_{M,\tau}^{\varpi,Q}\Gamma_{M,\tau}^\varpi
        =
        (w^\top v)^2 E_{M,\tau}^{\varpi,2}.
\]
Taking norms and using Proposition~\ref{prop:lambda-shadow-benchmark} gives the final identity.
\end{proof}

\begin{assumption}[Leading-generator dominance of the centered \(\varpi\)-benchmark]\label{ass:lambda-benchmark-generator}
In the nondegenerate regime
\[
        b_{M,\tau}^{\varpi,Q}\neq0,
\]
the centered \(\varpi\)-benchmark generator defect from
Proposition~\ref{prop:lambda-shadow-benchmark} satisfies
\[
        \mathfrak d_{M,\tau}^{\varpi,B}
        =
        o\!\left(|b_{M,\tau}^{\varpi,Q}|\right)
        \qquad
        \text{as }\tau\downarrow0.
\]
In the exact branch, assume instead
\[
        \mathfrak d_{M,\tau}^{\varpi,B}=0
        \qquad
        \text{for all sufficiently small }\tau.
\]
\end{assumption}

\begin{corollary}[Residual-variance screened reduction of the benchmark-generator defect]\label{cor:lambda-shadow-lambda-square}
Assume the hypotheses of Proposition~\ref{prop:lambda-shadow-lambda-square-coeff}, and suppose
\[
        B_{M,\eta}^{\varpi,Q}\in L^2(\QQ,\cG_T),
        \qquad
        w^\top v\neq0.
\]
Assume further that
\[
        b_{M,\tau}^{\varpi,2}\neq0
        \qquad
        \text{for all sufficiently small }\tau.
\]
then Assumption~\ref{ass:lambda-benchmark-generator} follows from
\[
        \|E_{M,\tau}^{\varpi,2}\|_{L^2(\QQ)}
        =
        o\!\left(|b_{M,\tau}^{\varpi,2}|\right)
        \qquad
        \text{as }\tau\downarrow0.
\]
In the exact branch, it is enough that
\[
        E_{M,\tau}^{\varpi,2}=0
        \qquad
        \text{for all sufficiently small }\tau.
\]
\end{corollary}

\begin{proof}
Under \(B_{M,\eta}^{\varpi,Q}\in L^2(\QQ,\cG_T)\),
Proposition~\ref{prop:lambda-shadow-lambda-square-coeff} gives
\[
        \mathfrak d_{M,\tau}^{\varpi,B}
        =
        |w^\top v|^2\|E_{M,\tau}^{\varpi,2}\|_{L^2(\QQ)}
        \qquad
        \text{and}
        \qquad
        b_{M,\tau}^{\varpi,Q}
        =
        (w^\top v)^2 b_{M,\tau}^{\varpi,2}.
\]
Since \(w^\top v\neq0\), the nondegenerate regimes \(b_{M,\tau}^{\varpi,Q}\neq0\) and
\(b_{M,\tau}^{\varpi,2}\neq0\) coincide.  Therefore
\[
        \|E_{M,\tau}^{\varpi,2}\|_{L^2(\QQ)}
        =
        o\!\left(|b_{M,\tau}^{\varpi,2}|\right)
        \quad\Longrightarrow\quad
        \mathfrak d_{M,\tau}^{\varpi,B}
        =
        o\!\left(|b_{M,\tau}^{\varpi,Q}|\right),
\]
which is exactly Assumption~\ref{ass:lambda-benchmark-generator}.  If
\(E_{M,\tau}^{\varpi,2}=0\), then \(\mathfrak d_{M,\tau}^{\varpi,B}=0\), giving the exact
branch.
\end{proof}


\subsection{Scalar Gaussian seam and Companion~A source-family verification}\label{subsec:scalar-gaussian-seam}

The primitive Gaussian-Volterra source-family verification is long and mechanical.  It is kept in
Companion~A rather than in the main body.  This appendix records only the scalar
seam used by the option-surface branch.  Companion~A verifies this seam from raw Volterra
kernel primitives, Hilbert--Schmidt one-mode structure, leading-coordinate control, centered-square
comparison, and natural-scale lower-order defect bounds.

\begin{assumption}[Gaussian scalar seam inputs]\label{ass:gaussian-scalar-seam-package}
For the benchmark maturity family \(M\), assume the primitive Gaussian-Volterra hypotheses and
source-family estimates verified in Companion~A.  In scalar form, this means that
there exist common channel weights \(w=(w_j)_j\), a nonzero benchmark coefficient
\(B_M^{\mathrm{quad}}\), and a natural scale \(a_{M,\tau}^{\mathrm{quad}}>0\) such that the centered
quadratic option block has the shrinking-triangle expansion
\[
        \frac{1}{a_{M,\tau}^{\mathrm{quad}}}
        \left(
        \Gamma_{M,\tau}^{\QQ}-B_M^{\mathrm{quad}}
        \sum_j w_j Z_{j,\tau}
        \right)
        \longrightarrow 0
        \quad\text{in }L^2(\QQ)
\]
uniformly along the local maturity triangle.  The benchmark coefficient is nonzero and the
synchronized channel weights are the same weights selected by the Perron branch.
\end{assumption}

\begin{proposition}[Companion-verified scalar seam reduction]\label{thm:main-benchmark-model-to-seam}
Under Assumption~\ref{ass:gaussian-scalar-seam-package}, the Gaussian-Volterra benchmark produces
the scalar natural-scale seam needed by the local option-surface branch.  Equivalently, the
normalized centered quadratic profile has a nonzero common-channel limit in the Perron
coordinate, and every singular-value, scalar-comparison, and option-block defect is lower order at
that scale.
\end{proposition}

\begin{proof}
Companion~A proves the three primitive tasks.  First, the synchronized-amplitude coefficient,
the scalar quadratic comparison coefficient, and the benchmark-generated profile coefficient are
nonzero.  Second, the singular-value tail defect, scalar comparison defect, and combined option
block defect are all \(o(a_{M,\tau}^{\mathrm{quad}})\).  Third, the leading-coordinate map agrees
with the common Perron channel weights.  Substituting those three estimates into the displayed
scalar expansion in Assumption~\ref{ass:gaussian-scalar-seam-package} gives the claimed seam.  No
additional theorem-side wrapper is used in the main body.
\end{proof}

\begin{remark}[Companion source-family verification]\label{rem:companion-source-family-ledger}
The detailed source-family arguments that verify the primitive estimates live in the companion
technical note.  The appendix summary uses Assumption~\ref{ass:gaussian-scalar-seam-package},
Proposition~\ref{thm:main-benchmark-model-to-seam}, and the compact source-family verification below.
\end{remark}

\subsection{Gaussian--Volterra source-family verification}
\label{subsec:gaussian-public-route-ledger}

The companion Gaussian source-family note enters the compression argument through the compact assumption
and verification proposition below, followed by
Propositions~\ref{thm:main-common-channel-headline} and~\ref{thm:main-gaussian-option-factor}.  The
branch-by-branch Volterra estimates remain in Companion~A rather than in this article.

\begin{assumption}[Gaussian public source-family inputs]\label{ass:gaussian-lambda-public-frontier}
All derived results that invoke these inputs are read under the standing Gaussian-Volterra,
Perron-proxy, benchmark-generator, scalar-comparison, and option-defect hypotheses needed to
define the common-channel-weight structure.  In particular, the following named objects and
consequences are available:
\begin{enumerate}[label=(\roman*)]
\item the synchronized-amplitude coefficient, the scalar quadratic comparison coefficient, and
the benchmark-generated profile coefficient are the three leading coefficients for the
reduction;
\item the singular-value tail defect, the scalar comparison defect, and the combined option-block
defect are the three defect families whose lower-order estimates are measured at their natural
scales;
\item exact-branch implication templates use the corresponding exact vanishing of those defect
families, while asymptotic-branch templates use the displayed lower-order estimates stated in the
corollary itself or in the named source-family verification result being applied;
\item the non-temporal, temporal singular-value, combined source-family, visible-origin, lower-cell,
and denominator local-flatness branches are available only through their named verification or
terminal estimates;
\item the non-temporal, temporal singular-value, and combined source-family verification results are
available through Companion~A whenever their own hypotheses are separately
verified.
\end{enumerate}
This assumption records the notation and finite dependency list used by the appendix.
It does not by itself assert nonvanishing of
the three leading coefficients, lower-order smallness of the three defects, or exact-branch zero
defect identities.  Those conditions must be stated in the result that uses these inputs, or supplied by the companion
source-family verification.
\end{assumption}

\begin{remark}[Gaussian public source-family argument]\label{rem:gaussian-lambda-public-frontier}
The common-channel-weight reduction is used here as the input list for
Proposition~\ref{thm:main-common-channel-headline}.  Its three coefficients are the synchronized
amplitude coefficient, the scalar quadratic comparison coefficient, and the benchmark-generated
profile coefficient.  Its three defects are the singular-value tail defect, the scalar comparison
defect, and the combined option-block defect.  Once the stated coefficient lower bounds and
lower-order defect estimates are supplied, the public Gaussian reduction follows.

The source-family verifications are proved in Companion~A.  Here the grouped
argument is represented only by the compact verification proposition
Proposition~\ref{prop:gaussian-lambda-public-frontier-source-pair-dispatcher}.

On the model side, Companion~A verifies the raw Volterra singular-value chain: kernel
primitives, Hilbert--Schmidt approximation, one-mode and leading-coordinate structure, centered
square lift, scalar comparison profile control, and the natural-scale scalar seam.  The
visible-origin/lower-cell branch and the denominator local-flatness branch remain part of that
companion verification; in particular, the rescaled first-order fiber \(H^1\)-flat profile and the
\(L^2\)-variation predecessor are terminal estimates recorded there.
\end{remark}

\begin{proposition}[Companion Gaussian source-family verification]\label{prop:gaussian-lambda-public-frontier-source-pair-dispatcher}
Assume the source-family inputs in Assumption~\ref{ass:gaussian-lambda-public-frontier} and the scalar
seam of Proposition~\ref{thm:main-benchmark-model-to-seam}.  If a source-family argument in the companion
technical note verifies the source separation, lower-order defect, exact-branch, and quote-transfer
conditions required by those inputs, then it supplies the common-channel input required by the
Gaussian public argument.
\end{proposition}

\begin{proof}
The proposition is a dispatch statement from the companion verification to the present Gaussian
public argument.  The companion verification checks the primitive source-family conditions branch
by branch.  It proves the required nonzero source contribution, the natural-scale temporal estimate,
the non-cancellation condition, the public-compression coefficient, and the quote-transfer
remainder.  These are the inputs listed in
Assumption~\ref{ass:gaussian-lambda-public-frontier}, so the verified branch contributes the
common-channel input asserted here.
\end{proof}


\begin{proposition}[Companion-verified common-channel reduction]\label{thm:main-common-channel-headline}
Under Assumption~\ref{ass:gaussian-lambda-public-frontier}, the scalar seam of
Proposition~\ref{thm:main-benchmark-model-to-seam}, and a successful Companion~A source-family verification in
Proposition~\ref{prop:gaussian-lambda-public-frontier-source-pair-dispatcher}, the Gaussian public
argument has a nonzero common-channel coefficient at the local maturity anchor.  In particular, the
public option-facing coordinate is aligned with the Perron direction up to lower-order
residuals at the stated maturity scale.
\end{proposition}

\begin{proof}
The scalar seam supplies the normalized common-channel profile.  Assumption~\ref{ass:gaussian-lambda-public-frontier} fixes the
public objects and nonzero coefficient convention.  The companion verification supplies the
branch-specific estimates and defect controls.  The common-channel conclusion is the intersection of these three inputs.
\end{proof}


\begin{proposition}[Gaussian option-factor consequence]\label{thm:main-gaussian-option-factor}
In the Gaussian-Volterra/Bachelier specialization, suppose Assumption~\ref{ass:gaussian-lambda-public-frontier}
holds and the common-channel reduction of Proposition~\ref{thm:main-common-channel-headline}
has been verified by the companion source-family argument.  Then the option-facing factor coefficient inherits the same
Perron-aligned leading channel as the scalar seam.  Consequently the local Bachelier surface
coefficient is governed, to first order, by the compressed Perron stress coordinate, while
branch-specific singular-value tails, scalar-comparison errors, and option-block defects are lower
order.
\end{proposition}

\begin{proof}
Proposition~\ref{thm:main-common-channel-headline} identifies the nonzero common-channel coefficient.
The option-facing transfer map is linear at the local maturity anchor under the Gaussian/Bachelier
source-family hypotheses and its remainder is part of the lower-order option-block defect.  Applying that
transfer to the scalar seam preserves the leading Perron-aligned coefficient and absorbs all
controlled defects into the lower-order term.
\end{proof}


\begin{remark}[Companion location of option-defect source arguments]\label{rem:gaussian-option-defect-companion}
The combined option-block defect, direct option-side maturity-scale hypotheses, and Gaussian public
source-family checks are verified in Companion~A.  Only the consequence needed for quoted
surfaces is kept here.  Once the companion verification applies, the Bachelier/Gaussian
surface coefficient is governed by the common-channel Perron factor.
\end{remark}

\begin{corollary}[Gaussian-Volterra/Bachelier specialization of the non-Gaussian skew theorem]\label{cor:gaussian-market-skew}
Under the assumptions of Corollary~\ref{cor:bachelier-transfer}, let
\[
        \Gamma_{M,\tau}^\QQ
        :=
        C_{\cH_{M,\tau}^{\mathrm{opt}}}^\QQ(L_{M,\tau}^\QQ)
\]
be the compressed option-facing priced structural loading from
Proposition~\ref{prop:option-compression}.  Because
\(\cH_{M,\tau}^{\mathrm{opt}}\) is trivial, \(\Gamma_{M,\tau}^\QQ\) is deterministic.  Then, as
\(\tau\downarrow0\),
\[
        \text{if }\cE_M^{\mathrm{ret}}=\varnothing,
        \qquad
        \partial_\varepsilon
        \mathfrak S_{M,N}(\tau;\varepsilon)\big|_{\varepsilon=0}
        =
        \Gamma_{M,\tau}^\QQ
        =
        0;
\]
If \(\cE_M^{\mathrm{ret}}\neq\varnothing\) but
\(\mathrm{supp}^{\mathrm{comp}}(\Gamma_{M,\tau}^\QQ)=\varnothing\), then the Gaussian coefficient
system selects no nontrivial leading public power.  In the nontrivial case
\[
        \cE_M^{\mathrm{ret}}\neq\varnothing
        \quad\text{and}\quad
        \mathrm{supp}^{\mathrm{comp}}(\Gamma_{M,\tau}^\QQ)\neq\varnothing,
\]
assume in addition that the option-compression remainder from
Proposition~\ref{prop:option-compression} is negligible at the selected cumulant scale:
\[
        r_{M,\tau}=o\!\left(\tau^{\beta_M^{\mathrm{comp}}-1/2}\right).
\]
Then
\[
        \partial_\varepsilon
        \mathfrak S_{M,N}(\tau;\varepsilon)\big|_{\varepsilon=0}
        =
        \Gamma_{M,\tau}^\QQ
        +
        o\!\left(\tau^{\beta_M^{\mathrm{comp}}-1/2}\right),
\]
and
\[
        \Gamma_{M,\tau}^\QQ
        =
        \sum_{\beta\in\mathrm{supp}^{\mathrm{comp}}(\Gamma_{M,\tau}^\QQ)}
        A_M(\beta)\,\tau^{\beta-1/2}
        +
        o\!\left(\tau^{\beta_M^{\mathrm{comp}}-1/2}\right).
\]
Equivalently, when \(\alpha_{M,\tau}^\QQ>0\),
\[
        \alpha_{M,\tau}^\QQ
        C_{\cH_{M,\tau}^{\mathrm{opt}}}^\QQ(\Lambda_{M,\tau})
        =
        \sum_{\beta\in\mathrm{supp}^{\mathrm{comp}}(\Gamma_{M,\tau}^\QQ)}
        A_M(\beta)\,\tau^{\beta-1/2}
        +
        o\!\left(\tau^{\beta_M^{\mathrm{comp}}-1/2}\right).
\]
\end{corollary}

\begin{proof}
For every exponent value \(\beta\), define the aggregated compressed coefficient
\[
        A_M(\beta)
        :=
        \sum_{(i,j)\in\cE_M^{\mathrm{ret}}:\,H_{ij}=\beta} c_{ij}^M.
\]
Under Proposition~\ref{prop:market-kappa3} and Corollary~\ref{cor:bachelier-transfer}, the
Gaussian-Volterra block verifies Assumption~\ref{ass:cumulant-class} with these aggregated
coefficients and \(G_M(\tau)\to1\).  Together with the displayed scale-compatible
option-remainder condition, the conclusion is Theorem~\ref{thm:market-skew} specialized to that
closed-form coefficient system.  Grouping channels with the same exponent produces the aggregated
coefficients \(A_M(\beta)\), so same-exponent cancellation is already absorbed into
\(\mathrm{supp}^{\mathrm{comp}}(\Gamma_{M,\tau}^\QQ)\).
\end{proof}

\begin{theorem}[Market-scale skew visibility and screening for the perturbation-induced option-facing loading]\label{thm:market-visibility}
Under Corollary~\ref{cor:gaussian-market-skew}, the perturbation-induced contribution carried by the option-facing priced
structural loading falls into one of the following regimes
for the visible summary family \(\mathfrak S_{M,N}(\tau;\varepsilon)\), equivalently for the
compressed option-facing priced structural loading
\[
        \Gamma_{M,\tau}^\QQ
        =
        \alpha_{M,\tau}^\QQ
        C_{\cH_{M,\tau}^{\mathrm{opt}}}^\QQ(\Lambda_{M,\tau})
\]
whenever \(\alpha_{M,\tau}^\QQ>0\):
\begin{enumerate}[label=(\roman*)]
\item if \(\mathrm{supp}^{\mathrm{comp}}(\Gamma_{M,\tau}^\QQ)=\varnothing\), then the present
Gaussian coefficient system selects no nontrivial leading public power for the perturbation-induced
loading;
\item if \(\beta_M^{\mathrm{off}}<\beta_M^{\mathrm{diag}}\) and
\[
        A_M^{\mathrm{off}}
        =
        \sum_{i\neq j:\,H_{ij}=\beta_M^{\mathrm{off}}} c_{ij}^M
        \neq0,
\]
then the perturbation-induced structural contribution is \(\QQ\)-visible at the leading scale
\[
        a_\tau=\tau^{\beta_M^{\mathrm{off}}-1/2};
\]
\item if \(\beta_M^{\mathrm{diag}}<\beta_M^{\mathrm{off}}\) and \(A_M^{\mathrm{diag}}\neq0\), then
the public leading order is diagonal, and the off-diagonal contagion contribution is screened at
the public scale
\[
        a_\tau=\tau^{\beta_M^{\mathrm{diag}}-1/2};
\]
it can appear only at the lower-order scale \(\tau^{\beta_M^{\mathrm{off}}-1/2}\) if the
corresponding off-diagonal coefficient is nonzero;
\item if \(\beta_M^{\mathrm{diag}}=\beta_M^{\mathrm{off}}<\infty\), then leading-order visibility
depends on the total coefficient on the tied exponent
\[
        A_M^{\mathrm{tie}}
        :=
        \sum_{(i,j)\in\cE_M^{\mathrm{ret}}:\,H_{ij}=\beta_M}
        c_{ij}^M.
\]
\end{enumerate}
Equivalently, the leading maturity power is \(\beta_M^{\mathrm{comp}}\), the smallest exponent in
the aggregated compressed exponent support of
\(
        \alpha_{M,\tau}^\QQ C_{\cH_{M,\tau}^{\mathrm{opt}}}^\QQ(\Lambda_{M,\tau})
\),
and visibility at that power is decided by non-cancellation of the corresponding compressed
coefficient.  In the one-source/one-target limit, this reduces to the threshold logic of
\cite{vidal2026latent}.  Since \(\cH_{M,\tau}^{\mathrm{opt}}\) is trivial, this is equivalently the
statement that \(|\Gamma_{M,\tau}^\QQ|>0\) at the leading scale selected by the compressed priced
structural loading.
\end{theorem}

\begin{proof}
If \(\mathrm{supp}^{\mathrm{comp}}(\Gamma_{M,\tau}^\QQ)=\varnothing\), this is exactly the first
case.  Assume therefore that the aggregated compressed support is nonempty.
Corollary~\ref{cor:gaussian-market-skew} expresses the compressed priced structural loading
\(\Gamma_{M,\tau}^\QQ\) as a finite sum of powers of \(\tau\) grouped by exponent.  The smallest
active exponent in \(\mathrm{supp}^{\mathrm{comp}}(\Gamma_{M,\tau}^\QQ)\), equivalently the
smallest \(\beta\) with \(A_M(\beta)\neq0\), determines the public leading order.  If the smallest
active exponent is off-diagonal, visibility at that order requires non-cancellation of the
aggregate coefficient \(A_M^{\mathrm{off}}\).  If the smallest active exponent is diagonal and its
aggregate coefficient is nonzero, then the public summary is dominated by diagonal channels, so the
off-diagonal contagion contribution is screened at that leading scale.  If channels with the same
candidate leading exponent cancel, the effective public scale lifts to the next exponent in the
aggregated compressed support.  Restricting to a single incoming off-diagonal channel and a single
diagonal benchmark recovers the threshold mechanism of \cite{vidal2026latent}.
\end{proof}

\begin{corollary}[Market-index option summaries are visible but generally incomplete]\label{cor:market-index-option-summaries}
The market-index ATM normal skew family is a legitimate visible \(\QQ\)-summary for the
option-facing claim \(\Xi_{M,\tau}\), but it need not span the full perturbation-induced
structural contribution.  In
particular, under the hypotheses of Theorem~\ref{thm:market-visibility}(iii), the option-facing
structural loading \(L_{M,\tau}^\QQ\) is nonzero while its off-diagonal contagion contribution is
screened at the public leading scale.
\end{corollary}

\begin{proof}
Corollary~\ref{cor:gaussian-market-skew} gives the nontrivial option-facing priced loading, and
Theorem~\ref{thm:market-visibility} identifies when the leading public scale is diagonal rather
than off-diagonal.
\end{proof}

\section{Black and implied-volatility conventions}\label{sec:black}

This section records that the visibility classification proved for the normal/Bachelier quote is
unchanged by a smooth local passage to Black-implied-volatility conventions.  The first result
derives that passage directly from the option price surface by implicit differentiation; the second
records the abstract conversion-map version used by derived visibility summaries.  The local surface
language follows the Black--Scholes/Black convention and local-volatility surface tradition
\cite{blackScholes1973,black1976,dupire1994}.

\begin{theorem}[Black ATM skew conversion by implicit differentiation]\label{thm:self-contained-black}
Assume \(S_0^M>0\).  Let the same call price surface \(C_M(\tau,k;\varepsilon)\) be quoted both in
Bachelier volatility \(\sigma_N^M\) and in Black implied volatility \(\sigma_B^M\), with strike
\(K=S_0^M+k\).  Assume the baseline \(\varepsilon=0\) surface is regular near \(k=0\) and that the
Black vega is nonzero.  Suppose also that the perturbation is a pure skew perturbation at the
money, in the sense that
\[
        \partial_\varepsilon C_M(\tau,0;\varepsilon)\big|_{\varepsilon=0}=0 .
\]
Then
\[
        \partial_\varepsilon\partial_k\sigma_B^M(\tau,k;\varepsilon)
        \big|_{k=0,\varepsilon=0}
        =
        M_{B,N}(\tau)
        \partial_\varepsilon\partial_k\sigma_N^M(\tau,k;\varepsilon)
        \big|_{k=0,\varepsilon=0},
\]
where
\[
        M_{B,N}(\tau)
        :=
        \frac{\phi(0)}{S_0^M\phi(w_\tau/2)}
        =
        \frac{1}{S_0^M}\exp\!\left(\frac{w_\tau^2}{8}\right),
        \qquad
        w_\tau:=\sigma_B^M(\tau,0;0)\sqrt\tau .
\]
Here \(\phi\) is the standard normal density.
In particular, if \(w_\tau\to0\) as \(\tau\downarrow0\), then
\[
        M_{B,N}(\tau)\to \frac1{S_0^M}>0.
\]
Thus Black quoting preserves the leading exponent support and all visibility/screening
classifications carried by the Bachelier ATM skew derivative.
\end{theorem}

\begin{proof}
For fixed \(k\), implicit differentiation of the Bachelier and Black pricing equations gives
\[
        \partial_\varepsilon\sigma_N^M(\tau,k;\varepsilon)|_{\varepsilon=0}
        =
        \frac{\partial_\varepsilon C_M(\tau,k;\varepsilon)|_{\varepsilon=0}}{V_N(\tau,k)},
        \qquad
        \partial_\varepsilon\sigma_B^M(\tau,k;\varepsilon)|_{\varepsilon=0}
        =
        \frac{\partial_\varepsilon C_M(\tau,k;\varepsilon)|_{\varepsilon=0}}{V_B(\tau,k)},
\]
where, at the baseline,
\[
        V_N(\tau,k)=\sqrt\tau\,\phi\!\left(
        \frac{k}{\sigma_N^M(\tau,k;0)\sqrt\tau}\right),
\]
and
\[
        V_B(\tau,k)=S_0^M\sqrt\tau\,\phi(d_+(\tau,k)),
        \qquad
        d_+(\tau,0)=\frac{w_\tau}{2}.
\]
Therefore
\[
        \partial_\varepsilon\sigma_B^M(\tau,k;\varepsilon)|_{\varepsilon=0}
        =
        Q_\tau(k)\,
        \partial_\varepsilon\sigma_N^M(\tau,k;\varepsilon)|_{\varepsilon=0},
\]
with
\[
        Q_\tau(k)=\frac{V_N(\tau,k)}{V_B(\tau,k)}.
\]
At \(k=0\),
\[
        Q_\tau(0)=\frac{\phi(0)}{S_0^M\phi(w_\tau/2)}=M_{B,N}(\tau).
\]
The pure-skew condition implies
\[
        \partial_\varepsilon\sigma_N^M(\tau,0;\varepsilon)|_{\varepsilon=0}=0,
\]
because the Bachelier vega is nonzero.  Differentiating the relation
\(\partial_\varepsilon\sigma_B^M=Q_\tau\partial_\varepsilon\sigma_N^M\) in \(k\) at zero therefore
removes the possible \(Q_\tau'(0)\partial_\varepsilon\sigma_N^M(\tau,0)|_{\varepsilon=0}\) term and
yields the stated formula.  Positivity of the Bachelier and Black vegas gives
\(M_{B,N}(\tau)>0\); the assumed short-maturity limits of the quoting map give the stated positive
limit.
\end{proof}

\begin{corollary}[Black expansion for the primitive non-Gaussian Volterra class]\label{cor:black-ng-volterra}
Consider an option price surface whose return perturbation satisfies
Assumption~\ref{ass:ng-volterra-primitives} and whose Bachelier/Black quoting maps satisfy the
regularity hypotheses of Theorem~\ref{thm:self-contained-black}.  Suppose the associated Bachelier
quote satisfies the pure-skew condition at the money.  Then, in the nonempty support case of
Theorem~\ref{thm:ng-volterra-cumulants},
\[
        \partial_\varepsilon\partial_k\sigma_B^M(\tau,k;\varepsilon)
        \big|_{k=0,\varepsilon=0}
        =
        M_{B,N}(\tau)
        \left(
                \sum_{\beta\in\mathcal B_J}A_J(\beta)\tau^{\beta-1/2}
                +O(\tau^{\beta_0-1/2+\eta})
        \right).
\]
Whenever the selected-support condition of Theorem~\ref{thm:ng-volterra-cumulants} holds, the
Black expansion has the same selected exponent as the Bachelier expansion.
\end{corollary}

\begin{proof}
Corollary~\ref{cor:self-contained-bachelier-ng-volterra} supplies the Bachelier slope expansion on
the selected support.  Theorem~\ref{thm:self-contained-black} applies the local Black conversion
with multiplier \(G_B(\tau)\) and a lower-order Taylor remainder.  Multiplication by \(G_B(\tau)\)
therefore preserves the selected exponent and multiplies the coefficient family by the displayed
Black conversion factor.
\end{proof}

\begin{proposition}[Black-implied visibility is equivalent at leading order]\label{prop:black}
Assume in addition that \(S_0^M>0\) and that the market quotes admit a local normal-to-Black ATM
slope conversion map \(\Phi_\tau\) near the baseline, so that
\[
        \mathfrak S_{M,B}(\tau;\varepsilon)
        =
        \Phi_\tau\bigl(\mathfrak S_{M,N}(\tau;\varepsilon)\bigr).
\]
Assume that \(\Phi_\tau\) is differentiable at
\(\mathfrak S_{M,N}(\tau;0)\) with derivative \(G_M(\tau)\), where
\[
        G_M(\tau)\to G_M(0)>0,
\]
and that, in the perturbation regime considered here, its Taylor remainder satisfies
\[
        \Phi_\tau\bigl(\mathfrak S_{M,N}(\tau;\varepsilon)\bigr)
        -
        \Phi_\tau\bigl(\mathfrak S_{M,N}(\tau;0)\bigr)
        -
        G_M(\tau)
        \bigl(
                \mathfrak S_{M,N}(\tau;\varepsilon)
                -
                \mathfrak S_{M,N}(\tau;0)
        \bigr)
        =
        o\!\left(\varepsilon\,\tau^{\beta_M^{\mathrm{comp}}-1/2}\right).
\]
Assume also that
\[
        \mathrm{supp}^{\mathrm{comp}}(\Gamma_{M,\tau}^\QQ)\neq\varnothing,
        \qquad
        \text{so }
        \beta_M^{\mathrm{comp}}<\infty.
\]
Then
\[
        \mathfrak S_{M,B}(\tau;\varepsilon)-\mathfrak S_{M,B}(\tau;0)
        =
        G_M(\tau)\,
        \bigl(
                \mathfrak S_{M,N}(\tau;\varepsilon)-\mathfrak S_{M,N}(\tau;0)
        \bigr)
        +
        o\!\left(\varepsilon\,\tau^{\beta_M^{\mathrm{comp}}-1/2}\right),
\]
where \(\mathfrak S_{M,B}\) denotes the public ATM Black-implied volatility slope summary.
Consequently,
Theorem~\ref{thm:market-visibility} remains unchanged under Black quoting conventions.
\end{proposition}

\begin{proof}
The displayed Taylor expansion for \(\Phi_\tau\) gives the leading-order relation between the
Black and normal ATM slope increments.  The limiting positivity of \(G_M\) preserves visibility,
screening, and cancellation classifications.
\end{proof}

\begin{assumption}[Regular local Black surface window]\label{ass:black-window}
Let \(\cD\subset(0,T]\times\R\) be a compact maturity--moneyness window.  Let
\(C^\varepsilon(\tau,k)\) be an arbitrage-regular price surface, and let
\(\sigma_B^\varepsilon(\tau,k)\) be the corresponding Black implied-volatility surface.  Assume
\[
        C^\varepsilon=C^0+\varepsilon\dot C+\rho^\varepsilon,
        \qquad
        \|\rho^\varepsilon\|_{C^m(\cD)}=o(\varepsilon a_\tau),
\]
for the leading scale \(a_\tau\) being tracked.  Assume the baseline Black surface is regular on
\(\cD\):
\[
        V_B(\tau,k)
        :=
        \partial_\sigma C_{BS}(\tau,k,\sigma_B^0(\tau,k))
        >0,
\]
and the inverse Black map has uniformly bounded derivatives up to order \(m+1\) on a neighborhood
of \(C^0(\cD)\).  Equivalently, the window excludes local vega degeneracy.
\end{assumption}

\begin{theorem}[Local Black implied-volatility surface transfer]\label{thm:black-surface-transfer}
Under Assumption~\ref{ass:black-window},
\[
        \sigma_B^\varepsilon
        =
        \sigma_B^0+
        \varepsilon\dot\sigma_B+o_{C^m(\cD)}(\varepsilon a_\tau),
\]
where
\[
        \dot\sigma_B(\tau,k)
        =
        \frac{\dot C(\tau,k)}{V_B(\tau,k)}.
\]
Consequently, for every continuous linear surface-jet functional \(\ell:C^m(\cD)\to\R\),
\[
        \ell(\sigma_B^\varepsilon-\sigma_B^0)
        =
        \varepsilon\,\ell\!\left(\frac{\dot C}{V_B}\right)
        +
        o(\varepsilon a_\tau).
\]
If
\[
        \ell\!\left(\frac{\dot C}{V_B}\right)
        =
        c_\ell\tau^\beta+o(\tau^\beta),
        \qquad
        c_\ell\ne0,
\]
then the Black surface jet has the same leading visibility exponent \(\beta\).  If this coefficient
vanishes for a given \(\ell\), that Black jet does not select the exponent.
\end{theorem}

\begin{proof}
The Black price equation is
\[
        C^\varepsilon(\tau,k)=C_{BS}(\tau,k,\sigma_B^\varepsilon(\tau,k)).
\]
The Banach-space implicit function theorem applies on the regular window because vega is positive
and the inverse map has bounded derivatives.  Differentiating the identity at
\(\varepsilon=0\) gives
\[
        \dot C(\tau,k)=V_B(\tau,k)\dot\sigma_B(\tau,k),
\]
which proves the formula for \(\dot\sigma_B\).  The remainder estimate follows from the Taylor
remainder of the inverse map in \(C^m(\cD)\).  Applying \(\ell\) gives the jet statement and the
visibility conclusion.
\end{proof}

\begin{corollary}[Normal-to-Black surface transfer]\label{cor:normal-black-surface}
Suppose, on the same regular window, the price perturbation is represented through a normal implied
surface by
\[
        C^\varepsilon=C_N(\sigma_N^\varepsilon),
        \qquad
        \sigma_N^\varepsilon=
        \sigma_N^0+
        \varepsilon\dot\sigma_N+o_{C^m(\cD)}(\varepsilon a_\tau).
\]
Let
\[
        V_N(\tau,k):=\partial_\sigma C_N(\tau,k,\sigma_N^0(\tau,k)).
\]
Then
\[
        \dot\sigma_B
        =
        \frac{V_N}{V_B}\dot\sigma_N.
\]
Thus Black and normal visibility agree for every surface jet \(\ell\) for which the linear
operator
\[
        f\mapsto \ell\!\left(\frac{V_N}{V_B}f\right)
\]
does not annihilate the leading normal-surface coefficient.
\end{corollary}

\begin{proof}
Differentiating \(C_N(\sigma_N^\varepsilon)\) at \(\varepsilon=0\) gives
\(\dot C=V_N\dot\sigma_N\).  Substitute this into
Theorem~\ref{thm:black-surface-transfer}.
\end{proof}

\begin{remark}[Relation to the ATM-skew proposition]\label{rem:black-surface-atm}
Proposition~\ref{prop:black} is the one-dimensional jet
\(\ell(f)=\partial_k f(\tau,0)\) applied to Theorem~\ref{thm:black-surface-transfer}, after the
small-maturity vega ratio is expanded.  The pure-skew hypothesis is needed only when one wants a
single scalar exponent to be carried by that one jet.  The surface theorem transfers the whole
local jet vector on any regular pricing window.
\end{remark}

\bibliographystyle{unsrt}
\bibliography{references}

\end{document}
